% Generated mechanically from frozen input p04/ko/divisors.tex.
% Do not edit this generated file; edit the bound locale source instead.

% OK, start here.
%

\setcounter{chapter}{30}
\chapter{인자}



\phantomsection
\label{divisors-section-phantom}


\section{서론}
\label{divisors-section-introduction}

\noindent
이 장에서는 인자의 정의 등과 관련된 몇 가지 매우 기본적인 문제를
연구한다. 기본 참고문헌은 \cite{EGA}이다.



\section{연관점}
\label{divisors-section-associated}

\noindent
$R$을 환이라 하고 $M$을 $R$-가군이라 하자.
소 아이디얼 $\mathfrak p \subset R$가 $M$에 {\it 연관}된다는 것은
$M$의 어떤 원소의 소멸자가 $\mathfrak p$라는 뜻임을 상기하자.
Algebra, Definition \ref{algebra-definition-associated}을 보라.
다음은 \cite[IV Definition 3.1.1]{EGA}에 주어진 스킴 위의 준연접층에
대한 연관점의 정의이다.

\begin{definition}
\label{divisors-definition-associated}
$X$를 스킴이라 하자.
$\mathcal{F}$를 $X$ 위의 준연접층이라 하자.
\begin{enumerate}
\item $x \in X$가 $\mathcal{F}$에 {\it 연관}되어 있다고 함은 극대
아이디얼 $\mathfrak m_x$가 $\mathcal{O}_{X, x}$-가군
$\mathcal{F}_x$에 연관되어 있다는 뜻이다.
\item $\text{Ass}(\mathcal{F})$ 또는 $\text{Ass}_X(\mathcal{F})$로
$\mathcal{F}$의 연관점들의 집합을 나타낸다.
\item {\it $X$의 연관점}은 $\mathcal{O}_X$의 연관점이다.
\end{enumerate}
\end{definition}

\noindent
이 정의들은 $X$가 국소 뇌터이고 $\mathcal{F}$가 유한형일 때 가장
유용하다. 예를 들어 $X$의 한 기약 성분의 일반점이 $X$에 연관되지 않을
수도 있다. Example \ref{divisors-example-no-associated-prime}을 보라.
비뇌터인 경우에는 약연관점을 사용하는 편이 더 편리할 수 있다.
Section \ref{divisors-section-weakly-associated}를 보라.
스킴론적 개념과 아핀 열린집합 위의 대수적 개념을 연결해 보자. 이 대응은
국소 뇌터 스킴에 대해서만 완전하게 작동함을 주의하자.

\begin{lemma}
\label{divisors-lemma-associated-affine-open}
$X$를 스킴이라 하자. $\mathcal{F}$를 $X$ 위의 준연접층이라 하자.
$\Spec(A) = U \subset X$를 아핀 열린집합이라 하고
$M = \Gamma(U, \mathcal{F})$로 놓자.
$x \in U$라 하고 $\mathfrak p \subset A$를 이에 대응하는 소 아이디얼이라
하자.
\begin{enumerate}
\item $\mathfrak p$가 $M$에 연관되어 있으면 $x$는 $\mathcal{F}$에
연관되어 있다.
\item $\mathfrak p$가 유한 생성이면 그 역도 성립한다.
\end{enumerate}
특히 $X$가 국소 뇌터이면 위와 같은 모든 쌍 $(\mathfrak p, x)$에 대해
동치
$$
\mathfrak p \in \text{Ass}(M) \Leftrightarrow x \in \text{Ass}(\mathcal{F})
$$
가 성립한다.
\end{lemma}

\begin{proof}
이는 Algebra, Lemma \ref{algebra-lemma-associated-primes-localize}에서
따른다. 다음과 같이 직접 논증할 수도 있다.
$\mathfrak p$가 $M$에 연관되어 있다고 하자.
그러면 $m \in M$ 가운데 소멸자가 $\mathfrak p$인 것이 존재한다.
국소화가 완전하므로
$\mathfrak pA_{\mathfrak p}$는
$m/1 \in M_{\mathfrak p}$의 소멸자이다. 또한
$M_{\mathfrak p} = \mathcal{F}_x$이므로
(Schemes, Lemma \ref{schemes-lemma-spec-sheaves})
$x$가 $\mathcal{F}$에 연관되어 있다는 결론을 얻는다.

\medskip\noindent
역으로 $x$가 $\mathcal{F}$에 연관되어 있고
$\mathfrak p$가 유한 생성이라고 가정하자.
$x$가 $\mathcal{F}$에 연관되어 있으므로 소멸자가
$m' \in M_{\mathfrak p}$인 원소가 존재하고, 그 소멸자는
$\mathfrak pA_{\mathfrak p}$이다.
$m' = m/f$로 쓰되, 여기서 $f \in A$이고
$f \not \in \mathfrak p$이다. $I$를 $m$의 소멸자라 하자. 그러면 이는
$A$의 아이디얼이며 $IA_{\mathfrak p} = \mathfrak pA_{\mathfrak p}$를
만족한다. 따라서 $I \subset \mathfrak p$이고
$(\mathfrak p/I)_{\mathfrak p} = 0$이다.
$\mathfrak p$가 유한 생성이므로 어떤 $g \in A$,
$g \not \in \mathfrak p$가 존재하여 $g(\mathfrak p/I) = 0$이다.
따라서 $gm$의 소멸자는 $\mathfrak p$이고 원하는 결론을 얻는다.

\medskip\noindent
$X$가 국소 뇌터이면 $A$는 뇌터 환이고
(Properties, Lemma \ref{properties-lemma-locally-Noetherian})
$\mathfrak p$는 항상 유한 생성이다.
\end{proof}

\begin{lemma}
\label{divisors-lemma-ass-support}
$X$를 스킴이라 하자.
$\mathcal{F}$를 준연접 $\mathcal{O}_X$-가군이라 하자.
그러면 $\text{Ass}(\mathcal{F}) \subset \text{Supp}(\mathcal{F})$이다.
\end{lemma}

\begin{proof}
이는 정의에서 곧바로 따른다.
\end{proof}

\begin{lemma}
\label{divisors-lemma-ses-ass}
$X$를 스킴이라 하자.
$0 \to \mathcal{F}_1 \to \mathcal{F}_2 \to \mathcal{F}_3 \to 0$을
$X$ 위의 준연접층들의 짧은 완전열이라 하자.
그러면
$\text{Ass}(\mathcal{F}_2) \subset
\text{Ass}(\mathcal{F}_1) \cup \text{Ass}(\mathcal{F}_3)$
이고
$\text{Ass}(\mathcal{F}_1) \subset \text{Ass}(\mathcal{F}_2)$이다.
\end{lemma}

\begin{proof}
모든 점 $x \in X$에 대해 줄기들의 열
$0 \to \mathcal{F}_{1, x} \to \mathcal{F}_{2, x} \to \mathcal{F}_{3, x} \to 0$
은 $\mathcal{O}_{X, x}$-가군들의 짧은 완전열이다.
따라서 이 보조정리는 Algebra, Lemma \ref{algebra-lemma-ass}에서 따른다.
\end{proof}

\begin{lemma}
\label{divisors-lemma-finite-ass}
$X$를 국소 뇌터 스킴이라 하자.
$\mathcal{F}$를 연접 $\mathcal{O}_X$-가군이라 하자.
그러면 $\text{Ass}(\mathcal{F}) \cap U$는 모든 준콤팩트 열린집합
$U \subset X$에 대해 유한이다.
\end{lemma}

\begin{proof}
뇌터 환 위의 유한 가군의 연관 소 아이디얼들의 집합은 유한하므로 참이다.
Algebra, Lemma \ref{algebra-lemma-finite-ass}를 보라.
스킴의 언어에서 대수의 언어로 옮기려면 $U$가 아핀 열린집합들의 유한
합집합이고, 각 열린집합은 뇌터 환의 스펙트럼이며
(Properties, Lemma \ref{properties-lemma-locally-Noetherian}),
$\mathcal{F}$는 이 환 위의 유한 가군에 대응한다는 사실
(Cohomology of Schemes, Lemma \ref{coherent-lemma-coherent-Noetherian})과
마지막으로 Lemma \ref{divisors-lemma-associated-affine-open}을 사용한다.
\end{proof}

\begin{lemma}
\label{divisors-lemma-ass-zero}
$X$를 국소 뇌터 스킴이라 하자. $\mathcal{F}$를 준연접
$\mathcal{O}_X$-가군이라 하자. 그러면
$$
\mathcal{F} = 0 \Leftrightarrow \text{Ass}(\mathcal{F}) = \emptyset.
$$
\end{lemma}

\begin{proof}
$\mathcal{F} = 0$이면 정의에 의해
$\text{Ass}(\mathcal{F}) = \emptyset$이다. 역으로
$\text{Ass}(\mathcal{F}) = \emptyset$이면 Algebra, Lemma
\ref{algebra-lemma-ass-zero}에 의해 $\mathcal{F} = 0$이다.
스킴의 언어에서 대수의 언어로 옮기려면 임의의 아핀 열린집합에 제한한 뒤
Lemma \ref{divisors-lemma-associated-affine-open}을 사용한다.
\end{proof}

\begin{example}
\label{divisors-example-no-associated-prime}
$k$를 체라 하자. 환 $R = k[x_1, x_2, x_3, \ldots]/(x_i^2)$는 국소
멱영 극대 아이디얼 $\mathfrak m$을 갖는 국소환이다.
$R$에는 소멸자가 $\mathfrak m$인 원소가 존재하지 않는다.
따라서 $\text{Ass}(R) = \emptyset$이고, $X = \Spec(R)$은 연관점이 없는
스킴의 예이다.
\end{example}

\begin{lemma}
\label{divisors-lemma-restriction-injective-open-contains-ass}
$X$를 국소 뇌터 스킴이라 하자. $\mathcal{F}$를 준연접
$\mathcal{O}_X$-가군이라 하자. $U \subset X$가 열린집합이고
$\text{Ass}(\mathcal{F}) \subset U$이면
$\Gamma(X, \mathcal{F}) \to \Gamma(U, \mathcal{F})$는 단사이다.
\end{lemma}

\begin{proof}
$s \in \Gamma(X, \mathcal{F})$를 $U$에 제한하면 영이 되는 절단이라 하자.
$\mathcal{F}' \subset \mathcal{F}$를 사상
$\mathcal{O}_X \to \mathcal{F}$의 상이라 하자. 이 사상은 $s$가
정의한다. 그러면
$\text{Supp}(\mathcal{F}') \cap U = \emptyset$이다. 한편 Lemma
\ref{divisors-lemma-ses-ass}에 의해
$\text{Ass}(\mathcal{F}') \subset \text{Ass}(\mathcal{F})$이다. 또한
$\text{Ass}(\mathcal{F}') \subset \text{Supp}(\mathcal{F}')$이므로
(Lemma \ref{divisors-lemma-ass-support})
$\text{Ass}(\mathcal{F}') = \emptyset$이라는 결론을 얻는다.
따라서 Lemma \ref{divisors-lemma-ass-zero}에 의해 $\mathcal{F}' = 0$이다.
\end{proof}

\begin{lemma}
\label{divisors-lemma-minimal-support-in-ass}
$X$를 국소 뇌터 스킴이라 하자.
$\mathcal{F}$를 준연접 $\mathcal{O}_X$-가군이라 하자.
$x \in \text{Supp}(\mathcal{F})$를 $\mathcal{F}$의 지지에 속하면서
$\text{Supp}(\mathcal{F})$의 다른 점의 특수화가 아닌 점이라 하자.
그러면 $x \in \text{Ass}(\mathcal{F})$이다.
특히 $X$의 한 기약 성분의 모든 일반점은 $X$의 연관점이다.
\end{lemma}

\begin{proof}
$x \in \text{Supp}(\mathcal{F})$이므로 가군 $\mathcal{F}_x$는 영이
아니다. 따라서 Algebra, Lemma \ref{algebra-lemma-ass-zero}에 의해
$\text{Ass}(\mathcal{F}_x) \subset \Spec(\mathcal{O}_{X, x})$는
공집합이 아니다. 한편 가정에 의해
$\text{Supp}(\mathcal{F}_x) = \{\mathfrak m_x\}$이다.
또
$\text{Ass}(\mathcal{F}_x) \subset \text{Supp}(\mathcal{F}_x)$이므로
(Algebra, Lemma \ref{algebra-lemma-ass-support})
$\mathfrak m_x$가 $\mathcal{F}_x$에 연관되어 있음을 알 수 있고 원하는
결론을 얻는다.
\end{proof}

\noindent
다음 보조정리는 More on Algebra, Lemma
\ref{more-algebra-lemma-check-injective-on-ass}의 유사물이다.

\begin{lemma}
\label{divisors-lemma-check-injective-on-ass}
$X$를 국소 뇌터 스킴이라 하자.
$\varphi : \mathcal{F} \to \mathcal{G}$를 준연접
$\mathcal{O}_X$-가군들의 사상이라 하자.
모든 $x \in X$에 대해 다음 중 적어도 하나가 성립한다고 가정하자.
\begin{enumerate}
\item $\mathcal{F}_x \to \mathcal{G}_x$가 단사이거나,
\item $x \not \in \text{Ass}(\mathcal{F})$이다.
\end{enumerate}
그러면 $\varphi$는 단사이다.
\end{lemma}

\begin{proof}
가정에 의해 $\text{Ass}(\Ker(\varphi)) = \emptyset$이고, 따라서 Lemma
\ref{divisors-lemma-ass-zero}에 의해 $\Ker(\varphi) = 0$이다.
\end{proof}

\begin{lemma}
\label{divisors-lemma-check-isomorphism-via-depth-and-ass}
$X$를 국소 뇌터 스킴이라 하자.
$\varphi : \mathcal{F} \to \mathcal{G}$를 준연접
$\mathcal{O}_X$-가군들의 사상이라 하자. $\mathcal{F}$가 연접이고 모든
$x \in X$에 대해 다음 중 하나가 성립한다고 가정하자.
\begin{enumerate}
\item $\mathcal{F}_x \to \mathcal{G}_x$가 동형사상이거나,
\item $\text{depth}(\mathcal{F}_x) \geq 2$이고
$x \not \in \text{Ass}(\mathcal{G})$이다.
\end{enumerate}
그러면 $\varphi$는 동형사상이다.
\end{lemma}

\begin{proof}
이는 More on Algebra, Lemma
\ref{more-algebra-lemma-check-isomorphism-via-depth-and-ass}를 스킴의 언어로
옮긴 것이다.
\end{proof}



\section{사상과 연관점}
\label{divisors-section-morphisms-associated}

\noindent
$f : X \to S$를 스킴의 사상이라 하자.
$\mathcal{F}$를 $\mathcal{O}_X$-가군층이라 하자.
$s \in S$가 점이면, $\mathcal{F}_s$로 $\mathcal{O}_{X_s}$-가군을
나타내는 것이 흔히 편리하다. 이는 $\mathcal{F}$를 사상
$i_s : X_s \to X$에 따라 당겨 얻는다. 여기서 $X_s$는 $f$의 $s$ 위
스킴론적 올이다. 식으로 쓰면
$$
\mathcal{F}_s = i_s^*\mathcal{F}
$$
이다. 물론 이 표기는 $\mathcal{F}$의 점 $x \in X$에서의 줄기에 이미
쓰이는 표기와 충돌하는데, 이는 $f = \text{id}_X$일 때 발생한다. 그러나 다음
보조정리의 서술에서처럼 이 표기가 편리한 경우가 많다.

\begin{lemma}
\label{divisors-lemma-bourbaki}
$f : X \to S$를 스킴의 사상이라 하자.
$\mathcal{F}$를 $X$ 위의 준연접층으로서 $S$ 위 평탄인 것이라 하자.
$\mathcal{G}$를 $S$ 위의 준연접층이라 하자. 그러면
$$
\text{Ass}_X(\mathcal{F} \otimes_{\mathcal{O}_X} f^*\mathcal{G})
\supset
\bigcup\nolimits_{s \in \text{Ass}_S(\mathcal{G})}
\text{Ass}_{X_s}(\mathcal{F}_s)
$$
이고, $S$가 국소 뇌터이면 등호가 성립한다($\mathcal{F}_s$의 표기는
위에서 설명하였다).
\end{lemma}

\begin{proof}
$x \in X$라 하고 $s = f(x) \in S$라 하자.
$B = \mathcal{O}_{X, x}$, $A = \mathcal{O}_{S, s}$,
$N = \mathcal{F}_x$, $M = \mathcal{G}_s$로 놓자.
$\mathcal{F} \otimes_{\mathcal{O}_X} f^*\mathcal{G}$의 $x$에서의 줄기는
$B$-가군 $M \otimes_A N$과 같음을 주의하자. 따라서
$x \in \text{Ass}_X(\mathcal{F} \otimes_{\mathcal{O}_X} f^*\mathcal{G})$
일 필요충분조건은
$\mathfrak m_B$가 $\text{Ass}_B(M \otimes_A N)$에 속하는 것이다.
마찬가지로 $s \in \text{Ass}_S(\mathcal{G})$이고
$x \in \text{Ass}_{X_s}(\mathcal{F}_s)$일 필요충분조건은
$\mathfrak m_A \in \text{Ass}_A(M)$이고
$\mathfrak m_B/\mathfrak m_A B \in
\text{Ass}_{B \otimes \kappa(\mathfrak m_A)}(N \otimes \kappa(\mathfrak m_A))$
인 것이다. 그러므로 이 보조정리는 Algebra, Lemma
\ref{algebra-lemma-bourbaki-fibres}에서 따른다.
\end{proof}



\section{매장점}
\label{divisors-section-embedded}

\noindent
$R$을 환이라 하고 $M$을 $R$-가군이라 하자.
소 아이디얼 $\mathfrak p \subset R$가 $M$의 {\it 매장된 연관 소아이디얼}이라는
것은 그것이 $M$의 연관 소아이디얼이지만 $M$의 연관 소아이디얼들 가운데
극소가 아니라는 뜻임을 상기하자. Algebra, Definition
\ref{algebra-definition-embedded-primes}을 보라. 다음은
\cite[IV Definition 3.1.1]{EGA}에 주어진 스킴 위의 준연접층에 대한
매장된 연관점의 정의이다.

\begin{definition}
\label{divisors-definition-embedded}
$X$를 스킴이라 하자.
$\mathcal{F}$를 $X$ 위의 준연접층이라 하자.
\begin{enumerate}
\item $\mathcal{F}$의 {\it 매장된 연관점}이란 $\mathcal{F}$의 연관점들
가운데 극대가 아닌 연관점, 즉 $\mathcal{F}$의 다른 연관점의 특수화인
점을 말한다.
\item $x$를 $X$의 점이라 하자. $x$가 $\mathcal{O}_X$의 매장된
연관점이면 {\it 매장점}이라 한다.
\item $X$의 {\it 매장 성분}이란 기약 닫힌집합
$Z = \overline{\{x\}}$인데, 여기서 $x$는 $X$의 매장점인 것을 말한다.
\end{enumerate}
\end{definition}

\noindent
뇌터인 경우 $\mathcal{F}$가 연접이면 다음이 성립한다.

\begin{lemma}
\label{divisors-lemma-embedded}
$X$를 국소 뇌터 스킴이라 하자.
$\mathcal{F}$를 연접 $\mathcal{O}_X$-가군이라 하자. 그러면
\begin{enumerate}
\item $\text{Supp}(\mathcal{F})$의 기약 성분들의 일반점은
$\mathcal{F}$의 연관점이고,
\item $\mathcal{F}$의 한 연관점이 매장되어 있을 필요충분조건은 그것이
$\text{Supp}(\mathcal{F})$의 한 기약 성분의 일반점이 아닌 것이다.
\end{enumerate}
특히 $X$의 매장점은 $X$의 한 기약 성분의 일반점이 아닌 $X$의
연관점이다.
\end{lemma}

\begin{proof}
이 경우 $Z = \text{Supp}(\mathcal{F})$가 닫혀 있음을 상기하자. Morphisms,
Lemma \ref{morphisms-lemma-support-finite-type}을 보라. 또한 $Z$의 기약
성분들의 일반점은 $\mathcal{F}$의 연관점이다. Lemma
\ref{divisors-lemma-minimal-support-in-ass}을 보라. 끝으로 Lemma
\ref{divisors-lemma-ass-support}에 의해 $\text{Ass}(\mathcal{F}) \subset Z$이다.
이 결과들을 $Z$가 소버 위상공간이고 따라서 $Z$의 모든 점이 $Z$의 한
일반점의 특수화라는 사실과 결합하면 (1)과 (2)를 얻는다.
\end{proof}

\begin{lemma}
\label{divisors-lemma-S1-no-embedded}
$X$를 국소 뇌터 스킴이라 하자.
$\mathcal{F}$를 $X$ 위의 연접층이라 하자. 다음은 동치이다.
\begin{enumerate}
\item $\mathcal{F}$에는 매장된 연관점이 없고,
\item $\mathcal{F}$는 성질 $(S_1)$을 갖는다.
\end{enumerate}
\end{lemma}

\begin{proof}
이는 Algebra, Lemma \ref{algebra-lemma-criterion-no-embedded-primes}와 위의
Lemma \ref{divisors-lemma-associated-affine-open}을 결합한 것이다.
\end{proof}

\begin{lemma}
\label{divisors-lemma-noetherian-dim-1-CM-no-embedded-points}
$X$를 차원이 $\leq 1$인 국소 뇌터 스킴이라 하자. 다음은 동치이다.
\begin{enumerate}
\item $X$는 Cohen--Macaulay이고,
\item $X$에는 매장점이 없다.
\end{enumerate}
\end{lemma}

\begin{proof}
Lemma \ref{divisors-lemma-S1-no-embedded}과 정의들에서 따른다.
\end{proof}

\begin{lemma}
\label{divisors-lemma-scheme-theoretically-dense-contain-embedded-points}
$X$를 국소 뇌터 스킴이라 하자. $U \subset X$를 열린 부분스킴이라 하자.
다음은 동치이다.
\begin{enumerate}
\item $U$는 $X$에서 스킴론적으로 조밀하고
(Morphisms, Definition
\ref{morphisms-definition-scheme-theoretically-dense}),
\item $U$는 $X$에서 조밀하며 $U$는 $X$의 모든 매장점을 포함하고,
\item $U$는 $X$의 모든 연관점을 포함한다.
\end{enumerate}
\end{lemma}

\begin{proof}
문제는 $X$ 위에서 국소적이므로 $X = \Spec(A)$이고 $A$가 뇌터 환이라고
가정해도 된다. 그러면 $U$는 준콤팩트이고
(Properties, Lemma \ref{properties-lemma-immersion-into-noetherian}), 따라서
$U = D(f_1) \cup \ldots \cup D(f_n)$이다
(Algebra, Lemma \ref{algebra-lemma-qc-open}). 이 상황에서 $U$가 $X$에서
스킴론적으로 조밀할 필요충분조건은
$A \to A_{f_1} \times \ldots \times A_{f_n}$이 단사인 것이다. Morphisms,
Example \ref{morphisms-example-scheme-theoretic-closure}를 보라. 조건 (3)을
대수로 옮기면 $\mathfrak p$를 $A$의 임의 연관 소아이디얼이라 할 때
$i$가 존재하여 $f_i \not \in \mathfrak p$라는 뜻이다.

\medskip\noindent
(1), 즉 $A \to A_{f_1} \times \ldots \times A_{f_n}$이 단사라고 가정하자.
$a \in A$의 소멸자가 소 아이디얼 $\mathfrak p$라고 하자. $a$가 어떤
$A_{f_i}$의 영이 아닌 원소로 가게 되는 $i$가 존재하므로
$f_i \not \in \mathfrak p$이다. 따라서 (3)이 성립한다.

\medskip\noindent
$U$가 $\Spec(A)$에서 조밀할 필요충분조건은 모든 극소 소아이디얼을
포함하는 것임을 주의하자. $A$의 연관 소아이디얼들의 집합은 극소
소아이디얼들의 집합과 매장된 연관 소아이디얼들의 집합의 합집합이므로
(Algebra, Proposition
\ref{algebra-proposition-minimal-primes-associated-primes}), (2)와 (3)은
동치이다.

\medskip\noindent
(3)을 가정하자. 이는 임의의 연관 소아이디얼 $\mathfrak p$가 $A$의
것일 때 $A_{f_i}$의 한 소아이디얼에 대응하게 되는 $i$가 존재한다는
뜻이다. 그러면 $A \to A_{f_1} \times \ldots \times A_{f_n}$은 단사이다.
실제로 $A \to \prod_{\mathfrak p \in \text{Ass}(A)} A_\mathfrak p$가
Algebra, Lemma \ref{algebra-lemma-zero-at-ass-zero}에 의해 단사이다. 이로써 (3)이
(1)을 함의함을 알 수 있다.
\end{proof}

\begin{lemma}
\label{divisors-lemma-remove-embedded-points}
$X$를 국소 뇌터 스킴이라 하자.
$\mathcal{F}$를 $X$ 위의 연접층이라 하자. 다음 연접 부분층들의 집합은
$$
\{
\mathcal{K} \subset \mathcal{F}
\mid
\text{Supp}(\mathcal{K})\text{ is nowhere dense in }\text{Supp}(\mathcal{F})
\}
$$
극대 원소 $\mathcal{K}$를 갖는다. $\mathcal{F}' = \mathcal{F}/\mathcal{K}$로
놓으면 다음이 성립한다.
\begin{enumerate}
\item $\text{Supp}(\mathcal{F}') = \text{Supp}(\mathcal{F})$,
\item $\mathcal{F}'$에는 매장된 연관점이 없고,
\item 조밀한 열린집합 $U \subset X$가 존재하여
$U \cap \text{Supp}(\mathcal{F})$는 $\text{Supp}(\mathcal{F})$에서 조밀하고
$\mathcal{F}'|_U \cong \mathcal{F}|_U$이다.
\end{enumerate}
\end{lemma}

\begin{proof}
이는 Algebra, Lemmas \ref{algebra-lemma-remove-embedded-primes}와
\ref{algebra-lemma-remove-embedded-primes-localize}에서 따른다.
$U$는 $\mathcal{F}$의 매장된 연관점들의 집합의 폐포의 여집합으로 택할 수
있음을 주의하자.
\end{proof}

\begin{lemma}
\label{divisors-lemma-no-embedded-points-endos}
$X$를 국소 뇌터 스킴이라 하자.
$\mathcal{F}$를 매장된 연관점이 없는 연접 $\mathcal{O}_X$-가군이라 하자.
다음과 같이 놓자.
$$
\mathcal{I}
=
\Ker(\mathcal{O}_X
\longrightarrow
\SheafHom_{\mathcal{O}_X}(\mathcal{F}, \mathcal{F})).
$$
이는 매장점이 없는 닫힌 부분스킴 $Z \subset X$를 정의하는 연접
아이디얼층이다. 또한 연접층 $\mathcal{G}$가 $Z$ 위에 존재하여
(a) $\mathcal{F} = (Z \to X)_*\mathcal{G}$이고,
(b) $\mathcal{G}$에는 매장된 연관점이 없으며,
(c) $\text{Supp}(\mathcal{G}) = Z$이다(집합으로서).
\end{lemma}

\begin{proof}
몇몇 명제는 Cohomology of Schemes, Lemma
\ref{coherent-lemma-coherent-support-closed}의 증명에서 이미 보았다.
나머지는 Algebra, Lemma
\ref{algebra-lemma-no-embedded-primes-endos}에서 따른다.
\end{proof}



\section{약연관점}
\label{divisors-section-weakly-associated}

\noindent
$R$을 환이라 하고 $M$을 $R$-가군이라 하자.
소 아이디얼 $\mathfrak p \subset R$가 $M$에 {\it 약연관}된다는 것은
어떤 원소 $m$이 $M$ 안에 존재하여, $\mathfrak p$가 $m$의 소멸자를
포함하는 소 아이디얼들 가운데 극소라는 뜻임을 상기하자. Algebra,
Definition \ref{algebra-definition-weakly-associated}을 보라.
$R$이 극대 아이디얼 $\mathfrak m$을 갖는 국소환이면,
$\mathfrak m$이 $M$에 약연관될 필요충분조건은 원소 $m \in M$이
존재하여 그 소멸자의 근기가 $\mathfrak m$인 것이다. Algebra, Lemma
\ref{algebra-lemma-weakly-ass-local}을 보라.

\begin{definition}
\label{divisors-definition-weakly-associated}
$X$를 스킴이라 하자.
$\mathcal{F}$를 $X$ 위의 준연접층이라 하자.
\begin{enumerate}
\item $x \in X$가 $\mathcal{F}$에 {\it 약연관}되어 있다고 함은 극대
아이디얼 $\mathfrak m_x$가 $\mathcal{O}_{X, x}$-가군
$\mathcal{F}_x$에 약연관되어 있다는 뜻이다.
\item $\text{WeakAss}(\mathcal{F})$로 $\mathcal{F}$의 약연관점들의
집합을 나타낸다.
\item {\it $X$의 약연관점}은 $\mathcal{O}_X$의 약연관점이다.
\end{enumerate}
\end{definition}

\noindent
이 경우 임의의 아핀 열린집합 위에서 위 개념은 앞서 정의한 약연관
소아이디얼과 정확히 대응한다. 정확한 명제는 다음과 같다.

\begin{lemma}
\label{divisors-lemma-weakly-associated-affine-open}
$X$를 스킴이라 하자. $\mathcal{F}$를 $X$ 위의 준연접층이라 하자.
$\Spec(A) = U \subset X$를 아핀 열린집합이라 하고
$M = \Gamma(U, \mathcal{F})$로 놓자.
$x \in U$라 하고 $\mathfrak p \subset A$를 이에 대응하는 소
아이디얼이라 하자. 다음은 동치이다.
\begin{enumerate}
\item $\mathfrak p$는 $M$에 약연관되고,
\item $x$는 $\mathcal{F}$에 약연관된다.
\end{enumerate}
\end{lemma}

\begin{proof}
Algebra, Lemma \ref{algebra-lemma-weakly-ass-local}에서 따른다.
\end{proof}

\begin{lemma}
\label{divisors-lemma-weakly-ass-support}
$X$를 스킴이라 하자.
$\mathcal{F}$를 준연접 $\mathcal{O}_X$-가군이라 하자. 그러면
$$
\text{Ass}(\mathcal{F}) \subset \text{WeakAss}(\mathcal{F}) \subset
\text{Supp}(\mathcal{F}).
$$
\end{lemma}

\begin{proof}
이는 정의들에서 곧바로 따른다.
\end{proof}

\begin{lemma}
\label{divisors-lemma-ses-weakly-ass}
$X$를 스킴이라 하자.
$0 \to \mathcal{F}_1 \to \mathcal{F}_2 \to \mathcal{F}_3 \to 0$을
$X$ 위의 준연접층들의 짧은 완전열이라 하자. 그러면
$\text{WeakAss}(\mathcal{F}_2) \subset
\text{WeakAss}(\mathcal{F}_1) \cup \text{WeakAss}(\mathcal{F}_3)$
이고
$\text{WeakAss}(\mathcal{F}_1) \subset \text{WeakAss}(\mathcal{F}_2)$이다.
\end{lemma}

\begin{proof}
모든 점 $x \in X$에 대해 줄기들의 열
$0 \to \mathcal{F}_{1, x} \to \mathcal{F}_{2, x} \to \mathcal{F}_{3, x} \to 0$
은 $\mathcal{O}_{X, x}$-가군들의 짧은 완전열이다. 따라서 이 보조정리는
Algebra, Lemma \ref{algebra-lemma-weakly-ass}에서 따른다.
\end{proof}

\begin{lemma}
\label{divisors-lemma-weakly-ass-zero}
$X$를 스킴이라 하자.
$\mathcal{F}$를 준연접 $\mathcal{O}_X$-가군이라 하자. 그러면
$$
\mathcal{F} = (0) \Leftrightarrow \text{WeakAss}(\mathcal{F}) = \emptyset
$$
이다.
\end{lemma}

\begin{proof}
Lemma \ref{divisors-lemma-weakly-associated-affine-open}과 Algebra, Lemma
\ref{algebra-lemma-weakly-ass-zero}에서 따른다.
\end{proof}

\begin{lemma}
\label{divisors-lemma-restriction-injective-open-contains-weakly-ass}
$X$를 스킴이라 하자. $\mathcal{F}$를 준연접
$\mathcal{O}_X$-가군이라 하자. $U \subset X$가 열린집합이고
$\text{WeakAss}(\mathcal{F}) \subset U$이면
$\Gamma(X, \mathcal{F}) \to \Gamma(U, \mathcal{F})$는 단사이다.
\end{lemma}

\begin{proof}
$s \in \Gamma(X, \mathcal{F})$를 $U$에 제한하면 영이 되는 절단이라 하자.
$\mathcal{F}' \subset \mathcal{F}$를 사상
$\mathcal{O}_X \to \mathcal{F}$의 상이라 하자. 이 사상은 $s$가 정의한다.
그러면 $\text{Supp}(\mathcal{F}') \cap U = \emptyset$이다. 한편 Lemma
\ref{divisors-lemma-ses-weakly-ass}에 의해
$\text{WeakAss}(\mathcal{F}') \subset \text{WeakAss}(\mathcal{F})$이다.
또한 $\text{WeakAss}(\mathcal{F}') \subset \text{Supp}(\mathcal{F}')$이므로
(Lemma \ref{divisors-lemma-weakly-ass-support})
$\text{WeakAss}(\mathcal{F}') = \emptyset$이라는 결론을 얻는다. 따라서
Lemma \ref{divisors-lemma-weakly-ass-zero}에 의해 $\mathcal{F}' = 0$이다.
\end{proof}

\begin{lemma}
\label{divisors-lemma-minimal-support-in-weakly-ass}
$X$를 스킴이라 하자.
$\mathcal{F}$를 준연접 $\mathcal{O}_X$-가군이라 하자.
$x \in \text{Supp}(\mathcal{F})$를 $\mathcal{F}$의 지지에 속하면서
$\text{Supp}(\mathcal{F})$의 다른 점의 특수화가 아닌 점이라 하자. 그러면
$x \in \text{WeakAss}(\mathcal{F})$이다. 특히 $X$의 한 기약 성분의
임의의 일반점은 $\mathcal{O}_X$에 약연관된다.
\end{lemma}

\begin{proof}
$x \in \text{Supp}(\mathcal{F})$이므로 가군 $\mathcal{F}_x$는 영이 아니다.
따라서 Algebra, Lemma \ref{algebra-lemma-weakly-ass-zero}에 의해
$\text{WeakAss}(\mathcal{F}_x) \subset \Spec(\mathcal{O}_{X, x})$는
공집합이 아니다. 한편 가정에 의해
$\text{Supp}(\mathcal{F}_x) = \{\mathfrak m_x\}$이다. 또한
$\text{WeakAss}(\mathcal{F}_x) \subset \text{Supp}(\mathcal{F}_x)$이므로
(Algebra, Lemma \ref{algebra-lemma-weakly-ass-support})
$\mathfrak m_x$가 $\mathcal{F}_x$에 약연관되어 있음을 알 수 있고 원하는
결론을 얻는다.
\end{proof}

\begin{lemma}
\label{divisors-lemma-ass-weakly-ass}
$X$를 스킴이라 하자.
$\mathcal{F}$를 준연접 $\mathcal{O}_X$-가군이라 하자.
$\mathfrak m_x$가 $\mathcal{O}_{X, x}$의 유한 생성 아이디얼이면
$$
x \in \text{Ass}(\mathcal{F}) \Leftrightarrow
x \in \text{WeakAss}(\mathcal{F}).
$$
이다. 특히 $X$가 국소 뇌터이면
$\text{Ass}(\mathcal{F}) = \text{WeakAss}(\mathcal{F})$이다.
\end{lemma}

\begin{proof}
Algebra, Lemma \ref{algebra-lemma-ass-weakly-ass}을 보라.
\end{proof}

\begin{lemma}
\label{divisors-lemma-weakass-pushforward}
$f : X \to S$를 스킴의 준콤팩트 준분리 사상이라 하자.
$\mathcal{F}$를 준연접 $\mathcal{O}_X$-가군이라 하자.
$s \in S$를 $f$의 상에 속하지 않는 점이라 하자. 그러면 $s$는
$f_*\mathcal{F}$에 약연관되지 않는다.
\end{lemma}

\begin{proof}
$f' : X' \to \Spec(\mathcal{O}_{S, s})$를 $f$를 사상
$g : \Spec(\mathcal{O}_{S, s}) \to S$에 따라 바꾼 밑이라 하고, 다른 사영을
$g' : X' \to X$로 나타내자. 그러면
$$
(f_*\mathcal{F})_s = (g^*f_*\mathcal{F})_s = (f'_*(g')^*\mathcal{F})_s
$$
이다. 첫째 등식은 $g$가 $s$에서 국소환의 동형사상을 유도하기 때문이고,
둘째 등식은 평탄 밑변환 때문이다(Cohomology of Schemes, Lemma
\ref{coherent-lemma-flat-base-change-cohomology}). 물론
$s \in \Spec(\mathcal{O}_{S, s})$는 $f'$의 상에 속하지 않는다. 따라서
$S$가 국소환 $(A, \mathfrak m)$의 스펙트럼이고 $s$가 $\mathfrak m$에
대응한다고 가정해도 된다. Schemes, Lemma
\ref{schemes-lemma-push-forward-quasi-coherent}에 의해 층
$f_*\mathcal{F}$는 준연접이므로, $A$-가군 $M$에 대응한다고 하자.
$s$가 $f$의 상에 속하지 않으므로
$X = \bigcup_{a \in \mathfrak m} f^{-1}D(a)$는 열린 덮개이다.
$X$가 준콤팩트이므로 $a_1, \ldots, a_n \in \mathfrak m$을 택하여
$X = f^{-1}D(a_1) \cup \ldots \cup f^{-1}D(a_n)$으로 만들 수 있다. 따라서
$$
M \to M_{a_1} \oplus \ldots \oplus M_{a_r}
$$
는 단사이다. 그러므로 영이 아닌 임의의 원소 $m$을 줄기
$M_\mathfrak p$에서 택하면 $i$가 존재하여 $a_i^n m$은 모든
$n \geq 0$에 대해 영이 아니다. 따라서 $\mathfrak m$은 $M$에
약연관되지 않는다.
\end{proof}

\begin{lemma}
\label{divisors-lemma-check-injective-on-weakass}
$X$를 스킴이라 하자. $\varphi : \mathcal{F} \to \mathcal{G}$를 준연접
$\mathcal{O}_X$-가군들의 사상이라 하자. 모든 $x \in X$에 대해 다음 중
적어도 하나가 성립한다고 가정하자.
\begin{enumerate}
\item $\mathcal{F}_x \to \mathcal{G}_x$가 단사이거나,
\item $x \not \in \text{WeakAss}(\mathcal{F})$이다.
\end{enumerate}
그러면 $\varphi$는 단사이다.
\end{lemma}

\begin{proof}
가정에 의해 $\text{WeakAss}(\Ker(\varphi)) = \emptyset$이고, 따라서 Lemma
\ref{divisors-lemma-weakly-ass-zero}에 의해 $\Ker(\varphi) = 0$이다.
\end{proof}

\begin{lemma}
\label{divisors-lemma-depth-2-hartog}
$X$를 국소 뇌터 스킴이라 하자. $\mathcal{F}$를 연접
$\mathcal{O}_X$-가군이라 하자. $j : U \to X$를 열린 부분스킴이라 하되,
$x \in X \setminus U$이면 $\text{depth}(\mathcal{F}_x) \geq 2$라고 하자.
그러면
$$
\mathcal{F} \longrightarrow j_*(\mathcal{F}|_U)
$$
는 동형사상이고, 따라서
$\Gamma(X, \mathcal{F}) \to \Gamma(U, \mathcal{F})$도 동형사상이다.
\end{lemma}

\begin{proof}
Lemma \ref{divisors-lemma-check-isomorphism-via-depth-and-ass}을 보조정리에 표시된
사상에 적용할 수 있다고 주장한다. $x \in X$라 하자. $x \in U$이면
$j_*(\mathcal{F}|_U)|_U = \mathcal{F}|_U$이므로 이 사상은 줄기 위에서
동형사상이다. $x \in X \setminus U$이면
$x \not \in \text{Ass}(j_*(\mathcal{F}|_U))$이다
(Lemmas \ref{divisors-lemma-weakass-pushforward}와
\ref{divisors-lemma-weakly-ass-support}). $\text{depth}(\mathcal{F}_x) \geq 2$를
가정했으므로 증명이 끝난다.
\end{proof}

\begin{lemma}
\label{divisors-lemma-weakass-reduced}
$X$를 축약 스킴이라 하자. 그러면 $X$의 약연관점은 정확히 $X$의 기약
성분들의 일반점이다.
\end{lemma}

\begin{proof}
Algebra, Lemma \ref{algebra-lemma-reduced-weakly-ass-minimal}에서 따른다.
\end{proof}



\section{사상과 약연관점}
\label{divisors-section-morphisms-weakly-associated}

\begin{lemma}
\label{divisors-lemma-weakly-ass-reverse-functorial}
$f : X \to S$를 스킴의 아핀 사상이라 하자.
$\mathcal{F}$를 준연접 $\mathcal{O}_X$-가군이라 하자. 그러면
$$
\text{WeakAss}_S(f_*\mathcal{F}) \subset f(\text{WeakAss}_X(\mathcal{F}))
$$
이다.
\end{lemma}

\begin{proof}
$X$와 $S$가 아핀이라고 가정해도 되며, 이때 $X \to S$는 환 준동형
$A \to B$에서 온다. 그러면 $\mathcal{F} = \widetilde M$로 쓸 수 있는데,
여기서 어떤 $B$-가군을 $M$이라 하였다. Lemma
\ref{divisors-lemma-weakly-associated-affine-open}에 의해 $\mathcal{F}$의 약연관점은
$M$의 약연관 소아이디얼과 정확히 대응한다. 마찬가지로
$f_*\mathcal{F}$의 약연관점은 $M$을 $A$-가군으로 보았을 때의 약연관
소아이디얼과 정확히 대응한다. 따라서 이 보조정리는 Algebra, Lemma
\ref{algebra-lemma-weakly-ass-reverse-functorial}에서 따른다.
\end{proof}

\begin{lemma}
\label{divisors-lemma-ass-functorial-equal}
$f : X \to S$를 스킴의 아핀 사상이라 하자.
$\mathcal{F}$를 준연접 $\mathcal{O}_X$-가군이라 하자.
$X$가 국소 뇌터이면
$$
f(\text{Ass}_X(\mathcal{F})) =
\text{Ass}_S(f_*\mathcal{F}) =
\text{WeakAss}_S(f_*\mathcal{F}) =
f(\text{WeakAss}_X(\mathcal{F}))
$$
이다.
\end{lemma}

\begin{proof}
$X$와 $S$가 아핀이라고 가정해도 되며, 이때 $X \to S$는 환 준동형
$A \to B$에서 온다. $X$가 국소 뇌터이므로 환 $B$는 뇌터이다. Properties,
Lemma \ref{properties-lemma-locally-Noetherian}을 보라.
$\mathcal{F} = \widetilde M$로 쓰자. 여기서 어떤 $B$-가군을 $M$이라 하였다. Lemma
\ref{divisors-lemma-associated-affine-open}에 의해 $\mathcal{F}$의 연관점은
$M$의 연관 소아이디얼과 정확히 대응하고, $M$을 $A$-가군으로 보았을 때의
모든 연관 소아이디얼은 $f_*\mathcal{F}$의 연관점이다. 따라서 포함관계
$$
f(\text{Ass}_X(\mathcal{F})) \subset \text{Ass}_S(f_*\mathcal{F})
$$
는 Algebra, Lemma \ref{algebra-lemma-ass-functorial-Noetherian}에서 따른다.
또 포함관계
$$
\text{Ass}_S(f_*\mathcal{F}) \subset \text{WeakAss}_S(f_*\mathcal{F})
$$
는 Lemma \ref{divisors-lemma-weakly-ass-support}에서 따른다. 포함관계
$$
\text{WeakAss}_S(f_*\mathcal{F}) \subset f(\text{WeakAss}_X(\mathcal{F}))
$$
는 Lemma \ref{divisors-lemma-weakly-ass-reverse-functorial}에서 따른다. 양끝 집합은
Lemma \ref{divisors-lemma-ass-weakly-ass}에 의해 같으므로 모든 곳에서 등호를 얻는다.
\end{proof}

\begin{lemma}
\label{divisors-lemma-weakly-associated-finite}
$f : X \to S$를 스킴의 유한 사상이라 하자.
$\mathcal{F}$를 준연접 $\mathcal{O}_X$-가군이라 하자. 그러면
$\text{WeakAss}(f_*\mathcal{F}) = f(\text{WeakAss}(\mathcal{F}))$이다.
\end{lemma}

\begin{proof}
$X$와 $S$가 아핀이라고 가정해도 되며, 이때 $X \to S$는 유한 환 준동형
$A \to B$에서 온다. $\mathcal{F} = \widetilde M$로 쓰자. 여기서 어떤
$B$-가군을 $M$이라 하였다. Lemma
\ref{divisors-lemma-weakly-associated-affine-open}에 의해 $\mathcal{F}$의 약연관점은
$M$의 약연관 소아이디얼과 정확히 대응한다. 마찬가지로
$f_*\mathcal{F}$의 약연관점은 $M$을 $A$-가군으로 보았을 때의 약연관
소아이디얼과 정확히 대응한다. 따라서 이 보조정리는 Algebra, Lemma
\ref{algebra-lemma-weakly-ass-finite-ring-map}에서 따른다.
\end{proof}

\begin{lemma}
\label{divisors-lemma-weakly-ass-pullback}
$f : X \to S$를 스킴의 사상이라 하자. $\mathcal{G}$를 준연접
$\mathcal{O}_S$-가군이라 하자. $x \in X$이고 $s = f(x)$라 하자.
$f$가 $x$에서 평탄하고, 점 $x$가 올 $X_s$의 일반점이며,
$s \in \text{WeakAss}_S(\mathcal{G})$이면
$x \in \text{WeakAss}(f^*\mathcal{G})$이다.
\end{lemma}

\begin{proof}
$A = \mathcal{O}_{S, s}$, $B = \mathcal{O}_{X, x}$,
$M = \mathcal{G}_s$로 놓자. $m \in M$을 그 소멸자
$I = \{a \in A \mid am = 0\}$의 근기가 $\mathfrak m_A$인 원소라 하자.
그러면 $m \otimes 1$의 소멸자는 $I B$인데, 이는 $A \to B$가 충실히
평탄하기 때문이다. 따라서 $\sqrt{I B} = \mathfrak m_B$임을 보이면 충분하다.
이는 $B/\mathfrak m_AB$의 극대 아이디얼이 국소 멱영이라는 사실
(Algebra, Lemma \ref{algebra-lemma-minimal-prime-reduced-ring})과 가정
$\sqrt{I} = \mathfrak m_A$에서 따른다. 몇몇 세부사항은 생략한다.
\end{proof}

\begin{lemma}
\label{divisors-lemma-weakly-ass-change-fields}
$K/k$를 체 확대라 하자. $X$를 $k$ 위의 스킴이라 하자.
$\mathcal{F}$를 준연접 $\mathcal{O}_X$-가군이라 하자.
$y \in X_K$의 상이 $x \in X$라고 하자. $y$가 역상 $\mathcal{F}_K$의
약연관점이면 $x$는 $\mathcal{F}$의 약연관점이다.
\end{lemma}

\begin{proof}
이는 Algebra, Lemma \ref{algebra-lemma-weakly-ass-change-fields}를 스킴의
언어로 옮긴 것이다.
\end{proof}

\noindent
다음은 직상이 흔히 깊이 적어도 2를 가짐을 보여 주는 간단한
보조정리이다.

\begin{lemma}
\label{divisors-lemma-depth-pushforward}
$f : X \to S$를 스킴의 준콤팩트 준분리 사상이라 하자.
$\mathcal{F}$를 준연접 $\mathcal{O}_X$-가군이라 하자.
$s \in S$라 하자.
\begin{enumerate}
\item $s \not \in f(X)$이면 $s$는 $f_*\mathcal{F}$에 약연관되지 않는다.
\item $s \not \in f(X)$이고 $\mathcal{O}_{S, s}$가 뇌터이면 $s$는
$f_*\mathcal{F}$에 연관되지 않는다.
\item $s \not \in f(X)$이고 $(f_*\mathcal{F})_s$가 유한
$\mathcal{O}_{S, s}$-가군이며 $\mathcal{O}_{S, s}$가 뇌터이면
$\text{depth}((f_*\mathcal{F})_s) \geq 2$이다.
\item $\mathcal{F}$가 $S$ 위에서 평탄하고 $a \in \mathfrak m_s$가
영인자가 아니면 $a$는 $(f_*\mathcal{F})_s$ 위에서도 영인자가 아니다.
\item $\mathcal{F}$가 $S$ 위에서 평탄하고 $a, b \in \mathfrak m_s$가
정칙열이면 $a$는 $(f_*\mathcal{F})_s$ 위에서 영인자가 아니고 $b$는
$(f_*\mathcal{F})_s/a(f_*\mathcal{F})_s$ 위에서 영인자가 아니다.
\item $\mathcal{F}$가 $S$ 위에서 평탄하고 $(f_*\mathcal{F})_s$가 유한
$\mathcal{O}_{S, s}$-가군이면
$\text{depth}((f_*\mathcal{F})_s) \geq
\min(2, \text{depth}(\mathcal{O}_{S, s}))$이다.
\end{enumerate}
\end{lemma}

\begin{proof}
(1)은 Lemma \ref{divisors-lemma-weakass-pushforward}이다. (2)는 (1)과 Lemma
\ref{divisors-lemma-ass-weakly-ass}에서 따른다.

\medskip\noindent
(3)의 증명. 깊이가 $\geq 2$임을 보이려면
$\Hom_{\mathcal{O}_{S, s}}(\kappa(s), (f_*\mathcal{F})_s) = 0$이고
$\Ext^1_{\mathcal{O}_{S, s}}(\kappa(s), (f_*\mathcal{F})_s) = 0$임을
보이면 충분하다. Algebra, Lemma \ref{algebra-lemma-depth-ext}를 보라.
완전열
$0 \to \mathfrak m_s \to \mathcal{O}_{S, s} \to \kappa(s) \to 0$을
사용하면 사상
$$
\Hom_{\mathcal{O}_{S, s}}(\mathcal{O}_{S, s}, (f_*\mathcal{F})_s)
\to
\Hom_{\mathcal{O}_{S, s}}(\mathfrak m_s, (f_*\mathcal{F})_s)
$$
이 동형사상임을 증명하면 충분하다. 평탄 밑변환에 의해(Cohomology of
Schemes, Lemma \ref{coherent-lemma-flat-base-change-cohomology}) $S$를
$\Spec(\mathcal{O}_{S, s})$로, $X$를
$\Spec(\mathcal{O}_{S, s}) \times_S X$로 바꾸어도 된다.
$\mathfrak m \subset \mathcal{O}_S$로 $s$의 아이디얼층을 나타내자. 그러면
$$
\Hom_{\mathcal{O}_{S, s}}(\mathfrak m_s, (f_*\mathcal{F})_s) =
\Hom_{\mathcal{O}_S}(\mathfrak m, f_*\mathcal{F}) =
\Hom_{\mathcal{O}_X}(f^*\mathfrak m, \mathcal{F})
$$
이다. 첫째 등식은 $S$가 닫힌점 $s$를 갖는 국소 스킴이기 때문이고, 둘째
등식은 준연접 가군 위의 $f^*, f_*$에 대한 수반에서 따른다. 그런데
$s \not \in f(X)$이므로 $f^*\mathfrak m = \mathcal{O}_X$이다. 논증을
거꾸로 따라가면 원하는 등식을 얻는다.

\medskip\noindent
(4), (5), (6)을 증명하기 위해 평탄 밑변환을 사용하여(Cohomology of
Schemes, Lemma \ref{coherent-lemma-flat-base-change-cohomology}) $S$가
$\mathcal{O}_{S, s}$의 스펙트럼인 경우로 환원한다. 그러면 영인자가 아닌
$a \in \mathcal{O}_{S, s}$는 짧은 완전열
$$
0 \to \mathcal{O}_S \xrightarrow{a} \mathcal{O}_S \to
\mathcal{O}_S/a \mathcal{O}_S \to 0
$$
을 결정한다. $\mathcal{F}$가 $S$ 위에서 평탄하므로 완전열
$$
0 \to \mathcal{F} \xrightarrow{a} \mathcal{F} \to
\mathcal{F}/a\mathcal{F} \to 0
$$
을 얻는다. 이를 직상하면 완전열
$$
0 \to f_*\mathcal{F} \xrightarrow{a} f_*\mathcal{F} \to
f_*(\mathcal{F}/a\mathcal{F})
$$
을 얻는다. 이는 (4)를 증명하고
$f_*\mathcal{F}/ af_*\mathcal{F} \subset f_*(\mathcal{F}/a\mathcal{F})$임을
보여 준다. $b$가 $\mathcal{O}_{S, s}/a\mathcal{O}_{S, s}$ 위에서
영인자가 아니면, 똑같은 논증으로
$b : \mathcal{F}/a\mathcal{F} \to \mathcal{F}/a\mathcal{F}$가 단사임을
알 수 있다. 이를 직상하면
$$
b : f_*(\mathcal{F}/a\mathcal{F}) \to f_*(\mathcal{F}/a\mathcal{F})
$$
가 단사이고, 따라서
$b : f_*\mathcal{F}/ af_*\mathcal{F} \to f_*\mathcal{F}/ af_*\mathcal{F}$도
단사이다. 이는 (5)를 증명한다. (6)은 (4), (5), 그리고 정의들에서
따른다.
\end{proof}



\section{상대 연관점 집합}
\label{divisors-section-relative-assassin}

\noindent
$A \to B$를 환 준동형이라 하자. $N$을 $B$-가군이라 하자.
소 아이디얼 $\mathfrak q \subset B$가 $N$의, $B/A$ 위 상대 연관점
집합에 속한다는 것은 $\mathfrak q$가
$N \otimes_A \kappa(\mathfrak p)$의 연관 소아이디얼이라는 뜻임을
상기하자. 여기서 $\mathfrak p = A \cap \mathfrak q$이다. Algebra,
Definition \ref{algebra-definition-relative-assassin}을 보라. 다음은 스킴의
사상 위에서 준연접층의 상대 연관점 집합을 정의한다.

\begin{definition}
\label{divisors-definition-relative-assassin}
$f : X \to S$를 스킴의 사상이라 하자.
$\mathcal{F}$를 준연접 $\mathcal{O}_X$-가군이라 하자.
$\mathcal{F}$의, $X$ 안에서 $S$ 위의 {\it 상대 연관점 집합}은
다음 집합이다.
$$
\text{Ass}_{X/S}(\mathcal{F}) =
\bigcup\nolimits_{s \in S} \text{Ass}_{X_s}(\mathcal{F}_s)
$$
여기서 $\mathcal{F}_s = (X_s \to X)^*\mathcal{F}$는 $\mathcal{F}$를
$f$의 $s$ 위 올에 제한한 것이다.
\end{definition}

\noindent
다시 말하지만 이는 $f$의 올들이 국소 뇌터이고 $\mathcal{F}$가 유한형일
때 사용하는 것이 가장 좋다는 단서가 있다. 일반적인 경우에는 아마 다음
절에서 정의할 상대 약연관점 집합을 사용해야 할 것이다. 스킴론적 개념과
아핀 열린집합 위의 대수적 개념을 연결해 보자. 이 대응은 올들이 국소
뇌터인 스킴의 사상에 대해서만 완전하게 작동함을 주의하자.

\begin{lemma}
\label{divisors-lemma-relative-assassin-affine-open}
$f : X \to S$를 스킴의 사상이라 하자.
$\mathcal{F}$를 $X$ 위의 준연접층이라 하자.
$U \subset X$와 $V \subset S$를 $f(U) \subset V$인 아핀 열린집합들이라
하자. $U = \Spec(A)$, $V = \Spec(R)$로 쓰고
$M = \Gamma(U, \mathcal{F})$로 놓자.
$x \in U$라 하고 $\mathfrak p \subset A$를 이에 대응하는 소
아이디얼이라 하자. 그러면
$$
\mathfrak p \in \text{Ass}_{A/R}(M) \Rightarrow
x \in \text{Ass}_{X/S}(\mathcal{F})
$$
이다. 모든 올 $X_s$가 $f$의 올로서 국소 뇌터이면
$\mathfrak p \in \text{Ass}_{A/R}(M) \Leftrightarrow
x \in \text{Ass}_{X/S}(\mathcal{F})$가 위와 같은 모든 쌍
$(\mathfrak p, x)$에 대해 성립한다.
\end{lemma}

\begin{proof}
집합 $\text{Ass}_{A/R}(M)$은 Algebra, Definition
\ref{algebra-definition-relative-assassin}에서 정의된다.
한 쌍 $(\mathfrak p, x)$를 택하자. $s = f(x)$라 하자.
$\mathfrak r \subset R$를 $\mathfrak p$ 아래에 있는 소 아이디얼, 즉
$s$에 대응하는 소 아이디얼이라 하자.
$\mathfrak p' \subset A \otimes_R \kappa(\mathfrak r)$를 그 역상이
$\mathfrak p$인 소 아이디얼, 즉 $x$를 그 올 $X_s$의 점으로 보았을 때
이에 대응하는 소 아이디얼이라 하자. 그러면
$\mathfrak p \in \text{Ass}_{A/R}(M)$일 필요충분조건은
$\mathfrak p'$가 $M \otimes_R \kappa(\mathfrak r)$의 연관
소아이디얼인 것이다. Algebra, Lemma
\ref{algebra-lemma-compare-relative-assassins}를 보라. 환
$A \otimes_R \kappa(\mathfrak r)$은 $U_s$에 대응하고, 가군
$M \otimes_R \kappa(\mathfrak r)$은 준연접층
$\mathcal{F}_s|_{U_s}$에 대응함을 주의하자. 따라서 Lemma
\ref{divisors-lemma-associated-affine-open}에 의해 $x$는 $\mathcal{F}_s$의
연관점이다. $\mathfrak p'$가 유한 생성이면 역함의도 성립하며, 이로써
마지막 문장이 참임을 알 수 있다.
\end{proof}

\begin{lemma}
\label{divisors-lemma-base-change-relative-assassin}
$f : X \to S$를 스킴의 사상이라 하자.
$\mathcal{F}$를 준연접 $\mathcal{O}_X$-가군이라 하자.
$g : S' \to S$를 스킴의 사상이라 하자. 밑변환 도식
$$
\xymatrix{
X' \ar[d] \ar[r]_{g'} & X \ar[d] \\
S' \ar[r]^g & S
}
$$
을 생각하고 $\mathcal{F}' = (g')^*\mathcal{F}$로 놓자.
$x' \in X'$의 상들이 $x \in X$, $s' \in S'$, $s \in S$라고 하자.
$f$가 국소 유한형이라고 가정하자. 그러면
$x' \in \text{Ass}_{X'/S'}(\mathcal{F}')$일 필요충분조건은
$x \in \text{Ass}_{X/S}(\mathcal{F})$이고 $x'$가
$\Spec(\kappa(s') \otimes_{\kappa(s)} \kappa(x))$의 한 기약 성분의
일반점에 대응하는 것이다.
\end{lemma}

\begin{proof}
올들의 사상 $X'_{s'} \to X_s$를 생각하자.
$X_{s'} = X_s \times_{\Spec(\kappa(s))} \Spec(\kappa(s'))$이므로 이는
평탄 사상이다. 또한 $\mathcal{F}'_{s'}$는 이 사상을 따른
$\mathcal{F}_s$의 역상이다. $X_s$는 뇌터 스킴
$\Spec(\kappa(s))$ 위에서 국소 유한형이므로 $X_s$는 국소 뇌터이다. Morphisms,
Lemma \ref{morphisms-lemma-finite-type-noetherian}을 보라. 따라서 Lemma
\ref{divisors-lemma-bourbaki}를 적용하면
$$
\text{Ass}_{X'_{s'}}(\mathcal{F}'_{s'}) =
\bigcup\nolimits_{x \in \text{Ass}(\mathcal{F}_s)} \text{Ass}((X'_{s'})_x).
$$
를 얻는다. 그러므로 보조정리를 증명하려면 올 $(X'_{s'})_x$, 즉 사상
$X'_{s'} \to X_s$의 $x$ 위 올의 연관점들이 그 일반점들임을
보이면 충분하다. 스킴으로서
$(X'_{s'})_x = \Spec(\kappa(s') \otimes_{\kappa(s)} \kappa(x))$임을
주의하자. Algebra, Lemma \ref{algebra-lemma-tensor-fields-CM}에 의해 환
$\kappa(s') \otimes_{\kappa(s)} \kappa(x)$은 뇌터 Cohen--Macaulay
환이다. 따라서 그 연관 소아이디얼들은 극소 소아이디얼들이다. Algebra,
Proposition \ref{algebra-proposition-minimal-primes-associated-primes}
(극소 소아이디얼은 연관된다)와 Algebra, Lemma
\ref{algebra-lemma-criterion-no-embedded-primes}(매장된 소아이디얼이 없다)를
보라.
\end{proof}

\begin{remark}
\label{divisors-remark-base-change-relative-assassin}
Lemma \ref{divisors-lemma-base-change-relative-assassin}과 같은 표기와 가정 아래에서
항상
$(g')^{-1}(\text{Ass}_{X/S}(\mathcal{F})) \supset
\text{Ass}_{X'/S'}(\mathcal{F}')$임을 알 수 있다.
사상 $S' \to S$가 국소 준유한이면 실제로
$$
(g')^{-1}(\text{Ass}_{X/S}(\mathcal{F}))
=
\text{Ass}_{X'/S'}(\mathcal{F}')
$$
인데, 이 경우 체 확대 $\kappa(s')/\kappa(s)$들이 항상 유한하기 때문이다.
사실 이는 임의의 사상 $g : S' \to S$에 대해 더 일반적으로 성립하는데,
모든 체 확대 $\kappa(s')/\kappa(s)$가 대수적인 경우가 그렇다. 이 경우
$\kappa(s') \otimes_{\kappa(s)} \kappa(x)$의 모든 소아이디얼은 극대
(그리고 극소) 소아이디얼이기 때문이다. Algebra, Lemma
\ref{algebra-lemma-integral-over-field}를 보라.
\end{remark}



\section{상대 약연관점 집합}
\label{divisors-section-relative-weak-assassin}

\begin{definition}
\label{divisors-definition-relative-weak-assassin}
$f : X \to S$를 스킴의 사상이라 하자.
$\mathcal{F}$를 준연접 $\mathcal{O}_X$-가군이라 하자.
$\mathcal{F}$의, $X$ 안에서 $S$ 위의 {\it 상대 약연관점 집합}은
다음 집합이다.
$$
\text{WeakAss}_{X/S}(\mathcal{F}) =
\bigcup\nolimits_{s \in S} \text{WeakAss}(\mathcal{F}_s)
$$
여기서 $\mathcal{F}_s = (X_s \to X)^*\mathcal{F}$는 $\mathcal{F}$를
$f$의 $s$ 위 올에 제한한 것이다.
\end{definition}

\begin{lemma}
\label{divisors-lemma-relative-weak-assassin-assassin-finite-type}
$f : X \to S$를 국소 유한형인 스킴의 사상이라 하자.
$\mathcal{F}$를 준연접 $\mathcal{O}_X$-가군이라 하자. 그러면
$\text{WeakAss}_{X/S}(\mathcal{F}) = \text{Ass}_{X/S}(\mathcal{F})$이다.
\end{lemma}

\begin{proof}
$f$의 올들은 국소 뇌터 스킴이고, 국소 뇌터 스킴에서는 연관점과
약연관점이 일치하므로 참이다. Lemma \ref{divisors-lemma-ass-weakly-ass}을 보라.
\end{proof}

\begin{lemma}
\label{divisors-lemma-relative-weak-assassin-finite}
$f : X \to S$를 스킴의 사상이라 하자.
$i : Z \to X$를 유한 사상이라 하자.
$\mathcal{F}$를 준연접 $\mathcal{O}_Z$-가군이라 하자. 그러면
$\text{WeakAss}_{X/S}(i_*\mathcal{F}) =
i(\text{WeakAss}_{Z/S}(\mathcal{F}))$이다.
\end{lemma}

\begin{proof}
$i_s : Z_s \to X_s$를 올 사이의 유도된 사상이라 하자. 그러면
$(i_*\mathcal{F})_s = i_{s, *}(\mathcal{F}_s)$이다. 이는 Cohomology of
Schemes, Lemma \ref{coherent-lemma-affine-base-change}과 $i$가 아핀이라는
사실에서 따른다. 따라서
Lemma \ref{divisors-lemma-weakly-associated-finite}를 적용하여 결론을 얻는다.
\end{proof}



\section{피팅 아이디얼}
\label{divisors-section-fitting-ideals}

\noindent
이 절은 More on Algebra, Section
\ref{more-algebra-section-fitting-ideals}의 논의를 이어 간다.
$S$를 스킴이라 하자. $\mathcal{F}$를 유한형 준연접
$\mathcal{O}_S$-가군이라 하자. 이 상황에서 피팅 아이디얼
$$
0 = \text{Fit}_{-1}(\mathcal{F}) \subset \text{Fit}_0(\mathcal{F}) \subset
\text{Fit}_1(\mathcal{F}) \subset \ldots \subset \mathcal{O}_S
$$
을 다음 성질로 특징지어지는 준연접 아이디얼들의 열로 구성할 수 있다.
모든 아핀 열린집합 $U = \Spec(A)$가 $S$에 있을 때 $\mathcal{F}|_U$가
$A$-가군 $M$에 대응하면 $\text{Fit}_i(\mathcal{F})|_U$는 아이디얼
$\text{Fit}_i(M) \subset A$에 대응한다. 이는 잘 정의되며 준연접
아이디얼층이다. 실제로 $f \in A$이면 $i$번째 피팅 아이디얼, 즉
$M_f$의 $A_f$ 위 피팅 아이디얼은 More on Algebra, Lemma
\ref{more-algebra-lemma-fitting-ideal-basics}에 의해
$\text{Fit}_i(M) A_f$와 같다.

\medskip\noindent
또는 $\mathcal{F}$의 국소 표시를 이용하여 피팅 아이디얼을 구성할 수 있다.
구체적으로 $U \subset X$가 열려 있고
$$
\bigoplus\nolimits_{i \in I} \mathcal{O}_U \to
\mathcal{O}_U^{\oplus n} \to \mathcal{F}|_U \to 0
$$
이 $\mathcal{F}$의 표시로서 $U$ 위에 놓이면,
$\text{Fit}_r(\mathcal{F})|_U$는 이 표시의 첫째 화살표를 정의하는 행렬의
$(n - r) \times (n - r)$ 소행렬식들로 생성된다. 이것은 환 위 가군의
피팅 아이디얼을 실제로 정의하는 방법이므로 위 구성과 양립한다. 몇몇
세부사항은 생략한다.

\begin{lemma}
\label{divisors-lemma-base-change-fitting-ideal}
$f : T \to S$를 스킴의 사상이라 하자.
$\mathcal{F}$를 유한형 준연접 $\mathcal{O}_S$-가군이라 하자. 그러면
$f^{-1}\text{Fit}_i(\mathcal{F}) \cdot \mathcal{O}_T =
\text{Fit}_i(f^*\mathcal{F})$이다.
\end{lemma}

\begin{proof}
More on Algebra, Lemma \ref{more-algebra-lemma-fitting-ideal-basics}의 (3)에서
곧바로 따른다.
\end{proof}

\begin{lemma}
\label{divisors-lemma-fitting-ideal-of-finitely-presented}
$S$를 스킴이라 하자.
$\mathcal{F}$를 유한 표시 $\mathcal{O}_S$-가군이라 하자. 그러면
$\text{Fit}_r(\mathcal{F})$는 유한형 준연접 아이디얼이다.
\end{lemma}

\begin{proof}
More on Algebra, Lemma \ref{more-algebra-lemma-fitting-ideal-basics}의 (4)에서
곧바로 따른다.
\end{proof}

\begin{lemma}
\label{divisors-lemma-on-subscheme-cut-out-by-Fit-0}
$S$를 스킴이라 하자.
$\mathcal{F}$를 유한형 준연접 $\mathcal{O}_S$-가군이라 하자.
$Z_0 \subset S$를 $\text{Fit}_0(\mathcal{F})$가 정의하는 닫힌
부분스킴이라 하자. $Z \subset S$를 $\mathcal{F}$의 스킴론적 지지라 하자.
그러면
\begin{enumerate}
\item 닫힌 부분스킴으로서 $Z \subset Z_0 \subset S$이고,
\item 닫힌 부분집합으로서 $Z = Z_0 = \text{Supp}(\mathcal{F})$이며,
\item 유한형 준연접 $\mathcal{O}_{Z_0}$-가군 $\mathcal{G}_0$가 존재하여
$$
(Z_0 \to X)_*\mathcal{G}_0 = \mathcal{F}.
$$
\end{enumerate}
\end{lemma}

\begin{proof}
$Z$는 국소적으로 $\mathcal{F}$의 소멸자가 정의함을 상기하자. Morphisms,
Definition \ref{morphisms-definition-scheme-theoretic-support}를 보라(여기서는
Morphisms, Lemma \ref{morphisms-lemma-scheme-theoretic-support}을 사용하여
$Z$를 정의한다). 따라서 More on Algebra, Lemma
\ref{more-algebra-lemma-fitting-ideal-basics}의 (6)에 의해 스킴론적으로
$Z \subset Z_0$임을 알 수 있다. 한편 Morphisms, Lemma
\ref{morphisms-lemma-scheme-theoretic-support}에 의해 집합론적으로
$Z = \text{Supp}(\mathcal{F})$이고, More on Algebra, Lemma
\ref{more-algebra-lemma-fitting-ideal-basics}의 (7)에 의해 집합론적으로
$Z_0 = Z$이다. 끝으로 (3)의 $\mathcal{G}_0$를 얻기 위해 Morphisms,
Lemma \ref{morphisms-lemma-scheme-theoretic-support}에서와 같이
$\mathcal{G}$를 $Z$ 위에서 사용하여
$\mathcal{G}_0 = (Z \to Z_0)_*\mathcal{G}$로 놓아도 되고, Morphisms,
Lemma \ref{morphisms-lemma-i-star-equivalence} 및 More on Algebra, Lemma
\ref{more-algebra-lemma-fitting-ideal-basics}의 (6)에 의해
$\text{Fit}_0(\mathcal{F})$가 $\mathcal{F}$를 소멸시킨다는 사실을
사용해도 된다.
\end{proof}

\begin{lemma}
\label{divisors-lemma-fitting-ideal-generate-locally}
$S$를 스킴이라 하자. $\mathcal{F}$를 유한형 준연접
$\mathcal{O}_S$-가군이라 하자. $s \in S$라 하자. $\mathcal{F}$가
$r$개의 원소로 $s$의 한 근방에서 생성될 필요충분조건은
$\text{Fit}_r(\mathcal{F})_s = \mathcal{O}_{S, s}$인 것이다.
\end{lemma}

\begin{proof}
More on Algebra, Lemma
\ref{more-algebra-lemma-fitting-ideal-generate-locally}에서 곧바로 따른다.
\end{proof}

\begin{lemma}
\label{divisors-lemma-fitting-ideal-finite-locally-free}
$S$를 스킴이라 하자. $\mathcal{F}$를 유한형 준연접
$\mathcal{O}_S$-가군이라 하자. $r \geq 0$이라 하자. 다음은 동치이다.
\begin{enumerate}
\item $\mathcal{F}$는 계수 $r$인 유한 국소 자유이고,
\item $\text{Fit}_{r - 1}(\mathcal{F}) = 0$이고
$\text{Fit}_r(\mathcal{F}) = \mathcal{O}_S$이며,
\item $\text{Fit}_k(\mathcal{F}) = 0$은 $k < r$일 때 성립하고,
$\text{Fit}_k(\mathcal{F}) = \mathcal{O}_S$는 $k \geq r$일 때 성립한다.
\end{enumerate}
\end{lemma}

\begin{proof}
More on Algebra, Lemma
\ref{more-algebra-lemma-fitting-ideal-finite-locally-free}에서 곧바로 따른다.
\end{proof}

\begin{lemma}
\label{divisors-lemma-locally-free-rank-r-pullback}
$S$를 스킴이라 하자. $\mathcal{F}$를 유한형 준연접
$\mathcal{O}_S$-가군이라 하자. $\mathcal{F}$의 피팅 아이디얼들이 정의하는
닫힌 부분스킴들
$$
S = Z_{-1} \supset Z_0 \supset Z_1 \supset Z_2 \ldots
$$
은 다음 성질들을 갖는다.
\begin{enumerate}
\item 교집합 $\bigcap Z_r$은 공집합이다.
\item 규칙
$$
T \longmapsto
\left\{
\begin{matrix}
\{*\} & \text{if }\mathcal{F}_T\text{ is locally generated by }
\leq r\text{ sections} \\
\emptyset & \text{otherwise}
\end{matrix}
\right.
$$
으로 정의되는 함자 $(\Sch/S)^{opp} \to \textit{Sets}$는 열린 부분스킴
$S \setminus Z_r$로 표현된다.
\item 규칙
$$
T \longmapsto
\left\{
\begin{matrix}
\{*\} & \text{if }\mathcal{F}_T\text{ locally free rank }r\\
\emptyset & \text{otherwise}
\end{matrix}
\right.
$$
으로 정의되는 함자 $F_r : (\Sch/S)^{opp} \to \textit{Sets}$는 국소 닫힌
부분스킴 $Z_{r - 1} \setminus Z_r$로 표현되며, 이는 $S$의 부분스킴이다.
\end{enumerate}
$\mathcal{F}$가 유한 표시이면 $Z_r \to S$, $S \setminus Z_r \to S$,
$Z_{r - 1} \setminus Z_r \to S$는 유한 표시이다.
\end{lemma}

\begin{proof}
(1)은 모든 아핀 열린집합 $U$ 위에서 어떤 정수 $n$이 존재하여
$\text{Fit}_n(\mathcal{F})|_U = \mathcal{O}_U$이므로 참이다. 실제로
$n$을 $\mathcal{F}$의, $U$ 위 생성원 수로 택할 수 있다. More on
Algebra, Section \ref{more-algebra-section-fitting-ideals}를 보라.

\medskip\noindent
임의의 사상 $g : T \to S$에 대해 Lemmas
\ref{divisors-lemma-base-change-fitting-ideal}과
\ref{divisors-lemma-fitting-ideal-generate-locally}에서 다음을 알 수 있다.
$\mathcal{F}_T$가 국소적으로 $\leq r$개의 절단으로 생성될 필요충분조건은
$\text{Fit}_r(\mathcal{F}) \cdot \mathcal{O}_T = \mathcal{O}_T$인 것이다.
이는 (2)를 증명한다.

\medskip\noindent
임의의 사상 $g : T \to S$에 대해 Lemmas
\ref{divisors-lemma-base-change-fitting-ideal}과
\ref{divisors-lemma-fitting-ideal-finite-locally-free}에서 다음을 알 수 있다.
$\mathcal{F}_T$가 계수 $r$인 자유 가군일 필요충분조건은
$\text{Fit}_r(\mathcal{F}) \cdot \mathcal{O}_T = \mathcal{O}_T$이고
$\text{Fit}_{r - 1}(\mathcal{F}) \cdot \mathcal{O}_T = 0$인 것이다.
이는 (3)을 증명한다.

\medskip\noindent
$\mathcal{F}$가 유한 표시라고 가정하자. 그러면 각 사상 $Z_r \to S$는
$\text{Fit}_r(\mathcal{F})$가 유한형이므로 유한 표시이다(Lemma
\ref{divisors-lemma-fitting-ideal-of-finitely-presented}와 Morphisms, Lemma
\ref{morphisms-lemma-closed-immersion-finite-presentation}). 이는
$Z_{r - 1} \setminus Z_r$가 $Z_r$ 안에서 역콤팩트 열린집합임을 함의하고
(Properties, Lemma
\ref{properties-lemma-quasi-coherent-finite-type-ideals}), 따라서 사상
$Z_{r - 1} \setminus Z_r \to Z_r$도 유한 표시이다.
\end{proof}

\noindent
Lemma \ref{divisors-lemma-locally-free-rank-r-pullback}에도 불구하고 다음 보조정리는
$\mathcal{F}$가 유한형 준연접 가군일 때에는 성립하지 않는다. 즉 층화는
여전히 존재하지만 일반적으로 함자 $F_{flat}$을 표현하지는 않는다.

\begin{lemma}
\label{divisors-lemma-finite-presentation-module}
$S$를 스킴이라 하자. $\mathcal{F}$를 유한 표시
$\mathcal{O}_S$-가군이라 하자. Lemma
\ref{divisors-lemma-locally-free-rank-r-pullback}에서와 같이
$S = Z_{-1} \supset Z_0 \supset Z_1 \supset \ldots$라 하자.
$S_r = Z_{r - 1} \setminus Z_r$로 놓자. 그러면
$S' = \coprod_{r \geq 0} S_r$는 함자
$$
F_{flat} : \Sch/S \longrightarrow \textit{Sets},\quad\quad
T \longmapsto
\left\{
\begin{matrix}
\{*\} & \text{if }\mathcal{F}_T\text{ flat over }T\\
\emptyset & \text{otherwise}
\end{matrix}
\right.
$$
를 표현한다. 또한 $\mathcal{F}|_{S_r}$는 계수 $r$인 국소 자유이고,
사상 $S_r \to S$와 $S' \to S$는 유한 표시이다.
\end{lemma}

\begin{proof}
$g : T \to S$를 역상 $\mathcal{F}_T = g^*\mathcal{F}$가 평탄인 스킴의
사상이라 하자. 그러면 $\mathcal{F}_T$는 유한 표시 평탄
$\mathcal{O}_T$-가군이다. 따라서 $\mathcal{F}_T$는 유한 국소 자유이다.
Properties, Lemma \ref{properties-lemma-finite-locally-free}를 보라.
그러므로 $T = \coprod_{r \geq 0} T_r$이고, 여기서
$\mathcal{F}_T|_{T_r}$는 계수 $r$인 국소 자유이다. 이는
$$
F_{flat} = \coprod\nolimits_{r \geq 0} F_r
$$
임을 $\Sch/S$ 위의 Zariski 층들의 범주에서 함의한다. 여기서 $F_r$은
Lemma \ref{divisors-lemma-locally-free-rank-r-pullback}에서와 같다. 따라서
$F_{flat}$은 $\coprod_{r \geq 0} (Z_{r - 1} \setminus Z_r)$로 표현된다.
여기서 $Z_r$은 Lemma \ref{divisors-lemma-locally-free-rank-r-pullback}에서와 같다.
나머지 명제들도 그 보조정리에서 따른다.
\end{proof}

\begin{example}
\label{divisors-example-not-fp-MB}
$R = \prod_{n \in \mathbf{N}} \mathbf{F}_2$라 하자. $I \subset R$을 성분이
거의 모두 영인 원소 $a = (a_n)_{n \in \mathbf{N}}$들의 아이디얼이라 하자.
$\mathfrak m$을 $I$를 포함하는 극대 아이디얼이라 하자. 그러면
$M = R/\mathfrak m$은 유한 평탄 $R$-가군이다. 실제로 $R$은 절대
평탄이다(More on Algebra, Lemma
\ref{more-algebra-lemma-product-fields-absolutely-flat}).
$S = \Spec(R)$와 $\mathcal{F} = \widetilde{M}$로 놓자. Lemma
\ref{divisors-lemma-locally-free-rank-r-pullback}의 닫힌 부분스킴들은
$S = Z_{-1}$, $Z_0 = \Spec(R/\mathfrak m)$, 그리고
$Z_i = \emptyset$이며 이는 $i > 0$일 때 성립한다. 그러나
$\text{id} : S \to S$는 $(S \setminus Z_0) \amalg Z_0$를 통해 분해되지
않는데, $\mathfrak m$이 $S$의 고립점이 아니기 때문이다. 따라서 Lemma
\ref{divisors-lemma-finite-presentation-module}은 유한형
가군에 대해 성립하지 않는다.
\end{example}



\section{사상의 특이 궤적}
\label{divisors-section-singular-locus-morphism}

\noindent
$f : X \to S$를 스킴의 유한형 사상이라 하자. 점들의 집합 $U$를
$f$가 매끄러운 점들의 집합이라 하자. 이는 $X$의 열린집합이다(Morphisms, Definition
\ref{morphisms-definition-smooth}에 의해). 여러 상황에서 표준적인 닫힌
부분스킴 $\text{Sing}(f) \subset X$를 갖는 것이 유용하다. 그 여집합은
$U$이고 그 형성은 임의의 밑변환과 가환한다.

\medskip\noindent
$f$가 유한 표시이면 한 선택으로 아핀 국소적으로 Smoothing Ring Maps,
Definition \ref{smoothing-definition-strictly-standard}의 의미에서
``엄밀히 표준적인'' 함수들이 정의하는 닫힌 부분스킴 $Z$를 생각할 수 있다.
Smoothing Ring Maps, Lemma
\ref{smoothing-lemma-strictly-standard-base-change}에 의해 다음을 알 수 있다.
$f' : X' \to S'$가 $f$의 밑변환이고 이 밑변환이 사상 $S' \to S$에
따른 것이면
$Z' \subset S' \times_S Z$이다. 여기서 $Z'$은 $X'$ 위에서 아핀
국소적으로 엄밀히 표준적인 함수들이 정의하는 닫힌 부분스킴이다. 하지만
등호인지는 분명하지 않다. 엄밀히 표준적인 원소라는 개념은 Popescu의
정리에 관한 장에서 유용했다. 이 원소들이 정의하는 닫힌 부분스킴은
(우리가 아는 한) 문헌에서 사용되지 않는다\footnote{$f$가 국소 완전
교차 사상이면(More on Morphisms, Definition
\ref{more-morphisms-definition-lci}), 국소적으로 엄밀히 표준적인 원소들이
정의하는 닫힌 부분스킴이 살펴보아야 할 올바른 대상이다.}.

\medskip\noindent
$f$가 평탄하고 유한 표시이며 $f$의 올들이 모두 차원 $d$인 등차원이면,
$d$번째 피팅 아이디얼을 사용하는데, 이는 $\Omega_{X/S}$의 피팅
아이디얼이며 좋은 닫힌 부분스킴을 준다. 임의의 유한형 사상에 대해
$X$의 닫힌 부분스킴들을 생각하자. 이들은 $\Omega_{X/S}$의 피팅
아이디얼들이 정의하며, $X$의 층화를 $\Omega_{X/S}$의 계수에 따라
정의한다. 그 형성은 밑변환과 가환한다. 이는 유용할 수 있으며
올들의 매장 차원과 관련된다. Varieties, Section
\ref{varieties-section-embedding-dimension}을 보라.

\begin{lemma}
\label{divisors-lemma-base-change-and-fitting-ideal-omega}
$f : X \to S$를 국소 유한형인 스킴의 사상이라 하자.
$X = Z_{-1} \supset Z_0 \supset Z_1 \supset \ldots$를
$\Omega_{X/S}$의 피팅 아이디얼들이 정의하는 닫힌 부분스킴들이라 하자.
그러면 $Z_i$의 형성은 임의의 밑변환과 가환한다.
\end{lemma}

\begin{proof}
$\Omega_{X/S}$는 유한형 준연접 $\mathcal{O}_X$-가군이므로(Morphisms,
Lemma \ref{morphisms-lemma-finite-type-differentials}) 피팅 아이디얼들이
정의된다. $f' : X' \to S'$가 $f$의 밑변환이고 그 밑변환 사상이
$g : S' \to S$이면,
$\Omega_{X'/S'} = (g')^*\Omega_{X/S}$이다. 여기서 $g' : X' \to X$는
사영이다(Morphisms, Lemma
\ref{morphisms-lemma-base-change-differentials}). 따라서
$(g')^{-1}\text{Fit}_i(\Omega_{X/S}) \cdot \mathcal{O}_{X'} =
\text{Fit}_i(\Omega_{X'/S'})$이다. 이는 스킴론적으로
$$
Z'_i = (g')^{-1}(Z_i) = Z_i \times_X X'
$$
라는 뜻이며, 이것이 보조정리의 명제의 의미이다.
\end{proof}

\noindent
$0$번째 피팅 아이디얼, 즉 $\Omega$의 그것은 사상의 ``분기 궤적''을
정의한다.

\begin{lemma}
\label{divisors-lemma-zero-fitting-ideal-omega-unramified}
$f : X \to S$를 국소 유한형인 스킴의 사상이라 하자. 닫힌 부분스킴
$Z \subset X$를 택하되, $0$번째 피팅 아이디얼, 즉 $\Omega_{X/S}$의
그것이 정의한다고 하자. 이는 정확히 $f$가 비분기이지 않은 점들의 집합이다.
\end{lemma}

\begin{proof}
Lemma \ref{divisors-lemma-on-subscheme-cut-out-by-Fit-0}에 의해 $Z$의 여집합은 정확히
$\Omega_{X/S}$가 영인 궤적이다. 이는 Morphisms, Lemma
\ref{morphisms-lemma-unramified-omega-zero}에 의해 정확히 $f$가 비분기인
점들의 집합이다.
\end{proof}

\begin{lemma}
\label{divisors-lemma-d-fitting-ideal-omega-smooth}
$f : X \to S$를 스킴의 사상이라 하자. $d \geq 0$을 정수라 하자. 다음을
가정하자.
\begin{enumerate}
\item $f$는 평탄하고,
\item $f$는 국소 유한 표시이며,
\item $f$의 모든 공집합이 아닌 올은 차원 $d$인 등차원이다.
\end{enumerate}
$Z \subset X$를 $d$번째 피팅 아이디얼이 정의하는 닫힌 부분스킴이라
하자. 이 아이디얼은 $\Omega_{X/S}$의 피팅 아이디얼이다. 그러면 $Z$는
정확히 $f$가 매끄럽지 않은 점들의
집합이다.
\end{lemma}

\begin{proof}
Lemma \ref{divisors-lemma-locally-free-rank-r-pullback}에 의해 $Z$의 여집합은 정확히
$\Omega_{X/S}$가 많아야 $d$개의 원소로 생성될 수 있는 궤적이다. 따라서
이 보조정리는 Morphisms, Lemma \ref{morphisms-lemma-smooth-at-point}에서
따른다.
\end{proof}


\section{꼬임 없는 가군}
\label{divisors-section-torsion-free}

\noindent
이 절은 준연접 가군에 대한 More on Algebra, Section
\ref{more-algebra-section-torsion-flat}의 대응물이다.

\begin{lemma}
\label{divisors-lemma-torsion-sections}
$X$를 일반점이 $\eta$인 정역 스킴이라 하자. $\mathcal{F}$를 준연접
$\mathcal{O}_X$-가군이라 하자. $U \subset X$를 공집합이 아닌 열린집합이라
하고 $s \in \mathcal{F}(U)$라 하자. 다음은 서로 동치이다.
\begin{enumerate}
\item 어떤 $x \in U$에 대해 $s$의 상이 $\mathcal{F}_x$에서 꼬임이다.
\item 모든 $x \in U$에 대해 $s$의 상이 $\mathcal{F}_x$에서 꼬임이다.
\item $s$의 상이 $\mathcal{F}_\eta$에서 영이다.
\item $s$의 상이 $j_*\mathcal{F}_\eta$에서 영이다. 여기서
$j : \eta \to X$는 포함 사상이다.
\end{enumerate}
\end{lemma}

\begin{proof}
생략한다.
\end{proof}

\begin{definition}
\label{divisors-definition-torsion}
$X$를 정역 스킴이라 하자. $\mathcal{F}$를 준연접
$\mathcal{O}_X$-가군이라 하자.
\begin{enumerate}
\item $\mathcal{F}$의 국소 절단이 Lemma
\ref{divisors-lemma-torsion-sections}의 서로 동치인 조건들을 만족하면 이를
{\it 꼬임}이라고 한다.
\item $\mathcal{F}$가 {\it 꼬임이 없다}고 하는 것은 $\mathcal{F}$의 모든
꼬임 절단이 $0$일 때이다.
\end{enumerate}
\end{definition}

\noindent
통상적인 대수적 개념과 이를 비교하는 필수 보조정리는 다음과 같다.

\begin{lemma}
\label{divisors-lemma-check-torsion-on-affines}
$X$를 정역 스킴이라 하자. $\mathcal{F}$를 준연접
$\mathcal{O}_X$-가군이라 하자. 다음은 서로 동치이다.
\begin{enumerate}
\item $\mathcal{F}$에는 꼬임이 없다.
\item 아핀 열린집합 $U \subset X$에 대해 $\mathcal{F}(U)$는 꼬임이 없는
$\mathcal{O}(U)$-가군이다.
\end{enumerate}
\end{lemma}

\begin{proof}
생략한다.
\end{proof}

\begin{lemma}
\label{divisors-lemma-torsion}
$X$를 정역 스킴이라 하자. $\mathcal{F}$를 준연접
$\mathcal{O}_X$-가군이라 하자. $\mathcal{F}$의 꼬임 절단들은 준연접
$\mathcal{O}_X$-부분가군
$\mathcal{F}_{tors} \subset \mathcal{F}$
을 이룬다. 몫가군 $\mathcal{F}/\mathcal{F}_{tors}$에는 꼬임이 없다.
\end{lemma}

\begin{proof}
생략한다. 대수적 대응물은 More on Algebra, Lemma
\ref{more-algebra-lemma-torsion}을 보라.
\end{proof}

\begin{lemma}
\label{divisors-lemma-flat-torsion-free}
$X$를 정역 스킴이라 하자. 임의의 평탄 준연접 $\mathcal{O}_X$-가군에는
꼬임이 없다.
\end{lemma}

\begin{proof}
생략한다. More on Algebra, Lemma
\ref{more-algebra-lemma-flat-torsion-free}를 보라.
\end{proof}

\begin{lemma}
\label{divisors-lemma-flat-pullback-torsion}
$f : X \to Y$를 정역 스킴들의 평탄 사상이라 하자. $\mathcal{G}$를 꼬임이
없는 준연접 $\mathcal{O}_Y$-가군이라 하자. 그러면 $f^*\mathcal{G}$는
꼬임이 없는 $\mathcal{O}_X$-가군이다.
\end{lemma}

\begin{proof}
생략한다. 대수적 대응물은 More on Algebra, Lemma
\ref{more-algebra-lemma-flat-pullback-torsion}을 보라.
\end{proof}

\begin{lemma}
\label{divisors-lemma-flat-over-integral-integral-fibre}
$f : X \to Y$를 스킴의 평탄 사상이라 하자. $Y$가 정역이고 $f$의 일반
올이 정역이면 $X$는 정역이다.
\end{lemma}

\begin{proof}
대수적 대응물은 다음과 같다. $A$를 분수체가 $K$인 정역이라 하고 $B$를
평탄 $A$-대수라 하되 $B \otimes_A K$가 정역이라고 하자. 그러면 $B$는
정역이다. 실제로 More on Algebra, Lemma
\ref{more-algebra-lemma-flat-torsion-free}에 의해 $B$에는 꼬임이 없고,
따라서 $B \subset B \otimes_A K$이기 때문이다.
\end{proof}

\begin{lemma}
\label{divisors-lemma-check-torsion}
$X$를 정역 스킴이라 하자. $\mathcal{F}$를 준연접
$\mathcal{O}_X$-가군이라 하자. 그러면 $\mathcal{F}$에 꼬임이 없는 것과
$\mathcal{F}_x$가 꼬임이 없는 $\mathcal{O}_{X, x}$-가군인 것은 모든
$x \in X$에 대해 서로 동치이다.
\end{lemma}

\begin{proof}
생략한다. More on Algebra, Lemma
\ref{more-algebra-lemma-check-torsion}을 보라.
\end{proof}

\begin{lemma}
\label{divisors-lemma-extension-torsion-free}
$X$를 정역 스킴이라 하자.
$0 \to \mathcal{F} \to \mathcal{F}' \to \mathcal{F}'' \to 0$
을 준연접 $\mathcal{O}_X$-가군들의 짧은 완전열이라 하자.
$\mathcal{F}$와 $\mathcal{F}''$에 꼬임이 없으면 $\mathcal{F}'$에도
꼬임이 없다.
\end{lemma}

\begin{proof}
생략한다. 대수적 대응물은 More on Algebra, Lemma
\ref{more-algebra-lemma-extension-torsion-free}를 보라.
\end{proof}

\begin{lemma}
\label{divisors-lemma-torsion-free-finite-noetherian-domain}
$X$를 일반점이 $\eta$인 국소 뇌터 정역 스킴이라 하자. $\mathcal{F}$를
영이 아닌 응집 $\mathcal{O}_X$-가군이라 하자. 다음은 서로 동치이다.
\begin{enumerate}
\item $\mathcal{F}$에는 꼬임이 없다.
\item $\eta$는 $\mathcal{F}$의 유일한 연관 소아이디얼이다.
\item $\eta$는 $\mathcal{F}$의 지지집합에 속하고 $\mathcal{F}$는
성질 $(S_1)$을 가진다.
\item $\eta$는 $\mathcal{F}$의 지지집합에 속하고 $\mathcal{F}$는 매장된
연관 소아이디얼을 갖지 않는다.
\end{enumerate}
\end{lemma}

\begin{proof}
이는 More on Algebra, Lemma
\ref{more-algebra-lemma-torsion-free-finite-noetherian-domain}을 스킴의
언어로 옮긴 것이다. 그 번역은 생략한다.
\end{proof}

\begin{lemma}
\label{divisors-lemma-torsion-free-over-regular-dim-1}
$X$를 차원이 $\leq 1$인 정칙 정역 스킴이라 하자. $\mathcal{F}$를 응집
$\mathcal{O}_X$-가군이라 하자. 다음은 서로 동치이다.
\begin{enumerate}
\item $\mathcal{F}$에는 꼬임이 없다.
\item $\mathcal{F}$는 유한 국소 자유이다.
\end{enumerate}
\end{lemma}

\begin{proof}
유한 국소 자유 가군에 꼬임이 없다는 것은 자명하다. 역으로,
$\mathcal{F}$에 꼬임이 없으면 $\mathcal{F}_x$가 자유
$\mathcal{O}_{X, x}$-가군임을 모든 $x \in X$에 대해 보이겠다. 이는
Algebra, Lemma \ref{algebra-lemma-finite-projective}와 $\mathcal{F}$가
응집이라는 사실에 의해 충분하다. $\dim(\mathcal{O}_{X, x}) = 0$이면
$\mathcal{O}_{X, x}$는 체이므로 명제는 자명하다.
$\dim(\mathcal{O}_{X, x}) = 1$이면 $\mathcal{O}_{X, x}$는 이산 부치환이다
(Algebra, Lemma \ref{algebra-lemma-characterize-dvr}). 또한
$\mathcal{F}_x$에는 꼬임이 없다. 따라서 More on Algebra, Lemma
\ref{more-algebra-lemma-dedekind-torsion-free-flat}에 의해
$\mathcal{F}_x$는 자유이다.
\end{proof}

\begin{lemma}
\label{divisors-lemma-hom-into-torsion-free}
$X$를 정역 스킴이라 하자. $\mathcal{F}$, $\mathcal{G}$를 준연접
$\mathcal{O}_X$-가군이라 하자. $\mathcal{G}$에 꼬임이 없고
$\mathcal{F}$가 유한 표시이면
$\SheafHom_{\mathcal{O}_X}(\mathcal{F}, \mathcal{G})$에는 꼬임이 없다.
\end{lemma}

\begin{proof}
$\SheafHom_{\mathcal{O}_X}(\mathcal{F}, \mathcal{G})$가 Schemes,
Section \ref{schemes-section-quasi-coherent}에 의해 준연접이므로 명제는
의미가 있다. 명제가 참임을 보려면 More on Algebra, Lemma
\ref{more-algebra-lemma-hom-into-torsion-free}를 보라. 일부 세부사항은
생략한다.
\end{proof}

\begin{lemma}
\label{divisors-lemma-isom-depth-2-torsion-free}
$X$를 정역 국소 뇌터 스킴이라 하자.
$\varphi : \mathcal{F} \to \mathcal{G}$를 준연접
$\mathcal{O}_X$-가군의 사상이라 하자. $\mathcal{F}$는 응집이고,
$\mathcal{G}$에는 꼬임이 없으며, 모든 $x \in X$에 대해 다음 중 하나가
성립한다고 가정하자.
\begin{enumerate}
\item $\mathcal{F}_x \to \mathcal{G}_x$는 동형이다.
\item $\text{depth}(\mathcal{F}_x) \geq 2$이다.
\end{enumerate}
그러면 $\varphi$는 동형이다.
\end{lemma}

\begin{proof}
이는 More on Algebra, Lemma
\ref{more-algebra-lemma-isom-depth-2-torsion-free}를 스킴의 언어로 옮긴
것이다.
\end{proof}


\section{가군의 계수}
\label{divisors-section-rank}

\noindent
이 절은 준연접 가군에 대한 More on Algebra, Section
\ref{more-algebra-section-rank}의 대응물이다. $X$가 일반점이 $\eta$인
정역 스킴이면 국소환 $\mathcal{O}_{X, \eta}$는 잉여체
$\kappa(\eta)$와 같다는 것을 상기하자. Morphisms, Lemma
\ref{morphisms-lemma-integral-scheme-rational-functions}를 보라.

\begin{definition}
\label{divisors-definition-rank}
$X$를 일반점이 $\eta$인 정역 스킴이라 하자. $\mathcal{F}$를 준연접
$\mathcal{O}_X$-가군이라 하자. $\mathcal{F}$의 {\it 계수}는
$\dim_{\kappa(\eta)} \mathcal{F}_\eta$이다.
\end{definition}

\noindent
이는 두 정의를 모두 적용할 수 있을 때 국소 자유 가군의 계수 정의와
충돌하지 않는다. 스킴 $X$가 기약이 아니거나 일반점에서의 국소환
$\mathcal{O}_{X, \eta}$가 체가 아닐 때 준연접 가군의 계수를 정의하는
것은 좋은 생각으로 보이지 않는다\footnote{$X$가 기약이고
$\mathcal{F}_\eta$가 $\mathcal{O}_{X, \eta}$ 위에서 자유이면 이 자유
가군의 계수를 사용하는 것이 자연스러워 보인다. 일부 문헌은
$\mathcal{F}_\eta$의 길이를 사용하여 계수를 정의하므로 주의하라.}

\begin{lemma}
\label{divisors-lemma-check-rank-on-affines}
$X$를 정역 스킴이라 하자. $\mathcal{F}$를 준연접
$\mathcal{O}_X$-가군이라 하자. $r$을 기수라 하자. 다음은 서로 동치이다.
\begin{enumerate}
\item $\mathcal{F}$의 계수는 $r$이다.
\item 모든 공집합이 아닌 아핀 열린집합 $U \subset X$에 대해
$\mathcal{F}(U)$는 $\mathcal{O}_X(U)$-가군이고 그 계수는 $r$이다.
\item 어떤 공집합이 아닌 아핀 열린집합 $U \subset X$에 대해
$\mathcal{F}(U)$는 $\mathcal{O}_X(U)$-가군이고 그 계수는 $r$이다.
\end{enumerate}
\end{lemma}

\begin{proof}
생략한다.
\end{proof}

\begin{lemma}
\label{divisors-lemma-dominant-pullback-rank}
$f : X \to Y$를 정역 스킴들의 우세 사상이라 하자. $\mathcal{G}$를
준연접 $\mathcal{O}_Y$-가군이라 하자. 그러면 $\mathcal{G}$가 $Y$ 위에서
갖는 계수는 $f^*\mathcal{G}$가 $X$ 위에서 갖는 계수와 같다.
\end{lemma}

\begin{proof}
Lemma \ref{divisors-lemma-check-rank-on-affines}, Morphisms, Lemma
\ref{morphisms-lemma-dominant-between-integral}, 그리고 More on Algebra,
Lemma \ref{more-algebra-lemma-dominant-pullback-rank}을 사용하라.
\end{proof}

\begin{lemma}
\label{divisors-lemma-rank-additive}
$X$를 정역 스킴이라 하자.
$0 \to \mathcal{F} \to \mathcal{F}' \to \mathcal{F}'' \to 0$
을 준연접 $\mathcal{O}_X$-가군들의 짧은 완전열이라 하자. 그러면
$\mathcal{F}'$의 계수는 $\mathcal{F}$와 $\mathcal{F}''$의 계수의 합이다.
\end{lemma}

\begin{proof}
생략한다. More on Algebra, Lemma \ref{more-algebra-lemma-rank-additive}을
보라.
\end{proof}

\begin{lemma}
\label{divisors-lemma-rank-tensor}
$X$를 정역 스킴이라 하자. $\mathcal{F}$와 $\mathcal{G}$를 준연접
$\mathcal{O}_X$-가군이라 하자.
\begin{enumerate}
\item $\mathcal{F} \otimes_{\mathcal{O}_X} \mathcal{G}$의 계수는
$\mathcal{F}$와 $\mathcal{G}$의 계수의 곱이다.
\item $\mathcal{F}$가 유한 표시이면
$\SheafHom_{\mathcal{O}_X}(\mathcal{F}, \mathcal{G})$의 계수는
$\mathcal{F}$와 $\mathcal{G}$의 계수의 곱이다.
\end{enumerate}
\end{lemma}

\begin{proof}
생략한다. (2)가 의미를 갖는 것은
$\SheafHom_{\mathcal{O}_X}(\mathcal{F}, \mathcal{G})$가 Schemes,
Section \ref{schemes-section-quasi-coherent}의 항목 (10)에 의해
준연접이기 때문이다. 이 보조정리의 대수적 판본은 More on Algebra,
Lemma \ref{more-algebra-lemma-rank-tensor}이다.
\end{proof}

\begin{lemma}
\label{divisors-lemma-finite-module-rank}
$X$를 정역 스킴이라 하자. $\mathcal{F}$를 유한형 준연접
$\mathcal{O}_X$-가군이라 하자. 그러면 다음이 성립한다.
\begin{enumerate}
\item $\mathcal{F}$는 유한 계수 $r \geq 0$을 갖는다.
\item 공집합이 아닌 열린집합 $U \subset X$가 존재하여
$\mathcal{F}|_U$는 계수 $r$인 자유 가군이다.
\end{enumerate}
\end{lemma}

\begin{proof}
대수적 판본은 More on Algebra, Lemma
\ref{more-algebra-lemma-finite-module-rank}을 보라.
\end{proof}


\section{반사적 가군}
\label{divisors-section-reflexive}

\noindent
이 절은 국소 뇌터 스킴 위의 응집 가군에 대한 More on Algebra, Section
\ref{more-algebra-section-reflexive}의 대응물이다. 응집 가군을 다루는
이유는 $\SheafHom_{\mathcal{O}_X}(\mathcal{F}, \mathcal{G})$가 임의의 응집
$\mathcal{O}_X$-가군 $\mathcal{F}, \mathcal{G}$에 대해 응집이기 때문이다.
Modules, Lemma
\ref{modules-lemma-internal-hom-locally-kernel-direct-sum}을 보라.

\begin{definition}
\label{divisors-definition-reflexive}
$X$를 정역 국소 뇌터 스킴이라 하자. $\mathcal{F}$를 응집
$\mathcal{O}_X$-가군이라 하자. $\mathcal{F}$의 {\it 반사적 포락}은
$\mathcal{O}_X$-가군
$$
\mathcal{F}^{**} = \SheafHom_{\mathcal{O}_X}(
\SheafHom_{\mathcal{O}_X}(\mathcal{F}, \mathcal{O}_X), \mathcal{O}_X)
$$
이다. $\mathcal{F}$가 {\it 반사적}이라고 하는 것은 자연스러운 사상
$j : \mathcal{F} \longrightarrow \mathcal{F}^{**}$이 동형일 때이다.
\end{definition}

\noindent
Lemma \ref{divisors-lemma-dual-reflexive}에 의해 반사적 포락은 반사적
$\mathcal{O}_X$-가군이다. 더 일반적인 상황에서 같은 정의를 사용하여
반사적 가군을 정의할 수도 있지만, 이는 그다지 유용해 보이지 않는다.
통상적인 대수적 개념과 이를 비교하는 필수 보조정리는 다음과 같다.

\begin{lemma}
\label{divisors-lemma-check-reflexive-on-affines}
$X$를 정역 국소 뇌터 스킴이라 하자. $\mathcal{F}$를 응집
$\mathcal{O}_X$-가군이라 하자. 다음은 서로 동치이다.
\begin{enumerate}
\item $\mathcal{F}$는 반사적이다.
\item 아핀 열린집합 $U \subset X$에 대해 $\mathcal{F}(U)$는 반사적
$\mathcal{O}(U)$-가군이다.
\end{enumerate}
\end{lemma}

\begin{proof}
생략한다.
\end{proof}

\begin{remark}
\label{divisors-remark-different-reflexive}
$X$가 체 위의 유한형 스킴이면 때때로 반사적 가군에 다른 개념을 사용한다
(예를 들어 \cite[bottom of page 5 and Definition 1.1.9]{HL}를 보라). 이 다른
개념은 $R\SheafHom$을 쌍대화 복합체 $\omega_X^\bullet$로 가도록
사용하며 $\mathcal{O}_X$로 가게 하지 않는다. $X$가 고렌슈타인이 아닐 때 두 개념은
다를 수 있으므로 아마 다른 이름을 붙여야 한다. 예를 들어
$X = \Spec(k[t^3, t^4, t^5])$이면, 계산에 의해 쌍대화 층 $\omega_X$는
우리의 의미에서는 반사적이 아니지만 다른 의미에서는 반사적이다. 실제로
$\omega_X \to \SheafHom(\SheafHom(\omega_X, \omega_X), \omega_X)$는
동형이다.
\end{remark}

\begin{lemma}
\label{divisors-lemma-reflexive-torsion-free}
$X$를 정역 국소 뇌터 스킴이라 하자. $\mathcal{F}$를 응집
$\mathcal{O}_X$-가군이라 하자.
\begin{enumerate}
\item $\mathcal{F}$가 반사적이면 $\mathcal{F}$에는 꼬임이 없다.
\item 사상 $j : \mathcal{F} \longrightarrow \mathcal{F}^{**}$이 단사일
필요충분조건은 $\mathcal{F}$에 꼬임이 없는 것이다.
\end{enumerate}
\end{lemma}

\begin{proof}
생략한다. More on Algebra, Lemma
\ref{more-algebra-lemma-reflexive-torsion-free}를 보라.
\end{proof}

\begin{lemma}
\label{divisors-lemma-check-reflexive}
$X$를 정역 국소 뇌터 스킴이라 하자. $\mathcal{F}$를 응집
$\mathcal{O}_X$-가군이라 하자. 다음은 서로 동치이다.
\begin{enumerate}
\item $\mathcal{F}$는 반사적이다.
\item $\mathcal{F}_x$는 반사적 $\mathcal{O}_{X, x}$-가군이다. 이는 모든
$x \in X$에 대해 성립한다.
\item $\mathcal{F}_x$는 반사적 $\mathcal{O}_{X, x}$-가군이다. 이는 모든
닫힌점 $x \in X$에 대해 성립한다.
\end{enumerate}
\end{lemma}

\begin{proof}
Modules, Lemma \ref{modules-lemma-stalk-internal-hom}에 의해 (1)과 (2)는
동치이다. $X$의 모든 점은 닫힌점으로 특수화되므로(Properties, Lemma
\ref{properties-lemma-locally-Noetherian-closed-point}) (2)와 (3)도
동치이다.
\end{proof}

\begin{lemma}
\label{divisors-lemma-flat-pullback-reflexive}
$f : X \to Y$를 정역 국소 뇌터 스킴들의 평탄 사상이라 하자.
$\mathcal{G}$를 응집 반사적 $\mathcal{O}_Y$-가군이라 하자. 그러면
$f^*\mathcal{G}$는 응집 반사적 $\mathcal{O}_X$-가군이다.
\end{lemma}

\begin{proof}
생략한다. 대수적 대응물은 More on Algebra, Lemma
\ref{more-algebra-lemma-flat-pullback-reflexive}를 보라.
\end{proof}

\begin{lemma}
\label{divisors-lemma-sequence-reflexive}
$X$를 정역 국소 뇌터 스킴이라 하자.
$0 \to \mathcal{F} \to \mathcal{F}' \to \mathcal{F}''$을 응집
$\mathcal{O}_X$-가군들의 완전열이라 하자. $\mathcal{F}'$가 반사적이고
$\mathcal{F}''$에 꼬임이 없으면 $\mathcal{F}$는 반사적이다.
\end{lemma}

\begin{proof}
생략한다. More on Algebra, Lemma
\ref{more-algebra-lemma-sequence-reflexive}를 보라.
\end{proof}

\begin{lemma}
\label{divisors-lemma-dual-reflexive}
$X$를 정역 국소 뇌터 스킴이라 하자. $\mathcal{F}$, $\mathcal{G}$를 응집
$\mathcal{O}_X$-가군이라 하자. $\mathcal{G}$가 반사적이면
$\SheafHom_{\mathcal{O}_X}(\mathcal{F}, \mathcal{G})$는 반사적이다.
\end{lemma}

\begin{proof}
$\SheafHom_{\mathcal{O}_X}(\mathcal{F}, \mathcal{G})$가 Cohomology of
Schemes, Lemma \ref{coherent-lemma-tensor-hom-coherent}에 의해
응집이므로 명제는 의미가 있다. 명제가 참임을 보려면 More on Algebra,
Lemma \ref{more-algebra-lemma-dual-reflexive}를 보라. 일부 세부사항은
생략한다.
\end{proof}

\begin{remark}
\label{divisors-remark-tensor}
$X$를 정역 국소 뇌터 스킴이라 하자. Lemma
\ref{divisors-lemma-dual-reflexive} 덕분에 반사적 포락 $\mathcal{F}^{**}$은 응집
$\mathcal{O}_X$-가군에 대해 응집 반사적임을 안다. 범주 $\mathcal{C}$를
응집 반사적 $\mathcal{O}_X$-가군들의 범주라 하자. 반사적 포락을
취하는 것은 포함 함자 $\mathcal{C} \to \textit{Coh}(\mathcal{O}_X)$의
왼쪽 수반을 준다. $\mathcal{C}$가 핵과 여핵을 갖는 가법 범주임에
유의하자. 실제로 $\varphi : \mathcal{F} \to \mathcal{G}$가
$\mathcal{C}$ 안에서 주어지면, 통상적인 핵
$\Ker(\varphi)$은 반사적이고(Lemma \ref{divisors-lemma-sequence-reflexive}),
통상적인 여핵의 반사적 포락 $\Coker(\varphi)^{**}$은 $\mathcal{C}$에서의
여핵이다. 더욱이 $\mathcal{C}$는 텐서곱
$$
\mathcal{F} \otimes_\mathcal{C} \mathcal{G} =
(\mathcal{F} \otimes_{\mathcal{O}_X} \mathcal{G})^{**}
$$
을 물려받으며, 이는 결합적이고 대칭적이다. 임의의 세 대상
$\mathcal{F}, \mathcal{G}, \mathcal{H}$가 $\mathcal{C}$에 속할 때 항등식
$$
\Hom_\mathcal{C}(\mathcal{F} \otimes_\mathcal{C} \mathcal{G}, \mathcal{H}) =
\Hom_\mathcal{C}(\mathcal{F},
\SheafHom_{\mathcal{O}_X}(\mathcal{G}, \mathcal{H}))
$$
을 갖는다는 의미에서 내부 Hom이 있다. Modules, Lemma
\ref{modules-lemma-internal-hom}을 보라. $\mathcal{C}$에서 모든 대상
$\mathcal{F}$는 {\it 쌍대 대상}
$\SheafHom_{\mathcal{O}_X}(\mathcal{F}, \mathcal{O}_X)$을 갖는다.
$X$에 추가 조건을 주지 않으면
$$
\SheafHom_{\mathcal{O}_X}(\mathcal{F}, \mathcal{G}) \not \cong
\SheafHom_{\mathcal{O}_X}(\mathcal{F}, \mathcal{O}_X)
\otimes_\mathcal{C} \mathcal{G}
\quad\text{and}\quad
\mathcal{F} \otimes_\mathcal{C}
\SheafHom_{\mathcal{O}_X}(\mathcal{F}, \mathcal{O}_X)
\not \cong \mathcal{O}_X
$$
가 $\mathcal{F}, \mathcal{G}$의 계수가 $1$이고 이들이 $\mathcal{C}$에
속할 때 성립할 수 있다. 예를 만들기 위해 $X = \Spec(R)$라 하자. 여기서 $R$은 More on
Algebra, Example \ref{more-algebra-example-ring-not-S2}에서와 같고,
$\mathcal{F}, \mathcal{G}$는 $M$에 대응하는 가군들이라 하자. 계산은
생략한다.
\end{remark}

\begin{lemma}
\label{divisors-lemma-reflexive-depth-2}
$X$를 정역 국소 뇌터 스킴이라 하자. $\mathcal{F}$를 응집
$\mathcal{O}_X$-가군이라 하자. 다음은 서로 동치이다.
\begin{enumerate}
\item $\mathcal{F}$는 반사적이다.
\item 각 $x \in X$에 대해 다음 중 하나가 성립한다.
\begin{enumerate}
\item $\mathcal{F}_x$는 반사적 $\mathcal{O}_{X, x}$-가군이다.
\item $\text{depth}(\mathcal{F}_x) \geq 2$이다.
\end{enumerate}
\end{enumerate}
\end{lemma}

\begin{proof}
생략한다. More on Algebra, Lemma
\ref{more-algebra-lemma-reflexive-depth-2}를 보라.
\end{proof}

\begin{lemma}
\label{divisors-lemma-reflexive-S2}
$X$를 정역 국소 뇌터 스킴이라 하자. $\mathcal{F}$를 응집 반사적
$\mathcal{O}_X$-가군이라 하자. $x \in X$라 하자.
\begin{enumerate}
\item $\text{depth}(\mathcal{O}_{X, x}) \geq 2$이면
$\text{depth}(\mathcal{F}_x) \geq 2$이다.
\item $X$가 $(S_2)$이면 $\mathcal{F}$도 $(S_2)$이다.
\end{enumerate}
\end{lemma}

\begin{proof}
생략한다. More on Algebra, Lemma \ref{more-algebra-lemma-reflexive-S2}를
보라.
\end{proof}

\begin{lemma}
\label{divisors-lemma-reflexive-S2-extend}
$X$를 정역 국소 뇌터 스킴이라 하자. $j : U \to X$를 여집합이 $Z$인 열린
부분스킴이라 하자. $\mathcal{O}_{X, z}$의 깊이가 $\geq 2$라고 모든
$z \in Z$에 대해 가정하자. 그러면 $j^*$와 $j_*$는 응집 반사적
$\mathcal{O}_X$-가군들의 범주와 응집 반사적 $\mathcal{O}_U$-가군들의
범주 사이의 범주 동치를 정의한다.
\end{lemma}

\begin{proof}
$\mathcal{F}$를 응집 반사적 $\mathcal{O}_X$-가군이라 하자. $z \in Z$에
대해 줄기 $\mathcal{F}_z$의 깊이는 Lemma \ref{divisors-lemma-reflexive-S2}에 의해
$\geq 2$이다. 따라서 Lemma \ref{divisors-lemma-depth-2-hartog}에 의해
$\mathcal{F} \to j_*j^*\mathcal{F}$는 동형이다. 역으로 $\mathcal{G}$를
응집 반사적 $\mathcal{O}_U$-가군이라 하자. $j_*\mathcal{G}$가 응집
반사적 $\mathcal{O}_X$-가군임을 보이면 충분하다. 이를 증명하기 위해
$X$가 아핀이라고 가정해도 된다. Properties, Lemma
\ref{properties-lemma-lift-finite-presentation}에 의해 응집
$\mathcal{O}_X$-가군 $\mathcal{F}$가 존재하여
$\mathcal{G} = j^*\mathcal{F}$이다. $\mathcal{F}$를 그 반사적 포락으로 바꾼 뒤
$\mathcal{F}$가 반사적이라고 가정해도 된다(위의 논의, 특히 Lemma
\ref{divisors-lemma-dual-reflexive}를 보라). 위에 의해
$j_*\mathcal{G} = j_*j^*\mathcal{F} = \mathcal{F}$이며, 원하는 바이다.
\end{proof}

\noindent
스킴이 정규이면 반사적인 것은 꼬임이 없고 $(S_2)$인 것과 같다.

\begin{lemma}
\label{divisors-lemma-reflexive-over-normal}
$X$를 정역 국소 뇌터 정규 스킴이라 하자. $\mathcal{F}$를 응집
$\mathcal{O}_X$-가군이라 하자. 다음은 서로 동치이다.
\begin{enumerate}
\item $\mathcal{F}$는 반사적이다.
\item $\mathcal{F}$에는 꼬임이 없고 성질 $(S_2)$를 가진다.
\item 다음을 만족하는 열린 부분스킴 $j : U \to X$가 존재한다.
\begin{enumerate}
\item $X \setminus U$의 모든 기약 성분은 여차원이 $\geq 2$이며, 이는
$X$ 안에서의 여차원이다.
\item $j^*\mathcal{F}$는 유한 국소 자유이다.
\item $\mathcal{F} = j_*j^*\mathcal{F}$이다.
\end{enumerate}
\end{enumerate}
\end{lemma}

\begin{proof}
Lemma \ref{divisors-lemma-check-reflexive-on-affines}을 사용하면 (1)과 (2)의 동치는
More on Algebra, Lemma \ref{more-algebra-lemma-reflexive-over-normal}에서
따른다. $U \subset X$가 (3)에서와 같다고 하자. Properties, Lemma
\ref{properties-lemma-criterion-normal}에 의해
$\text{depth}(\mathcal{O}_{X, x}) \geq 2$임을 모든 $x \not \in U$에 대해
안다. 유한 국소 자유 가군은
반사적이므로 Lemma \ref{divisors-lemma-reflexive-S2-extend}에 의해 (3)은 (1)을
함의한다.

\medskip\noindent
(1)을 가정하자. $U \subset X$를 $j^*\mathcal{F} = \mathcal{F}|_U$가 유한
국소 자유가 되게 하는 최대 열린 부분스킴이라 하자. 그러면 (3)(b)가
성립한다. $x \in X$를 한 점이라 하자. $\mathcal{F}_x$가 자유
$\mathcal{O}_{X, x}$-가군이면 $x \in U$이다. Modules, Lemma
\ref{modules-lemma-finite-presentation-stalk-free}를 보라.
$\dim(\mathcal{O}_{X, x}) \leq 1$이면 $\mathcal{O}_{X, x}$는 체이거나
이산 부치환이다(Properties, Lemma
\ref{properties-lemma-criterion-normal}). 따라서 $\mathcal{F}_x$는
자유이다(More on Algebra, Lemma
\ref{more-algebra-lemma-dedekind-torsion-free-flat}). 그러므로
$x \not \in U \Rightarrow \dim(\mathcal{O}_{X, x}) \geq 2$이다. 이제
Properties, Lemma \ref{properties-lemma-codimension-local-ring}에 의해
(3)(a)가 성립한다. 앞에서 이미 사용한 Properties, Lemma
\ref{properties-lemma-criterion-normal}에 의해
$\text{depth}(\mathcal{O}_{X, x}) \geq 2$임도 $x \not \in U$에 대해
알 수 있으므로 (3)(c)는
Lemma \ref{divisors-lemma-reflexive-S2-extend}에서 따른다.
\end{proof}

\begin{lemma}
\label{divisors-lemma-describe-reflexive-hull}
$X$를 일반점이 $\eta$인 정역 국소 뇌터 정규 스킴이라 하자.
$\mathcal{F}$, $\mathcal{G}$를 응집 $\mathcal{O}_X$-가군이라 하자.
$T : \mathcal{G}_\eta \to \mathcal{F}_\eta$를 선형 사상이라 하자. 그러면
$T$가 사상 $\mathcal{G} \to \mathcal{F}^{**}$, 즉 $\mathcal{O}_X$-가군의 사상으로
확장될 필요충분조건은 다음과 같다.
\begin{itemize}
\item[(*)] 모든 $x \in X$ 중 $\dim(\mathcal{O}_{X, x}) = 1$인 것에 대해
$$
T\left(\Im(\mathcal{G}_x \to \mathcal{G}_\eta)\right) \subset
\Im(\mathcal{F}_x \to \mathcal{F}_\eta).
$$
\end{itemize}
\end{lemma}

\begin{proof}
$\mathcal{F}^{**}$에는 꼬임이 없고
$\mathcal{F}_\eta = \mathcal{F}^{**}_\eta$이므로 확장이 존재한다면
유일하다. 따라서 $X$의 열린 덮개의 각 원소 위에서 보조정리를 증명하면
충분하다. 즉 $X$가 아핀이라고 가정해도 된다. 이 경우 다음 대수 문제를
묻는 것이다. $R$을 분수체가 $K$인 뇌터 정규 정역이라 하고 $M$, $N$을
유한 $R$-가군이라 하며,
$T : M \otimes_R K \to N \otimes_R K$를 $K$-선형 사상이라 하자. 언제
$T$가 사상 $N \to M^{**}$으로 확장되는가? More on Algebra, Lemma
\ref{more-algebra-lemma-describe-reflexive-hull}에 의해 이는 $N_\mathfrak p$가
$(M/M_{tors})_\mathfrak p$ 안으로 사상될 필요충분조건이다. 여기서 높이가
$1$인 모든 소아이디얼 $\mathfrak p$를 취하며 이들은 $R$의 소아이디얼이다. 이는 정확히
보조정리의 조건 $(*)$이다.
\end{proof}

\begin{lemma}
\label{divisors-lemma-reflexive-over-regular-dim-2}
$X$를 차원이 $\leq 2$인 정칙 스킴이라 하자. $\mathcal{F}$를 응집
$\mathcal{O}_X$-가군이라 하자. 다음은 서로 동치이다.
\begin{enumerate}
\item $\mathcal{F}$는 반사적이다.
\item $\mathcal{F}$는 유한 국소 자유이다.
\end{enumerate}
\end{lemma}

\begin{proof}
유한 국소 자유 가군이 반사적이라는 것은 자명하다. 역으로,
$\mathcal{F}$가 반사적이면 $\mathcal{F}_x$가 자유
$\mathcal{O}_{X, x}$-가군임을 모든 $x \in X$에 대해 보이겠다. 이는 Algebra, Lemma
\ref{algebra-lemma-finite-projective}와 $\mathcal{F}$가 응집이라는 사실에
의해 충분하다. $\dim(\mathcal{O}_{X, x}) = 0$이면
$\mathcal{O}_{X, x}$는 체이므로 명제는 자명하다.
$\dim(\mathcal{O}_{X, x}) = 1$이면 $\mathcal{O}_{X, x}$는 이산 부치환이다
(Algebra, Lemma \ref{algebra-lemma-characterize-dvr}). 또한
$\mathcal{F}_x$에는 꼬임이 없다. 따라서 More on Algebra, Lemma
\ref{more-algebra-lemma-dedekind-torsion-free-flat}에 의해
$\mathcal{F}_x$는 자유이다. $\dim(\mathcal{O}_{X, x}) = 2$이면
$\mathcal{O}_{X, x}$는 차원 $2$인 정칙 국소환이다. More on Algebra,
Lemma \ref{more-algebra-lemma-reflexive-over-normal}에 의해
$\mathcal{F}_x$의 깊이는 $\geq 2$이다. 따라서 Algebra, Lemma
\ref{algebra-lemma-regular-mcm-free}에 의해 $\mathcal{F}_x$는 자유이다.
\end{proof}


\section{유효 카르티에 인자}
\label{divisors-section-effective-Cartier-divisors}

\noindent
다른 유형의 인자를 정의하기 전에 유효 카르티에 인자의 개념을 정의한다.

\begin{definition}
\label{divisors-definition-effective-Cartier-divisor}
$S$를 스킴이라 하자.
\begin{enumerate}
\item $S$의 {\it 국소 주인 닫힌 부분스킴}은 그 아이디얼 층이 국소적으로
한 원소에 의해 생성되는 닫힌 부분스킴이다.
\item $S$ 위의 {\it 유효 카르티에 인자}는 닫힌 부분스킴 $D \subset S$ 중
그 아이디얼 층 $\mathcal{I}_D \subset \mathcal{O}_S$가 가역
$\mathcal{O}_S$-가군인 것이다.
\end{enumerate}
\end{definition}

\noindent
따라서 유효 카르티에 인자는 국소 주인 닫힌 부분스킴이지만 그 역은 항상
참은 아니다. 유효 카르티에 인자는 가능한 가장 강한 의미에서 순수
여차원 $1$인 닫힌 부분스킴이다. 즉 국소적으로 영인자가 아닌 하나의
원소에 의해 잘려 나온다. 특히 어디에서도 조밀하지 않다.

\begin{lemma}
\label{divisors-lemma-characterize-effective-Cartier-divisor}
$S$를 스킴이라 하자. $D \subset S$를 닫힌 부분스킴이라 하자. 다음은 서로
동치이다.
\begin{enumerate}
\item 부분스킴 $D$는 $S$ 위의 유효 카르티에 인자이다.
\item 모든 $x \in D$에 대해 아핀 열린집합 $\Spec(A) = U \subset S$가
$x$의 근방으로 존재하여, $U \cap D = \Spec(A/(f))$이다. 여기서
$f \in A$는 영인자가 아닌 원소이다.
\end{enumerate}
\end{lemma}

\begin{proof}
(1)을 가정하자. 모든 $x \in D$에 대해 아핀 열린집합
$\Spec(A) = U \subset S$가 $x$의 근방으로 존재하여
$\mathcal{I}_D|_U \cong \mathcal{O}_U$이다. 다시 말해 제한
$f \in \Gamma(U, \mathcal{I}_D)$가 존재하여 $\mathcal{I}_D|_U$를
자유롭게 생성한다. 따라서 $f \in A$이고 곱셈
사상 $f : A \to A$는 단사이다. 또한 $\mathcal{I}_D$가 준연접이므로
$D \cap U = \Spec(A/(f))$임을 안다.

\medskip\noindent
(2)를 가정하자. $x \in D$라 하자. 가정에 의해 아핀 열린집합
$\Spec(A) = U \subset S$가 $x$의 근방으로 존재하여,
$U \cap D = \Spec(A/(f))$이다. 여기서 $f \in A$는 영인자가 아닌
원소이다. 그러면
$\mathcal{I}_D|_U \cong \mathcal{O}_U$이다. 실제로 이는
$\widetilde{(f)} \cong \widetilde{A} \cong \mathcal{O}_U$와 같다. 물론
$\mathcal{I}_D$를 열린 부분스킴 $S \setminus D$로 제한한 것은
$\mathcal{O}_{S \setminus D}$와 동형이다. 따라서 $\mathcal{I}_D$는 가역
$\mathcal{O}_S$-가군이다.
\end{proof}

\begin{lemma}
\label{divisors-lemma-complement-locally-principal-closed-subscheme}
$S$를 스킴이라 하자. $Z \subset S$를 국소 주인 닫힌 부분스킴이라 하자.
$U = S \setminus Z$라 하자. 그러면 $U \to S$는 아핀 사상이다.
\end{lemma}

\begin{proof}
이 문제는 $S$ 위에서 국소적이다. Morphisms, Lemmas
\ref{morphisms-lemma-characterize-affine}를 보라. 따라서
$S = \Spec(A)$이고 $Z = V(f)$이며 $f \in A$라고 가정해도 된다. 이 경우
$U = D(f) = \Spec(A_f)$는 아핀이므로 $U \to S$는 아핀이다.
\end{proof}

\begin{lemma}
\label{divisors-lemma-complement-effective-Cartier-divisor}
$S$를 스킴이라 하자. $D \subset S$를 유효 카르티에 인자라 하자.
$U = S \setminus D$라 하자. 그러면 $U \to S$는 아핀 사상이고 $U$는
$S$에서 스킴론적으로 조밀하다.
\end{lemma}

\begin{proof}
아핀성은 Lemma
\ref{divisors-lemma-complement-locally-principal-closed-subscheme}이다. 조밀성 문제는
$S$ 위에서 국소적이다. Morphisms, Lemma
\ref{morphisms-lemma-characterize-scheme-theoretically-dense}를 보라. 따라서
$S = \Spec(A)$이고 $D$가 영인자가 아닌 원소 $f \in A$에 대응한다고
가정해도 된다. Lemma
\ref{divisors-lemma-characterize-effective-Cartier-divisor}를 보라. 따라서
$A \subset A_f$이고, 이는 $U \subset S$가 스킴론적으로 조밀함을
함의한다. Morphisms, Example
\ref{morphisms-example-scheme-theoretic-closure}를 보라.
\end{proof}

\begin{lemma}
\label{divisors-lemma-effective-Cartier-makes-dimension-drop}
$S$를 스킴이라 하자. $D \subset S$를 유효 카르티에 인자라 하자.
$s \in D$라 하자. $\dim_s(S) < \infty$이면
$\dim_s(D) < \dim_s(S)$이다.
\end{lemma}

\begin{proof}
$\dim_s(S) < \infty$라고 가정하자. 아핀 열린 근방
$U = \Spec(A) \subset S$를 $s$에 대해 택하되, $\dim(U) = \dim_s(S)$이고
$D = V(f)$가 되게 하자. 여기서 $f \in A$는 영인자가 아닌 원소이다(Lemma
\ref{divisors-lemma-characterize-effective-Cartier-divisor}를 보라). $\dim(U)$는
환 $A$의 크룰 차원이고 $\dim(U \cap D)$는 환 $A/(f)$의 크룰 차원임을
상기하자. 그러면 $f$는 $A$의 어떤 극소 소아이디얼에도 속하지 않는다.
따라서 $A/(f)$ 안의 소아이디얼들의 임의의 극대 사슬을 $A$ 안의
소아이디얼들의 사슬로 보면 극소 소아이디얼을 덧붙여 연장할 수 있다.
\end{proof}

\begin{definition}
\label{divisors-definition-sum-effective-Cartier-divisors}
$S$를 스킴이라 하자. 유효 카르티에 인자 $D_1$, $D_2$가 $S$ 위에
주어졌을 때 $D = D_1 + D_2$를 $S$의 닫힌 부분스킴으로 두되, 준연접
아이디얼 층 $\mathcal{I}_{D_1}\mathcal{I}_{D_2} \subset \mathcal{O}_S$에
대응하게 한다. 이를 {\it 유효 카르티에 인자 $D_1$과
$D_2$의 합}이라고 한다.
\end{definition}

\noindent
합 $\sum n_iD_i$를 정의할 수 있음은 분명하다. 여기서 유한 개의 유효
카르티에 인자 $D_i$는 $X$ 위에 있고 $n_i$는 음이 아닌 정수이다.

\begin{lemma}
\label{divisors-lemma-sum-effective-Cartier-divisors}
두 유효 카르티에 인자의 합은 유효 카르티에 인자이다.
\end{lemma}

\begin{proof}
생략한다. 국소적으로 $f_1, f_2 \in A$는 영인자가 아닌 원소들이며,
그러면 $f_1f_2 \in A$도 영인자가 아닌 원소이다.
\end{proof}

\begin{lemma}
\label{divisors-lemma-difference-effective-Cartier-divisors}
$X$를 스킴이라 하자. $D, D'$을 $X$ 위의 두 유효 카르티에 인자라 하자.
$D \subset D'$이면($X$의 닫힌 부분스킴으로서) 유효 카르티에 인자
$D''$이 존재하여 $D' = D + D''$이다.
\end{lemma}

\begin{proof}
생략한다.
\end{proof}

\begin{lemma}
\label{divisors-lemma-sum-closed-subschemes-effective-Cartier}
$X$를 스킴이라 하자. $Z, Y$를 $X$의 두 닫힌 부분스킴이라 하고 그
아이디얼 층을 $\mathcal{I}$와 $\mathcal{J}$라 하자. $\mathcal{I}\mathcal{J}$가
유효 카르티에 인자 $D \subset X$를 정의하면 $Z$와 $Y$는 유효 카르티에
인자이고 $D = Z + Y$이다.
\end{lemma}

\begin{proof}
Lemma \ref{divisors-lemma-characterize-effective-Cartier-divisor}를 적용하면 다음
대수적 상황을 얻는다. $A$는 환이고 $I, J \subset A$는 아이디얼들이며,
$f \in A$는 영인자가 아닌 원소로서 $IJ = (f)$를 만족한다. 따라서
결과는 Algebra, Lemma \ref{algebra-lemma-product-ideals-principal}에서
따른다.
\end{proof}

\begin{lemma}
\label{divisors-lemma-sum-effective-Cartier-divisors-union}
$X$를 스킴이라 하자. $D, D' \subset X$를 유효 카르티에 인자들이라 하되,
스킴론적 교집합 $D \cap D'$이 $D'$ 위의 유효 카르티에 인자라고 하자.
그러면 $D + D'$은 $D$와 $D'$의 스킴론적 합집합이다.
\end{lemma}

\begin{proof}
스킴론적 교집합과 합집합의 정의는 Morphisms, Definition
\ref{morphisms-definition-scheme-theoretic-intersection-union}을 보라. 이
보조정리를 국소적으로 증명하면(Lemma
\ref{divisors-lemma-characterize-effective-Cartier-divisor}를 사용하여) 다음 대수
문제를 얻는다. 환 $A$와 영인자가 아닌 원소들 $f_1, f_2 \in A$가
주어지고 $f_1$이 $A/f_2A$에서 영인자가 아닌 원소로 사상된다고 하자.
$f_1A \cap f_2A = f_1f_2A$임을 보여라. 직접적인 논증은 생략한다.
\end{proof}

\noindent
스킴의 임의의 사상 아래에서 닫힌 부분스킴의 역상을 Schemes, Definition
\ref{schemes-definition-inverse-image-closed-subscheme}에서 정의했음을
상기하자.

\begin{lemma}
\label{divisors-lemma-pullback-locally-principal}
$f : S' \to S$를 스킴의 사상이라 하자. $Z \subset S$를 국소 주인 닫힌
부분스킴이라 하자. 그러면 역상 $f^{-1}(Z)$는 $S'$의 국소 주인 닫힌
부분스킴이다.
\end{lemma}

\begin{proof}
생략한다.
\end{proof}

\begin{definition}
\label{divisors-definition-pullback-effective-Cartier-divisor}
$f : S' \to S$를 스킴의 사상이라 하자. $D \subset S$를 유효 카르티에
인자라 하자. {\it $D$의 $f$에 의한 당김이 정의된다}고 하는 것은 닫힌
부분스킴 $f^{-1}(D) \subset S'$이 유효 카르티에 인자일 때이다. 이 경우 이를
$f^*D$ 또는 $f^{-1}(D)$로 표기하고 {\it 유효 카르티에 인자의 당김}이라
한다.
\end{definition}

\noindent
$f^{-1}(D)$가 유효 카르티에 인자라는 조건은 실제로 흔히 만족된다. 한
예가 되는 보조정리는 다음과 같다.

\begin{lemma}
\label{divisors-lemma-pullback-effective-Cartier-defined}
$f : X \to Y$를 스킴의 사상이라 하자. $D \subset Y$를 유효 카르티에
인자라 하자. 다음 각 경우에 $D$의 $f$에 의한 당김이 정의된다.
\begin{enumerate}
\item $f(x) \not \in D$이다. 여기서 $x$는 $X$의 임의의 약연관점이다.
\item $X$, $Y$는 정역이고 $f$는 우세이다.
\item $X$는 축약이고 $f(\xi) \not \in D$이다. 여기서 $\xi$는 $X$의
임의의 기약 성분의 임의의 일반점이다.
\item $X$는 국소 뇌터이고 $f(x) \not \in D$이다. 여기서 $x$는 $X$의
임의의 연관점이다.
\item $X$는 국소 뇌터이고 매장점이 없으며 $f(\xi) \not \in D$이다.
여기서 $\xi$는 $X$의 기약 성분의 임의의 일반점이다.
\item $f$는 평탄하다.
\item 필요에 따라 여기에 더 추가한다.
\end{enumerate}
\end{lemma}

\begin{proof}
이 문제는 $X$ 위에서 국소적이므로 $X = \Spec(A)$, $Y = \Spec(R)$이고
$f$가 $\varphi : R \to A$에 의해 주어지며, $D = \Spec(R/(t))$이고
$t \in R$가 영인자가 아닌 원소인 경우로 환원한다. 각 경우의 목표는
$\varphi(t) \in A$가 영인자가 아닌 원소임을 보이는 것이다.

\medskip\noindent
경우 (1)은 Algebra, Lemma
\ref{algebra-lemma-weakly-ass-zero-divisors}에서 따른다. 경우 (4)는 Lemma
\ref{divisors-lemma-ass-weakly-ass}에 의해 (1)의 특수한 경우이다. 경우 (5)는
(4)와 정의들에서 따른다. 경우 (3)은 Lemma
\ref{divisors-lemma-weakass-reduced}에 의해 (1)의 특수한 경우이다. 경우 (2)는
(3)의 특수한 경우이다. $R \to A$가 평탄이면 $t : R \to R$가
단사라는 사실로부터 $t : A \to A$가 단사임을 얻는다. 이는 (6)을
증명한다.
\end{proof}

\begin{lemma}
\label{divisors-lemma-pullback-effective-Cartier-divisors-additive}
$f : S' \to S$를 스킴의 사상이라 하자. $D_1$, $D_2$를 $S$ 위의 유효
카르티에 인자들이라 하자. $D_1$과 $D_2$의 당김들이 정의되면
$D = D_1 + D_2$의 당김도 정의되고
$f^*D = f^*D_1 + f^*D_2$이다.
\end{lemma}

\begin{proof}
생략한다.
\end{proof}


\section{유효 카르티에 인자와 가역층}
\label{divisors-section-effective-Cartier-invertible}

\noindent
유효 카르티에 인자는 가역 아이디얼 층을 가지므로(Definition
\ref{divisors-definition-effective-Cartier-divisor}) 다음 정의는 의미가 있다.

\begin{definition}
\label{divisors-definition-invertible-sheaf-effective-Cartier-divisor}
$S$를 스킴이라 하자. $D \subset S$를 아이디얼 층이 $\mathcal{I}_D$인
유효 카르티에 인자라 하자.
\begin{enumerate}
\item {\it 가역층 $\mathcal{O}_S(D)$, 즉 $D$에 결부된 것}을 다음과 같이
정의한다.
$$
\mathcal{O}_S(D) =
\SheafHom_{\mathcal{O}_S}(\mathcal{I}_D, \mathcal{O}_S) =
\mathcal{I}_D^{\otimes -1}.
$$
\item 보통 $1$ 또는 $1_D$로 표기하는 {\it 표준 절단}은
$\mathcal{O}_S(D)$의 전역 절단 중 포함 사상
$\mathcal{I}_D \to \mathcal{O}_S$에 대응하는 것이다.
\item 다음과 같이 쓴다.
$\mathcal{O}_S(-D) = \mathcal{O}_S(D)^{\otimes -1} = \mathcal{I}_D$.
\item 두 번째 유효 카르티에 인자 $D' \subset S$가 주어지면 다음과 같이
정의한다.
$\mathcal{O}_S(D - D') =
\mathcal{O}_S(D) \otimes_{\mathcal{O}_S} \mathcal{O}_S(-D')$.
\end{enumerate}
\end{definition}

\noindent
몇 가지 설명을 덧붙이자. 대응 $D \mapsto \mathcal{O}_S(D)$은 유효
카르티에 인자들의 덧셈(Definition
\ref{divisors-definition-sum-effective-Cartier-divisors})을 $S$의 피카르 군에서의
덧셈으로 보낸다는 것을 아래에서 볼 것이다(Lemma
\ref{divisors-lemma-invertible-sheaf-sum-effective-Cartier-divisors}). 그러나 위
정의의 표현 $D - D'$은 기하학적 의미를 갖지 않는다. 더 정확히 말해,
$S$ 위의 유효 카르티에 인자들의 집합을 영원이 공인 유효 카르티에
인자인 가환 모노이드 $\text{EffCart}(S)$로 생각할 수 있다. 그러면 대응
$(D, D') \mapsto \mathcal{O}_S(D - D')$은 군 준동형
$$
\text{EffCart}(S)^{gp} \longrightarrow \Pic(S)
$$
을 정의하며, 여기서 왼쪽은 $\text{EffCart}(S)$의 군 완성이다. 다시 말해
$\mathcal{O}_S(D - D')$라고 쓸 때 $D - D'$을
$\text{EffCart}(S)^{gp}$의 원소로 생각할 수 있다.

\begin{lemma}
\label{divisors-lemma-conormal-effective-Cartier-divisor}
$S$를 스킴이라 하고 $D \subset S$를 유효 카르티에 인자라 하자. 그러면
여법선층은 $\mathcal{C}_{D/S} = \mathcal{I}_D|_D =
\mathcal{O}_S(-D)|_D$이고 법선층은
$\mathcal{N}_{D/S} = \mathcal{O}_S(D)|_D$이다.
\end{lemma}

\begin{proof}
이는 Morphisms, Lemma \ref{morphisms-lemma-affine-conormal}에서 따른다.
\end{proof}

\begin{lemma}
\label{divisors-lemma-ses-add-divisor}
$X$를 스킴이라 하자. $D, D' \subset X$를 $D \subset D'$인 유효
카르티에 인자들이라 하고, $C \subset X$를 $D' = D + C$가 되게 하는
유효 카르티에 인자라 하자(이는 Lemma
\ref{divisors-lemma-difference-effective-Cartier-divisors}에 의해 존재한다).
그러면 $\mathcal{O}_X$-가군들의 짧은 완전열
$$
0 \to \mathcal{O}_X(-D)|_C \to \mathcal{O}_{D'} \to \mathcal{O}_D \to 0
$$
이 존재한다.
\end{lemma}

\begin{proof}
보조정리의 명제와 증명에서 Morphisms, Lemma
\ref{morphisms-lemma-i-star-equivalence}의 동치를 사용하여 $X$의 닫힌
부분스킴 위의 준연접 가군을 $X$ 위의 준연접 가군으로 생각한다.
$\mathcal{I}$를 $D$의 $D'$ 안에서의 아이디얼 층이라 하고
$\mathcal{I}_D = \mathcal{O}_X(-D)$를 $D$의 $X$ 안에서의 아이디얼 층이라
하자. 그러면 표준 전사 $\mathcal{I}_D \to \mathcal{I}$가 존재한다.
따라서 이 사상의 핵이 $\mathcal{I}_C\mathcal{I}_D$와 같음을 보이면
충분하다(표기는 자명하다). 이는 국소적으로 확인할 수 있으므로,
$X = \Spec(A)$, $D = \Spec(A/fA)$, $C = \Spec(A/gA)$이고
$f, g \in A$가 영인자가 아닌 원소들이라고 가정해도 된다(Lemma
\ref{divisors-lemma-characterize-effective-Cartier-divisor}). 그러면
$\mathcal{I}_D$는 $fA$에 대응하고, 영인자가 아닌 원소 $fg \in A$는
$D'$의 아이디얼을 생성하며, $\mathcal{I}$는 $I = fA/fgA$에 대응한다.
따라서 결과는 분명하다.
\end{proof}

\begin{lemma}
\label{divisors-lemma-invertible-sheaf-sum-effective-Cartier-divisors}
$S$를 스킴이라 하자. $D_1$, $D_2$를 $S$ 위의 유효 카르티에 인자들이라
하자. $D = D_1 + D_2$라 하자. 그러면 $1_{D_1} \otimes 1_{D_2}$를
$1_D$로 보내는 유일한 동형
$$
\mathcal{O}_S(D_1) \otimes_{\mathcal{O}_S} \mathcal{O}_S(D_2)
\longrightarrow
\mathcal{O}_S(D)
$$
이 존재한다.
\end{lemma}

\begin{proof}
생략한다.
\end{proof}

\begin{lemma}
\label{divisors-lemma-pullback-effective-Cartier-divisors}
$f : S' \to S$를 스킴의 사상이라 하자. $D$를 $S$ 위의 유효 카르티에
인자라 하자. $D$의 당김이 정의되면
$f^*\mathcal{O}_S(D) = \mathcal{O}_{S'}(f^*D)$이고 표준 절단 $1_D$은
표준 절단 $1_{f^*D}$로 당겨진다.
\end{lemma}

\begin{proof}
생략한다.
\end{proof}

\begin{definition}
\label{divisors-definition-regular-section}
$(X, \mathcal{O}_X)$를 국소 환 달린 공간이라 하자. $\mathcal{L}$을 $X$
위의 가역층이라 하자. 전역 절단 $s \in \Gamma(X, \mathcal{L})$이
{\it 정칙 절단}이라고 하는 것은 사상 $\mathcal{O}_X \to \mathcal{L}$,
$f \mapsto fs$가 단사일 때이다.
\end{definition}

\begin{lemma}
\label{divisors-lemma-regular-section-structure-sheaf}
$X$를 국소 환 달린 공간이라 하자. $f \in \Gamma(X, \mathcal{O}_X)$라
하자. 다음은 서로 동치이다.
\begin{enumerate}
\item $f$는 정칙 절단이다.
\item 임의의 $x \in X$에 대해 상 $f \in \mathcal{O}_{X, x}$는 영인자가
아닌 원소이다.
\end{enumerate}
$X$가 스킴이면 이들은 다음과도 동치이다.
\begin{enumerate}
\item[(3)] 임의의 아핀 열린집합 $\Spec(A) = U \subset X$에 대해 상
$f \in A$는 영인자가 아닌 원소이다.
\item[(4)] 아핀 열린 덮개 $X = \bigcup \Spec(A_i)$가 존재하여 $f$의 상은
$A_i$에서 영인자가 아닌 원소이다. 이는 모든 $i$에 대해 성립한다.
\end{enumerate}
\end{lemma}

\begin{proof}
생략한다.
\end{proof}

\noindent
$s$가 가역 $\mathcal{O}_X$-가군 $\mathcal{L}$의 전역 절단이면 이를
$\mathcal{O}_X$-가군 사상 $s : \mathcal{O}_X \to \mathcal{L}$로 볼 수
있음에 유의하자. 따라서 그 쌍대는 사상
$s : \mathcal{L}^{\otimes -1} \to \mathcal{O}_X$이다. 쌍대 가역층의
정의는 Modules, Definition \ref{modules-definition-powers}를 보라.

\begin{definition}
\label{divisors-definition-zero-scheme-s}
$X$를 스킴이라 하자. $\mathcal{L}$을 가역층이라 하자. 전역 절단
$s \in \Gamma(X, \mathcal{L})$를 택하자. $s$의 {\it 영점 스킴}은 닫힌
부분스킴 $Z(s) \subset X$이며, 준연접 아이디얼 층
$\mathcal{I} \subset \mathcal{O}_X$, 즉 사상
$s : \mathcal{L}^{\otimes -1} \to \mathcal{O}_X$의 상에 의해 정의된다.
\end{definition}

\begin{lemma}
\label{divisors-lemma-zero-scheme}
$X$를 스킴이라 하자. $\mathcal{L}$을 가역층이라 하자.
$s \in \Gamma(X, \mathcal{L})$라 하자.
\begin{enumerate}
\item 닫힌 몰입 $i : Z \to X$ 중
$i^*s \in \Gamma(Z, i^*\mathcal{L})$가 영이 되게 하는 것들을
포함관계로 순서 짓자. 영점 스킴 $Z(s)$은 이
순서집합의 최대 원소이다.
\item 스킴의 임의의 사상 $f : Y \to X$에 대해, $f^*s = 0$이
$\Gamma(Y, f^*\mathcal{L})$에서 성립할 필요충분조건은 $f$가
$Z(s)$를 경유하는 것이다.
\item 영점 스킴 $Z(s)$은 국소 주인 닫힌 부분스킴이다.
\item 영점 스킴 $Z(s)$이 유효 카르티에 인자일 필요충분조건은 $s$가
$\mathcal{L}$의 정칙 절단인 것이다.
\end{enumerate}
\end{lemma}

\begin{proof}
생략한다.
\end{proof}

\begin{lemma}
\label{divisors-lemma-characterize-OD}
\begin{slogan}
스킴 위의 유효 카르티에 인자는 고정된 정칙 전역 절단을 지닌 가역층과
같은 것이다.
\end{slogan}
$X$를 스킴이라 하자.
\begin{enumerate}
\item $D \subset X$가 유효 카르티에 인자이면 표준 절단 $1_D$, 즉
$\mathcal{O}_X(D)$의 표준 절단은 정칙이다.
\item 역으로 $s$가 가역층 $\mathcal{L}$의 정칙 절단이면, 유일한 유효
카르티에 인자 $D = Z(s) \subset X$와 유일한 동형
$\mathcal{O}_X(D) \to \mathcal{L}$이 존재하며, 그 동형은 $1_D$를 $s$로
보낸다.
\end{enumerate}
구성 $D \mapsto (\mathcal{O}_X(D), 1_D)$와
$(\mathcal{L}, s) \mapsto Z(s)$은 서로 역인 사상
$$
\left\{
\begin{matrix}
\text{effective Cartier divisors on }X
\end{matrix}
\right\}
\leftrightarrow
\left\{
\begin{matrix}
\text{isomorphism classes of pairs }(\mathcal{L}, s)\\
\text{consisting of an invertible }
\mathcal{O}_X\text{-module}\\
\mathcal{L}\text{ and a regular global section }s
\end{matrix}
\right\}
$$
을 준다.
\end{lemma}

\begin{proof}
생략한다.
\end{proof}

\begin{remark}
\label{divisors-remark-ses-regular-section}
$X$를 스킴, $\mathcal{L}$을 가역 $\mathcal{O}_X$-가군, $s$를
$\mathcal{L}$의 정칙 절단이라 하자. 그러면 영점 스킴 $D = Z(s)$은 $X$
위의 유효 카르티에 인자이고 짧은 완전열들
$$
0 \to \mathcal{O}_X \to \mathcal{L} \to i_*(\mathcal{L}|_D) \to 0
\quad\text{and}\quad
0 \to \mathcal{L}^{\otimes -1} \to \mathcal{O}_X \to i_*\mathcal{O}_D \to 0.
$$
이 존재한다. 유효 카르티에 인자 $D \subset X$가 주어지면 Lemmas
\ref{divisors-lemma-characterize-OD}와
\ref{divisors-lemma-conormal-effective-Cartier-divisor}를 사용하여
$$
0 \to \mathcal{O}_X \to \mathcal{O}_X(D) \to i_*(\mathcal{N}_{D/X}) \to 0
\quad\text{and}\quad
0 \to \mathcal{O}_X(-D) \to \mathcal{O}_X \to i_*(\mathcal{O}_D) \to 0
$$
을 얻는다.
\end{remark}


\section{뇌터 스킴 위의 유효 카르티에 인자}
\label{divisors-section-Noetherian-effective-Cartier}

\noindent
국소 뇌터인 상황에서는 유효 카르티에 인자와 정칙 절단에 관한 논의의
대부분이 다소 단순해진다.

\begin{lemma}
\label{divisors-lemma-regular-section-associated-points}
$X$를 국소 뇌터 스킴이라 하자. $\mathcal{L}$을 가역
$\mathcal{O}_X$-가군이라 하자. $s \in \Gamma(X, \mathcal{L})$라 하자.
그러면 $s$가 정칙 절단일 필요충분조건은 $s$가 $X$의 연관점들에서
사라지지 않는 것이다.
\end{lemma}

\begin{proof}
생략한다. 힌트: $\mathcal{L}$이 자명한 아핀 경우로 환원한 뒤 Lemma
\ref{divisors-lemma-regular-section-structure-sheaf}와 Algebra, Lemma
\ref{algebra-lemma-ass-zero-divisors}를 사용하라.
\end{proof}

\begin{lemma}
\label{divisors-lemma-effective-Cartier-in-points}
$X$를 국소 뇌터 스킴이라 하자. $D \subset X$를 준연접 아이디얼 층
$\mathcal{I} \subset \mathcal{O}_X$에 대응하는 닫힌 부분스킴이라 하자.
\begin{enumerate}
\item 모든 $x \in D$에 대해 아이디얼
$\mathcal{I}_x \subset \mathcal{O}_{X, x}$이 한 원소에 의해 생성될 수
있으면 $D$는 국소 주이다.
\item 모든 $x \in D$에 대해 아이디얼
$\mathcal{I}_x \subset \mathcal{O}_{X, x}$이 영인자가 아닌 하나의
원소에 의해 생성될 수 있으면 $D$는 유효 카르티에 인자이다.
\end{enumerate}
\end{lemma}

\begin{proof}
$\Spec(A)$를 한 점 $x \in D$의 아핀 근방이라 하자. 소아이디얼
$\mathfrak p \subset A$가 $x$에 대응한다고 하자. 아이디얼 $I \subset A$가
$D$의 $\Spec(A)$ 위 자취를 정의한다고 하자. $A$는 뇌터이므로
($X$가 국소 뇌터이기 때문에) 아이디얼 $I$는 유한 개의 원소에 의해
생성된다. 이를 $I = (f_1, \ldots, f_r)$라 하자. (1)의 가정 아래
$I_\mathfrak p = (f)$이며 어떤 $f \in A_\mathfrak p$가 이를 생성한다.
따라서 $f_i = g_i f$이며 어떤 $g_i \in A_\mathfrak p$가 존재한다.
$g_i = a_i/h_i$ 및 $f = f'/h$로 쓰자. 여기서
$a_i, h_i, f', h \in A$이고 $h_i, h \not \in \mathfrak p$이다. 그러면
$I_{h_1 \ldots h_r h} \subset A_{h_1 \ldots h_r h}$는 주아이디얼이다.
실제로 $f'$에 의해 생성된다. 이는 (1)을 증명한다. (2)를 위해
$I = (f)$라고 가정해도 된다. 가정은 $f$의 상이 $A_\mathfrak p$에서
영인자가 아닌 원소임을 함의한다. 그러면 Algebra, Lemma
\ref{algebra-lemma-regular-sequence-in-neighbourhood}에 의해 $f$는 $x$의
한 근방에서 영인자가 아닌 원소이다. 이는 (2)를 증명한다.
\end{proof}

\begin{lemma}
\label{divisors-lemma-effective-Cartier-codimension-1}
$X$를 국소 뇌터 스킴이라 하자.
\begin{enumerate}
\item $D \subset X$를 국소 주인 닫힌 부분스킴이라 하자. $\xi \in D$를
$D$의 한 기약 성분의 일반점이라 하자. 그러면
$\dim(\mathcal{O}_{X, \xi}) \leq 1$이다.
\item $D \subset X$를 유효 카르티에 인자라 하자. $\xi \in D$를 $D$의
한 기약 성분의 일반점이라 하자. 그러면
$\dim(\mathcal{O}_{X, \xi}) = 1$이다.
\end{enumerate}
\end{lemma}

\begin{proof}
(1)의 증명. 가정에 의해 $X = \Spec(A)$이고
$D = \Spec(A/(f))$이며 $A$는 뇌터 환이고 $f \in A$라고 가정해도 된다.
$\xi$에 대응하는 소아이디얼을 $\mathfrak p \subset A$라 하자. $\xi$가
$D$의 한 기약 성분의 일반점이라는 가정은 $\mathfrak p$가 $(f)$ 위에서
극소임을 뜻한다. 따라서 Algebra, Lemma
\ref{algebra-lemma-minimal-over-1}에 의해
$\dim(A_\mathfrak p) \leq 1$이다.

\medskip\noindent
(2)의 증명. (1)에 의해 $\dim(\mathcal{O}_{X, \xi}) \leq 1$임을 안다.
한편 국소 방정식 $f$는 Lemma
\ref{divisors-lemma-characterize-effective-Cartier-divisor}에 의해
$A_\mathfrak p$에서 영인자가 아닌 원소이다. 이는 차원이 적어도 $1$임을
함의한다(초등적인 Algebra, Lemma
\ref{algebra-lemma-Zariski-topology}에 의해 $A_\mathfrak p$ 안에는 $f$를
포함하지 않는 소아이디얼이 있어야 하기 때문이다).
\end{proof}

\begin{lemma}
\label{divisors-lemma-integral-effective-Cartier-divisor-dvr}
$X$를 뇌터 스킴이라 하자. $D \subset X$를 정역인 닫힌 부분스킴이면서
유효 카르티에 인자인 것이라 하자. 그러면 $X$의 국소환은 $D$의 일반점에서
국소환은 이산 부치환이다.
\end{lemma}

\begin{proof}
Lemma \ref{divisors-lemma-characterize-effective-Cartier-divisor}에 의해
$X = \Spec(A)$이고 $D = \Spec(A/(f))$이며, $A$는 뇌터 환이고
$f \in A$는 영인자가 아닌 원소라고 가정해도 된다. $D$가 정역이라는
가정은 $(f)$가 소아이디얼임을 뜻한다. 따라서 일반점에서 $X$의 국소환은
$A_{(f)}$이고, 이는 극대 아이디얼이 영인자가 아닌 원소 하나에 의해
생성되는 뇌터 국소환이다. 그러므로 Algebra, Lemma
\ref{algebra-lemma-characterize-dvr}에 의해 이산 부치환이다.
\end{proof}

\begin{lemma}
\label{divisors-lemma-effective-Cartier-divisor-Sk}
$X$를 국소 뇌터 스킴이라 하자. $D \subset X$를 유효 카르티에 인자라
하자. $X$가 $(S_k)$이면 $D$는 $(S_{k - 1})$이다.
\end{lemma}

\begin{proof}
$x \in D$라 하자. 그러면
$\mathcal{O}_{D, x} = \mathcal{O}_{X, x}/(f)$이며,
$f \in \mathcal{O}_{X, x}$는 영인자가 아닌 원소이다. 가정에 의해
$\text{depth}(\mathcal{O}_{X, x}) \geq \min(\dim(\mathcal{O}_{X, x}), k)$이다.
Algebra, Lemma \ref{algebra-lemma-depth-drops-by-one}에 의해
$\text{depth}(\mathcal{O}_{D, x}) = \text{depth}(\mathcal{O}_{X, x}) - 1$이고,
Algebra, Lemma \ref{algebra-lemma-one-equation}에 의해
$\dim(\mathcal{O}_{D, x}) = \dim(\mathcal{O}_{X, x}) - 1$이다. 따라서
$\text{depth}(\mathcal{O}_{D, x}) \geq \min(\dim(\mathcal{O}_{D, x}), k - 1)$이며,
이는 원하는 바이다.
\end{proof}

\begin{lemma}
\label{divisors-lemma-normal-effective-Cartier-divisor-S1}
$X$를 국소 뇌터 정규 스킴이라 하자. $D \subset X$를 유효 카르티에
인자라 하자. 그러면 $D$는 $(S_1)$이다.
\end{lemma}

\begin{proof}
Properties, Lemma \ref{properties-lemma-criterion-normal}에 의해 $X$는
$(S_2)$이다. 따라서 Lemma \ref{divisors-lemma-effective-Cartier-divisor-Sk}로
결론을 얻는다.
\end{proof}

\begin{lemma}
\label{divisors-lemma-weil-divisor-is-cartier-UFD}
$X$를 국소 뇌터 스킴이라 하자. $D \subset X$를 정역인 닫힌
부분스킴이라 하자. 다음을 가정하자.
\begin{enumerate}
\item $D$의 여차원은 $1$이고, 이는 $X$ 안에서의 여차원이다.
\item $\mathcal{O}_{X, x}$는 모든 $x \in D$에 대해 UFD이다.
\end{enumerate}
그러면 $D$는 유효 카르티에 인자이다.
\end{lemma}

\begin{proof}
$x \in D$라 하고 $A = \mathcal{O}_{X, x}$로 두자. $\mathfrak p \subset A$를
$D$의 일반점에 대응하는 것이라 하자. 가정 (1)에 의해
$A_\mathfrak p$의 차원은 $1$이다. 따라서 $\mathfrak p$는 높이 $1$인
소아이디얼이다. $A$가 UFD이므로 $\mathfrak p = (f)$이고 어떤
$f \in A$가 존재한다. 물론 $f$는 영인자가 아닌 원소이고 Lemma
\ref{divisors-lemma-effective-Cartier-in-points}에 의해 결론을 얻는다.
\end{proof}

\begin{lemma}
\label{divisors-lemma-codim-1-part}
$X$를 뇌터 스킴이라 하자. $Z \subset X$를 닫힌 부분스킴이라 하자. 다음이
존재한다고 가정하자.
\begin{enumerate}
\item 정역인 유효 카르티에 인자들의 모음 $D_i \subset X$, $i \in I$,
\item 닫힌 부분집합 $Z' \subset X$로서 그 모든 기약 성분은 여차원
$\geq 2$이며, 이는 $X$ 안에서의 여차원이다(Topology, Definition
\ref{topology-definition-codimension}).
\end{enumerate}
또한 집합론적으로 $Z \subset Z' \cup \bigcup_{i \in I} D_i$라고 하자.
그러면 $a_i \geq 0$이고 거의 모든 경우 $a_i = 0$이며, 여기서 $i \in I$인 정수들이
존재하여 유효 카르티에 인자
$$
D = \sum a_i D_i
$$
는 $Z$에 포함되고, 포함 사상 $D \to Z$는 여차원 $2$인 부분을 $X$에서
부분을 제외하면 동형이다(즉 열린집합 $U \subset Z$가 존재하여
$D \cap U \to Z \cap U$가 동형이고 $Z \setminus U$의 모든 기약 성분은
여차원 $\geq 2$이며, 이는 $X$ 안에서의 여차원이다). $Z$가 $X$의
어디에서도 조밀하지 않을 때 $D_i$, $i \in I$와 $Z'$의 존재가 보장된다.
그 조건은 $\mathcal{O}_{X, x}$가 모든 $x \in Z$에 대해 UFD이거나 $X$가
정칙인 것이다.
\end{lemma}

\begin{proof}
$\xi_i \in D_i$를 일반점이라 하고 $\mathcal{O}_i = \mathcal{O}_{X, \xi_i}$를
국소환이라 하자. 이는 Lemma
\ref{divisors-lemma-integral-effective-Cartier-divisor-dvr}에 의해 이산
부치환이다. $a_i \geq 0$을
$\mathcal{I}_{Z, \xi_i} \subset \mathcal{O}_i$의 원소가 갖는 최소
부치로 두자. 물론 $a_i > 0$일 수 있는 것은 $D_i$가 $Z$의 기약 성분일
때뿐이므로, $a_i > 0$인 것은 유한 개의 $i \in I$에 대해서뿐이다. 유효
카르티에 인자 $D = \sum a_i D_i$가 원하는 것이라고 주장한다.

\medskip\noindent
실제로 $x \in X$라고 하자. $A = \mathcal{O}_{X, x}$로 두자.
$D_1, \ldots, D_n$을 서로 다른 인자 $D_i$ 중 $x \in D_i$이고
$a_i > 0$인 것들이라 하자. $1 \leq i \leq n$에 대해 $f_i \in A$를 $D_i$의
국소 방정식이라 하자. 그러면 $f_i$는 $A$의 소원소이고
$\mathcal{O}_i = A_{(f_i)}$이다. $I = \mathcal{I}_{Z, x} \subset A$를
$Z$의 아이디얼 층의 줄기라 하자. $a_i$의 선택에 의해
$I A_{(f_i)} = f_i^{a_i}A_{(f_i)}$이다. 다음을 주장한다.
$I \subset (\prod_{i = 1, \ldots, n} f_i^{a_i})$.

\medskip\noindent
주장의 증명. 국소화 사상
$\varphi : A/(f_i) \to A_{(f_i)}/f_iA_{(f_i)}$는 단사이다. 실제로
소아이디얼 $(f_i)$는 극대 아이디얼 $f_iA_{(f_i)}$의 역상이다. $n$에
대한 귀납법에 의해
$\varphi_n : A/(f_i^n)\to A_{(f_i)}/f_i^nA_{(f_i)}$도 단사임을 얻는다.
$\varphi_{a_i}(I) = 0$이므로 $I \subset (f_i^{a_i})$이다. 따라서 임의의
$x \in I$에 대해 $x = f_1^{a_1}x_1$로 쓸 수 있으며, 여기서 어떤
$x_1 \in A$가 존재한다. $D_1, \ldots, D_n$은 서로 다르므로
$f_i$는 $A_{(f_j)}$에서 단위원이다. 이는 $i \not = j$일 때이다. $x$와
$x_1$을 $A_{(f_i)}$에서, $n \geq i > 1$인 경우에 비교하면 여전히
$x_1 \in (f_i^{a_i})$이다. 앞의 과정을 반복하면 귀납적으로
$x_i = f_{i + 1}^{a_{i + 1}}x_{i + 1}$로 쓸 수 있다. 이는 임의의
$n > i \geq 1$에 대해 성립한다. 결론적으로
$x \in (\prod_{i = 1, \ldots n} f_i^{a_i})$이며, 이는 임의의 $x \in I$에
대해 성립하고,
이는 원하는 바이다.

\medskip\noindent
이 주장은 $\mathcal{I}_Z \subset \mathcal{I}_D$, 즉 $D \subset Z$임을
보인다. 더욱이 $D$와 $Z$가 $\xi_i$에서 일치함도 알 수 있으므로
$D \to Z$는 여차원 $2$인 부분을 $X$ 위에서 제외하면 동형이다.

\medskip\noindent
마지막 명제들을 보이기 위해 다음과 같이 논증한다. 정칙 국소환은
UFD이므로(More on Algebra, Lemma
\ref{more-algebra-lemma-regular-local-UFD}) UFD인 경우를 논증하면 충분하다.
그 경우 $D_i$를 $Z$의 기약 성분들 중 여차원이 $1$이고 이것이 $X$
안에서의 여차원인 것들이라 하자.
Lemma \ref{divisors-lemma-weil-divisor-is-cartier-UFD}에 의해 각 $D_i$는 유효
카르티에 인자이다.
\end{proof}

\begin{lemma}
\label{divisors-lemma-codimension-1-is-effective-Cartier}
$Z \subset X$를 뇌터 스킴의 닫힌 부분스킴이라 하자. 다음을 가정하자.
\begin{enumerate}
\item $Z$는 매장점을 갖지 않는다.
\item $Z$의 모든 기약 성분의 여차원은 $1$이고, 이는 $X$ 안에서의
여차원이다.
\item 각 국소환 $\mathcal{O}_{X, x}$, $x \in Z$는 UFD이거나 $X$가
정칙이다.
\end{enumerate}
그러면 $Z$는 유효 카르티에 인자이다.
\end{lemma}

\begin{proof}
Lemma \ref{divisors-lemma-codim-1-part}에서와 같이 $D = \sum a_i D_i$라 하자.
여기서 $D_i \subset Z$는 $Z$의 기약 성분들이다.
$D \to Z$가 동형이 아니면 $\mathcal{O}_Z \to \mathcal{O}_D$는 여차원
$\geq 2$인 부분에 놓이는 영이 아닌 핵을 갖는다. 이는 $Z$가 매장점을
갖는다는 뜻인데, 가정 (1)에 어긋난다. 따라서 원하는 대로
$D \cong Z$이다.
\end{proof}

\begin{lemma}
\label{divisors-lemma-UFD-one-equation-CM}
$R$을 뇌터 UFD라 하자. $I \subset R$을 아이디얼이라 하되, $R/I$는
매장된 소아이디얼을 갖지 않고 $I$ 위의 모든 극소 소아이디얼은 높이
$1$이라고 하자. 그러면 $I = (f)$이고 어떤 $f \in R$가 존재한다.
\end{lemma}

\begin{proof}
Lemma \ref{divisors-lemma-codimension-1-is-effective-Cartier}에 의해 아이디얼 층
$\tilde I$는 $\Spec(R)$ 위에서 가역이다. More on Algebra, Lemma
\ref{more-algebra-lemma-UFD-Pic-trivial}에 의해 이는 한 원소로 생성된다.
\end{proof}

\begin{lemma}
\label{divisors-lemma-effective-Cartier-divisor-is-a-sum}
$X$를 뇌터 스킴이라 하자. $D \subset X$를 유효 카르티에 인자라 하자.
정역인 유효 카르티에 인자들 $D_i \subset X$가 존재하여 집합론적으로
$D \subset \bigcup D_i$라고 가정하자. 그러면
$D = \sum a_i D_i$이며 어떤 $a_i \geq 0$가 존재한다. $D_i$의 존재는
$\mathcal{O}_{X, x}$가 모든 $x \in D$에 대해 UFD이거나 $X$가 정칙이면
보장된다.
\end{lemma}

\begin{proof}
Lemma \ref{divisors-lemma-codim-1-part}에서처럼 $a_i$를 택하고
$D' = \sum a_i D_i$로 두자. 그러면 $D' \to D$는 유효 카르티에 인자들의
포함으로, 여차원 $2$인 부분을 $X$ 위에서 제외하면 동형이다.
$x \in X$를 택하자. $A = \mathcal{O}_{X, x}$로 두고 $f, f' \in A$를
$D, D'$의 아이디얼을 $A$에서 생성하는 영인자가 아닌 원소들이라 하자.
그러면 $f = gf'$이며 어떤 $g \in A$가 존재한다. 더욱이 모든 소아이디얼
$\mathfrak p$ 중 높이가 $\leq 1$이고 $A$의 소아이디얼인 것에 대해 $g$는
$A_\mathfrak p$의 단위원으로 사상된다. 이는 $g$가 단위원임을 함의한다.
실제로 $(g)$ 위의 극소 소아이디얼들의 높이는 $1$이다(Algebra, Lemma
\ref{algebra-lemma-minimal-over-1}).
\end{proof}

\begin{lemma}
\label{divisors-lemma-quasi-projective-Noetherian-pic-effective-Cartier}
\begin{slogan}
사영 스킴 위에서 모든 선다발은 정칙 유리형 절단을 갖는다.
\end{slogan}
$X$를 풍부한 가역층을 갖는 뇌터 스킴이라 하자. 그러면 모든 가역
$\mathcal{O}_X$-가군은 어떤 유효 카르티에 인자 $D, D'$에 대해, 이들이
$X$ 안에 있을 때,
$$
\mathcal{O}_X(D - D') =
\mathcal{O}_X(D) \otimes_{\mathcal{O}_X} \mathcal{O}_X(D')^{\otimes -1}
$$
과 동형이다. 더욱이 유한 부분집합 $E \subset X$가 주어지면 $D, D'$을
택하여 $E \cap D = \emptyset$이고 $E \cap D' = \emptyset$이 되게 할 수
있다. $X$가 준아핀이면 $D' = \emptyset$으로 택할 수 있다.
\end{lemma}

\begin{proof}
$x_1, \ldots, x_n$을 $X$의 연관점들이라 하자(Lemma
\ref{divisors-lemma-finite-ass}).

\medskip\noindent
$X$가 준아핀이고 $\mathcal{N}$이 임의의 가역 $\mathcal{O}_X$-가군이면,
절단 $t$를 $\mathcal{N}$에서 택하되
$E \cup \{x_1, \ldots, x_n\}$의 어떤 점에서도 사라지지 않게 할 수 있다. Properties, Lemma
\ref{properties-lemma-quasi-affine-invertible-nonvanishing-section}을 보라.
그러면 Lemma \ref{divisors-lemma-regular-section-associated-points}에 의해 $t$는
$\mathcal{N}$의 정칙 절단이다. 따라서
$\mathcal{N} \cong \mathcal{O}_X(D)$이며, 여기서 $D = Z(t)$는 $t$에
대응하는 유효 카르티에 인자이다. Lemma \ref{divisors-lemma-characterize-OD}를
보라. 구성에 의해 $E \cap D = \emptyset$이므로 이 경우에는 끝났다.

\medskip\noindent
일반적인 경우로 돌아가서 $\mathcal{L}$을 $X$ 위의 풍부한 가역층이라
하자. 어떤 $m > 0$과 절단 $s \in \Gamma(X, \mathcal{L}^{\otimes m})$가
존재하여 $X_s$는 아핀이고
$E \cup \{x_1, \ldots, x_n\} \subset X_s$이다(Properties, Lemma
\ref{properties-lemma-ample-finite-set-in-principal-affine}).

\medskip\noindent
$\mathcal{N}$을 임의의 가역 $\mathcal{O}_X$-가군이라 하자. 준아핀인
경우에 의해 절단 $t \in \mathcal{N}(X_s)$를 찾되
$E \cup \{x_1, \ldots, x_n\}$의 어떤 점에서도 사라지지 않게 할 수 있다. Properties, Lemma
\ref{properties-lemma-invert-s-sections}에 의해 어떤 $e \geq 0$에 대해
절단 $s^e|_{X_s} t$는 전역 절단 $\tau$, 즉
$\mathcal{L}^{\otimes me} \otimes \mathcal{N}$의 전역 절단으로
확장된다. 따라서 $\mathcal{L}^{\otimes me} \otimes \mathcal{N}$과
$\mathcal{L}^{\otimes me}$은 모두 가역층이고, 각각
$E \cup \{x_1, \ldots, x_n\}$의 어떤 점에서도 사라지지 않는 전역
절단을 갖는다. 따라서 이들은 Lemma
\ref{divisors-lemma-regular-section-associated-points}에 의해 정칙 절단이다.
그러므로
$\mathcal{L}^{\otimes me} \otimes \mathcal{N} \cong \mathcal{O}_X(D)$이고
$\mathcal{L}^{\otimes me} \cong \mathcal{O}_X(D')$이다. 여기서 어떤 유효
카르티에 인자 $D$와 $D'$이 존재한다. Lemma
\ref{divisors-lemma-characterize-OD}를 보라. 구성에 의해
$E \cap D = \emptyset$이고 $E \cap D' = \emptyset$이며, 증명이 끝난다.
\end{proof}

\begin{lemma}
\label{divisors-lemma-wedge-product-ses}
$X$를 차원 $2$인 정역 정칙 스킴이라 하자. $i : D \to X$를 유효
카르티에 인자의 몰입이라 하자.
$\mathcal{F} \to \mathcal{F}' \to i_*\mathcal{G} \to 0$을 응집
$\mathcal{O}_X$-가군들의 완전열이라 하자. 다음을 가정하자.
\begin{enumerate}
\item $\mathcal{F}, \mathcal{F}'$은 계수 $r$인 국소 자유이고, 이는 $X$의
공집합이 아닌 열린집합 위에서 성립한다.
\item $D$는 정역 스킴이다.
\item $\mathcal{G}$는 유한 국소 자유 $\mathcal{O}_D$-가군이고 그 계수는
$s$이다.
\end{enumerate}
그러면 $\mathcal{L} = (\wedge^r\mathcal{F})^{**}$와
$\mathcal{L}' = (\wedge^r \mathcal{F}')^{**}$은 가역
$\mathcal{O}_X$-가군이고, $\mathcal{L}' \cong \mathcal{L}(k D)$이며 어떤
$k \in \{0, \ldots, \min(s, r)\}$가 존재한다.
\end{lemma}

\begin{proof}
첫 명제는 Lemma \ref{divisors-lemma-reflexive-over-regular-dim-2}에서 따른다.
실제로 가정 (1)은 $\mathcal{L}$과 $\mathcal{L}'$의 계수가 $1$임을
함의한다. $\wedge^r$와 이중 쌍대를 취하는 것은 함자들이므로 표준 사상
$\sigma : \mathcal{L} \to \mathcal{L}'$을 얻는다. 이는 (1)의 공집합이
아닌 열린집합 위에서 동형이므로 영이 아니다. 증명을 마치려면 $\sigma$를
$\mathcal{L}' \otimes \mathcal{L}^{\otimes -1}$의 전역 절단으로 볼 때,
$X$의 어떤 여차원점에서도 사라지지 않음을 보이면 충분하다. 단 $D$의
일반점에서는 예외이고, 거기에서의 영점 차수는 많아야
$\min(s, r)$이어야 한다.

\medskip\noindent
대수로 옮기면 다음 문제에 이른다. $(A, \mathfrak m, \kappa)$를 분수체가
$K$인 이산 부치환이라 하자. $M \to M' \to N \to 0$을 유한
$A$-가군들의 완전열이라 하되,
$\dim_K(M \otimes K) = \dim_K(M' \otimes K) = r$이고
$N \cong \kappa^{\oplus s}$라고 하자. 유도된 사상
$L = \wedge^r(M)^{**} \to L' = \wedge^r(M')^{**}$의 영점 차수가 많아야
$\min(s, r)$임을 보여야 한다. $A$ 위 가군의 구조 정리를 사용하겠다.
More on Algebra, Lemma
\ref{more-algebra-lemma-generalized-valuation-ring-modules} 또는
\ref{more-algebra-lemma-modules-PID}를 보라. 유한 $A$-가군을 꼬임
부분가군으로 나누어도 이중 쌍대는 변하지 않는다. 따라서 $M$을
$M/M_{tors}$로, $M'$을 $M'/\Im(M_{tors} \to M')$로 바꾸어 $M$에 꼬임이
없다고 가정해도 된다. 그러면 $M \to M'$은 단사이고
$M'_{tors} \to N$도 단사이다. 따라서 $M'$을 $M'/M'_{tors}$로, $N$을
$N/M'_{tors}$로 바꾸어도 된다. 이로써 $M$과 $M'$이 계수 $r$인 자유
가군이고 $N \cong \kappa^{\oplus s}$인 경우로 환원된다. 이 경우
$\sigma$는 $M \to M'$의 행렬식이고, 그 영점 차수는 $s$이다. 예를 들어
Algebra, Lemma \ref{algebra-lemma-order-vanishing-determinant}를 보라.
\end{proof}


\section{아핀 열린집합의 여집합}
\label{divisors-section-complement-affine-open}

\noindent
이 절에서는 다형체의 아핀 열린집합의 여집합이 순수 여차원 $1$을 갖는다는
결과를 논의한다.

\begin{lemma}
\label{divisors-lemma-affine-punctured-spec}
$(A, \mathfrak m)$을 뇌터 국소환이라 하자. 뚫린 스펙트럼
$U = \Spec(A) \setminus \{\mathfrak m\}$, 즉 $A$의 뚫린 스펙트럼이
아핀일 필요충분조건은 $\dim(A) \leq 1$인 것이다.
\end{lemma}

\begin{proof}
$\dim(A) = 0$이면 $U$는 공집합이므로 아핀이다($0$인 환의 스펙트럼과
같다). $\dim(A) = 1$이면 원소 $f \in \mathfrak m$을 택하되, 그것이
$A$의 유한 개의 극소 소아이디얼 중 어느 것에도 속하지 않게 할 수 있다(Algebra,
Lemmas \ref{algebra-lemma-Noetherian-irreducible-components}와
\ref{algebra-lemma-silly}). 그러면 $U = \Spec(A_f)$는 아핀이다.

\medskip\noindent
역은 더 흥미롭다. 다소 비표준적인 증명을 제시하고 아래의 remark에서
표준적인 논증을 논의하겠다. $U = \Spec(B)$가 아핀이라고 가정하자.
축약으로 바꾸어도 아핀성과 차원은 변하지 않으므로, $A$를 그 멱영원들의
아이디얼에 의한 몫으로 바꾸어 $A$가 축약이라고 가정해도 된다.
$Q = B/A$로 두고 이를 $A$-가군으로 보자. $Q$의 지지집합은
$\{\mathfrak m\}$이다. 실제로 $A_\mathfrak p = B_\mathfrak p$는 모든
극대가 아닌 소아이디얼 $\mathfrak p$와 환 $A$에 대해 성립한다.
$\dim(A) \geq 1$이라고 가정해도 되므로, 위와 같이 원소
$f \in \mathfrak m$을 택하되 그것이 $A$의 어느 극소 아이디얼에도
속하지 않게 할 수 있다. $A$가 축약이므로 이는 $f$가
$A$에서 영인자가 아닌 원소임을 함의한다. 같은 논증은
$f : B \to B$가 단사임을 보인다. 실제로 $\Spec(B)$는
$\Spec(A)$에서 열려 있다. Algebra, Lemma \ref{algebra-lemma-one-equation}에 의해
$\dim(A/fA) = \dim(A) - 1$임에 유의하자. 짧은 완전열
$f$에 의한 곱셈을 짧은 완전열 $0 \to A \to B \to Q \to 0$에
적용하고 뱀 보조정리를 쓰면
$$
0 \to Q[f] \to A/fA \to B/fB \to Q/fQ \to 0
$$
을 얻는다. 여기서 $Q[f] = \Ker(f : Q \to Q)$이다. 이는 $Q[f]$가 유한
$A$-가군임을 함의한다. $Q[f]$의 지지집합이 $\{\mathfrak m\}$이므로
$l = \text{length}_A(Q[f]) < \infty$임을 안다(Algebra, Lemma
\ref{algebra-lemma-support-point}). $l_n = \text{length}_A(Q[f^n])$으로
두자. 완전열
$$
0 \to Q[f^n] \to Q[f^{n + 1}] \xrightarrow{f^n} Q[f]
$$
은 귀납적으로 $l_n < \infty$이고 $l_n \leq l_{n + 1}$임을 보인다.
완전열
$$
0 \to Q[f] \to Q[f^{n + 1}] \xrightarrow{f} Q[f^n] \to Q/fQ
$$
을 생각하면 $Q[f^n]$이 $Q/fQ$에서 갖는 상의 길이는
$l_n - l_{n + 1} + l \leq l$을 가짐을 안다. $Q = \bigcup Q[f^n]$이므로
$Q/fQ$의 길이는 많아야 $l$, 즉 유계임을 얻는다. 따라서 $Q/fQ$는 유한
$A$-가군이다. 그러므로 $A/fA \to B/fB$는 유한 환 사상이고, 특히
스펙트럼 사이의 닫힌 사상을 유도한다(Algebra, Lemmas
\ref{algebra-lemma-integral-going-up}과
\ref{algebra-lemma-going-up-closed}). 한편 $\Spec(B/fB)$는
$\Spec(A/fA)$의 뚫린 스펙트럼이다. $\Spec(B/fB) = \emptyset$이 아니면
이는 모순이다. 따라서 원하는 대로 $\dim(A/fA) = 0$이다.
\end{proof}

\begin{remark}
\label{divisors-remark-affine-punctured-spectrum-standard-proof}
$(A, \mathfrak m)$이 차원 $\geq 2$인 뇌터 국소 정규 정역이고 $U$가 $A$의
뚫린 스펙트럼이면 $\Gamma(U, \mathcal{O}_U) = A$이다. 하르톡스 정리의
이 대수적 판본은 Algebra, Lemma
\ref{algebra-lemma-normal-domain-intersection-localizations-height-1}에서 본
사실
$A = \bigcap_{\text{height}(\mathfrak p) = 1} A_\mathfrak p$에서 따른다.
따라서 이 경우 $U$는 아핀일 수 없다(아핀이면 $\mathfrak m$이 $U$의
점이어야 하기 때문이다). 이는 Lemma \ref{divisors-lemma-affine-punctured-spec}의
증명의 출발점으로 흔히 사용된다. 일반적인 뇌터 국소환의 경우를 이 경우로
환원하려면 먼저 완비화하고(나가타 국소환을 얻기 위해), 적절한 극소
소아이디얼을 택해 $A$를 $A/\mathfrak q$로 바꾼 뒤
정규화한다. 이 각 단계는 차원을 바꾸지 않으며 모순을 얻는다. 완비화
단계를 생략할 수도 있지만 그러면 일반적으로 정규화는 뇌터 정역이 아니다.
그럼에도 이는 같은 차원의 크룰 정역이다(이는 크룰--아키즈키를 사용하여
증명된다). 따라서 같은 논증을 적용할 수 있다.
\end{remark}

\begin{remark}
\label{divisors-remark-affine-puctured-spectrum-general}
뚫린 스펙트럼이 아핀인 비뇌터 국소환 $(A, \mathfrak m)$을 어떻게
특징지어야 할지는 분명하지 않다. 이러한 환은
$I$라는 유한 생성 아이디얼을 가지며 $\mathfrak m = \sqrt{I}$이다. 물론 $I$를
$1$개의 원소로 생성되도록 택할 수 있으면 $A$는 아핀인 뚫린 스펙트럼을
갖는다. 이는 비뇌터 예를 많이 준다. 역으로 Lemma
\ref{divisors-lemma-affine-punctured-spec}의 증명에 있는 논증에 의해, 이러한 환은
원소 $f \in \mathfrak m$ 중 영인자가 아니면서
$H^0_I(A/fA) = 0$을 만족하는 것을 가질 수 없다(따라서 $A$는 길이 $2$인 정칙열을
가질 수 없다). 더욱이 같은 것은 임의의 환 $A'$에 대해서도 성립한다.
여기서 $A \to A'$는 국소환들의 국소 준동형이고
$\mathfrak m_{A'} = \sqrt{\mathfrak mA'}$라고 가정한다.
\end{remark}

\begin{lemma}
\label{divisors-lemma-complement-affine-open-immersion}
\begin{reference}
\cite[EGA IV, Corollaire 21.12.7]{EGA4}
\end{reference}
$X$를 국소 뇌터 스킴이라 하자. $U \subset X$를 포함 사상
$U \to X$가 아핀인 열린 부분스킴이라 하자. 모든 일반점 $\xi$, 즉
$X \setminus U$의 한 기약 성분의 일반점에 대해 국소환
$\mathcal{O}_{X, \xi}$의 차원은
$\leq 1$이다. $U$가 조밀하거나 $\xi$가 $U$의 폐포에 속하면
$\dim(\mathcal{O}_{X, \xi}) = 1$이다.
\end{lemma}

\begin{proof}
$\xi$가 $X \setminus U$의 일반점이므로
$$
U_\xi = U \times_X \Spec(\mathcal{O}_{X, \xi}) \subset
\Spec(\mathcal{O}_{X, \xi})
$$
는 $\mathcal{O}_{X, \xi}$의 뚫린 스펙트럼이다(힌트: Schemes, Lemma
\ref{schemes-lemma-specialize-points}를 사용하라). $U \to X$가 아핀이므로
$U_\xi \to \Spec(\mathcal{O}_{X, \xi})$도 아핀이고(Morphisms, Lemma
\ref{morphisms-lemma-base-change-affine}), 따라서 $U_\xi$는 아핀이다.
그러므로 Lemma \ref{divisors-lemma-affine-punctured-spec}에 의해
$\dim(\mathcal{O}_{X, \xi}) \leq 1$이다. $\xi \in \overline{U}$이면
특수화 $\eta \to \xi$가 존재하고 여기서 $\eta \in U$이다($\eta$를
$\overline{U}$의 한 기약 성분 중 $\xi$를 포함하는 성분의 일반점으로
택하면 된다. $\overline{U}$가 국소 뇌터이므로 국소적으로 유한 개의
기약 성분을 가지며, 따라서 $\eta \in U$임을 안다). 그러면
$\eta \in \Spec(\mathcal{O}_{X, \xi})$이고 차원은 $0$일 수 없다.
\end{proof}

\begin{lemma}
\label{divisors-lemma-complement-affine-open}
$X$를 분리된 국소 뇌터 스킴이라 하자. $U \subset X$를 아핀
열린집합이라 하자. 모든 일반점 $\xi$, 즉 $X \setminus U$의 한 기약
성분의 일반점에 대해 국소환 $\mathcal{O}_{X, \xi}$의 차원은
$\leq 1$이다.
$U$가 조밀하거나 $\xi$가 $U$의 폐포에 속하면
$\dim(\mathcal{O}_{X, \xi}) = 1$이다.
\end{lemma}

\begin{proof}
이는 Lemma \ref{divisors-lemma-complement-affine-open-immersion}에서 따른다.
실제로 사상 $U \to X$는 Morphisms, Lemma
\ref{morphisms-lemma-affine-permanence}에 의해 아핀이다.
\end{proof}

\noindent
다음 보조정리는 때때로 유효 카르티에 인자를 만드는 데 사용할 수 있다.

\begin{lemma}
\label{divisors-lemma-complement-open-affine-effective-cartier-divisor}
$X$를 뇌터 분리 스킴이라 하자. $U \subset X$를 조밀한 아핀 열린집합이라
하자. $\mathcal{O}_{X, x}$가 모든 $x \in X \setminus U$에 대해 UFD이면,
유효 카르티에 인자 $D \subset X$가 존재하여 $U = X \setminus D$를
만족한다.
\end{lemma}

\begin{proof}
$X$가 뇌터이므로 여집합 $X \setminus U$는 유한 개의 기약 성분
$D_1, \ldots, D_r$을 갖는다(Properties, Lemma
\ref{properties-lemma-Noetherian-irreducible-components}를
$X \setminus U$의 축약 유도 부분스킴 구조에 적용하라). 각
$D_i \subset X$는 Lemma \ref{divisors-lemma-complement-affine-open}에 의해 여차원
$1$이다(Properties, Lemma
\ref{properties-lemma-codimension-local-ring}도 보라). 따라서 Lemma
\ref{divisors-lemma-weil-divisor-is-cartier-UFD}에 의해 $D_i$는 유효 카르티에
인자이다. 그러므로 $D = D_1 + \ldots + D_r$로 택할 수 있다.
\end{proof}

\begin{lemma}
\label{divisors-lemma-complement-open-affine-effective-cartier-divisor-bis}
$X$를 대각선이 아핀인 뇌터 스킴이라 하자. $U \subset X$를 조밀한 아핀
열린집합이라 하자. $\mathcal{O}_{X, x}$가 모든 $x \in X \setminus U$에
대해 UFD이면, 유효 카르티에 인자 $D \subset X$가 존재하여
$U = X \setminus D$를 만족한다.
\end{lemma}

\begin{proof}
$X$가 뇌터이므로 여집합 $X \setminus U$는 유한 개의 기약 성분
$D_1, \ldots, D_r$을 갖는다(Properties, Lemma
\ref{properties-lemma-Noetherian-irreducible-components}를
$X \setminus U$의 축약 유도 부분스킴 구조에 적용하라). $D_i$를 $X$의
축약 닫힌 부분스킴으로 보자. $X = \bigcup_{j \in J} X_j$를 $X$의 아핀
열린 덮개라 하자. $j$가 $J$를 달릴 때마다 $U_j = U \cap X_j$로 두자.
$X$의 대각선이 아핀이므로 스킴
$$
U_j = X \times_{(X \times X)} (U \times X_j)
$$
는 아핀이다. 따라서 $X_j$가 분리되어 있으므로 Lemma
\ref{divisors-lemma-complement-open-affine-effective-cartier-divisor}와 그 증명에
의해 모든 $j \in J$ 및 $1 \leq i \leq r$에 대해 교집합
$D_i \cap X_j$는 공집합이거나 $X_j$ 안의 유효 카르티에 인자이다.
따라서 $D_i \subset X$는 유효 카르티에 인자이다(이는 국소적 성질이기
때문이다). 그러므로 $D = D_1 + \ldots + D_r$로 택할 수 있다.
\end{proof}

\begin{lemma}
\label{divisors-lemma-regular-ample-family}
$X$를 준콤팩트 정칙 스킴이라 하되 대각선은 아핀이라고 하자. 그러면
$X$는 풍부한 가역 가군족을 갖는다(Morphisms, Definition
\ref{morphisms-definition-family-ample-invertible-modules}).
\end{lemma}

\begin{proof}
$X$가 정역 스킴들의 유한 서로소 합집합임에 유의하자(Properties, Lemmas
\ref{properties-lemma-regular-normal}과
\ref{properties-lemma-normal-Noetherian}). 따라서 $X$가 정역이면서 뇌터,
정칙이고 대각선이 아핀이라고 가정해도 된다. $x \in X$라 하자.
아핀 열린 근방 $U \subset X$를 택하되 $x$를 포함하게 하자. $X$가
정역이므로 $U$는 $X$에서
조밀하다. More on Algebra, Lemma
\ref{more-algebra-lemma-regular-local-UFD}에 의해 $X$의 국소환들은 UFD이다.
따라서 Lemma
\ref{divisors-lemma-complement-open-affine-effective-cartier-divisor-bis}에 의해
유효 카르티에 인자 $D \subset X$를 찾을 수 있고 그 여집합은 $U$이다.
그러면 $s = 1_D$는 $\mathcal{L} = \mathcal{O}_X(D)$의 표준 절단이고(Definition
\ref{divisors-definition-invertible-sheaf-effective-Cartier-divisor}를 보라) 정확히
$D$를 따라 사라지므로 $U = X_s$이다. 따라서 Morphisms, Definition
\ref{morphisms-definition-family-ample-invertible-modules}의 두 조건이 모두
성립하며 끝난다.
\end{proof}


\section{노름}
\label{divisors-section-norms}

\noindent
$\pi : X \to Y$를 스킴의 유한 사상이라 하고 $d \geq 1$을 정수라 하자.
다음 조건을 만족하는 곱셈적 사상이 존재할 때,
{\it 차수 $d$이고 $\pi$에 대한 노름}
\footnote{이는 표준적이지 않은 표기이다.}이 존재한다고 하자.
$$
\text{Norm}_\pi : \pi_*\mathcal{O}_X \to \mathcal{O}_Y
$$
그 조건은 다음과 같다.
\begin{enumerate}
\item 합성
$\mathcal{O}_Y \xrightarrow{\pi^\sharp} \pi_*\mathcal{O}_X
\xrightarrow{\text{Norm}_\pi} \mathcal{O}_Y$는 $g \mapsto g^d$와 같고,
\item 열린집합 $V \subset Y$에 대해 $f \in \mathcal{O}_X(\pi^{-1}V)$가
$x \in \pi^{-1}(V)$에서 영이면, $\text{Norm}_\pi(f)$는
$\pi(x)$에서 영이다.
\end{enumerate}
조건 (1)은 $\pi$가 전사이도록 강제함에 유의하자.
$\text{Norm}_\pi$는 곱셈적이므로 단원을 단원으로 보낸다. 따라서
$y \in Y$가 주어졌을 때, $f$가 $X$ 위의 정칙 함수이고 정의되어 있으며,
어떤 $x \in X$에서도 사라지지 않되 그 점이 $\pi(x) = y$를 만족하면
$\text{Norm}_\pi(f)$는 정의되어 있고 $y$에서 사라지지 않는다. 이는
(2)를 요구하지 않고도 성립한다. 실제로 이 절의 구성은 조건 (1)만을
요구하며, 특정한 소멸 성질들만이 성질 (2)를 요구한다(특히 Lemma
\ref{divisors-lemma-norm-ample}의 증명에서 사용된다).

\begin{lemma}
\label{divisors-lemma-finite-trivialize-invertible-upstairs}
$\pi : X \to Y$를 스킴의 유한 사상이라 하자. $\mathcal{L}$을 가역
$\mathcal{O}_X$-가군이라 하고 $y \in Y$라 하자. 열린 근방
$V \subset Y$가 존재하여 $y$를 포함하고
$\mathcal{L}|_{\pi^{-1}(V)}$는 자명하다.
\end{lemma}

\begin{proof}
$Y$와 따라서 $X$가 아핀이라고 가정해도 됨은 분명하다. $\pi$가
유한이므로 올 $\pi^{-1}(\{y\})$는 $y$ 위에서 유한하다. $X$가 아핀이므로
$s \in \Gamma(X, \mathcal{L})$를 택하되 $\pi^{-1}(\{y\})$의 어느
점에서도 사라지지 않게 할 수 있다. 이는 Properties, Lemma
\ref{properties-lemma-quasi-affine-invertible-nonvanishing-section}에서
따르지만 직접적인 논증도 제시하겠다. 즉, 닫힌점들의 유한 집합
$E \subset X$를 택하여 모든 $x \in \pi^{-1}(\{y\})$가 $E$의 어떤 점으로
특수화되게 할 수 있다. $x \in E$에 대해 닫힌 몰입을
$i_x : x \to X$로 나타내자. 그러면
$\mathcal{L} \to \bigoplus_{x \in E} i_{x, *}i_x^*\mathcal{L}$은
준연접 $\mathcal{O}_X$-가군의 전사 사상이고, 따라서 사상
$$
\Gamma(X, \mathcal{L}) \to
\bigoplus\nolimits_{x \in E} \mathcal{L}_x/\mathfrak m_x\mathcal{L}_x
$$
도 전사이다(준연접 $\mathcal{O}_X$-가군의 범주에서 대역 절단을 취하는
함자는 완전하기 때문이다. Schemes, Lemma
\ref{schemes-lemma-equivalence-quasi-coherent}를 보라). 따라서
$s \in \Gamma(X, \mathcal{L})$를 찾아 $E$의 한 점으로 특수화되는 어느
점에서도 사라지지 않게 할 수 있다. 그러면 $X_s \subset X$는
$\pi^{-1}(\{y\})$의 열린 근방이다. $\pi$는 유한이고 따라서 닫힌
사상이므로, 열린 근방 $V \subset Y$가 존재하여 $y$를 포함하고 그 역상이
원하는 대로 $X_s$에 포함됨을 얻는다.
\end{proof}

\begin{lemma}
\label{divisors-lemma-norm-invertible}
$\pi : X \to Y$를 스킴의 유한 사상이라 하자. 차수 $d$이고 $\pi$에 대한
노름이 존재하면 아벨 군의 준동형
$$
\text{Norm}_\pi : \Pic(X) \to \Pic(Y)
$$
이 존재하며 $\text{Norm}_\pi(\pi^*\mathcal{N}) \cong \mathcal{N}^{\otimes d}$가
성립한다. 이는 모든 가역 $\mathcal{O}_Y$-가군 $\mathcal{N}$에 대한 말이다.
\end{lemma}

\begin{proof}
Cohomology, Lemma \ref{cohomology-lemma-h1-invertible}에서 주어진 가역
$\mathcal{O}_X$-가군의 동형류와 $H^1(X, \mathcal{O}_X^*)$의 원소 사이의
대응을 더 언급하지 않고 사용하겠다. 가역 $\mathcal{O}_X$-가군
$\mathcal{L}$의 노름을 취하는 법을 설명하자. Lemma
\ref{divisors-lemma-finite-trivialize-invertible-upstairs}에 의해 열린 덮개
$Y = \bigcup V_j$가 존재하여 $\mathcal{L}|_{\pi^{-1}V_j}$는 자명하다.
생성 절단 $s_j \in \mathcal{L}(\pi^{-1}V_j)$를 각 $j$에 대해 택하자.
겹침 $\pi^{-1}V_j \cap \pi^{-1}V_{j'}$ 위에서
$$
s_j = u_{jj'} s_{j'}
$$
로 쓸 수 있으며, 여기서
$u_{jj'} \in \mathcal{O}^*_X(\pi^{-1}V_j \cap \pi^{-1}V_{j'})$는 유일하다.
따라서 원소
$$
v_{jj'} = \text{Norm}_\pi(u_{jj'}) \in \mathcal{O}_Y^*(V_j \cap V_{j'})
$$
를 생각할 수 있다. 이 원소들은 코사이클 조건을 만족하고($u_{jj'}$가
그렇고 $\text{Norm}_\pi$가 곱셈적이기 때문이다), 따라서 가역
$\mathcal{O}_Y$-가군을 정의한다. 이것이 잘 정의되고, 피카르 군 위에서
가법적이며, $\text{Norm}_\pi(\pi^*\mathcal{N}) \cong \mathcal{N}^{\otimes d}$를
만족한다는 확인은 생략한다. 마지막 성질은 모든 가역 $\mathcal{O}_Y$-가군
$\mathcal{N}$에 대한 것이다.
\end{proof}

\begin{lemma}
\label{divisors-lemma-norm-map-invertible}
$\pi : X \to Y$를 스킴의 유한 사상이라 하자. 차수 $d$이고 $\pi$에 대한
노름이 존재한다고 가정하자. 임의의 $\mathcal{O}_X$-선형 사상
$\varphi : \mathcal{L} \to \mathcal{L}'$가 가역 $\mathcal{O}_X$-가군
사이에 주어졌다고 하자. 이에 대해
$\mathcal{O}_Y$-선형 사상
$$
\text{Norm}_\pi(\varphi) :
\text{Norm}_\pi(\mathcal{L})
\longrightarrow
\text{Norm}_\pi(\mathcal{L}')
$$
이 존재한다. 여기서 $\text{Norm}_\pi(\mathcal{L})$과
$\text{Norm}_\pi(\mathcal{L}')$은 Lemma \ref{divisors-lemma-norm-invertible}에서와
같다. 더욱이 $y \in Y$에 대해 다음은 동치이다.
\begin{enumerate}
\item $\varphi$는 어떤 점 $x \in X$에서 영이고 그 점은 $\pi(x) = y$를 만족하며,
\item $\text{Norm}_\pi(\varphi)$는 $y$에서 영이다.
\end{enumerate}
\end{lemma}

\begin{proof}
열린 덮개 $Y = \bigcup V_j$를 택하여 $\mathcal{L}$과 $\mathcal{L}'$이
열린집합 $\pi^{-1}V_j$ 위에서 자명하게 하자. 이는 Lemma
\ref{divisors-lemma-finite-trivialize-invertible-upstairs}에 의해 가능하다.
생성 절단 $s_j$와 $s'_j$를 각각 $\mathcal{L}$과 $\mathcal{L}'$에 대해
$\pi^{-1}V_j$ 위에서 택하자. 그러면 $\varphi(s_j) = f_js'_j$이고, 여기서
$f_j \in \mathcal{O}_X(\pi^{-1}V_j)$이다.
$\text{Norm}_\pi(\varphi)$를 $\text{Norm}_\pi(f_j)$에 의한 곱셈으로
$V_j$ 위에서 정의하자.
Lemma \ref{divisors-lemma-norm-invertible}의 증명에서
$\text{Norm}_\pi(\mathcal{L})$, $\text{Norm}_\pi(\mathcal{L}')$을 구성하는
데 사용한 코사이클에 대한 간단한 계산으로 이것이 보조정리에 진술된
사상을 정의함을 알 수 있다. 마지막 명제는 차수 $d$인 노름 사상의
정의에 있는 조건 (2)에서 따른다. 몇몇 세부사항은 생략한다.
\end{proof}

\begin{lemma}
\label{divisors-lemma-norm-ample}
$\pi : X \to Y$를 스킴의 유한 사상이라 하자. $X$가 풍부한 가역층을
가지며 차수 $d$이고 $\pi$에 대한 노름이 존재한다고 가정하자. 그러면
$Y$는 풍부한 가역층을 갖는다.
\end{lemma}

\begin{proof}
가정으로 주어진 풍부한 가역층을 $\mathcal{L}$이라 하며 이는 $X$ 위에 있다.
$\mathcal{N} = \text{Norm}_\pi(\mathcal{L})$이 $Y$ 위에서 풍부함을
보이겠다.

\medskip\noindent
$X$는 준콤팩트이고(Properties, Definition
\ref{properties-definition-ample}) $X \to Y$는 전사이므로
($\text{Norm}_\pi$의 존재에 의해) $Y$는 준콤팩트이다. 점 $y \in Y$를
택하자. 증명을 끝내기 위해 절단 $t$가 $\mathcal{N}$의 어떤 양의 텐서
거듭제곱에 존재하여 $y$에서 사라지지 않고 $Y_t$가 아핀임을 보이겠다.
이를 위해 아핀 열린 근방 $V \subset Y$를 택하여 $y$를 포함하게 하자.
$n \gg 0$과 절단
$s \in \Gamma(X, \mathcal{L}^{\otimes n})$을 택하여
$$
\pi^{-1}(\{y\}) \subset X_s \subset \pi^{-1}V
$$
가 되게 할 수 있다. 이는 Properties, Lemma
\ref{properties-lemma-ample-finite-set-in-principal-affine}에서 따른다.
그러면 $t = \text{Norm}_\pi(s)$는 $\mathcal{N}^{\otimes n}$의 절단이며
$x$에서 사라지지 않고 $Y_t \subset V$를 만족한다. Lemma
\ref{divisors-lemma-norm-map-invertible}를 보라. 그러면 $Y_t$는 Properties, Lemma
\ref{properties-lemma-affine-cap-s-open}에 의해 아핀이다.
\end{proof}

\begin{lemma}
\label{divisors-lemma-norm-quasi-affine}
$\pi : X \to Y$를 스킴의 유한 사상이라 하자. $X$가 준아핀이고
차수 $d$이고 $\pi$에 대한 노름이 존재한다고 가정하자. 그러면 $Y$는
준아핀이다.
\end{lemma}

\begin{proof}
Properties, Lemma \ref{properties-lemma-quasi-affine-O-ample}에 의해
$\mathcal{O}_X$는 $X$ 위의 풍부한 가역층이다. Lemma
\ref{divisors-lemma-norm-ample}의 증명은
$\text{Norm}_\pi(\mathcal{O}_X) = \mathcal{O}_Y$가 풍부한 가역
$\mathcal{O}_Y$-가군임을 보인다. 따라서 Properties, Lemma
\ref{properties-lemma-quasi-affine-O-ample}에 의해 $Y$는 준아핀이다.
\end{proof}

\begin{lemma}
\label{divisors-lemma-finite-locally-free-has-norm}
$\pi : X \to Y$를 차수 $d \geq 1$인 유한 국소 자유 사상이라 하자.
그러면 임의의 밑변환과 가환하는 차수 $d$의 표준 노름이 존재한다.
\end{lemma}

\begin{proof}
$V \subset Y$를 아핀 열린집합이라 하되 $(\pi_*\mathcal{O}_X)|_V$가 계수
$d$인 유한 자유 가군이게 하자. 기저를 택하면 동형
$$
\mathcal{O}_V^{\oplus d} \cong (\pi_*\mathcal{O}_X)|_V
$$
을 얻는다. 모든
$f \in \pi_*\mathcal{O}_X(V) = \mathcal{O}_X(\pi^{-1}(V))$에 대해
$f$에 의한 곱셈은 표시된 자유 벡터 다발의 $\mathcal{O}_V$-선형
자기준동형 $m_f$를 정의한다. 따라서 $d \times d$ 행렬
$M_f \in \text{Mat}(d \times d, \mathcal{O}_Y(V))$를 얻고 다음과 같이
둘 수 있다.
$$
\text{Norm}_\pi(f) = \det(M_f)
$$
행렬의 행렬식은 택한 기저와 무관하므로 이는 잘 정의된다. 이는 또한 이
구성이 서로 붙어 원하는 대역 사상을 이룸을 뜻한다. 밑변환과의 가환성은
구성에서 곧바로 따른다.

\medskip\noindent
성질 (1)은 $d \times d$ 대각행렬의 대각 성분이 $g, g, \ldots, g$일 때
행렬식이 $g^d$라는 사실에서 따른다. 성질 (2)를 보이기 위해 밑변환하여
$Y$가 체 $k$의 스펙트럼이라고 가정해도 된다. 그러면
$X = \Spec(A)$이고 $A$는 $k$-대수이며 $\dim_k(A) = d$이다.
$x \in X$가 존재하여 $f \in A$가 $x$에서 사라지면, 어떤 체로 가는 사상
$A \to \kappa$가 존재하여 $f$를 $\kappa$의 영으로 보낸다. 따라서
$f : A \to A$는 전사일 수 없고, 그러므로 원하는 대로
$\det(f : A \to A) = 0$이다.
\end{proof}

\begin{lemma}
\label{divisors-lemma-norm-in-normal-case}
$\pi : X \to Y$를 유한 전사 사상이라 하되 $X$와 $Y$는 정역이고
$Y$는 정규라고 하자. 그러면 차수 $[R(X) : R(Y)]$이고 $\pi$에 대한 노름이
존재한다.
\end{lemma}

\begin{proof}
$\Spec(B) \subset Y$를 아핀 열린 부분집합이라 하고 그 역상을
$\Spec(A) \subset X$라 하자. 그러면 $A$와 $B$는 정역이다. $K$를 $A$의
분수체라 하고 $L$을 $B$의 분수체라 하자. 그림은 다음과 같다.
$$
\xymatrix{
L \ar[r] & K \\
B \ar[u] \ar[r] & A \ar[u]
}
$$
$K/L$은 유한 확대이므로 노름 사상
$\text{Norm}_{K/L} : K^* \to L^*$가 존재하며 그 차수는 $d = [K : L]$이다. 이는 Lemma
\ref{divisors-lemma-finite-locally-free-has-norm}의 증명에서처럼 $f \in K$를
$\det_L(f : K \to K)$로 보냄으로써 주어진다. $f : K \to K$의 특성
다항식은 $f$의 $L$ 위 극소 다항식의 거듭제곱임에 유의하자. 특히
$\text{Norm}_{K/L}(f)$는 $f$의 $L$ 위 극소 다항식의 상수항의
거듭제곱이다. 따라서 Algebra, Lemma
\ref{algebra-lemma-minimal-polynomial-normal-domain}에 의해
$\text{Norm}_{K/L}$는 $A$를 $B$ 안으로 보낸다. 이는 아핀 열린집합 위의
절단들에 대한 가환하는 사상계를 결정하며, 따라서 대역 노름 사상
$\text{Norm}_\pi$를 결정하고 그 차수는 $d$이다.

\medskip\noindent
성질 (1)은 구성에서 바로 따른다. 성질 (2)를 보기 위해 $f \in A$가
소아이디얼 $\mathfrak p \subset A$에 속한다고 하자.
$f^m + b_1 f^{m - 1} + \ldots + b_m$을 $f$의 $L$ 위 극소 다항식이라
하자. Algebra, Lemma \ref{algebra-lemma-minimal-polynomial-normal-domain}에
의해 $b_i \in B$이다. 따라서 $b_0 \in B \cap \mathfrak p$이다.
$\text{Norm}_{K/L}(f) = b_0^{d/m}$이므로(위 참조), 노름은
$\mathfrak p$의 상인 점에서 사라진다고 결론짓는다.
\end{proof}

\begin{lemma}
\label{divisors-lemma-Frobenius-gives-norm-for-reduction}
$X$를 뇌터 스킴이라 하자. $p$를 $p\mathcal{O}_X = 0$을 만족하는
소수라 하자. 그러면 어떤 $e > 0$에 대해 차수 $p^e$이고
$X_{red} \to X$에 대한 노름이 존재하며, 여기서 $X_{red}$는 $X$의 축약이다.
\end{lemma}

\begin{proof}
$A$를 $pA = 0$인 뇌터 환이라 하자. $I \subset A$를 멱영원들의
아이디얼이라 하자. 그러면 $I^n = 0$인 어떤 $n$이 존재한다(Algebra,
Lemma \ref{algebra-lemma-Noetherian-power}). $e$를 택하여 $p^e \geq n$이
되게 하자. 그러면
$$
A/I \longrightarrow A,\quad
f \bmod I \longmapsto f^{p^e}
$$
는 잘 정의된다. 이는 차수 $p^e$이고 $\Spec(A/I) \to \Spec(A)$에 대한
노름을 준다. 이제 $X$가 아핀 스킴들 $\Spec(A_i)$을 붙여서 얻어진다면,
어떤 $e \gg 0$에 대해 이 사상들이 서로 붙어 $X_{red} \to X$에 대한
노름 사상을 이룬다. 세부사항은 생략한다.
\end{proof}

\begin{proposition}
\label{divisors-proposition-push-down-ample}
$\pi : X \to Y$를 스킴의 유한 전사 사상이라 하자. $X$가 풍부한 가역
$\mathcal{O}_X$-가군을 갖는다고 가정하자. 다음 중 하나가 성립한다고 하자.
\begin{enumerate}
\item $\pi$가 유한 국소 자유이거나,
\item $Y$가 정역 정규 스킴이거나,
\item $Y$가 뇌터이고 $p\mathcal{O}_Y = 0$이며 $X = Y_{red}$이다.
\end{enumerate}
그러면 $Y$는 풍부한 가역 $\mathcal{O}_Y$-가군을 갖는다.
\end{proposition}

\begin{proof}
경우 (1)은 Lemmas \ref{divisors-lemma-finite-locally-free-has-norm}와
\ref{divisors-lemma-norm-ample}의 결합에서 따른다. 경우 (3)은 Lemmas
\ref{divisors-lemma-Frobenius-gives-norm-for-reduction}와
\ref{divisors-lemma-norm-ample}의 결합에서 따른다. 경우 (2)에서는 먼저 $X$를
$X$의 한 기약 성분 중 $Y$를 지배하는 것으로 바꾼다(이를 $X$의 축약
닫힌 부분스킴으로 본다). 그런 다음 Lemma
\ref{divisors-lemma-norm-in-normal-case}를 적용할 수 있다.
\end{proof}

\begin{lemma}
\label{divisors-lemma-push-down-quasi-affine}
$\pi : X \to Y$를 스킴의 유한 전사 사상이라 하자. $X$가 준아핀이라고
가정하자. 다음 중 하나가 성립한다고 하자.
\begin{enumerate}
\item $\pi$가 유한 국소 자유이거나,
\item $Y$가 정역 정규 스킴이다.
\end{enumerate}
그러면 $Y$는 준아핀이다.
\end{lemma}

\begin{proof}
경우 (1)은 Lemmas \ref{divisors-lemma-finite-locally-free-has-norm}와
\ref{divisors-lemma-norm-quasi-affine}의 결합에서 따른다. 경우 (2)에서는 먼저
$X$를 $X$의 한 기약 성분 중 $Y$를 지배하는 것으로 바꾼다(이를 $X$의
축약 닫힌 부분스킴으로 본다). 그런 다음 Lemma
\ref{divisors-lemma-norm-in-normal-case}를 적용할 수 있다.
\end{proof}


\section{상대 유효 카르티에 인자}
\label{divisors-section-effective-Cartier-morphisms}

\noindent
다음 보조정리는 밑 위에서 평탄한 유효 카르티에 인자가 실제로 밑 위의
``유효 카르티에 인자들의 족''임을 보인다. 예를 들어 임의의 올에 제한하면
유효 카르티에 인자가 된다.

\begin{lemma}
\label{divisors-lemma-relative-Cartier}
$f : X \to S$를 스킴의 사상이라 하고 $D \subset X$를 닫힌 부분스킴이라
하자. 다음을 가정하자.
\begin{enumerate}
\item $D$는 유효 카르티에 인자이고,
\item $D \to S$는 평탄 사상이다.
\end{enumerate}
그러면 스킴의 모든 사상 $g : S' \to S$에 대해 당김
$(g')^{-1}D$는 $X' = S' \times_S X$ 위의 유효 카르티에 인자이다. 여기서
$g' : X' \to X$는 사영이다.
\end{lemma}

\begin{proof}
Lemma \ref{divisors-lemma-characterize-effective-Cartier-divisor}를 사용하여 이를
다음 대수 명제로 옮긴다. $A \to B$를 환 사상이라 하고 $h \in B$라 하자.
$h$가 영인자가 아닌 원소이고 $B/hB$가 $A$ 위에서 평탄하다고 가정하자.
그러면
$$
0 \to B \xrightarrow{h} B \to B/hB \to 0
$$
은 $A$-가군의 짧은 완전열이고 $B/hB$는 $A$ 위에서 평탄하다. Algebra,
Lemma \ref{algebra-lemma-flat-tor-zero}에 의해 이 열은 임의의 가군과
$A$ 위에서 텐서한 뒤에도 완전하며, 특히 임의의 $A$-대수 $A'$과 텐서한
뒤에도 그러하다.
\end{proof}

\noindent
이 보조정리는 다음 정의의 동기이다.

\begin{definition}
\label{divisors-definition-relative-effective-Cartier-divisor}
$f : X \to S$를 스킴의 사상이라 하자. $X/S$ 위의
{\it 상대 유효 카르티에 인자}란 유효 카르티에 인자 $D \subset X$ 중
$D \to S$가 스킴의 평탄 사상인 것을 말한다.
\end{definition}

\noindent
이 표기는 표준적이지 않을 수 있음을 독자에게 경고한다. 특히
\cite[IV, Section 21.15]{EGA}에서는 상대 인자의 개념을 $X \to S$가
평탄하고 국소 유한 표시인 경우에만 논의한다. 우리의 정의는 조금 더
일반적이다. 그러나 많은 경우(항상은 아니지만) $x \in D$이면
$X \to S$는 $x$에서 평탄임이 드러난다.

\begin{lemma}
\label{divisors-lemma-sum-relative-effective-Cartier-divisor}
$f : X \to S$를 스킴의 사상이라 하자. $D_1, D_2 \subset X$가 $X/S$ 위의
상대 유효 카르티에 인자이면 $D_1 + D_2$도 그러하다(Definition
\ref{divisors-definition-sum-effective-Cartier-divisors}).
\end{lemma}

\begin{proof}
이는 다음 대수 명제로 옮겨진다. $A \to B$를 환 사상이라 하고
$h_1, h_2 \in B$라 하자. $h_i$들이 영인자가 아닌 원소이고
$B/h_iB$가 $A$ 위에서 평탄하다고 가정하자. 그러면 $h_1h_2$는 영인자가
아닌 원소이고 $B/h_1h_2B$는 $A$ 위에서 평탄하다. 그 이유는 짧은 완전열
$$
0 \to B/h_1B \to B/h_1h_2B \to B/h_2B \to 0
$$
이 있고 첫째 화살표가 $h_2$에 의한 곱셈으로 주어지기 때문이다. 양 끝의
두 가군이 $A$ 위에서 평탄하므로 가운데 것도 그러하다. Algebra, Lemma
\ref{algebra-lemma-flat-ses}를 보라.
\end{proof}

\begin{lemma}
\label{divisors-lemma-difference-relative-effective-Cartier-divisor}
$f : X \to S$를 스킴의 사상이라 하자. $D_1, D_2 \subset X$가 $X/S$ 위의
상대 유효 카르티에 인자이고 닫힌 부분스킴으로서 $D_1 \subset D_2$이면,
$D$라는 유효 카르티에 인자 중 $D_2 = D_1 + D$를 만족하는 것은(Lemma
\ref{divisors-lemma-difference-effective-Cartier-divisors}) $X/S$ 위의 상대 유효
카르티에 인자이다.
\end{lemma}

\begin{proof}
이는 다음 대수 명제로 옮겨진다. $A \to B$를 환 사상이라 하고
$h_1, h_2 \in B$라 하자. $h_i$들이 영인자가 아닌 원소이고
$B/h_iB$가 $A$ 위에서 평탄하며 $(h_2) \subset (h_1)$이라고 가정하자.
그러면 $h_2 = h h_1$로 쓸 수 있고 여기서 $h \in B$는 영인자가 아닌
원소이다. 짧은 완전열
$$
0 \to B/hB \to B/h_2B \to B/h_1B \to 0
$$
을 얻으며, 첫째 화살표는 $h_1$에 의한 곱셈으로 주어진다. 오른쪽 두
가군이 $A$ 위에서 평탄하므로 가운데 것도 그러하다. Algebra, Lemma
\ref{algebra-lemma-flat-ses}를 보라.
\end{proof}

\begin{lemma}
\label{divisors-lemma-flat-at-x}
$f : X \to S$를 스킴의 사상이라 하자. $D \subset X$를 $X/S$ 위의 상대
유효 카르티에 인자라 하자. $x \in D$이고 $\mathcal{O}_{X, x}$가
뇌터이면 $f$는 $x$에서 평탄하다.
\end{lemma}

\begin{proof}
$A = \mathcal{O}_{S, f(x)}$와 $B = \mathcal{O}_{X, x}$로 두자.
$h \in B$를 택하여 $D$의 아이디얼을 생성하게 하자. 그러면 $h$는 $B$에서
영인자가 아닌 원소이고 $B/hB$는 평탄한 국소 $A$-대수이다. $I \subset A$를
유한 생성 아이디얼이라 하자. 다음 가환 그림을 생각하자.
$$
\xymatrix{
0 \ar[r] &
B \ar[r]_h &
B \ar[r] &
B/hB \ar[r] & 0 \\
0 \ar[r] &
B \otimes_A I \ar[r]^h \ar[u] &
B \otimes_A I \ar[r] \ar[u] &
B/hB \otimes_A I \ar[r] \ar[u] & 0
}
$$
$B/hB$가 $A$ 위에서 평탄하므로 아래 열은 짧은 완전열이다. Algebra,
Lemma \ref{algebra-lemma-flat-tor-zero}를 보라. 오른쪽 수직 화살표는
$B/hB$가 $A$ 위에서 평탄하므로 단사이다. Algebra, Lemma
\ref{algebra-lemma-flat}을 보라. 따라서 $h$에 의한 곱셈은 가운데 수직
화살표의 핵 $K$ 위에서 전사이다. 나카야마 보조정리, 즉 Algebra, Lemma
\ref{algebra-lemma-NAK}에 의해 $K= 0$임을 얻는다. 따라서 $B$는 $A$ 위에서
평탄하다. Algebra, Lemma \ref{algebra-lemma-flat}을 보라.
\end{proof}

\noindent
다음 보조정리는 평탄 궤적의 열린성의 대수적 판본에 의존한다. 스킴론적
판본은 More on Morphisms, Section \ref{more-morphisms-section-open-flat}에
있다.

\begin{lemma}
\label{divisors-lemma-flat-relative-Cartier-divisor}
$f : X \to S$를 스킴의 사상이라 하자. $D \subset X$를 상대 유효 카르티에
인자라 하자. $f$가 국소 유한 표시이면 열린 부분스킴 $U \subset X$가
존재하여 $D \subset U$이고 $f|_U : U \to S$는 평탄하다.
\end{lemma}

\begin{proof}
$x \in D$를 택하자. 열린 근방 $U \subset X$ 중 $x$를 포함하고 $f|_U$가 평탄한
것을 찾으면 충분하다. 따라서 보조정리는 $X = \Spec(B)$와
$S = \Spec(A)$가 아핀이고 $D$가 영인자가 아닌 원소 $h \in B$로 주어진
경우로 환원된다. 가정에 의해 $B$는 유한 표시 $A$-대수이고 $B/hB$는
평탄한 $A$-대수이다. 절대 뇌터 근사를 사용하겠다.

\medskip\noindent
$B = A[x_1, \ldots, x_n]/(g_1, \ldots, g_m)$로 쓰자. $h$가
$h' \in A[x_1, \ldots, x_n]$의 상이라고 가정하자. 유한형
$\mathbf{Z}$-부분대수 $A_0 \subset A$를 택하여 다항식
$h', g_1, \ldots, g_m$의 모든 계수가 $A_0$에 속하게 하자. 그러면
$B_0 = A_0[x_1, \ldots, x_n]/(g_1, \ldots, g_m)$로 둘 수 있고, $h_0$를
$h'$의 $B_0$ 안의 상으로 둘 수 있다. 그러면
$B = B_0 \otimes_{A_0} A$이고
$B/hB = B_0/h_0B_0 \otimes_{A_0} A$이다. Algebra, Lemma
\ref{algebra-lemma-flat-finite-presentation-limit-flat}에 의해 $A_0$를
확대하고 나면 $B_0/h_0B_0$가 $A_0$ 위에서 평탄하다고 가정해도 된다.
$K_0 = \Ker(h_0 : B_0 \to B_0)$라 하자. $B_0$가 $\mathbf{Z}$ 위에서
유한형이므로 $K_0$는 유한 생성 아이디얼이다. $A_1 \subset A$를 유한형
$\mathbf{Z}$-부분대수라 하되 $A_0$를 포함하게 하고, 대응하는 대상들을
$B_1$, $h_1$, $K_1$으로 나타내되 모두 $A_1$ 위에서 취한다. More on Algebra, Lemma
\ref{more-algebra-lemma-base-change-H1-regular}에 의해 사상
$K_0 \otimes_{A_0} A_1 \to K_1$은 전사이다. 한편 가정에 의해
$h : B \to B$의 핵은 영이다. 따라서 $K_0$의 모든 원소는 $K_1$에서
영으로 가며 이는 충분히 큰 부분환 $A_1 \subset A$에 대해 성립한다.
$K_0$가 유한 생성이므로 $K_1 = 0$이 되는 적절한 $A_1$을 택할 수 있다.

\medskip\noindent
$f_1 : X_1 \to S_1$을 $\Spec$과 같게 두되, 여기서 환 사상은
$A_1 \to B_1$이다.
$D_1 = \Spec(B_1/h_1B_1)$로 두자.
$B = B_1 \otimes_{A_1} A$, 즉 $X = X_1 \times_{S_1} S$이므로, 이제
$X_1 \to S_1$과 상대 유효 카르티에 인자 $D_1$에 대해 보조정리를 보이면
충분하다. Morphisms, Lemma
\ref{morphisms-lemma-base-change-module-flat}를 보라. 따라서 $A$가 뇌터
환인 경우로 환원했다. 이 경우 환 사상 $A \to B$는 모든 소아이디얼
$\mathfrak q$, 즉 $V(h)$의 소아이디얼에서 평탄임을 Lemma
\ref{divisors-lemma-flat-at-x}에 의해
안다. 이 경우 평탄 궤적이 열린집합이라는 사실, 즉 Algebra, Theorem
\ref{algebra-theorem-openness-flatness}과 결합하면 끝난다.
\end{proof}

\noindent
아무런 유한성 가정도 필요로 하지 않는 다음 보조정리도 있다. 그 아이디어는
Michael Artin에게서 나온 것으로 보인다. \cite{Nobile}을 보라.

\begin{lemma}
\label{divisors-lemma-michael-artin}
$f : X \to S$를 스킴의 사상이라 하자. $D \subset X$를 $X/S$ 위의 상대
유효 카르티에 인자라 하자. $f$가 $X \setminus D$의 모든 점에서 평탄하면
$f$는 평탄하다.
\end{lemma}

\begin{proof}
이는 다음 대수 명제로 옮겨진다. $A \to B$를 환 사상이라 하고
$h \in B$라 하자. $h$가 영인자가 아닌 원소이고 $B/hB$가 $A$ 위에서
평탄하며 국소화 $B_h$가 $A$ 위에서 평탄하다고 가정하자. 그러면 $B$는
$A$ 위에서 평탄하다. 그 이유는 짧은 완전열
$$
0 \to B \to B_h \to \colim_n (1/h^n)B/B \to 0
$$
이 있고 둘째와 셋째 항이 $A$ 위에서 평탄하여 $B$도 $A$ 위에서 평탄하기
때문이다(Algebra, Lemma \ref{algebra-lemma-flat-ses}를 보라). 평탄 가군의
여과 여극한은 평탄하고(Algebra, Lemma
\ref{algebra-lemma-colimit-flat}를 보라), $n$에 관한 귀납법으로 각
$(1/h^n)B/B \cong B/h^nB$는 짧은 완전열
$$
0 \to B/h^{n - 1}B \xrightarrow{h} B/h^nB \to B/hB \to 0
$$
에 들어가므로 $A$ 위에서 평탄함에 유의하자. 세부사항은 생략한다.
\end{proof}

\begin{example}
\label{divisors-example-relative-cartier-ambient-space-not-flat}
다음은 상대 유효 카르티에 인자 $D$의 예인데, 주위 스킴은 $D$의 한 근방에서
평탄하지 않다. 즉, $A = k[t]$라 하고
$$
B = k[t, x, y, x^{-1}y, x^{-2}y, \ldots]/(ty, tx^{-1}y, tx^{-2}y, \ldots)
$$
라 하자. 그러면 $B$는 $A$ 위에서 평탄하지 않지만 $B/xB \cong A$는 $A$
위에서 평탄하다. 또한 $x$는 영인자가 아닌 원소이므로 $\Spec(B)$ 안에서
$\Spec(A)$ 위의 상대 유효 카르티에 인자를 정의한다.
\end{example}

\noindent
주위 스킴이 밑 위에서 평탄하고 국소 유한 표시이면, 올을 사용하여 상대
유효 카르티에 인자를 특징지을 수 있다. 이 보조정리에 대한 조금 다른
관점은 More on Morphisms, Lemma
\ref{more-morphisms-lemma-slice-given-element}도 보라.

\begin{lemma}
\label{divisors-lemma-fibre-Cartier}
$\varphi : X \to S$를 국소 유한 표시인 평탄 사상이라 하자.
$Z \subset X$를 닫힌 부분스킴이라 하고 $x \in Z$를 택하자. 그 상을
$s \in S$라 하자.
\begin{enumerate}
\item $Z_s \subset X_s$가 $x$의 한 근방에서 카르티에 인자이면, 열린집합
$U \subset X$와 상대 유효 카르티에 인자 $D \subset U$가 존재하여
$Z \cap U \subset D$이고 $Z_s \cap U = D_s$이다.
\item $Z_s \subset X_s$가 $x$의 한 근방에서 카르티에 인자이고, 사상
$Z \to X$가 유한 표시이며 $Z \to S$가 $x$에서 평탄하면, $U$와 $D$를
택하여 $Z \cap U = D$가 되게 할 수 있다.
\item $Z_s \subset X_s$가 $x$의 한 근방에서 카르티에 인자이고 $Z$가
$X$의 국소 주인 닫힌 부분스킴이며 이 성질이 $x$의 한 근방에서 성립하면,
$U$와 $D$를 택하여
$Z \cap U = D$가 되게 할 수 있다.
\end{enumerate}
특히 $Z \to S$가 국소 유한 표시이고 평탄하며 모든 올
$Z_s \subset X_s$가 유효 카르티에 인자이면 $Z$는 상대 유효 카르티에
인자이다. 마찬가지로 $Z$가 $X$의 국소 주인 닫힌 부분스킴이고 모든 올
$Z_s \subset X_s$가 유효 카르티에 인자이면 $Z$는 상대 유효 카르티에
인자이다.
\end{lemma}

\begin{proof}
아핀 열린 근방 $\Spec(A)$를 택하여 $s$를 포함하게 하고, 아핀 열린 근방
$\Spec(B)$를 택하여 $x$를 포함하게 하며
$\varphi(\Spec(B)) \subset \Spec(A)$가 되게 하자. 소아이디얼
$\mathfrak p \subset A$를 택하여 $s$에 대응하게 하고, 소아이디얼
$\mathfrak q \subset B$를 택하여 $x$에 대응하게 하자. 아이디얼
$I \subset B$를 택하여 $Z$에 대응하게 하자. 보조정리의 첫 가정에 의해
$A \to B$가 평탄하고 유한 표시임을 안다.
(1)의 가정은 $\Spec(B)$를 줄인 뒤
$I(B \otimes_A \kappa(\mathfrak p))$가 $B \otimes_A \kappa(\mathfrak p)$에서
영인자가 아닌 하나의 원소로 생성된다고 가정해도 된다는 뜻이다.
$f \in I$가 이 생성원으로 간다고 하자. 어떤 원소 $g \in B$를 역으로
만들되 $g \not \in \mathfrak q$가 되게 하면, 닫힌 부분스킴
$D = V(f) \subset \Spec(B_g)$가 상대 유효 카르티에 인자라고 주장한다.

\medskip\noindent
Algebra, Lemma \ref{algebra-lemma-flat-finite-presentation-limit-flat}에 의해
뇌터 환의 평탄 유한형 환 사상 $A_0 \to B_0$, 원소 $f_0 \in B_0$, 환 사상
$A_0 \to A$, 그리고 동형 $A \otimes_{A_0} B_0 \cong B$를 찾을 수 있다.
$\mathfrak p_0 = A_0 \cap \mathfrak p$이면
$$
B \otimes_A \kappa(\mathfrak p) =
\left(B_0 \otimes_{A_0} \kappa(\mathfrak p_0)\right)
\otimes_{\kappa(\mathfrak p_0))} \kappa(\mathfrak p)
$$
임을 안다. 따라서 $f_0$는 $B_0 \otimes_{A_0} \kappa(\mathfrak p_0)$에서
영인자가 아닌 원소이다. Algebra, Lemma
\ref{algebra-lemma-grothendieck}에 의해 $f_0$는
$(B_0)_{\mathfrak q_0}$에서 영인자가 아닌 원소이고, 여기서
$\mathfrak q_0 = B_0 \cap \mathfrak q$이다. 또한
$(B_0/f_0B_0)_{\mathfrak q_0}$는 $A_0$ 위에서 평탄하다. 따라서 Algebra,
Lemma \ref{algebra-lemma-regular-sequence-in-neighbourhood}와 Algebra,
Theorem \ref{algebra-theorem-openness-flatness}에 의해
$g_0 \in B_0$, $g_0 \not \in \mathfrak q_0$가 존재하여 $f_0$는
$(B_0)_{g_0}$에서 영인자가 아닌 원소이고
$(B_0/f_0B_0)_{g_0}$는 $A_0$ 위에서 평탄하다. 따라서
$D_0 = V(f_0) \subset \Spec((B_0)_{g_0})$는 상대 유효 카르티에 인자이다.
이 성질은 밑변환 아래 보존되므로(Lemma \ref{divisors-lemma-relative-Cartier} 참조),
위 주장은 $g$를 $g_0$의 $B$ 안의 상으로 택하면 성립한다.

\medskip\noindent
이제 (1)을 증명했다. (2)를 보기 위해 닫힌 몰입 $Z \to D$를 생각하자.
전사 환 사상 $u : \mathcal{O}_{D, x} \to \mathcal{O}_{Z, x}$는
본질적으로 유한 표시인 평탄한 국소 $\mathcal{O}_{S, s}$-대수 사이의
사상이고, $\mathfrak m_s$로 나눈 뒤 동형이 된다. 따라서 이는 동형이다.
Algebra, Lemma \ref{algebra-lemma-mod-injective-general}를 보라. 그러므로
$Z \to D$는 $x$의 한 근방에서 동형이다. Algebra, Lemma
\ref{algebra-lemma-local-isomorphism}를 보라. (3)을 보기 위해 필요하면
$U$를 줄여 $D$의 아이디얼이 영인자가 아닌 한 원소 $f$로 생성되고 $Z$의
아이디얼이 원소 $g$로 생성된다고 가정해도 된다. 그러면 $f = gh$이다.
그러나 $g|_{U_s}$와 $f|_{U_s}$는 $x$의 한 근방에서 같은 유효 카르티에
인자를 잘라낸다. 따라서 $h|_{X_s}$는 $\mathcal{O}_{X_s, x}$의 단원이고,
따라서 $h$는 $\mathcal{O}_{X, x}$의 단원이며, 따라서 $h$는 $x$의 열린
근방에서 단원이다. 즉, $Z \cap U = D$이고 이는 $U$를 줄인 뒤 성립한다.

\medskip\noindent
보조정리의 마지막 명제들은 (2), (3)과 다음 사실을 결합하면 바로 따른다.
$Z \to S$가 국소 유한 표시일 필요충분조건은 $Z \to X$가 유한 표시인
것이다. Morphisms, Lemmas
\ref{morphisms-lemma-composition-finite-presentation}와
\ref{morphisms-lemma-finite-presentation-permanence}를 보라.
\end{proof}


\section{몰입의 법선 원뿔}
\label{divisors-section-normal-cone}

\noindent
$i : Z \to X$를 닫힌 몰입이라 하자. 대응하는 준연접 아이디얼층을
$\mathcal{I} \subset \mathcal{O}_X$라 하자. 등급 $\mathcal{O}_X$-대수의
준연접층
$\bigoplus_{n \geq 0} \mathcal{I}^n/\mathcal{I}^{n + 1}$을 생각하자.
각 층 $\mathcal{I}^n/\mathcal{I}^{n + 1}$은 $\mathcal{I}$에 의해 영이
되므로 이 등급 대수는 Morphisms, Lemma
\ref{morphisms-lemma-i-star-equivalence}에 의해 등급 $\mathcal{O}_Z$-대수의
준연접층에 대응한다. 이 준연접 등급 $\mathcal{O}_Z$-대수를
{\it $Z$의 $X$ 안에서의 여법선 대수}라 부른다. Morphisms, Section
\ref{morphisms-section-closed-immersions-quasi-coherent}에서 언급한 표기의
남용으로 흔히 단순히
$\bigoplus_{n \geq 0} \mathcal{I}^n/\mathcal{I}^{n + 1}$로 나타낸다.

\medskip\noindent
$f : Z \to X$를 몰입이라 하자. $f$의 여법선 대수를 닫힌 몰입
$i : Z \to X \setminus \partial Z$의 여법선 대수로 정의한다. 여기서
$\partial Z = \overline{Z} \setminus Z$이다. 이를 흔히
$\bigoplus_{n \geq 0} \mathcal{I}^n/\mathcal{I}^{n + 1}$로 나타내며,
$\mathcal{I}$는 닫힌 몰입 $i : Z \to X \setminus \partial Z$의 아이디얼층이다.

\begin{definition}
\label{divisors-definition-conormal-sheaf}
$f : Z \to X$를 몰입이라 하자. {\it 여법선 대수
$\mathcal{C}_{Z/X, *}$, 즉 $Z$의 $X$ 안에서의 여법선 대수} 또는
{\it $f$의 여법선 대수}란 위에서 설명한
등급 $\mathcal{O}_Z$-대수의 준연접층
$\bigoplus_{n \geq 0} \mathcal{I}^n/\mathcal{I}^{n + 1}$을 말한다.
\end{definition}

\noindent
따라서 $\mathcal{C}_{Z/X, 1} = \mathcal{C}_{Z/X}$은 몰입의 여법선층이다.
또한 $\mathcal{C}_{Z/X, 0} = \mathcal{O}_Z$이고,
$\mathcal{C}_{Z/X, n}$은 다음 성질로 특징지어지는 준연접
$\mathcal{O}_Z$-가군이다.
\begin{equation}
\label{divisors-equation-conormal-in-degree-n}
i_*\mathcal{C}_{Z/X, n} = \mathcal{I}^n/\mathcal{I}^{n + 1}
\end{equation}
여기서 $i : Z \to X \setminus \partial Z$이고 $\mathcal{I}$는 위와 같은
$i$의 아이디얼층이다. 끝으로, 준연접 등급 $\mathcal{O}_Z$-대수의 표준
전사 사상
\begin{equation}
\label{divisors-equation-conormal-algebra-quotient}
\text{Sym}^*(\mathcal{C}_{Z/X}) \longrightarrow \mathcal{C}_{Z/X, *}
\end{equation}
이 있으며 차수 $0$과 $1$에서는 동형임에 유의하자.

\begin{lemma}
\label{divisors-lemma-affine-conormal-sheaf}
$i : Z \to X$를 몰입이라 하자. $i$의 여법선 대수는 다음 성질을 갖는다.
\begin{enumerate}
\item $U \subset X$를 $i(Z)$가 $U$의 닫힌 부분집합이 되는 임의의 열린집합이라
하자. 아이디얼층을 $\mathcal{I} \subset \mathcal{O}_U$라 하되 대응하는
닫힌 부분스킴이 $i(Z) \subset U$가 되게 하자. 그러면
$$
\mathcal{C}_{Z/X, *} =
i^*\left(\bigoplus\nolimits_{n \geq 0} \mathcal{I}^n\right) =
i^{-1}\left(
\bigoplus\nolimits_{n \geq 0} \mathcal{I}^n/\mathcal{I}^{n + 1}
\right)
$$
이다.
\item 임의의 아핀 열린집합 $\Spec(R) = U \subset X$가
$Z \cap U = \Spec(R/I)$를 만족하면 표준 동형
$\Gamma(Z \cap U, \mathcal{C}_{Z/X, *}) = \bigoplus_{n \geq 0} I^n/I^{n + 1}$
이 있다.
\end{enumerate}
\end{lemma}

\begin{proof}
대부분 정의에서 분명하다. 환 $R$과 아이디얼 $I$가 주어지고 이 아이디얼이
$R$의 것이면
$I^n/I^{n + 1} = I^n \otimes_R R/I$임에 유의하자. 세부사항은 생략한다.
\end{proof}

\begin{lemma}
\label{divisors-lemma-conormal-algebra-functorial}
스킴의 범주에서 다음 가환 그림이 주어졌다고 하자.
$$
\xymatrix{
Z \ar[r]_i \ar[d]_f & X \ar[d]^g \\
Z' \ar[r]^{i'} & X'
}
$$
$i$, $i'$이 몰입이라고 가정하자. 그러면 등급 $\mathcal{O}_Z$-대수의 표준
사상
$$
f^*\mathcal{C}_{Z'/X', *}
\longrightarrow
\mathcal{C}_{Z/X, *}
$$
이 존재하며 다음 성질로 특징지어진다. 아핀 열린집합의 모든 쌍
$(\Spec(R) = U \subset X, \Spec(R') = U' \subset X')$가
$g(U) \subset U'$를 만족하고
$Z \cap U = \Spec(R/I)$ 및 $Z' \cap U' = \Spec(R'/I')$를 만족하면, 유도 사상
$$
\Gamma(Z' \cap U', \mathcal{C}_{Z'/X', *}) =
\bigoplus\nolimits (I')^n/(I')^{n + 1}
\longrightarrow
\bigoplus\nolimits_{n \geq 0} I^n/I^{n + 1} =
\Gamma(Z \cap U, \mathcal{C}_{Z/X, *})
$$
은 환 사상 $f^\sharp : R' \to R$에서 유도된 것이며, 이 환 사상은
$f^\sharp(I') \subset I$를 만족한다.
\end{lemma}

\begin{proof}
$\partial Z' = \overline{Z'} \setminus Z'$ 및
$\partial Z = \overline{Z} \setminus Z$라 하자. 이들은 각각 $X'$와 $X$의
닫힌 부분집합이다. $X'$를 $X' \setminus \partial Z'$로 바꾸고 $X$를
$X \setminus \big(g^{-1}(\partial Z') \cup \partial Z\big)$로 바꾸면
$i$와 $i'$이 닫힌 몰입이라고 가정해도 된다.

\medskip\noindent
$g \circ i$가 $i'$을 통해 분해된다는 사실은 $g^*\mathcal{I}'$이
$\mathcal{I}$ 안으로 사상됨을 함의하며, 이는 표준 사상
$g^*\mathcal{I}' \to \mathcal{O}_X$ 아래에서의 말이다. Schemes, Lemmas
\ref{schemes-lemma-characterize-closed-subspace}와
\ref{schemes-lemma-restrict-map-to-closed}를 보라. 따라서 준연접층의 유도
사상
$g^*((\mathcal{I}')^n/(\mathcal{I}')^{n + 1}) \to
\mathcal{I}^n/\mathcal{I}^{n + 1}$을 얻는다. $i$로 당기면
$i^*g^*((\mathcal{I}')^n/(\mathcal{I}')^{n + 1}) \to
i^*(\mathcal{I}^n/\mathcal{I}^{n + 1})$를 얻는다. 다음에 유의하자.
$i^*(\mathcal{I}^n/\mathcal{I}^{n + 1}) = \mathcal{C}_{Z/X, n}$.
한편
$i^*g^*((\mathcal{I}')^n/(\mathcal{I}')^{n + 1}) =
f^*(i')^*((\mathcal{I}')^n/(\mathcal{I}')^{n + 1}) =
f^*\mathcal{C}_{Z'/X', n}$이다. 이것이 원하는 사상을 준다.

\medskip\noindent
그 사상이 국소적으로 주어진 사상
$(I')^n/(I')^{n + 1} \to I^n/I^{n + 1}$로 기술됨을 확인하려면 정의를
풀어 쓰면 되므로 생략한다. 또 하나 유의할 점은 임의의 $x \in i(Z)$에
대해 $U$, $U'$라는 아핀 열린 근방들이 실제로 존재하여
$f(U) \subset U'$이고, $Z \cap U$와 $U' \cap Z'$가 모두 닫히며,
$x \in U$가 된다는 것이다. 증명은 생략한다. 따라서 보조정리의 요건은
실제로 사상을 특징짓고, 그 요건을 정의로 사용할 수도 있었다.
\end{proof}

\begin{lemma}
\label{divisors-lemma-conormal-algebra-functorial-flat}
다음 그림이 스킴의 범주에서의 올곱 그림이라고 하자.
$$
\xymatrix{
Z \ar[r]_i \ar[d]_f & X \ar[d]^g \\
Z' \ar[r]^{i'} & X'
}
$$
$i$, $i'$은 몰입이라 하자. 그러면 Lemma
\ref{divisors-lemma-conormal-algebra-functorial}의 표준 사상
$f^*\mathcal{C}_{Z'/X', *} \to \mathcal{C}_{Z/X, *}$은 전사이다.
$g$가 평탄하면 이는 동형이다.
\end{lemma}

\begin{proof}
$R' \to R$을 환 사상이라 하고 $I' \subset R'$을 아이디얼이라 하자.
$I = I'R$로 두자. 그러면
$(I')^n/(I')^{n + 1} \otimes_{R'} R \to I^n/I^{n + 1}$은 전사이다.
$R' \to R$이 평탄하면 $I^n = (I')^n \otimes_{R'} R$이고 이 사상은
동형임을 안다.
\end{proof}

\begin{definition}
\label{divisors-definition-normal-cone}
$i : Z \to X$를 스킴의 몰입이라 하자. {\it 법선 원뿔 $C_ZX$, 즉
$Z$의 $X$ 안에서의 법선 원뿔}은
$$
C_ZX = \underline{\Spec}_Z(\mathcal{C}_{Z/X, *})
$$
이다. Constructions, Definitions \ref{constructions-definition-cone}와
\ref{constructions-definition-abstract-cone}를 보라. $Z$의 $X$ 안에서의
{\it 법선 다발}은 벡터 다발
$$
N_ZX = \underline{\Spec}_Z(\text{Sym}(\mathcal{C}_{Z/X}))
$$
이다. Constructions, Definitions
\ref{constructions-definition-vector-bundle}와
\ref{constructions-definition-abstract-vector-bundle}을 보라.
\end{definition}

\noindent
따라서 $C_ZX \to Z$는 $Z$ 위의 원뿔이고 $N_ZX \to Z$는 $Z$ 위의 벡터
다발이다(우리의 용어에서 이것은 여법선층이 유한 국소 자유층임을 함의하지
않음을 상기하라). 또한 등급 대수의 표준 전사
(\ref{divisors-equation-conormal-algebra-quotient})는 $Z$ 위의 원뿔의 표준 닫힌 몰입
\begin{equation}
\label{divisors-equation-normal-cone-in-normal-bundle}
C_ZX \longrightarrow N_ZX
\end{equation}
을 정의한다.


\section{정칙 아이디얼층}
\label{divisors-section-regular-ideal-sheaves}

\noindent
이 절에서는 유효 카르티에 인자의 개념을 더 높은 여차원으로 일반화한다.
원소열 $f_1, \ldots, f_r$이 환 $R$에서 {\it 정칙열}이라는 것은 각
$i = 1, \ldots, r$에 대해 원소 $f_i$가
$R/(f_1, \ldots, f_{i - 1})$ 위에서 영인자가 아니고
$R/(f_1, \ldots, f_r) \not = 0$인 것을 뜻한다. Algebra, Definition
\ref{algebra-definition-regular-sequence}을 보라. 이와 밀접하게 관련된
더 약한 조건 세 가지를 부과할 수 있다. 첫째 조건은
$f_1, \ldots, f_r$이 {\it 코쥘 정칙열}, 즉
$H_i(K_\bullet(f_1, \ldots, f_r)) = 0$이 $i > 0$에 대해 성립한다고
가정하는 것이다. More on Algebra, Definition
\ref{more-algebra-definition-koszul-regular-sequence}을 보라. 이 열을
{\it $H_1$-정칙열}이라 부르는 것은
$H_1(K_\bullet(f_1, \ldots, f_r)) = 0$일 때이다. 또 다른 조건은
$J = (f_1, \ldots, f_r)$로 둘 때 사상
$$
R/J[T_1, \ldots, T_r]
\longrightarrow
\bigoplus\nolimits_{n \geq 0}
J^n/J^{n + 1}
$$
이 $T_i$를 $f_i \bmod J^2$로 보내며 동형인 것이다. 이 경우
$f_1, \ldots, f_r$을 {\it 준정칙열}이라 한다. Algebra, Definition
\ref{algebra-definition-quasi-regular-sequence}을 보라. $R$-가군 $M$이
주어지면 $M$-정칙열과 $M$-준정칙열이라는 개념도 있다.

\medskip\noindent
이를 환 달린 공간의 경우에 다음과 같이 일반화할 수 있다. $X$를 환 달린
공간이라 하고 $f_1, \ldots, f_r \in \Gamma(X, \mathcal{O}_X)$라 하자.
$f_1, \ldots, f_r$이 {\it 정칙열}이라는 것은 각 $i = 1, \ldots, r$에 대해
사상
\begin{equation}
\label{divisors-equation-map-regular}
f_i :
\mathcal{O}_X/(f_1, \ldots, f_{i - 1})
\longrightarrow
\mathcal{O}_X/(f_1, \ldots, f_{i - 1})
\end{equation}
이 층의 단사 사상인 것을 뜻한다. $f_1, \ldots, f_r$이
{\it 코쥘 정칙열}이라는 것은 코쥘 복합체
\begin{equation}
\label{divisors-equation-koszul}
K_\bullet(\mathcal{O}_X, f_\bullet),
\end{equation}
즉 Modules, Definition \ref{modules-definition-koszul-complex}의 복합체가
차수 $> 0$에서 비순환인 것을 뜻한다. $f_1, \ldots, f_r$이
{\it $H_1$-정칙열}이라는 것은 코쥘 복합체
$K_\bullet(\mathcal{O}_X, f_\bullet)$가 차수 $1$에서 완전한 것을 뜻한다.
끝으로 $f_1, \ldots, f_r$이 {\it 준정칙열}이라는 것은 사상
\begin{equation}
\label{divisors-equation-map-quasi-regular}
\mathcal{O}_X/\mathcal{J}[T_1, \ldots, T_r]
\longrightarrow
\bigoplus\nolimits_{d \geq 0} \mathcal{J}^d/\mathcal{J}^{d + 1}
\end{equation}
이 층의 동형인 것을 뜻한다. 여기서 $\mathcal{J} \subset \mathcal{O}_X$는
$f_1, \ldots, f_r$로 생성되는 아이디얼층이다. ($\mathcal{F}$-정칙열과
$\mathcal{F}$-준정칙열의 개념도 있으며, 이는 주어진 $\mathcal{O}_X$-가군
$\mathcal{F}$에 대한 것이다. 필요해지면 여기서 도입할 것이다.)

\begin{lemma}
\label{divisors-lemma-types-regular-sequences-implications}
$X$를 환 달린 공간이라 하고
$f_1, \ldots, f_r \in \Gamma(X, \mathcal{O}_X)$라 하자. 다음 함의들이
성립한다.
$f_1, \ldots, f_r$이 정칙열 $\Rightarrow$
$f_1, \ldots, f_r$이 코쥘 정칙열 $\Rightarrow$
$f_1, \ldots, f_r$이 $H_1$-정칙열 $\Rightarrow$
$f_1, \ldots, f_r$이 준정칙열.
\end{lemma}

\begin{proof}
완전성은 줄기에서 확인할 수 있으므로, 열 $f_1, \ldots, f_r$이 정칙열일
필요충분조건은 사상
$$
f_i :
\mathcal{O}_{X, x}/(f_1, \ldots, f_{i - 1})
\longrightarrow
\mathcal{O}_{X, x}/(f_1, \ldots, f_{i - 1})
$$
이 모든 $x \in X$에 대해 단사인 것이다. 다른 정칙성의 종류도 줄기별로
확인할 수 있다(세부사항은 생략한다). 따라서 함의들은 More on Algebra,
Lemmas \ref{more-algebra-lemma-regular-koszul-regular},
\ref{more-algebra-lemma-koszul-regular-H1-regular}, 및
\ref{more-algebra-lemma-H1-regular-quasi-regular}에서 따른다.
\end{proof}

\begin{definition}
\label{divisors-definition-regular-ideal-sheaf}
\begin{reference}
코쥘 정칙 아이디얼층의 개념은 \cite[Expose VII, Definition 1.4]{SGA6}에서
도입되었으며, 거기서는 정칙 아이디얼층이라 불렸다.
\end{reference}
$X$를 환 달린 공간이라 하고 $\mathcal{J} \subset \mathcal{O}_X$를
아이디얼층이라 하자.
\begin{enumerate}
\item $\mathcal{J}$가 {\it 정칙}이라는 것은 모든
$x \in \text{Supp}(\mathcal{O}_X/\mathcal{J})$에 대해 열린 근방
$x \in U \subset X$와 정칙열 $f_1, \ldots, f_r \in \mathcal{O}_X(U)$가
존재하여 $\mathcal{J}|_U$가 $f_1, \ldots, f_r$로 생성되는 것이다.
\item $\mathcal{J}$가 {\it 코쥘 정칙}이라는 것은 모든
$x \in \text{Supp}(\mathcal{O}_X/\mathcal{J})$에 대해 열린 근방
$x \in U \subset X$와 코쥘 정칙열
$f_1, \ldots, f_r \in \mathcal{O}_X(U)$가 존재하여 $\mathcal{J}|_U$가
$f_1, \ldots, f_r$로 생성되는 것이다.
\item $\mathcal{J}$가 {\it $H_1$-정칙}이라는 것은 모든
$x \in \text{Supp}(\mathcal{O}_X/\mathcal{J})$에 대해 열린 근방
$x \in U \subset X$와 $H_1$-정칙열
$f_1, \ldots, f_r \in \mathcal{O}_X(U)$가 존재하여 $\mathcal{J}|_U$가
$f_1, \ldots, f_r$로 생성되는 것이다.
\item $\mathcal{J}$가 {\it 준정칙}이라는 것은 모든
$x \in \text{Supp}(\mathcal{O}_X/\mathcal{J})$에 대해 열린 근방
$x \in U \subset X$와 준정칙열 $f_1, \ldots, f_r \in \mathcal{O}_X(U)$가
존재하여 $\mathcal{J}|_U$가 $f_1, \ldots, f_r$로 생성되는 것이다.
\end{enumerate}
\end{definition}

\noindent
이 개념의 많은 성질은 환에서의 정칙열과 준정칙열에 대한 대응 개념으로부터
곧바로 따른다.

\begin{lemma}
\label{divisors-lemma-regular-quasi-regular-scheme}
$X$를 환 달린 공간이라 하고 $\mathcal{J}$를 아이디얼층이라 하자. 다음
함의들이 성립한다.
$\mathcal{J}$가 정칙 $\Rightarrow$
$\mathcal{J}$가 코쥘 정칙 $\Rightarrow$
$\mathcal{J}$가 $H_1$-정칙 $\Rightarrow$
$\mathcal{J}$가 준정칙.
\end{lemma}

\begin{proof}
이는 바로 Lemma \ref{divisors-lemma-types-regular-sequences-implications}로 환원된다.
\end{proof}

\begin{lemma}
\label{divisors-lemma-quasi-regular-ideal}
$X$를 국소 환 달린 공간이라 하고 $\mathcal{J} \subset \mathcal{O}_X$를
아이디얼층이라 하자. 그러면 $\mathcal{J}$가 준정칙일 필요충분조건은 다음
조건들이 성립하는 것이다.
\begin{enumerate}
\item $\mathcal{J}$는 유한형 $\mathcal{O}_X$-가군이고,
\item $\mathcal{J}/\mathcal{J}^2$는 유한 국소 자유
$\mathcal{O}_X/\mathcal{J}$-가군이며,
\item 표준 사상
$$
\text{Sym}^n_{\mathcal{O}_X/\mathcal{J}}(\mathcal{J}/\mathcal{J}^2)
\longrightarrow
\mathcal{J}^n/\mathcal{J}^{n + 1}
$$
들은 모든 $n \geq 0$에 대해 동형이다.
\end{enumerate}
\end{lemma}

\begin{proof}
열린집합 $U \subset X$에서 $\mathcal{J}|_U$가 준정칙열
$f_1, \ldots, f_r \in \mathcal{O}_X(U)$로 생성되면 $\mathcal{J}|_U$는
유한형이고, $\mathcal{J}|_U/\mathcal{J}^2|_U$는 기저
$f_1, \ldots, f_r$를 갖는 자유 가군이며, (3)의 사상들은 동형임이
분명하다. 마지막 사실은 이 사상들이
(\ref{divisors-equation-map-quasi-regular})의 차수 $n$ 부분을 좌표와 무관하게
표현한 것이기 때문이다. 따라서 준정칙이면 (1), (2), (3)이 성립함이
분명하다.

\medskip\noindent
역으로 (1), (2), (3)이 성립한다고 하자. 점
$x \in \text{Supp}(\mathcal{O}_X/\mathcal{J})$를 택하자. 그러면 근방
$U \subset X$가 존재하여 $x$를 포함하고,
$\mathcal{J}|_U/\mathcal{J}^2|_U$가 계수 $r$인 자유
$\mathcal{O}_U/\mathcal{J}|_U$-가군이 된다.
필요하면 $U$를 줄여 $f_1, \ldots, f_r \in \mathcal{J}(U)$가 존재하고
$\mathcal{J}|_U/\mathcal{J}^2|_U$의 기저로 가며 이를
$\mathcal{O}_U/\mathcal{J}|_U$-가군으로 보게 할 수 있다. 특히
$f_1, \ldots, f_r$의 상들이 $\mathcal{J}_x/\mathcal{J}^2_x$ 안에서 생성함을
안다. 따라서 나카야마 보조정리(Algebra, Lemma
\ref{algebra-lemma-NAK})에 의해 $f_1, \ldots, f_r$은 줄기
$\mathcal{J}_x$를 생성한다. $\mathcal{J}$가 유한형이므로 Modules, Lemma
\ref{modules-lemma-finite-type-surjective-on-stalk}에 의해 $U$를 줄여
$f_1, \ldots, f_r$이 $\mathcal{J}$를 생성한다고 가정해도 된다. 끝으로
(3)과 동형
$\mathcal{J}|_U/\mathcal{J}^2|_U = \bigoplus \mathcal{O}_U/\mathcal{J}|_U f_i$
로부터 $f_1, \ldots, f_r \in \mathcal{O}_X(U)$가 준정칙열임이 분명하다.
\end{proof}

\begin{lemma}
\label{divisors-lemma-generate-regular-ideal}
$(X, \mathcal{O}_X)$를 국소 환 달린 공간이라 하자.
$\mathcal{J} \subset \mathcal{O}_X$를 아이디얼층이라 하자. $x \in X$라
하고 $f_1, \ldots, f_r \in \mathcal{J}_x$의 상들이
$\kappa(x)$-벡터 공간 $\mathcal{J}_x/\mathfrak m_x\mathcal{J}_x$의 기저를
이룬다고 하자.
\begin{enumerate}
\item $\mathcal{J}$가 준정칙이면 열린 근방이 존재하여
$f_1, \ldots, f_r \in \mathcal{O}_X(U)$가 $\mathcal{J}|_U$를 생성하는
준정칙열을 이룬다.
\item $\mathcal{J}$가 $H_1$-정칙이면 열린 근방이 존재하여
$f_1, \ldots, f_r \in \mathcal{O}_X(U)$가 $H_1$-정칙열을 이루고
$\mathcal{J}|_U$를 생성한다.
\item $\mathcal{J}$가 코쥘 정칙이면 열린 근방이 존재하여
$f_1, \ldots, f_r \in \mathcal{O}_X(U)$가 $\mathcal{J}|_U$를 생성하는
코쥘 정칙열을 이룬다.
\end{enumerate}
\end{lemma}

\begin{proof}
먼저 $\mathcal{J}$가 준정칙이라고 가정하자. 열린 근방
$U \subset X$를 택하여 $x$를 포함하게 하고, 준정칙열
$g_1, \ldots, g_s \in \mathcal{O}_X(U)$를 택하여 $\mathcal{J}|_U$를
생성하게 할 수 있다. 이는 $\mathcal{J}/\mathcal{J}^2$가 계수 $s$인 자유
$\mathcal{O}_U/\mathcal{J}|_U$-가군임을 함의하므로(Lemma
\ref{divisors-lemma-quasi-regular-ideal}와
그 증명을 보라) $r = s$이다. $U$를 줄여
$f_1, \ldots, f_r \in \mathcal{J}(U)$라고 가정해도 된다. 따라서
$$
f_i = \sum a_{ij} g_j
$$
로 쓸 수 있고, 여기서 어떤 $a_{ij} \in \mathcal{O}_X(U)$들이 있다.
가정에 의해 행렬 $A = (a_{ij})$는 $\kappa(x)$ 위의 가역행렬로 간다.
따라서 $U$를 한 번 더 줄여 $(a_{ij})$가 가역이라고 가정해도 된다.
그러므로 $f_1, \ldots, f_r$이 $(\mathcal{J}/\mathcal{J}^2)|_U$의 기저를
주고, 이는 $f_1, \ldots, f_r$이 $U$ 위의 준정칙열임을 증명한다.

\medskip\noindent
(2)와 (3)의 가정은 (1)의 가정보다 강하므로, (2)와 (3)을 증명할 때에는
이미 $f_1, \ldots, f_r \in \mathcal{J}(U)$이고
$f_i = \sum a_{ij}g_j$이며 위와 같이 $(a_{ij})$가 가역이라고 가정해도
됨에 유의하자. 이제 $g_1, \ldots, g_r$은 $H_1$-정칙열 또는 코쥘
정칙열이다. $f_1, \ldots, f_r$ 위의 코쥘 복합체는
$g_1, \ldots, g_r$ 위의 코쥘 복합체와 동형이고, 이 동형은 행렬
$(a_{ij})$로 주어진다(More on Algebra, Lemma
\ref{more-algebra-lemma-change-basis}를 보라). 따라서 원하는 대로
$f_1, \ldots, f_r$은 $H_1$-정칙 또는 코쥘 정칙이다.
\end{proof}

\begin{lemma}
\label{divisors-lemma-regular-ideal-sheaf-quasi-coherent}
스킴 위의 정칙, 코쥘 정칙, $H_1$-정칙 또는 준정칙 아이디얼층은 모두
유한형 준연접 아이디얼층이다.
\end{lemma}

\begin{proof}
이러한 아이디얼층은 국소적으로 유한 개의 절단으로 생성되므로 결론이
따른다. 스킴 위에서 절단들로 국소 생성되는 모든 아이디얼층은 준연접이다.
Schemes, Lemma \ref{schemes-lemma-closed-subspace-scheme}를 보라.
\end{proof}

\begin{lemma}
\label{divisors-lemma-regular-ideal-sheaf-scheme}
$X$를 스킴이라 하고 $\mathcal{J}$를 아이디얼층이라 하자. 그러면
$\mathcal{J}$가 정칙(각각 코쥘 정칙, $H_1$-정칙, 준정칙)일 필요충분조건은
모든 $x \in \text{Supp}(\mathcal{O}_X/\mathcal{J})$에 대해 아핀 열린 근방
$x \in U \subset X$, $U = \Spec(A)$가 존재하여
$\mathcal{J}|_U = \widetilde{I}$이고 $I$가 정칙(각각 코쥘 정칙,
$H_1$-정칙, 준정칙)열 $f_1, \ldots, f_r \in A$로 생성되는 것이다.
\end{lemma}

\begin{proof}
가정에 의해 열린 근방 $U$를 찾아 $x$를 포함하게 하고, 그 위에서
$\mathcal{J}$가 정칙
(각각 코쥘 정칙, $H_1$-정칙, 준정칙)열
$f_1, \ldots, f_r \in \mathcal{O}_X(U)$로 생성되게 할 수 있다. $U$를
줄여 $U$가 아핀, 즉 $U = \Spec(A)$라고 가정해도 된다. $\mathcal{J}$는
Lemma \ref{divisors-lemma-regular-ideal-sheaf-quasi-coherent}에 의해 준연접이므로,
$\mathcal{J}|_U = \widetilde{I}$이고 $I \subset A$인 아이디얼을 택할 수 있다.
이제 함자
$$
\widetilde{\ } : \text{Mod}_A \longrightarrow \QCoh(\mathcal{O}_U)
$$
가 완전성을 보존하는 범주의 동치라는 사실을 사용할 수 있다. 예를 들어
함수들 $f_i$가 $\mathcal{J}$를 생성한다는 것은 $f_i$들을 $A$의 원소로 보면
$I$를 생성한다는 뜻이다. (\ref{divisors-equation-map-regular})가 단사라는 사실
(각각 (\ref{divisors-equation-koszul})가 완전, (\ref{divisors-equation-koszul})가 차수
$1$에서 완전, (\ref{divisors-equation-map-quasi-regular})가 동형)은 대응 사상
$A/(f_1, \ldots, f_{i - 1}) \to A/(f_1, \ldots, f_{i - 1})$(각각 복합체
$K_\bullet(A, f_1, \ldots, f_r)$, 사상
$A/I[T_1, \ldots, T_r] \to \bigoplus I^n/I^{n + 1}$)의 대응 성질을
함의한다. 따라서 $f_1, \ldots, f_r \in A$는 정칙(각각 코쥘 정칙,
$H_1$-정칙, 준정칙)열이며, 환 $A$의 원소열이다.
\end{proof}

\begin{lemma}
\label{divisors-lemma-Noetherian-scheme-regular-ideal}
$X$를 국소 뇌터 스킴이라 하고 $\mathcal{J} \subset \mathcal{O}_X$를 준연접
아이디얼층이라 하자. $x$를 $\mathcal{O}_X/\mathcal{J}$의 지지집합의 한
점이라 하자. 다음은 동치이다.
\begin{enumerate}
\item $\mathcal{J}_x$는 $\mathcal{O}_{X, x}$ 안의 정칙열로 생성되고,
\item $\mathcal{J}_x$는 $\mathcal{O}_{X, x}$ 안의 코쥘 정칙열로 생성되고,
\item $\mathcal{J}_x$는 $H_1$-정칙열로 생성되고 그 열은
$\mathcal{O}_{X, x}$ 안에 있으며,
\item $\mathcal{J}_x$는 $\mathcal{O}_{X, x}$ 안의 준정칙열로 생성되고,
\item 아핀 근방 $U = \Spec(A)$가 존재하여 $x$를 포함하고,
$\mathcal{J}|_U = \widetilde{I}$이고 $I$는 $A$ 안의 정칙열로 생성되며,
\item 아핀 근방 $U = \Spec(A)$가 존재하여 $x$를 포함하고,
$\mathcal{J}|_U = \widetilde{I}$이고 $I$는 $A$ 안의 코쥘 정칙열로 생성되며,
\item 아핀 근방 $U = \Spec(A)$가 존재하여 $x$를 포함하고,
$\mathcal{J}|_U = \widetilde{I}$이고 $I$는 $H_1$-정칙열로 생성되며 그 열은
$A$ 안에 있고,
\item 아핀 근방 $U = \Spec(A)$가 존재하여 $x$를 포함하고,
$\mathcal{J}|_U = \widetilde{I}$이고 $I$는 $A$ 안의 준정칙열로 생성되고,
\item 근방 $U$가 존재하여 $x$를 포함하고 $\mathcal{J}|_U$가 정칙이고,
\item 근방 $U$가 존재하여 $x$를 포함하고 $\mathcal{J}|_U$가 코쥘 정칙이고,
\item 근방 $U$가 존재하여 $x$를 포함하고 $\mathcal{J}|_U$가 $H_1$-정칙이고,
\item 근방 $U$가 존재하여 $x$를 포함하고 $\mathcal{J}|_U$가 준정칙이다.
\end{enumerate}
특히 국소 뇌터 스킴 위에서 정칙, 코쥘 정칙, $H_1$-정칙, 준정칙
아이디얼층의 개념은 모두 일치한다.
\end{lemma}

\begin{proof}
Lemma \ref{divisors-lemma-regular-ideal-sheaf-scheme}로부터 (5) $\Leftrightarrow$ (9),
(6) $\Leftrightarrow$ (10), (7) $\Leftrightarrow$ (11), (8)
$\Leftrightarrow$ (12)를 얻는다. (5) $\Rightarrow$ (1), (6)
$\Rightarrow$ (2), (7) $\Rightarrow$ (3), (8) $\Rightarrow$ (4)는 분명하다.
Algebra, Lemma \ref{algebra-lemma-regular-sequence-in-neighbourhood}에 의해
(1) $\Rightarrow$ (5)이다. Lemma
\ref{divisors-lemma-regular-quasi-regular-scheme}에 의해 (9) $\Rightarrow$ (10)
$\Rightarrow$ (11) $\Rightarrow$ (12)이다. 끝으로 Algebra, Lemma
\ref{algebra-lemma-quasi-regular-regular}에 의해 (4) $\Rightarrow$ (1)이다.
이제 열두 명제는 모두 동치이다.
\end{proof}



\section{정칙 몰입}
\label{divisors-section-regular-immersions}

\noindent
$i : Z \to X$를 스킴의 몰입이라 하자. 정의에 의해 이는 열린 부분스킴
$U \subset X$가 존재하여 $Z$가 $U$의 닫힌 부분스킴과 동일시된다는
뜻이다. 이에 대응하는 준연접 아이디얼층을
$\mathcal{I} \subset \mathcal{O}_U$라 하자. $U' \subset X$가 이러한 두
번째 열린 부분스킴이라고 하고, 이에 대응하는 준연접 아이디얼층을
$\mathcal{I}' \subset \mathcal{O}_{U'}$라 쓰자. 그러면
$\mathcal{I}|_{U \cap U'} = \mathcal{I}'|_{U \cap U'}$이다. 또한
$\mathcal{O}_U/\mathcal{I}$의 지지집합은 $Z$이고, 이는 $U \cap U'$에
포함되며 $\mathcal{O}_{U'}/\mathcal{I}'$의 지지집합이기도 하다. 따라서
Definition \ref{divisors-definition-regular-ideal-sheaf}로부터 $\mathcal{I}$가 정칙
아이디얼일 필요충분조건은 $\mathcal{I}'$가 정칙 아이디얼인 것임을 알 수
있다. 코쥘 정칙, $H_1$-정칙 또는 준정칙인 경우도 마찬가지이다.

\begin{definition}
\label{divisors-definition-regular-immersion}
\begin{reference}
코쥘 정칙 몰입의 개념은 \cite[Expose VII, Definition 1.4]{SGA6}에서
도입되었으며, 거기서는 정칙 몰입이라 불렸다.
\end{reference}
$i : Z \to X$를 스킴의 몰입이라 하자. 열린 부분스킴 $U \subset X$를
택하여 $i$가 $Z$를 $U$의 닫힌 부분스킴과 동일시하게 하고, 이에
대응하는 준연접 아이디얼층을 $\mathcal{I} \subset \mathcal{O}_U$라 하자.
\begin{enumerate}
\item $i$가 {\it 정칙 몰입}이라는 것은 $\mathcal{I}$가 정칙인 것이다.
\item $i$가 {\it 코쥘 정칙 몰입}이라는 것은 $\mathcal{I}$가 코쥘
정칙인 것이다.
\item $i$가 {\it $H_1$-정칙 몰입}이라는 것은 $\mathcal{I}$가
$H_1$-정칙인 것이다.
\item $i$가 {\it 준정칙 몰입}이라는 것은 $\mathcal{I}$가 준정칙인 것이다.
\end{enumerate}
\end{definition}

\noindent
위의 논의로부터 이 정의가 $U$의 선택과 무관함을 알 수 있다. 조건들은
강한 것부터 약한 것의 순서로 나열되어 있다. Lemma
\ref{divisors-lemma-regular-quasi-regular-immersion}를 보라. 코쥘 정칙 닫힌 몰입은
매끄러운 위상에서 국소적으로 정칙 몰입이다. Lemma
\ref{divisors-lemma-koszul-regular-smooth-locally-regular}를 보라. 국소 뇌터인
경우에는 네 개념이 모두 일치한다. Lemma
\ref{divisors-lemma-Noetherian-scheme-regular-ideal}을 보라.

\begin{lemma}
\label{divisors-lemma-regular-quasi-regular-immersion}
$i : Z \to X$를 스킴의 몰입이라 하자. 다음 함의들이 성립한다.
$i$가 정칙 $\Rightarrow$
$i$가 코쥘 정칙 $\Rightarrow$
$i$가 $H_1$-정칙 $\Rightarrow$
$i$가 준정칙.
\end{lemma}

\begin{proof}
이는 바로 Lemma \ref{divisors-lemma-regular-quasi-regular-scheme}로 환원된다.
\end{proof}

\begin{lemma}
\label{divisors-lemma-regular-immersion-noetherian}
$i : Z \to X$를 스킴의 몰입이라 하고 $X$가 국소 뇌터라고 가정하자. 그러면
$i$가 정칙 $\Leftrightarrow$
$i$가 코쥘 정칙 $\Leftrightarrow$
$i$가 $H_1$-정칙 $\Leftrightarrow$
$i$가 준정칙.
\end{lemma}

\begin{proof}
이는 Lemma \ref{divisors-lemma-regular-quasi-regular-immersion}와 Lemma
\ref{divisors-lemma-Noetherian-scheme-regular-ideal}에서 바로 따른다.
\end{proof}

\begin{lemma}
\label{divisors-lemma-flat-base-change-regular-immersion}
$i : Z \to X$를 정칙(각각 코쥘 정칙, $H_1$-정칙, 준정칙) 몰입이라 하자.
$X' \to X$를 평탄 사상이라 하자. 그러면 밑변환
$i' : Z \times_X X' \to X'$는 정칙(각각 코쥘 정칙, $H_1$-정칙,
준정칙) 몰입이다.
\end{lemma}

\begin{proof}
Lemma \ref{divisors-lemma-regular-ideal-sheaf-scheme}를 통해 이는 Algebra, Lemmas
\ref{algebra-lemma-flat-increases-depth}와
\ref{algebra-lemma-flat-base-change-quasi-regular}, 그리고 More on Algebra,
Lemma \ref{more-algebra-lemma-koszul-regular-flat-base-change}의 대수적
명제로 옮겨진다.
\end{proof}

\begin{lemma}
\label{divisors-lemma-quasi-regular-immersion}
$i : Z \to X$를 스킴의 몰입이라 하자. 그러면 $i$가 준정칙 몰입일
필요충분조건은 다음 조건들이 성립하는 것이다.
\begin{enumerate}
\item $i$는 국소 유한 표시이고,
\item 여법선층 $\mathcal{C}_{Z/X}$는 유한 국소 자유이며,
\item 사상 (\ref{divisors-equation-conormal-algebra-quotient})은 동형이다.
\end{enumerate}
\end{lemma}

\begin{proof}
열린 몰입은 국소 유한 표시이다. 따라서 $X$를 열린 부분스킴
$U \subset X$로 바꾸어 $i$가 $Z$를 $U$의 닫힌 부분스킴과 동일시하게
할 수 있다. 즉, $i$가 닫힌 몰입이라고 가정해도 된다. 이에 대응하는 준연접
아이디얼층을 $\mathcal{I} \subset \mathcal{O}_X$라 하자. Morphisms,
Lemma \ref{morphisms-lemma-closed-immersion-finite-presentation}를 떠올리면
$\mathcal{I}$가 유한형일 필요충분조건은 $i$가 국소 유한 표시인 것이다.
따라서 동치는 Lemma \ref{divisors-lemma-quasi-regular-ideal}와 정의를 풀어 쓰는
것으로부터 따른다.
\end{proof}

\begin{lemma}
\label{divisors-lemma-transitivity-conormal-quasi-regular}
$Z \to Y \to X$를 스킴의 몰입들이라 하자. $Z \to Y$가
$H_1$-정칙이라고 가정하자. 그러면 Morphisms, Lemma
\ref{morphisms-lemma-transitivity-conormal}의 표준 열
$$
0 \to i^*\mathcal{C}_{Y/X} \to
\mathcal{C}_{Z/X} \to
\mathcal{C}_{Z/Y} \to 0
$$
은 완전하며 국소적으로 분할된다.
\end{lemma}

\begin{proof}
$\mathcal{C}_{Z/Y}$는 유한 국소 자유이므로(Lemma
\ref{divisors-lemma-quasi-regular-immersion}와 Lemma
\ref{divisors-lemma-regular-quasi-regular-scheme}를 보라), 이 열이 완전임을
보이면 충분하다. Morphisms, Lemma
\ref{morphisms-lemma-transitivity-conormal}에서 증명된 바에 따르면 첫
번째 사상이 단사임을 보이면 충분하다. 아핀 국소적으로 작업하면 이는 다음
문제로 환원된다. 환 $A$와 아이디얼들 $I \subset J \subset A$가 있다고
하자. $J/I \subset A/I$가 $H_1$-정칙열로 생성된다고 가정하자. 그러면
$I/I^2 \otimes_A A/J \to J/J^2$가 단사인가? 여기서
$I/I^2 \otimes_A A/J = I/IJ$임에 유의하자. 따라서 우리가 증명하려는 것은
$I \cap J^2 = IJ$이다. 이는 More on Algebra, Lemma
\ref{more-algebra-lemma-conormal-sequence-H1-regular}의 결과이다.
\end{proof}

\noindent
준정칙 몰입의 합성은 준정칙이 아닐 수 있다. Algebra, Remark
\ref{algebra-remark-join-quasi-regular-sequences}를 보라. 다른 종류의
정칙 몰입들은 합성에 대해 보존된다.

\begin{lemma}
\label{divisors-lemma-composition-regular-immersion}
$i : Z \to Y$와 $j : Y \to X$를 스킴의 몰입들이라 하자.
\begin{enumerate}
\item $i$와 $j$가 정칙 몰입이면 $j \circ i$도 정칙 몰입이다.
\item $i$와 $j$가 코쥘 정칙 몰입이면 $j \circ i$도 코쥘 정칙 몰입이다.
\item $i$와 $j$가 $H_1$-정칙 몰입이면 $j \circ i$도 그러하다.
\item $i$가 $H_1$-정칙 몰입이고 $j$가 준정칙 몰입이면 $j \circ i$는
준정칙 몰입이다.
\end{enumerate}
\end{lemma}

\begin{proof}
(1)의 대수적 형태는 Algebra, Lemma
\ref{algebra-lemma-join-regular-sequences}이다. (2)의 대수적 형태는
More on Algebra, Lemma
\ref{more-algebra-lemma-join-koszul-regular-sequences}이다. (3)의 대수적
형태는 More on Algebra, Lemma
\ref{more-algebra-lemma-join-H1-regular-sequences}이다. (4)의 대수적 형태는
More on Algebra, Lemma
\ref{more-algebra-lemma-join-quasi-regular-H1-regular}이다.
\end{proof}

\begin{lemma}
\label{divisors-lemma-permanence-regular-immersion}
$i : Z \to Y$와 $j : Y \to X$를 스킴의 몰입들이라 하자. $j$가 국소 유한
표시이고 Morphisms, Lemma
\ref{morphisms-lemma-transitivity-conormal}의 열
$$
0 \to i^*\mathcal{C}_{Y/X} \to
\mathcal{C}_{Z/X} \to
\mathcal{C}_{Z/Y} \to 0
$$
이 완전하며 국소적으로 분할된다고 가정하자.
\begin{enumerate}
\item $j \circ i$가 준정칙 몰입이면 $i$도 준정칙 몰입이다.
\item $j \circ i$가 $H_1$-정칙 몰입이면 $i$도 그러하다.
\item $j$와 $j \circ i$가 모두 코쥘 정칙 몰입이면 $i$도 코쥘 정칙
몰입이다.
\end{enumerate}
\end{lemma}

\begin{proof}
$Y$와 $X$를 줄여 $i$와 $j$가 닫힌 몰입이라고 가정해도 된다. 아이디얼층
$\mathcal{I} \subset \mathcal{O}_X$를 $Y$의 것으로, 아이디얼층
$\mathcal{J} \subset \mathcal{O}_X$를 $Z$의 것으로 쓰자. 여법선열은
$0 \to \mathcal{I}/\mathcal{I}\mathcal{J}
\to \mathcal{J}/\mathcal{J}^2 \to
\mathcal{J}/(\mathcal{I} + \mathcal{J}^2) \to 0$이다.
$z \in Z$라 하고 $y = i(z)$, $x = j(y) = j(i(z))$라 두자.
$f_1, \ldots, f_n \in \mathcal{I}_x$를 택하여 그 상들이
$\mathcal{I}_x/\mathfrak m_z\mathcal{I}_x$의 기저를 이루게 하자. 이를
$f_1, \ldots, f_n, g_1, \ldots, g_m \in \mathcal{J}_x$로 확장하여 그
상들이 $\mathcal{J}_x/\mathfrak m_z\mathcal{J}_x$의 기저를 이루게 하자.
$z$의 한 근방에서 여법선층들의 열이 분할된다고 가정했으므로 이것은
가능하다. 실제로
$\mathcal{I}_x/\mathfrak m_x\mathcal{I}_x \to
\mathcal{J}_x/\mathfrak m_x\mathcal{J}_x$가 단사이다.

\medskip\noindent
(1)의 증명. Lemma \ref{divisors-lemma-generate-regular-ideal}에 의해 아핀
열린 근방 $U$를 찾아 $x$를 포함하게 하고,
$f_1, \ldots, f_n, g_1, \ldots, g_m$이 $\mathcal{J}$를 생성하는 준정칙열을
이루게 할 수 있다. $i$가 국소 유한 표시이므로 추가로
$\mathcal{I}$가 $f_1, \ldots, f_n$으로 생성된다고 가정해도 된다
(Morphisms, Lemma
\ref{morphisms-lemma-closed-immersion-finite-presentation}, 나카야마
보조정리, 그리고 Modules, Lemma
\ref{modules-lemma-finite-type-surjective-on-stalk}를 사용하라). 따라서
Algebra, Lemma \ref{algebra-lemma-truncate-quasi-regular}에 의해
$g_1, \ldots, g_m$은 $Y \cap U$ 위에서 $Z$를 잘라내는 준정칙열을
유도한다.

\medskip\noindent
(2)의 증명. More on Algebra, Lemma
\ref{more-algebra-lemma-truncate-H1-regular}를 사용한다는 점만 제외하면
(1)의 증명과 완전히 같다.

\medskip\noindent
(3)의 증명. Lemma \ref{divisors-lemma-generate-regular-ideal}를 두 번 적용하면
아핀 열린 근방 $U$를 찾아 $x$를 포함하게 하고 $f_1, \ldots, f_n$이
$\mathcal{I}$를 생성하는 코쥘 정칙열을 이루고,
$f_1, \ldots, f_n, g_1, \ldots, g_m$이 $\mathcal{J}$를 생성하는 코쥘
정칙열을 이루게 할 수 있다. 따라서 More on Algebra, Lemma
\ref{more-algebra-lemma-truncate-koszul-regular}에 의해
$g_1, \ldots, g_m$은 $Y \cap U$ 위에서 $Z$를 잘라내는 코쥘 정칙열을
유도한다.
\end{proof}

\begin{lemma}
\label{divisors-lemma-extra-permanence-regular-immersion-noetherian}
$i : Z \to Y$와 $j : Y \to X$를 스킴의 몰입들이라 하자. $z \in Z$를
택하고 이에 대응하는 점들을 $y \in Y$, $x \in X$라 쓰자. $X$가 국소
뇌터라고 가정하자. 다음은 동치이다.
\begin{enumerate}
\item $i$는 $z$의 한 근방에서 정칙 몰입이고 $j$는 $y$의 한 근방에서
정칙 몰입이다.
\item $i$와 $j \circ i$는 $z$의 한 근방에서 정칙 몰입이다.
\item $j \circ i$는 $z$의 한 근방에서 정칙 몰입이고, 여법선열
$$
0 \to i^*\mathcal{C}_{Y/X} \to
\mathcal{C}_{Z/X} \to
\mathcal{C}_{Z/Y} \to 0
$$
은 $z$의 한 근방에서 분할 완전이다.
\end{enumerate}
\end{lemma}

\begin{proof}
$X$(따라서 $Y$도)가 국소 뇌터이므로 네 종류의 정칙 몰입은 모두
일치한다. 또한 Lemma \ref{divisors-lemma-Noetherian-scheme-regular-ideal}에 의해
사상이 정칙 몰입인지 여부를 국소환의 수준에서 확인할 수 있다. 함의 (1)
$\Rightarrow$ (2)는 Lemma \ref{divisors-lemma-composition-regular-immersion}이고,
함의 (2) $\Rightarrow$ (3)은 Lemma
\ref{divisors-lemma-transitivity-conormal-quasi-regular}이다. 따라서 (3)이 (1)을
함의함을 보이면 충분하다.

\medskip\noindent
(3)을 가정하자. $A = \mathcal{O}_{X, x}$라 두자. $I \subset A$를 전사
사상 $\mathcal{O}_{X, x} \to \mathcal{O}_{Y, y}$의 핵이라 쓰고,
$J \subset A$를 전사 사상
$\mathcal{O}_{X, x} \to \mathcal{O}_{Z, z}$의 핵이라 쓰자. $J$를 $A$
안에서 생성하는 임의의 극소 원소열은
준정칙열이고 따라서 정칙열이다. Lemma
\ref{divisors-lemma-generate-regular-ideal}을 보라. 가정에 의해 여법선열
$$
0 \to I/IJ \to J/J^2 \to J/(I + J^2) \to 0
$$
은 $A/J$-가군의 열로서 분할 완전이다. 따라서 극소 생성계
$f_1, \ldots, f_n, g_1, \ldots, g_m$을 $J$에 대해 택하되
$f_1, \ldots, f_n \in I$가 $I$의 극소 생성계가 되게 할 수 있다. 위에서
지적했듯이 $f_1, \ldots, f_n, g_1, \ldots, g_m$은 $A$ 안의 정칙열이다.
정칙열의 정의로부터 $f_1, \ldots, f_n$은 $A$ 안의 정칙열이고
$\overline{g}_1, \ldots, \overline{g}_m$은 $A/I$ 안의 정칙열임이 바로
따른다. 따라서 $j$는 $y$에서 정칙 몰입이고 $i$는 $z$에서 정칙 몰입이다.
\end{proof}

\begin{remark}
\label{divisors-remark-not-always-extra-permanence}
Lemma \ref{divisors-lemma-extra-permanence-regular-immersion-noetherian}의 상황에서
(1), (2), (3)은 “$j \circ i$와 $j$가 $z$와 $y$에서 정칙 몰입이다”와
{\bf 동치가 아니다}. 한 예는
$X = \mathbf{A}^1_k = \Spec(k[x])$,
$Y = \Spec(k[x]/(x^2))$, $Z = \Spec(k[x]/(x))$이다.
\end{remark}

\begin{lemma}
\label{divisors-lemma-koszul-regular-smooth-locally-regular}
$i : Z \to X$를 코쥘 정칙 닫힌 몰입이라 하자. 그러면 전사인 매끄러운 사상
$X' \to X$가 존재하여 $i' : Z \times_X X' \to X'$가 $i$의 밑변환이며
정칙 몰입이 된다.
\end{lemma}

\begin{proof}
$X$가 아핀이고 $Z$의 아이디얼이 코쥘 정칙열로 생성된다고 가정해도
된다. 이를 위해 $X$를 적절한 아핀 열린 덮개의 원소들로 바꾼다(Lemma
\ref{divisors-lemma-regular-ideal-sheaf-scheme}에서와 같은 아핀 열린집합을
취한다). 아핀인 경우는 More on Algebra, Lemma
\ref{more-algebra-lemma-Koszul-regular-flat-locally-regular}이다.
\end{proof}

\begin{lemma}
\label{divisors-lemma-immersion-regular-regular-immersion}
$i : Z \to X$를 몰입이라 하자. $Z$와 $X$가 정칙 스킴이면 $i$는 정칙
몰입이다.
\end{lemma}

\begin{proof}
$z \in Z$라 하자. Lemma \ref{divisors-lemma-Noetherian-scheme-regular-ideal}에 의해
$\mathcal{O}_{X, z} \to \mathcal{O}_{Z, z}$의 핵이 정칙열로 생성됨을
보이면 충분하다. 이는 Algebra, Lemmas
\ref{algebra-lemma-regular-quotient-regular}와
\ref{algebra-lemma-regular-ring-CM}에서 따른다.
\end{proof}

\section{상대 정칙 몰입}
\label{divisors-section-relative-regular-immersion}

\noindent
이 절에서는 정칙 몰입의 밑변환 성질을 고찰한다. 다음 보조정리는 정칙
몰입이나 코쥘 몰입에 대해서는 성립하지 않는다. Examples, Lemma
\ref{examples-lemma-base-change-regular-sequence}를 보라.

\begin{lemma}
\label{divisors-lemma-relative-regular-immersion}
$f : X \to S$를 스킴의 사상이라 하고 $i : Z \subset X$를 몰입이라 하자.
다음을 가정하자.
\begin{enumerate}
\item $i$는 $H_1$-정칙(각각 준정칙) 몰입이고,
\item $Z \to S$는 평탄 사상이다.
\end{enumerate}
그러면 스킴의 모든 사상 $g : S' \to S$에 대해 밑변환
$Z' = S' \times_S Z \to X' = S' \times_S X$는 $H_1$-정칙(각각 준정칙)
몰입이다.
\end{lemma}

\begin{proof}
정의를 풀어 쓰고 Lemma \ref{divisors-lemma-regular-ideal-sheaf-scheme}를 사용하면
이는 More on Algebra, Lemma
\ref{more-algebra-lemma-relative-regular-immersion-algebra}로 옮겨진다.
\end{proof}

\noindent
이 보조정리가 다음 정의의 동기이다.

\begin{definition}
\label{divisors-definition-relative-H1-regular-immersion}
$f : X \to S$를 스킴의 사상이라 하고 $i : Z \to X$를 몰입이라 하자.
\begin{enumerate}
\item $i$를 {\it 상대 준정칙 몰입}이라 하는 것은 $Z \to S$가 평탄이고
$i$가 준정칙 몰입인 것이다.
\item $i$를 {\it 상대 $H_1$-정칙 몰입}이라 하는 것은 $Z \to S$가
평탄이고 $i$가 $H_1$-정칙 몰입인 것이다.
\end{enumerate}
\end{definition}

\noindent
이 용어법이 표준적이지 않을 수 있음을 독자에게 경고한다. Lemma
\ref{divisors-lemma-relative-regular-immersion}은 상대 준정칙(각각 $H_1$-정칙)
몰입이 임의의 밑변환에 대해 보존됨을 보장한다. 상대 $H_1$-정칙 몰입은
상대 준정칙 몰입이다. Lemma
\ref{divisors-lemma-regular-quasi-regular-immersion}를 보라. Lemma
\ref{divisors-lemma-flat-relative-H1-regular}(또는 Lemma
\ref{divisors-lemma-relative-regular-immersion-flat-in-neighbourhood})를 살펴보자.
이는 $Z \to X$가 상대 $H_1$-정칙(또는 준정칙) 몰입이고 주위 스킴이
$S$ 위에서 (평탄이며) 국소 유한 표시이면, $Z \to X$가 실제로 정칙
몰입이고 이 사실이 모든 밑변환 뒤에도 유지됨을 보인다.

\begin{lemma}
\label{divisors-lemma-quasi-regular-immersion-flat-at-x}
$f : X \to S$를 스킴의 사상이라 하고 $Z \to X$를 상대 준정칙 몰입이라
하자. $x \in Z$이고 $\mathcal{O}_{X, x}$가 뇌터이면 $f$는 $x$에서
평탄이다.
\end{lemma}

\begin{proof}
준정칙열 $f_1, \ldots, f_r \in \mathcal{O}_{X, x}$를 택하여 $Z$의
아이디얼을 $x$에서 잘라내게 하자. Algebra, Lemma
\ref{algebra-lemma-quasi-regular-regular}에 의해
$f_1, \ldots, f_r$은 정칙열이다. 따라서 $f_r$은
$\mathcal{O}_{X, x}/(f_1, \ldots, f_{r - 1})$ 위의 영인자가 아닌
원소이고, 그 몫은 평탄 $\mathcal{O}_{S, f(x)}$-가군이다. Lemma
\ref{divisors-lemma-flat-at-x}에 의해
$\mathcal{O}_{X, x}/(f_1, \ldots, f_{r - 1})$가 평탄
$\mathcal{O}_{S, f(x)}$-가군임을 얻는다. 귀납을 계속하면
$\mathcal{O}_{X, x}$가 평탄 $\mathcal{O}_{S, s}$-가군임을 얻는다.
\end{proof}

\begin{lemma}
\label{divisors-lemma-relative-regular-immersion-flat-in-neighbourhood}
$X \to S$를 스킴의 사상이라 하고 $Z \to X$를 몰입이라 하자. 다음을
가정하자.
\begin{enumerate}
\item $X \to S$는 평탄이며 국소 유한 표시이고,
\item $Z \to X$는 상대 준정칙 몰입이다.
\end{enumerate}
그러면 $Z \to X$는 정칙 몰입이며, 임의의 밑변환 뒤에도 그러하다.
\end{lemma}

\begin{proof}
$x \in Z$를 택하고 그 상을 $s \in S$라 하자. 이를 증명하려면 $x$의
아핀 근방을 찾아 $U$ 안에 포함되게 하고 그 아핀 열린집합 위에서 결론이
성립함을 보이면 충분하다. 따라서 $X$가 아핀이고 준정칙열
$f_1, \ldots, f_r \in \Gamma(X, \mathcal{O}_X)$가 존재하여
$Z = V(f_1, \ldots, f_r)$라고 가정해도 된다. More on Algebra, Lemma
\ref{more-algebra-lemma-relative-regular-immersion-algebra}에 의해 열
$f_1|_{X_s}, \ldots, f_r|_{X_s}$는
$\Gamma(X_s, \mathcal{O}_{X_s})$ 안의 준정칙열이다. $X_s$가 뇌터이므로,
필요하면 $X$를 조금 줄인 뒤 $f_1|_{X_s}, \ldots, f_r|_{X_s}$가 정칙열임을
얻는다. Algebra, Lemmas \ref{algebra-lemma-quasi-regular-regular}와
\ref{algebra-lemma-regular-sequence-in-neighbourhood}를 보라. Lemma
\ref{divisors-lemma-fibre-Cartier}에 의해 $Z_1 = V(f_1) \subset X$는 상대 유효
카르티에 인자이다. 여기서도 필요하면 $X$를 조금 줄인다. 같은 보조정리를
$Z_2 = V(f_1, f_2) \subset Z_1$에 다시 적용하면 $Z_2 \subset Z_1$이
상대 유효 카르티에 인자임을 알 수 있다. 이 과정을 계속하여
$Z = Z_n = V(f_1, \ldots, f_n)$에 이른다. 상대 유효 카르티에 인자라는
성질은 임의의 밑변환에 대해 보존되므로(Lemma
\ref{divisors-lemma-relative-Cartier}를 보라), 보조정리의 마지막 명제도 따른다.
\end{proof}

\begin{remark}
\label{divisors-remark-relative-regular-immersion-elements}
상대 준정칙 몰입의 여차원이 상수이면, 밑변환 뒤에도 바뀌지 않는다.
실제로 환 사상 $A \to B$와 준정칙열
$f_1, \ldots, f_r \in B$가 주어지고 $B/(f_1, \ldots, f_r)$가 $A$ 위에서
평탄하다고 하자. 그러면 모든 환 사상 $A \to A'$에 대해
$f_1 \otimes 1, \ldots, f_r \otimes 1$은
$B' = B \otimes_A A'$ 안의 준정칙열이다. 이는 More on Algebra, Lemma
\ref{more-algebra-lemma-relative-regular-immersion-algebra}로부터 따른다
(이 보조정리는 위의 Lemma \ref{divisors-lemma-relative-regular-immersion}의
증명에서 사용되었다). 이제 Lemma
\ref{divisors-lemma-relative-regular-immersion-flat-in-neighbourhood}의 증명은
$A \to B$가 평탄이며 국소 유한 표시이면 모든 소 아이디얼
$\mathfrak q' \subset B'$에 대해 열
$f_1 \otimes 1, \ldots, f_r \otimes 1$이 국소환
$B'_{\mathfrak q'}$ 안에서 실제로 정칙열임을 보인다.
\end{remark}

\begin{lemma}
\label{divisors-lemma-flat-relative-H1-regular}
$X \to S$를 스킴의 사상이라 하고 $Z \to X$를 상대 $H_1$-정칙 몰입이라
하자. $X \to S$가 국소 유한 표시라고 가정하자. 그러면
\begin{enumerate}
\item 열린 부분스킴 $U \subset X$가 존재하여 $Z \subset U$이고
$U \to S$가 평탄이며,
\item $Z \to X$는 정칙 몰입이고 이는 임의의 밑변환 뒤에도 성립한다.
\end{enumerate}
\end{lemma}

\begin{proof}
$x \in Z$를 택하자. (1)을 증명하려면 열린 근방 $U \subset X$를 찾아
$x$를 포함하게 하고 $U \to S$가 평탄임을 보이면 충분하다. 따라서
$X = \Spec(B)$와
$S = \Spec(A)$가 아핀이고 $Z$가 $H_1$-정칙열
$f_1, \ldots, f_r \in B$로 주어진 경우로 환원된다. 가정에 의해 $B$는
유한 표시 $A$-대수이고 $B/(f_1, \ldots, f_r)B$는 평탄 $A$-대수이다.
절대 뇌터 근사를 사용할 것이다.

\medskip\noindent
$B = A[x_1, \ldots, x_n]/(g_1, \ldots, g_m)$로 쓰자. $f_i$가
$f_i' \in A[x_1, \ldots, x_n]$의 상이라고 가정하자. 유한형
$\mathbf{Z}$-부분대수 $A_0 \subset A$를 택하여 다항식
$f_1', \ldots, f_r', g_1, \ldots, g_m$의 모든 계수가 $A_0$에 속하게
하자. $B_0 = A_0[x_1, \ldots, x_n]/(g_1, \ldots, g_m)$라 두고 그 상을
$f_{i, 0}$라 쓰자. 여기서 이는 $f_i'$의 $B_0$에서의 상이다. 그러면
$B = B_0 \otimes_{A_0} A$이고
$$
B/(f_1, \ldots, f_r) =
B_0/(f_{0, 1}, \ldots, f_{0, r}) \otimes_{A_0} A.
$$
Algebra, Lemma \ref{algebra-lemma-flat-finite-presentation-limit-flat}에
따라 $A_0$를 확대하여 $B_0/(f_{0, 1}, \ldots, f_{0, r})$가 $A_0$ 위에서
평탄이라고 가정해도 된다. 이 단계에서 코쥘 코호몰로지 군
$H_1(K_\bullet(B_0, f_{0, 1}, \ldots, f_{0, r}))$이 영일 필요는 없다.
한편 $B_0$가 뇌터이므로 이는 유한 생성 $B_0$-가군이다. 생성원들을
$\xi_1, \ldots, \xi_n \in H_1(K_\bullet(B_0, f_{0, 1}, \ldots, f_{0, r}))$
라 하자. $A_0 \subset A_1 \subset A$를 더 큰 유한형
$\mathbf{Z}$-부분대수라 하되 이는 $A$의 부분대수이다. 그 상을
$f_{1, i}$라 쓰자. 여기서 이는 $f_{0, i}$의
$B_1 = B_0 \otimes_{A_0} A_1$에서의 상이다. More on
Algebra, Lemma \ref{more-algebra-lemma-base-change-H1-regular}에 의해 사상
$$
H_1(K_\bullet(B_0, f_{0, 1}, \ldots, f_{0, r})) \otimes_{A_0} A_1
\longrightarrow
H_1(K_\bullet(B_1, f_{1, 1}, \ldots, f_{1, r}))
$$
은 전사이다. 또한 위와 같은 모든 $A_1$의 선택에 대한 복합체
$K_\bullet(B_1, f_{1, 1}, \ldots, f_{1, r})$의 여극한은 복합체
$K_\bullet(B, f_1, \ldots, f_r)$이고 이는 차수 $1$에서 비순환임이
분명하다. 따라서
Algebra, Lemma \ref{algebra-lemma-directed-colimit-exact}에 의해
$$
\colim_{A_0 \subset A_1 \subset A}
H_1(K_\bullet(B_1, f_{1, 1}, \ldots, f_{1, r}))
= 0
$$
이다. 그러므로 $A_1$을 택하여 $\xi_1, \ldots, \xi_n$이 모두
$H_1(K_\bullet(B_1, f_{1, 1}, \ldots, f_{1, r}))$에서 영으로 가게 할 수
있다. 다시 말해 코쥘 코호몰로지 군
$H_1(K_\bullet(B_1, f_{1, 1}, \ldots, f_{1, r}))$은 영이다.

\medskip\noindent
아핀 스킴의 사상 $X_1 \to S_1$을 고찰하자. 이는 $\Spec$을 환 사상
$A_1 \to B_1$에 취한 것이다. 또한
$Z_1 = \Spec(B_1/(f_{1, 1}, \ldots, f_{1, r}))$라 하자.
$B = B_1 \otimes_{A_1} A$, 즉 $X = X_1 \times_{S_1} S$이고 마찬가지로
$Z = Z_1 \times_S S_1$이므로, 이제 $X_1 \to S_1$과 상대
$H_1$-정칙 몰입 $Z_1 \to X_1$에 대해 (1)을 증명하면 충분하다.
Morphisms, Lemma \ref{morphisms-lemma-base-change-module-flat}를 보라.
따라서 $X \to S$가 뇌터 스킴들의 유한형 사상인 경우로 환원되었다.
이 경우 Lemma \ref{divisors-lemma-quasi-regular-immersion-flat-at-x}에 의해
$X \to S$는 $Z$의 모든 점에서 평탄하다. 이 경우 평탄 자취가
열린집합이라는 사실(Algebra, Theorem
\ref{algebra-theorem-openness-flatness}를 보라)과 결합하면 (1)을 얻는다.
이제 (2)는 Lemma
\ref{divisors-lemma-relative-regular-immersion-flat-in-neighbourhood}를 적용하여
따른다.
\end{proof}

\noindent
주위 스킴이 밑스킴 위에서 평탄이며 국소 유한 표시이면, 상대 준정칙 몰입을
그 올들로 특징지을 수 있다.

\begin{lemma}
\label{divisors-lemma-fibre-quasi-regular}
$\varphi : X \to S$를 국소 유한 표시인 평탄 사상이라 하자.
$T \subset X$를 닫힌 부분스킴이라 하고, $x \in T$의 상을 $s \in S$라 하자.
\begin{enumerate}
\item $T_s \subset X_s$가 $x$의 한 근방에서 준정칙 몰입이면, 열린집합
$U \subset X$와 상대 준정칙 몰입 $Z \subset U$가 존재하여
$Z_s = T_s \cap U_s$이고 $T \cap U \subset Z$이다.
\item $T_s \subset X_s$가 $x$의 한 근방에서 준정칙 몰입이고, 사상
$T \to X$가 유한 표시이며 $T \to S$가 $x$에서 평탄이면, (1)에서와 같이
$U$와 $Z$를 택하여 $T \cap U = Z$가 되게 할 수 있다.
\item $T_s \subset X_s$가 $x$의 한 근방에서 준정칙 몰입이고, $T$가
$c$개의 방정식으로 잘려 나오며 이는 $x$의 한 근방에서 성립하고,
$c = \dim_x(X_s) - \dim_x(T_s)$이면, (1)에서와 같이 $U$와 $Z$를 택하여
$T \cap U = Z$가 되게 할 수 있다.
\end{enumerate}
각 경우에 $Z \to U$는 Lemma
\ref{divisors-lemma-relative-regular-immersion-flat-in-neighbourhood}에 의해 정칙
몰입이다. 특히 $T \to S$가 국소 유한 표시이며 평탄이고 모든 올
$T_s \subset X_s$가 준정칙 몰입이면, $T \to X$는 상대 준정칙 몰입이다.
\end{lemma}

\begin{proof}
아핀 열린 근방 $\Spec(A)$를 택하여 $s$를 포함하게 하고, 아핀 열린 근방
$\Spec(B)$를 택하여 $x$를 포함하게 하며
$\varphi(\Spec(B)) \subset \Spec(A)$가 되게 하자. 소 아이디얼
$\mathfrak p \subset A$를 $s$에 대응하는 것으로, 소 아이디얼
$\mathfrak q \subset B$를 $x$에 대응하는 것으로, 아이디얼
$I \subset B$를 $T$에 대응하는 것으로 택하자. 보조정리의 처음 가정으로부터 $A \to B$가
평탄이며 유한 표시임을 안다. (1)의 가정은 $\Spec(B)$를 줄인 뒤
$I(B \otimes_A \kappa(\mathfrak p))$가 어떤 준정칙 원소열로 생성된다고
가정해도 된다는 뜻이다. 필요하면 $B$를 어떤
$g \in B$, $g \not \in \mathfrak q$에서 국소화하여, 원소들
$f_1, \ldots, f_r \in I$이 존재하고 그 상들이
$B \otimes_A \kappa(\mathfrak p)$ 안에서 준정칙열을 이루며
$I(B \otimes_A \kappa(\mathfrak p))$를 생성한다고 가정해도 된다.
Algebra, Lemmas \ref{algebra-lemma-quasi-regular-regular}와
\ref{algebra-lemma-regular-sequence-in-neighbourhood}에 의해 한 번 더
국소화하여 $f_1, \ldots, f_r \in I$이
$B \otimes_A \kappa(\mathfrak p)$ 안에서 정칙열을 이룬다고 가정해도
된다. Lemma \ref{divisors-lemma-fibre-Cartier}에 의해
$Z_1 = V(f_1) \subset \Spec(B)$는 상대 유효 카르티에 인자이다. 여기서도
필요하면 $B$를 국소화한다. 같은 보조정리를
$Z_2 = V(f_1, f_2) \subset Z_1$에 다시 적용하면 $Z_2 \subset Z_1$이
상대 유효 카르티에 인자임을 알 수 있다. 이 과정을 계속하여
$Z = Z_n = V(f_1, \ldots, f_n)$에 이른다. 그러면
$Z \to \Spec(B)$는 정칙 몰입이고 $Z$는 $S$ 위에서 평탄하다. 특히
$Z \to \Spec(B)$는 $\Spec(A)$ 위의 상대 준정칙 몰입이다. 이로써 (1)이
증명되었다.

\medskip\noindent
(2)를 보이기 위해 닫힌 몰입 $Z \to D$를 고찰하자. 전사 환 사상
$u : \mathcal{O}_{D, x} \to \mathcal{O}_{Z, x}$는 본질적 유한 표시인 평탄
국소 $\mathcal{O}_{S, s}$-대수들 사이의 사상이고, $\mathfrak m_s$로
나누면 동형이 된다. 따라서 이는 동형이다. Algebra, Lemma
\ref{algebra-lemma-mod-injective-general}을 보라. 이에 따라 $Z \to D$는
$x$의 한 근방에서 동형이다. Algebra, Lemma
\ref{algebra-lemma-local-isomorphism}을 보라.

\medskip\noindent
(3)을 보자. 필요하면 $U$를 줄여 $Z$의 아이디얼이 정칙열
$f_1, \ldots, f_r$로 생성되고(위의 $Z$의 구성을 보라), $T$의 아이디얼이
$g_1, \ldots, g_c$로 생성된다고 가정해도 된다. $c = r$임을 주장한다.
실제로
\begin{align*}
\dim_x(X_s) & = \dim(\mathcal{O}_{X_s, x}) +
\text{trdeg}_{\kappa(s)}(\kappa(x)), \\
\dim_x(T_s) & = \dim(\mathcal{O}_{T_s, x}) +
\text{trdeg}_{\kappa(s)}(\kappa(x)), \\
\dim(\mathcal{O}_{X_s, x}) & = \dim(\mathcal{O}_{T_s, x}) + r
\end{align*}
이다. 앞의 두 등식은 Algebra, Lemma
\ref{algebra-lemma-dimension-at-a-point-finite-type-field}에서 따르고,
두 번째 등식은 Algebra, Lemma
\ref{algebra-lemma-one-equation}을 $r$번 적용하여 따른다.
$T \subset Z$이므로 $f_i = \sum b_{ij} g_j$임을 안다. 그러나 $Z$와 $T$의
아이디얼들은 $X_s$의 같은 준정칙 닫힌 부분스킴을 잘라내며 이는 $x$의 한
근방에서 성립한다. 따라서 행렬 $(b_{ij}) \bmod \mathfrak m_x$는 가역이다
(일부 세부사항은 생략한다). 그러므로 $(b_{ij})$는 $x$의 한 열린 근방에서
가역이다. 다시 말해 $T \cap U = Z$이다. 여기서 $U$를 줄인다.

\medskip\noindent
보조정리의 마지막 명제들은 (2)와, $Z \to S$가 국소 유한 표시일
필요충분조건이 $Z \to X$가 유한 표시인 것이라는 사실을 결합하여 바로
따른다. Morphisms, Lemmas
\ref{morphisms-lemma-composition-finite-presentation}와
\ref{morphisms-lemma-finite-presentation-permanence}를 보라.
\end{proof}

\noindent
다음 보조정리는 Morphisms, Lemma
\ref{morphisms-lemma-section-smooth-morphism}을 강화한 것이다.

\begin{lemma}
\label{divisors-lemma-section-smooth-regular-immersion}
$f : X \to S$를 스킴의 매끄러운 사상이라 하고
$\sigma : S \to X$를 $f$의 절단이라 하자. 그러면 $\sigma$는 정칙
몰입이다.
\end{lemma}

\begin{proof}
Schemes, Lemma \ref{schemes-lemma-semi-diagonal}에 의해 사상 $\sigma$는
몰입이다. $X$를 $\sigma(S)$의 열린 근방으로 바꾸어 $\sigma$가 닫힌
몰입이라고 가정해도 된다. 이에 대응하는 닫힌 부분스킴을
$T = \sigma(S)$라 하자. 이는 $X$의 부분스킴이다. $T \to S$는
동형이므로 평탄하며 유한 표시이다.
또한 매끄러운 사상은 평탄하며 국소 유한 표시이다. Morphisms, Lemmas
\ref{morphisms-lemma-smooth-flat}와
\ref{morphisms-lemma-smooth-locally-finite-presentation}를 보라. 따라서
Lemma \ref{divisors-lemma-fibre-quasi-regular}에 의해 $T_s \subset X_s$가 준정칙
닫힌 부분스킴임을 보이면 충분하다. 이는 Morphisms, Lemma
\ref{morphisms-lemma-section-smooth-morphism}에서 바로 따르지만, 다음과
같이 직접 볼 수도 있다. $k$를 체라 하고 $A$를 매끄러운 $k$-대수라 하자.
$\mathfrak m \subset A$를 잉여체가 $k$인 극대 아이디얼이라 하자. 그러면
$\mathfrak m$은 준정칙열로 생성된다. 필요하면 $A$를 $A_g$로 바꾸는데,
여기서 어떤 $g \in A$, $g \not \in \mathfrak m$를 취한다. Algebra,
Lemma \ref{algebra-lemma-characterize-smooth-over-field}에서
$A_{\mathfrak m}$이 정칙 국소환임을 증명했으므로
$\mathfrak mA_{\mathfrak m}$은 정칙열로 생성된다. 이는 실제로
$\mathfrak m$이 정칙열로 생성됨을 함의한다. 여기서 필요하면 $A$를
$A_g$로 바꾸는데, 어떤 $g \in A$, $g \not \in \mathfrak m$를 취한다. Algebra,
Lemma \ref{algebra-lemma-regular-sequence-in-neighbourhood}를 보라.
\end{proof}

\noindent
다음 보조정리에는 일종의 역이 있다. Lemma
\ref{divisors-lemma-push-regular-immersion-thru-smooth}를 보라.

\begin{lemma}
\label{divisors-lemma-lift-regular-immersion-to-smooth}
스킴의 사상들의 가환 그림
$$
\xymatrix{
Y \ar[rd]_j \ar[rr]_i & & X \ar[ld] \\
& S
}
$$
이 주어졌다고 하자. $X \to S$가 매끄럽고 $i$, $j$가 몰입이라고
가정하자. $j$가 정칙(각각 코쥘 정칙, $H_1$-정칙, 준정칙) 몰입이면
$i$도 그러하다.
\end{lemma}

\begin{proof}
$i$를 합성
$$
Y \to Y \times_S X \to X
$$
으로 쓸 수 있다. Lemma \ref{divisors-lemma-section-smooth-regular-immersion}에 의해
첫 번째 화살표는 정칙 몰입이다. 두 번째 화살표는 $Y \to S$의 평탄
밑변환이므로 정칙(각각 코쥘 정칙, $H_1$-정칙, 준정칙) 몰입이다. Lemma
\ref{divisors-lemma-flat-base-change-regular-immersion}를 보라. 이제 Lemma
\ref{divisors-lemma-composition-regular-immersion}을 적용하면 결론을 얻는다.
\end{proof}

\begin{lemma}
\label{divisors-lemma-immersion-lci-into-smooth-regular-immersion}
스킴의 사상들의 가환 그림
$$
\xymatrix{
Y \ar[rd] \ar[rr]_i & & X \ar[ld] \\
& S
}
$$
이 주어졌다고 하자. $Y \to S$가 신토믹이고 $X \to S$가 매끄러우며
$i$가 몰입이라고 가정하자. 그러면 $i$는 정칙 몰입이다.
\end{lemma}

\begin{proof}
$X$를 $i(Y)$의 열린 근방으로 바꾸어 $i$가 닫힌 몰입이라고 가정해도
된다. 이에 대응하는 닫힌 부분스킴을 $T = i(Y)$라 하자. 이는 $X$의
부분스킴이다.
$T \cong Y$이므로 사상 $T \to S$는 평탄이며 유한 표시이다(Morphisms,
Lemmas \ref{morphisms-lemma-syntomic-locally-finite-presentation}와
\ref{morphisms-lemma-syntomic-flat}). 또한 매끄러운 사상은 평탄이며
국소 유한 표시이다(Morphisms, Lemmas
\ref{morphisms-lemma-smooth-flat}와
\ref{morphisms-lemma-smooth-locally-finite-presentation}). 따라서 Lemma
\ref{divisors-lemma-fibre-quasi-regular}에 의해 $T_s \subset X_s$가 준정칙 닫힌
부분스킴임을 보이면 충분하다. $X_s$는 체 위에서 국소 유한형이므로
뇌터이다(Morphisms, Lemma
\ref{morphisms-lemma-finite-type-noetherian}). 따라서 점별로
$T_s \subset X_s$가 준정칙 몰입인지 확인할 수 있다. Lemma
\ref{divisors-lemma-Noetherian-scheme-regular-ideal}을 보라. $t \in T_s$를 택하자.
Morphisms, Lemma
\ref{morphisms-lemma-local-complete-intersection}에 의해 국소환
$\mathcal{O}_{T_s, t}$는 $\kappa(s)$ 위의 국소 완전 교차이다. 국소환
$\mathcal{O}_{X_s, t}$는 정칙이다. Algebra, Lemma
\ref{algebra-lemma-characterize-smooth-over-field}를 보라. Algebra, Lemma
\ref{algebra-lemma-lci-local}에 의해 전사
$\mathcal{O}_{X_s, t} \to \mathcal{O}_{T_s, t}$의 핵이 정칙열로
생성됨을 알 수 있고, 이것이 증명할 내용이었다.
\end{proof}

\begin{lemma}
\label{divisors-lemma-immersion-smooth-into-smooth-regular-immersion}
스킴의 사상들의 가환 그림
$$
\xymatrix{
Y \ar[rd] \ar[rr]_i & & X \ar[ld] \\
& S
}
$$
이 주어졌다고 하자. $Y \to S$와 $X \to S$가 매끄럽고 $i$가
몰입이라고 가정하자. 그러면 $i$는 정칙 몰입이다.
\end{lemma}

\begin{proof}
이는 Lemma
\ref{divisors-lemma-immersion-lci-into-smooth-regular-immersion}의 특수한 경우이다.
실제로 매끄러운 사상은 신토믹이다. Morphisms, Lemma
\ref{morphisms-lemma-smooth-syntomic}을 보라.
\end{proof}

\begin{lemma}
\label{divisors-lemma-push-regular-immersion-thru-smooth}
스킴의 사상들의 가환 그림
$$
\xymatrix{
Y \ar[rd]_j \ar[rr]_i & & X \ar[ld] \\
& S
}
$$
이 주어졌다고 하자. $X \to S$가 매끄럽고 $i$와 $j$가 몰입이라고
가정하자. $i$가 코쥘 정칙(각각 $H_1$-정칙, 준정칙) 몰입이면 $j$도
그러하다.
\end{lemma}

\begin{proof}
이하에서 Lemma \ref{divisors-lemma-regular-quasi-regular-immersion}을 별도의 언급
없이 사용한다. 임의의 점 $y \in Y$를 택하고 $x = i(y)$,
$s = j(y)$라 두자. $X$와 $S$를 열린 근방들 $U$와 $V$로 바꾸되 각각
$x$와 $s$를 포함하게 하고, $Y$를 $y$의 열린 근방으로 바꾸되
$i^{-1}(U) \cap j^{-1}(V)$ 안에 놓이게 한 뒤 결론을 증명하면 충분하다.

\medskip\noindent
먼저 $X = \mathbf{A}^n_S$인 경우를 증명한다. $S$를 아핀 열린집합
$V$로, $Y$를 $j^{-1}(V)$로 바꾸어 $j$가 닫힌 몰입이고 $S$가 아핀이라고
가정해도 된다. $S = \Spec(A)$로 쓰자. 그러면
$j : Y \to S$는 $Y$와 닫힌 부분스킴 $\Spec(A/I)$ 사이의 동형을
정하며, 여기서 $I \subset A$는 어떤 아이디얼이다. 사상
$i : Y = \Spec(A/I) \to
\mathbf{A}^n_S = \Spec(A[x_1, \ldots, x_n])$는 $A$-대수 준동형
$i^\sharp : A[x_1, \ldots, x_n] \to A/I$에 대응한다.
$a_i \in A$를 택하여 그 상이 $i^\sharp(x_i)$가 되게 하자. 이 상은
$A/I$에서 취한다. 닫힌 몰입
$i$의 아이디얼은
$$
J = (x_1 - a_1, \ldots, x_n - a_n) + IA[x_1, \ldots, x_n].
$$
임을 관찰하자. $K = (x_1 - a_1, \ldots, x_n - a_n)$라 두자. 열
$$
0 \to K/KJ \to J/J^2 \to J/(K + J^2) \to 0
$$
이 분할 완전임을 주장한다. 이를 보려면 $K/K^2$가 기저
$x_i - a_i$를 갖는 자유 가군이고 그 바탕환이
$A[x_1, \ldots, x_n]/K \cong A$임에 유의하면 된다. 따라서 $K/KJ$는 환
$A[x_1, \ldots, x_n]/J \cong A/I$ 위에서 같은 기저를 갖는 자유
가군이다. 한편 미분을 취하면 사상
$$
\text{d}_{A[x_1, \ldots, x_n]/A} :
J/J^2
\longrightarrow
\Omega_{A[x_1, \ldots, x_n]/A} \otimes_{A[x_1, \ldots, x_n]}
A[x_1, \ldots, x_n]/J
$$
을 얻고, 이 사상은 생성원 $x_i - a_i$를 오른쪽 자유 가군의 기저 원소
$\text{d}x_i$로 보낸다. 따라서 주장이 따른다. 또한
$x_1 - a_1, \ldots, x_n - a_n$은
$A[x_1, \ldots, x_n]$ 안의 정칙열이고 그 몫환은
$A[x_1, \ldots, x_n]/(x_1 - a_1, \ldots, x_n - a_n) \cong A$이다.
따라서 닫힌 몰입 $i$는 다음과 같이 분해된다.
$$
Y \to V(x_1 - a_1, \ldots, x_n - a_n) \to \mathbf{A}^n_S
$$
여기서 합성은 코쥘 정칙(각각 $H_1$-정칙, 준정칙)이고, 두 번째 화살표는
정칙 몰입이며, 이에 딸린 여법선열은 분할된다. 이제 Lemma
\ref{divisors-lemma-permanence-regular-immersion}으로부터 결론이 따른다.

\medskip\noindent
다음으로 $i$가 $H_1$-정칙이면 결론이 성립함을 증명한다. 증명의 첫
문단에서처럼 줄여 $Y$, $X$, $S$가 아핀이라고 가정해도 된다. 이 경우
닫힌 몰입 $h : X \to \mathbf{A}^n_S$를 $S$ 위에서 택할 수 있으며,
이는 어떤 $n$에 대해 가능하다. Lemma
\ref{divisors-lemma-immersion-smooth-into-smooth-regular-immersion}에 의해 $h$는
정칙 몰입이다. 따라서 $h \circ i$는 $H_1$-정칙 몰입이다. Lemma
\ref{divisors-lemma-composition-regular-immersion}을 보라(이 단계는 “준정칙인
경우”에는 작동하지 않음에 유의하라). 이로써 위에서 증명한
$X = \mathbf{A}^n_S$이고 $S$가 아핀인 경우로 환원된다.

\medskip\noindent
끝으로 $i$가 준정칙이라고 가정하자. 증명의 첫 문단에서처럼 줄인 뒤
Morphisms, Lemma
\ref{morphisms-lemma-smooth-etale-over-affine-space}를 사용하여 $f$를
$X \to \mathbf{A}^n_S \to S$로 분해할 수 있다. 여기서 첫 번째 사상
$X \to \mathbf{A}^n_S$는 에탈이다. 따라서 문제는 (a)
$X = \mathbf{A}^n_S$인 경우와 (b) $f$가 에탈인 경우로 환원된다.
경우 (a)는 증명의 두 번째 문단에서 다루었다. 경우 (b)는 다음 문단에서
다룬다.

\medskip\noindent
$f$가 에탈이라고 가정하자. 줄인 뒤 $X$, $Y$, $S$가 아핀이고 $i$와
$j$가 닫힌 몰입이라고 가정해도 된다(작은 세부사항은 생략한다).
$S = \Spec(A)$, $X = \Spec(B)$,
$Y = \Spec(B/J) = \Spec(A/I)$라 하자. 더 줄여 $J$가 준정칙열로
생성된다고 가정해도 된다. 환 사상 $A \to B$는 에탈이므로 형식적
에탈이다(Algebra, Lemma \ref{algebra-lemma-formally-etale-etale}).
따라서 Algebra, Lemma
\ref{algebra-lemma-formally-etale-lift-infinitesimal}에 의해
$\bigoplus I^n/I^{n + 1} \cong \bigoplus J^n/J^{n + 1}$이다. $J$가
준정칙열로 생성되므로 $I$도 그러하다. 이로써 증명이 끝난다.
\end{proof}

\section{유리형 함수와 절단}
\label{divisors-section-meromorphic-functions}

\noindent
이 절에서는 일반적인 정의와 몇 가지 기초적인 결과만 다룬다.
생길 수 있는 몇 가지 함정에 대해서는
\cite{misconceptions}를 보라\footnote{위험해, 윌 로빈슨!}.

\medskip\noindent
$(X, \mathcal{O}_X)$를 국소환 달린 공간이라 하자.
임의의 열린집합 $U \subset X$에 대하여
$\mathcal{S}(U) \subset \mathcal{O}_X(U)$를
$\mathcal{O}_X$의 $U$ 위 정칙 절단들의 집합으로 정의했으며,
Definition \ref{divisors-definition-regular-section}을 보라. 정칙 절단을 더 작은
열린집합에 제한해도 정칙이다. 따라서
$\mathcal{S} : U \mapsto \mathcal{S}(U)$는 $\mathcal{O}_X$의
(집합들의) 부분층이다. 때때로 $\mathcal{S} = \mathcal{S}_X$로
쓰며, 이는 $X$에 대한 의존성을 나타내기 위해서이다.
더욱이 $\mathcal{S}(U)$는 환 $\mathcal{O}_X(U)$의 곱셈적
부분집합이며, 이는 각 $U$에 대해 성립한다. 따라서 환의 준층
$$
U \longmapsto \mathcal{S}(U)^{-1} \mathcal{O}_X(U),
$$
을 생각할 수 있으며, Modules, Lemma
\ref{modules-lemma-simple-invert}를 보라.

\begin{definition}
\label{divisors-definition-sheaf-meromorphic-functions}
$(X, \mathcal{O}_X)$를 국소환 달린 공간이라 하자. 위에 표시한
준층에 수반되는 층을 $X$ 위의 {\it 유리형 함수층}이라 하고
{\it $\mathcal{K}_X$}로 쓴다. $X$ 위의 {\it 유리형 함수}란
$\mathcal{K}_X$의 대역 절단을 말한다.
\end{definition}

\noindent
각 $\mathcal{S}(U)$의 모든 원소는 $\mathcal{O}_X(U)$의
영인자가 아닌 원소이므로, 환의 층의 자연스러운 사상
$\mathcal{O}_X \to \mathcal{K}_X$는 단사이다.

\begin{example}
\label{divisors-example-no-change}
$A = \mathbf{C}[x, \{y_\alpha\}_{\alpha \in \mathbf{C}}]/
((x - \alpha)y_\alpha, y_\alpha y_\beta)$라 하자. $A$의 모든 원소는
$f(x) + \sum \lambda_\alpha y_\alpha$의 꼴로 유일하게 쓸 수 있다.
여기서 $f(x) \in \mathbf{C}[x]$이고
$\lambda_\alpha \in \mathbf{C}$이다.
$X = \Spec(A)$라 하자. 이 경우 $\mathcal{O}_X = \mathcal{K}_X$이다.
실제로 모든 아핀 열린집합 $D(f)$에서 환 $A_f$의 영인자가 아닌
모든 원소는 가역원이다(증명은 생략한다).
\end{example}

\noindent
$(X, \mathcal{O}_X)$를 국소환 달린 공간이라 하고
$\mathcal{F}$를 $\mathcal{O}_X$-가군의 층이라 하자.
준층 $U \mapsto \mathcal{S}(U)^{-1}\mathcal{F}(U)$를 생각하자.
그 층화는 층 $\mathcal{F} \otimes_{\mathcal{O}_X} \mathcal{K}_X$이며,
Modules, Lemma \ref{modules-lemma-simple-invert-module}를 보라.

\begin{definition}
\label{divisors-definition-meromorphic-section}
$X$를 국소환 달린 공간이라 하자.
$\mathcal{F}$를 $\mathcal{O}_X$-가군의 층이라 하자.
\begin{enumerate}
\item $\mathcal{K}_X(\mathcal{F})$로 쓰는 층은
$\mathcal{K}_X$-가군의 층이며, 준층
$U \mapsto \mathcal{S}(U)^{-1}\mathcal{F}(U)$의 층화이다. 동치로
$\mathcal{K}_X(\mathcal{F}) =
\mathcal{F} \otimes_{\mathcal{O}_X} \mathcal{K}_X$이다(위를 보라).
\item $\mathcal{F}$의 {\it 유리형 절단}이란
$\mathcal{K}_X(\mathcal{F})$의 대역 절단을 말한다.
\end{enumerate}
\end{definition}

\noindent
특히 임의의 점 $x \in X$에 대하여
$$
\mathcal{K}_X(\mathcal{F})_x
=
\mathcal{F}_x \otimes_{\mathcal{O}_{X, x}} \mathcal{K}_{X, x}
=
\mathcal{S}_x^{-1}\mathcal{F}_x
$$
이다. 그러나 $\mathcal{S}_x$가 국소환
$\mathcal{O}_{X, x}$의 영인자가 아닌 원소들의 집합이 아닐 수도
있으므로 주의해야 한다. 실제로 전몫환으로 가는 단사 사상
$$
\mathcal{K}_{X, x} \longrightarrow Q(\mathcal{O}_{X, x})
$$
이 언제나 존재한다. 이 사상이 전사인 것은 $\mathcal{S}_x$가
$\mathcal{O}_{X, x}$의 영인자가 아닌 원소들의 집합일 때이며
그때에 한한다. 일반적으로 유리형 절단의 층은 준연접 가군이 아니지만,
준연접 가군과 공통으로 지니는 성질들이 있다.

\begin{definition}
\label{divisors-definition-pullback-meromorphic-sections}
$f : (X, \mathcal{O}_X) \to (Y, \mathcal{O}_Y)$를 국소환 달린
공간들의 사상이라 하자. $f$에 대해, 모든 열린집합
$U \subset X$, $V \subset Y$의 쌍이 $f(U) \subset V$를 만족하고
모든 절단
$s \in \Gamma(V, \mathcal{S}_Y)$에 대하여 당김
$f^\sharp(s) \in \Gamma(U, \mathcal{O}_X)$가
$\Gamma(U, \mathcal{S}_X)$의 원소이면,
{\it 유리형 함수의 당김이 정의된다}고 한다.
\end{definition}

\noindent
이 경우 유도된 사상
$f^\sharp : f^{-1}\mathcal{K}_Y \to \mathcal{K}_X$가 존재한다.
다시 말해, 환 달린 공간들의 사상으로 이루어진 가환 그림
$$
\xymatrix{
(X, \mathcal{K}_X) \ar[r] \ar[d]^f &
(X, \mathcal{O}_X) \ar[d]^f \\
(Y, \mathcal{K}_Y) \ar[r] &
(Y, \mathcal{O}_Y)
}
$$
을 얻는다. 때때로 $f^*(s) = f^\sharp(s)$로 쓰며,
여기서 $s \in \Gamma(Y, \mathcal{K}_Y)$는 절단이다.

\begin{lemma}
\label{divisors-lemma-pullback-meromorphic-sections-defined}
$f : X \to Y$를 스킴의 사상이라 하자.
다음 각 경우에는 유리형 함수의 당김이 정의된다.
\begin{enumerate}
\item $X$의 모든 약하게 수반된 점이 $Y$의 어떤 기약 성분의
생성점으로 사상된다.
\item $X$, $Y$가 정역 스킴이고 $f$가 우세하다.
\item $X$가 정역 스킴이고 $X$의 생성점이 $Y$의 어떤 기약 성분의
생성점으로 사상된다.
\item $X$가 기약원 없는 스킴이고, $X$의 모든 기약 성분의
모든 생성점이 $Y$의 어떤 기약 성분의 생성점으로 사상된다.
\item $X$가 국소 뇌터 스킴이고, $X$의 모든 수반점이
$Y$의 어떤 기약 성분의 생성점으로 사상된다.
\item $X$가 국소 뇌터 스킴이고 매장점이 없으며,
$X$의 기약 성분의 모든 생성점이 $Y$의 어떤 기약 성분의
생성점으로 사상된다.
\item $f$가 평탄하다.
\end{enumerate}
\end{lemma}

\begin{proof}
문제는 $X$와 $Y$ 위에서 국소적이다. 따라서
$X = \Spec(A)$, $Y = \Spec(R)$이고 $f$가 환 사상
$\varphi : R \to A$로 주어지는 경우로 줄일 수 있다.
구조층의 정칙 절단에 대한 Lemma
\ref{divisors-lemma-regular-section-structure-sheaf}의 특징짓기에 의해,
$R \to A$가 영인자가 아닌 원소를 영인자가 아닌 원소로
보냄을 보이면 된다. $t \in R$를 영인자가 아닌 원소라 하자.

\medskip\noindent
$R \to A$가 평탄하면 $t : R \to R$가 단사이므로
$t : A \to A$도 단사이다. 이것으로 (7)이 증명된다.

\medskip\noindent
나머지 경우에는 $t$가 $R$의 어떤 극소 소아이디얼에도 속하지
않음을 관찰하자(환의 극소 소아이디얼의 모든 원소는 영인자이기
때문이다). 따라서 (1)의 경우 $\varphi(t)$는 $A$의 어떤 약하게
수반된 소아이디얼에도 속하지 않는다. 그러므로 이 경우는
Algebra, Lemma \ref{algebra-lemma-weakly-ass-zero-divisors}에서
따른다. (5)는 Lemma \ref{divisors-lemma-ass-weakly-ass}에 의해 (1)의
특별한 경우이다. (6)은 (5)와 정의들에서 따른다. (4)는
Lemma \ref{divisors-lemma-weakass-reduced}에 의해 (1)의 특별한 경우이다.
(2)와 (3)은 (4)의 특별한 경우이다.
\end{proof}

\begin{lemma}
\label{divisors-lemma-meromorphic-weakass-finite}
스킴 $X$가 다음을 만족한다고 하자.
\begin{enumerate}
\item[(a)] $X$의 모든 약하게 수반된 점은 $X$의 어떤 기약 성분의
생성점이다.
\item[(b)] 모든 준콤팩트 열린집합은 유한 개의 기약 성분을 갖는다.
\end{enumerate}
$X^0$를 $X$의 기약 성분들의 생성점으로 이루어진 집합이라 하자.
그러면
$$
\mathcal{K}_X =
\bigoplus\nolimits_{\eta \in X^0} j_{\eta, *}\mathcal{O}_{X, \eta} =
\prod\nolimits_{\eta \in X^0} j_{\eta, *}\mathcal{O}_{X, \eta}
$$
이다. 여기서 $j_\eta : \Spec(\mathcal{O}_{X, \eta}) \to X$는
Schemes, Section \ref{schemes-section-points}의 표준 사상이다.
더욱이 다음이 성립한다.
\begin{enumerate}
\item $\mathcal{K}_X$는 $\mathcal{O}_X$-대수의 준연접층이다.
\item 모든 준연접 $\mathcal{O}_X$-가군 $\mathcal{F}$에 대하여
유리형 절단의 층
$$
\mathcal{K}_X(\mathcal{F}) =
\bigoplus\nolimits_{\eta \in X^0} j_{\eta, *}\mathcal{F}_\eta =
\prod\nolimits_{\eta \in X^0} j_{\eta, *}\mathcal{F}_\eta
$$
은 $\mathcal{F}$의 유리형 절단으로 이루어지며 준연접이다.
\item $\mathcal{S}_x \subset \mathcal{O}_{X, x}$는 임의의
$x \in X$에 대하여 영인자가 아닌 원소들의 집합이다.
\item $\mathcal{K}_{X, x}$는 $\mathcal{O}_{X, x}$의 전몫환이며,
이는 임의의 $x \in X$에 대해 성립한다.
\item $\mathcal{K}_X(U)$는 $\mathcal{O}_X(U)$의 전몫환과 같으며,
이는 임의의 아핀 열린집합 $U \subset X$에 대해 성립한다.
\item $X$의 유리 함수환
(Morphisms, Definition \ref{morphisms-definition-rational-function})은
$X$ 위의 유리형 함수환이며, 식으로 쓰면
$R(X) = \Gamma(X, \mathcal{K}_X)$이다.
\end{enumerate}
\end{lemma}

\begin{proof}
예를 들어 줄기에서 확인하면 국소 유한한 가군층의 직합은 곱과
같음을 알 수 있다. 또한
$\mathcal{K}_X(\mathcal{F}) =
\mathcal{F} \otimes_{\mathcal{O}_X} \mathcal{K}_X$이므로
(2)는 나머지 명제들에서 따른다. 또한 $X^0$가 유한할 때에는 (6)이
보조정리의 첫 명제와 Morphisms, Lemma
\ref{morphisms-lemma-integral-scheme-rational-functions}에서 따름을
관찰하자. (6)의 일반적인 경우는 이를 이어붙여서 얻는다(논증은
생략한다).

\medskip\noindent
$j : Y = \coprod\nolimits_{\eta \in X^0} \Spec(\mathcal{O}_{X, \eta}) \to X$
를 사상 $j_\eta$들의 쌍대합이라 하자. 우리는
$\mathcal{K}_X = j_*\mathcal{O}_Y$임을 보여야 한다.
먼저 $\mathcal{K}_Y = \mathcal{O}_Y$임을 유의하자. 이는 $Y$가
차원이 $0$인 국소환들의 스펙트럼의 서로소 합이기 때문이다. 차원이 영인
국소환에서는 영인자가 아닌 모든 원소가 가역원이다. 다음으로
Lemma \ref{divisors-lemma-pullback-meromorphic-sections-defined}에 의해
$j$에 대해 유리형 함수의 당김이 정의됨을 유의하자. 이로부터 사상
$$
\mathcal{K}_X \longrightarrow j_*\mathcal{O}_Y.
$$
을 얻는다. $\Spec(A) = U \subset X$를 아핀 열린집합이라 하자.
그러면 $A$는 유한 개의 극소 소아이디얼
$\mathfrak q_1, \ldots, \mathfrak q_t$를 가지며, $A$의 모든
약하게 수반된 소아이디얼은 $\mathfrak q_i$ 중 하나이다.
Algebra, Lemmas \ref{algebra-lemma-total-ring-fractions-no-embedded-points}
와 \ref{algebra-lemma-weakly-ass-zero-divisors}에 의해
$Q(A) = \prod A_{\mathfrak q_i}$를 얻는다. 다시 말해, 준층
$U \mapsto \mathcal{S}(U)^{-1}\mathcal{O}_X(U)$의 값은 이미
$j_*\mathcal{O}_Y(U)$와 같으며, 이는 우리의 아핀 열린집합
$U$ 위에서 성립한다.
따라서 위에 표시한 사상은 동형이며, 이것으로 보조정리의 명제에
나오는 첫 번째 표시 등식이 증명된다.

\medskip\noindent
끝으로 (1), (3), (4), (5)를 증명하자. (5)에 대해서는 증명 도중에
$\mathcal{K}_X = j_*\mathcal{O}_Y$임을 보았다. 사상 $j$는
$X$의 기약 성분들의 집합이 국소 유한하다는 가정에 의해 준콤팩트이다.
따라서 $j$는 준콤팩트이고 준분리이다($Y$가 분리이기 때문이다).
Schemes, Lemma \ref{schemes-lemma-push-forward-quasi-coherent}에 의해
$j_*\mathcal{O}_Y$는 준연접이다. 이것으로 (1)이 증명된다.
$x \in X$라 하자. 아핀 열린집합 $U = \Spec(A)$를 고르되, 이것이
$x$의 근방이고 모든 기약 성분이 $x$를 지나도록 할 수 있다.
그러면 $A \subset A_\mathfrak p$이다. 실제로 $A$의 모든 약하게
수반된 소아이디얼은 $\mathfrak p$에 포함되므로, Algebra, Lemma
\ref{algebra-lemma-weakly-ass-zero-divisors}에 의해
$A \setminus \mathfrak p$의 원소들은 영인자가 아니다.
$A_\mathfrak p$의 영인자가 아닌 모든 원소는 $x$의 (필요하면
더 작은) 아핀 열린 근방 위 영인자가 아닌 원소의 상임이 쉽게 따른다.
이것으로 (3)이 증명된다. (4)는 줄기를 계산하면 (3)에서 따른다.
\end{proof}

\begin{definition}
\label{divisors-definition-regular-meromorphic-section}
$X$를 국소환 달린 공간이라 하자.
$\mathcal{L}$을 가역 $\mathcal{O}_X$-가군이라 하자.
유리형 절단 $s$가 $\mathcal{L}$의 절단이라고 하고, 이 절단이
유도하는 사상
$\mathcal{K}_X \to \mathcal{K}_X(\mathcal{L})$이 단사이면
$s$를 {\it 정칙}이라고 한다. 다시 말해, 이는 가역
$\mathcal{K}_X$-가군 $\mathcal{K}_X(\mathcal{L})$의 정칙
절단이며, Definition \ref{divisors-definition-regular-section}을 보라.
\end{definition}

\noindent
(정칙) 유리형 절단을 언제 당길 수 있는지 구체적으로 적어 보자.

\begin{lemma}
\label{divisors-lemma-meromorphic-sections-pullback}
$f : X \to Y$를 국소환 달린 공간들의 사상이라 하자.
$f$에 대해 유리형 함수의 당김이 정의된다고 가정하자
(Definition \ref{divisors-definition-pullback-meromorphic-sections}을 보라).
\begin{enumerate}
\item $\mathcal{F}$를 $\mathcal{O}_Y$-가군의 층이라 하자.
표준 당김 사상
$f^* : \Gamma(Y, \mathcal{K}_Y(\mathcal{F})) \to
\Gamma(X, \mathcal{K}_X(f^*\mathcal{F}))$이 존재하며, 이는
$\mathcal{F}$의 유리형 절단에 대한 당김이다.
\item $\mathcal{L}$을 가역 $\mathcal{O}_Y$-가군이라 하자.
정칙 유리형 절단 $s$가 $\mathcal{L}$의 절단이면, 그 당김
$f^*s$는 $f^*\mathcal{L}$의 정칙 유리형 절단이다.
\end{enumerate}
\end{lemma}

\begin{proof}
생략한다.
\end{proof}

\begin{lemma}
\label{divisors-lemma-regular-meromorphic-ideal-denominators}
$X$를 스킴이라 하자.
$\mathcal{L}$을 가역 $\mathcal{O}_X$-가군이라 하자.
$s$를 $\mathcal{L}$의 정칙 유리형 절단이라 하자.
다음 규칙으로 정의되는 아이디얼층을
$\mathcal{I} \subset \mathcal{O}_X$로 쓰자.
$$
\mathcal{I}(V)
=
\{f \in \mathcal{O}_X(V) \mid fs \in \mathcal{L}(V)\}.
$$
$\mathcal{L}(V) \subset \mathcal{K}_X(\mathcal{L})(V)$이므로 이 식은
의미가 있다. 그러면 $\mathcal{I}$는 준연접 아이디얼층이고 단사 사상
$$
1 : \mathcal{I} \longrightarrow \mathcal{O}_X,
\quad
s : \mathcal{I} \longrightarrow \mathcal{L}
$$
이 존재하며, 그 여핵들은 $X$의 닫힌 무소조 부분집합 위에 지지된다.
\end{lemma}

\begin{proof}
문제는 $X$ 위에서 국소적이다. 따라서 $X = \Spec(A)$이고
$\mathcal{L} = \mathcal{O}_X$라고 가정해도 된다. 더 축소하여
$s = a/b$이고 $a, b \in A$가 $A$에서 {\it 둘 다} 영인자가 아닌
원소라고 가정해도 된다.
$I = \{x \in A \mid x(a/b) \in A\}$로 놓자.

\medskip\noindent
$\mathcal{I}$가 준연접임을 보이려면
$I_f = \{x \in A_f \mid x(a/b) \in A_f\}$임을 모든
$f \in A$에 대하여 보여야 한다.
$c/f^n \in A_f$, $(c/f^n)(a/b) \in A_f$라 하자. 그러면 어떤
$f^mc(a/b) \in A$이 성립하며, 여기서 $m$은 어떤 정수이다.
따라서 $c/f^n \in I_f$이다.
역으로 $I_f$가 $\{x \in A_f \mid x(a/b) \in A_f\}$에 포함됨은
쉽게 알 수 있다. 이것으로 준연접성이 증명된다.

\medskip\noindent
마지막 명제를 증명하자. $(b) \subset I$임은 명백하다. 따라서
$V(I) \subset V(b)$는 무소조 부분집합이다. 실제로 $b$는
영인자가 아니다.
그러므로 $1$의 여핵은 닫힌 무소조 집합 위에 지지된다.
같은 논증을 $s$의 여핵에도 적용할 수 있는데, 이는
$s(b) = (a) \subset sI \subset A$이기 때문이다.
\end{proof}

\begin{definition}
\label{divisors-definition-regular-meromorphic-ideal-denominators}
$X$를 스킴이라 하자.
$\mathcal{L}$을 가역 $\mathcal{O}_X$-가군이라 하자.
$s$를 $\mathcal{L}$의 정칙 유리형 절단이라 하자.
Lemma \ref{divisors-lemma-regular-meromorphic-ideal-denominators}에서 구성한
아이디얼층 $\mathcal{I}$를 {\it $s$의 분모 아이디얼층}이라 한다.
\end{definition}



\section{유리형 함수와 절단; 뇌터 경우}
\label{divisors-section-meromorphic-noetherian}

\noindent
국소 뇌터 스킴에 대해서는 유리형 함수층에 관한 몇 가지 결과를
증명할 수 있다. 그러나 \cite{misconceptions}에는
$\mathcal{K}_X$가 준연접일 필요가 없는 뇌터 스킴 $X$를 보이는
예가 있다.

\begin{lemma}
\label{divisors-lemma-meromorphic-section-restricts-to-zero}
$X$를 준콤팩트 스킴이라 하자. $h \in \Gamma(X, \mathcal{O}_X)$이고
$f \in \Gamma(X, \mathcal{K}_X)$이며, $f$를 $X_h$에 제한하면
영이라고 하자. 그러면 $h^n f = 0$이며, 여기서 어떤 $n \gg 0$이다.
\end{lemma}

\begin{proof}
$X$를 아핀 열린집합 $U$들로 덮되, 각 열린집합에서
$f|_U = s^{-1}a$이고 $a \in \mathcal{O}_X(U)$,
$s \in \mathcal{S}(U)$가 되도록 할 수 있다. $X$가 준콤팩트이므로
이 꼴의 유한 개 아핀 열린집합으로 덮을 수 있다. 따라서
$X = \Spec(A)$이고 $f = s^{-1}a$인 경우에 보조정리를 증명하면
충분하다. $s \in A$는 영인자가 아니므로 $f = a$인 경우에
결과를 증명하면 충분함을 유의하자. 조건 $f|_{X_h} = 0$은
$a$의 $A_h = \mathcal{O}_X(X_h)$에서의 상이 영임을 뜻한다.
이는 $\mathcal{O}_X \subset \mathcal{K}_X$이기 때문이다.
따라서 원하는 대로 $h^na = 0$이며, 여기서 어떤 $n > 0$이다.
\end{proof}

\begin{lemma}
\label{divisors-lemma-locally-Noetherian-K}
$X$를 국소 뇌터 스킴이라 하자.
\begin{enumerate}
\item 임의의 $x \in X$에 대하여
$\mathcal{S}_x \subset \mathcal{O}_{X, x}$는 영인자가 아닌
원소들의 집합이고, 따라서 $\mathcal{K}_{X, x}$는
$\mathcal{O}_{X, x}$의 전몫환이다.
\item 임의의 아핀 열린집합 $U \subset X$에 대하여 환
$\mathcal{K}_X(U)$는 $\mathcal{O}_X(U)$의 전몫환과 같다.
\end{enumerate}
\end{lemma}

\begin{proof}
이 보조정리를 증명할 때 $X$가 뇌터 환 $A$의 스펙트럼이라고
가정해도 된다. $x \in X$가 $\mathfrak p \subset A$에 대응한다고
하자.

\medskip\noindent
(1)의 증명. $\mathcal{S}_x$가
$\mathcal{O}_{X, x} = A_\mathfrak p$의 영인자가 아닌 원소들의
집합에 포함됨은 명백하다. 역으로 $f, g \in A$,
$g \not \in \mathfrak p$이고 $f/g$가 $A_{\mathfrak p}$에서
영인자가 아니라고 하자.
$I = \{a \in A \mid af = 0\}$로 놓자. 국소화의 완전성에 의해
$I_{\mathfrak p} = 0$이다. $A$가 뇌터이므로 $I$는 유한 생성이고,
따라서 $g'I = 0$이다. 여기서 어떤 $g' \in A$를 택할 수 있으며
$g' \not \in \mathfrak p$이다. 그러므로 $f$는 $A_{g'}$, 즉 $\mathfrak p$의
자리스키 열린 근방에서 영인자가 아니다. 따라서 $f/g$는
$\mathcal{S}_x$의 원소이다.

\medskip\noindent
(2)의 증명. $f \in \Gamma(X, \mathcal{K}_X)$를 유리형 함수라 하자.
$I = \{a \in A \mid af \in A\}$로 놓자. 소아이디얼
$\mathfrak p \subset A$를 고정하되, 이에 대응하는 점을
$x \in X$라 하자. (1)에 의해
$f$의 $\mathfrak p$에서의 줄기 속 상을 $a/b$로 쓸 수 있다.
여기서 $a, b \in A_{\mathfrak p}$이고
$b \in A_{\mathfrak p}$는 영인자가 아니다.
$b = c/d$로 쓰자. 여기서 $c, d \in A$이고
$d \not \in \mathfrak p$이다. 그러면 $ad - cf$는
$\mathcal{K}_X$의 절단으로서 $x$의 어떤 열린 근방에서
영이 된다. 이 근방을 $D(e)$라 하자. 여기서 $e \in A$이고
$e \not \in \mathfrak p$이다. Lemma
\ref{divisors-lemma-meromorphic-section-restricts-to-zero}에 의해
$e^n(ad - cf) = 0$이며, 여기서 어떤 $n \gg 0$이다.
따라서 $e^nc \in I$이고, $e^nc$의 $A_{\mathfrak p}$에서의 상은
영인자가 아니다.
$\text{Ass}(A) = \{\mathfrak q_1, \ldots, \mathfrak q_t\}$를
$A$의 수반 소아이디얼들의 집합이라 하자.
$IA_{\mathfrak q_i}$를 살펴보고 Algebra, Lemma
\ref{algebra-lemma-associated-primes-localize}를 사용하면 위 결과는
$I \not \subset \mathfrak q_i$임을 뜻하며, 이는 각 $i$에 대해 성립한다.
Algebra, Lemma \ref{algebra-lemma-silly}에 의해
$z \in I$, $z \not \in \bigcup \mathfrak q_i$인 원소가 존재한다.
Algebra, Lemma \ref{algebra-lemma-ass-zero-divisors}에 의해
$z$는 $A$에서 영인자가 아니다. 따라서
$f = (zf)/z$는 $A$의 전분수환의 원소이다. 이것으로 (2)가 증명된다.
\end{proof}

\begin{lemma}
\label{divisors-lemma-quasi-coherent-K}
$X$를 매장점이 없는 국소 뇌터 스킴이라 하자.
$X^0$를 $X$의 기약 성분들의 생성점으로 이루어진 집합이라 하자.
그러면
$$
\mathcal{K}_X =
\bigoplus\nolimits_{\eta \in X^0} j_{\eta, *}\mathcal{O}_{X, \eta} =
\prod\nolimits_{\eta \in X^0} j_{\eta, *}\mathcal{O}_{X, \eta}
$$
이다. 여기서 $j_\eta : \Spec(\mathcal{O}_{X, \eta}) \to X$는
Schemes, Section \ref{schemes-section-points}의 표준 사상이다.
더욱이 다음이 성립한다.
\begin{enumerate}
\item $\mathcal{K}_X$는 $\mathcal{O}_X$-대수의 준연접층이다.
\item 모든 준연접 $\mathcal{O}_X$-가군 $\mathcal{F}$에 대하여
유리형 절단의 층
$$
\mathcal{K}_X(\mathcal{F}) =
\bigoplus\nolimits_{\eta \in X^0} j_{\eta, *}\mathcal{F}_\eta =
\prod\nolimits_{\eta \in X^0} j_{\eta, *}\mathcal{F}_\eta
$$
은 $\mathcal{F}$의 유리형 절단으로 이루어지며 준연접이다.
\item $X$의 유리 함수환은 $X$ 위의 유리형 함수환이며,
식으로 쓰면 $R(X) = \Gamma(X, \mathcal{K}_X)$이다.
\end{enumerate}
\end{lemma}

\begin{proof}
이 보조정리는 Lemma \ref{divisors-lemma-meromorphic-weakass-finite}의
특별한 경우이다. 실제로 국소 뇌터 경우에는 약하게 수반된 점과
수반점이 같으며, 이는 Lemma \ref{divisors-lemma-ass-weakly-ass}에 따른다.
\end{proof}

\begin{lemma}
\label{divisors-lemma-regular-meromorphic-section-exists-noetherian}
$X$를 매장점이 없는 국소 뇌터 스킴이라 하자.
$\mathcal{L}$을 가역 $\mathcal{O}_X$-가군이라 하자.
그러면 $\mathcal{L}$은 정칙 유리형 절단을 갖는다.
\end{lemma}

\begin{proof}
각 생성점 $\eta$에 대하여, 이 점이 $X$의 점일 때 생성원
$s_\eta$를 고르자. 이는 계수 $1$ 자유 가군
$\mathcal{L}_\eta$의 생성원이며, 이 가군은 아르틴 국소환
$\mathcal{O}_{X, \eta}$ 위의 가군이다.
Lemma \ref{divisors-lemma-quasi-coherent-K}에 나오는
$\mathcal{K}_X$와 $\mathcal{K}_X(\mathcal{L})$의 서술에서 곧바로
$s = \prod s_\eta$가 $\mathcal{L}$의 정칙 유리형 절단임이 따른다.
\end{proof}

\begin{lemma}
\label{divisors-lemma-make-maps-regular-section}
다음이 주어졌다고 하자.
\begin{enumerate}
\item 국소 뇌터 스킴 $X$,
\item 가역 가군 $\mathcal{L}$로서 $\mathcal{O}_X$-가군인 것,
\item 정칙 유리형 절단 $s$로서 $\mathcal{L}$의 절단인 것,
\item 연접층 $\mathcal{F}$로서 $X$ 위에 있고 매장 수반점이 없으며
$\text{Supp}(\mathcal{F}) = X$인 것.
\end{enumerate}
$\mathcal{I} \subset \mathcal{O}_X$를 $s$의 분모 아이디얼이라 하자.
$T \subset X$를 $\mathcal{O}_X/\mathcal{I}$와
$\mathcal{L}/s(\mathcal{I})$의 지지집합들의 합집합이라 하자.
Lemma \ref{divisors-lemma-regular-meromorphic-ideal-denominators}에 의해 이는
무소조 닫힌 부분집합 $T \subset X$이다. 그러면 표준 단사 사상
$$
1 : \mathcal{I}\mathcal{F} \to \mathcal{F}, \quad
s : \mathcal{I}\mathcal{F} \to \mathcal{F} \otimes_{\mathcal{O}_X}\mathcal{L}
$$
이 존재하며, 그 여핵들은 $T$ 위에 지지된다.
\end{lemma}

\begin{proof}
$\mathcal{L} \cong \mathcal{O}_X$이고 $s = a/b$이며
$a, b \in A$가 둘 다 영인자가 아닌 아핀 경우로 줄인다.
환원 단계의 증명은 생략한다.
$\mathcal{F} = \widetilde{M}$로 쓰자.
$I = \{x \in A \mid x(a/b) \in A\}$로 놓으면
$\mathcal{I} = \widetilde{I}$이다(Lemma
\ref{divisors-lemma-regular-meromorphic-ideal-denominators}의 증명을 보라).
$T = V(I) \cup V((a/b)I)$임을 유의하자.
임의의 $A$-가군 $M$에 대하여 사상 $1 : IM \to M$을 생각하자.
이는 보조정리의 사상 $1$을 주는 사상이다. 한편 사상
$\sigma : IM \to M_b, x \mapsto ax/b$를 생각하자.
$b$는 $A$에서 영인자가 아니고, $M$의 지지집합은
$\Spec(A)$이며 매장 소아이디얼이 없으므로 $b$는 $M$에서도
영인자가 아니다. 따라서 $M \subset M_b$이다. $I$의 정의에 의해
$\sigma(IM) \subset M$이며, 이는 $M_b$의 부분가군으로서의
포함이다. 따라서 $A$-가군 사상 $s : IM \to M$을 얻는다
(즉, $s(z)/1 = \sigma(z)$를 $M_b$에서 모든 $z \in IM$에 대하여
만족하는 유일한 사상이다).
$a$도 영인자가 아니므로($A$와 $M$ 모두에서) 이 사상은 단사이다.
$M/IM$은 $I$에 의해 소멸하고 $M/s(IM)$은 $(a/b)I$에 의해
소멸함이 명백하다. 따라서 보조정리가 따른다.
\end{proof}





\section{유리형 함수와 절단; 기약원 없는 경우}
\label{divisors-section-meromorphic-reduced}

\noindent
기약원 없고 국소적으로 유한 개의 기약 성분을 갖는 스킴에서는
유리형 함수층이 준연접이다.

\begin{lemma}
\label{divisors-lemma-reduced-finite-irreducible}
$X$를 모든 준콤팩트 열린집합이 유한 개의 기약 성분을 갖는
기약원 없는 스킴이라 하자. $X^0$를 $X$의 기약 성분들의
생성점으로 이루어진 집합이라 하자. 그러면
$$
\mathcal{K}_X =
\bigoplus\nolimits_{\eta \in X^0} j_{\eta, *}\kappa(\eta) =
\prod\nolimits_{\eta \in X^0} j_{\eta, *}\kappa(\eta)
$$
이다. 여기서 $j_\eta : \Spec(\kappa(\eta)) \to X$는
Schemes, Section \ref{schemes-section-points}의 표준 사상이다.
더욱이 다음이 성립한다.
\begin{enumerate}
\item $\mathcal{K}_X$는 $\mathcal{O}_X$-대수의 준연접층이다.
\item 모든 준연접 $\mathcal{O}_X$-가군 $\mathcal{F}$에 대하여
유리형 절단의 층
$$
\mathcal{K}_X(\mathcal{F}) =
\bigoplus\nolimits_{\eta \in X^0} j_{\eta, *}\mathcal{F}_\eta =
\prod\nolimits_{\eta \in X^0} j_{\eta, *}\mathcal{F}_\eta
$$
은 $\mathcal{F}$의 유리형 절단으로 이루어지며 준연접이다.
\item $\mathcal{S}_x \subset \mathcal{O}_{X, x}$는 임의의
$x \in X$에 대하여 영인자가 아닌 원소들의 집합이다.
\item $\mathcal{K}_{X, x}$는 $\mathcal{O}_{X, x}$의 전몫환이며,
이는 임의의 $x \in X$에 대해 성립한다.
\item $\mathcal{K}_X(U)$는 $\mathcal{O}_X(U)$의 전몫환과 같으며,
이는 임의의 아핀 열린집합 $U \subset X$에 대해 성립한다.
\item $X$의 유리 함수환은 $X$ 위의 유리형 함수환이며,
식으로 쓰면 $R(X) = \Gamma(X, \mathcal{K}_X)$이다.
\end{enumerate}
\end{lemma}

\begin{proof}
이 보조정리는 Lemma \ref{divisors-lemma-meromorphic-weakass-finite}의
특별한 경우이다. 실제로 기약원 없는 스킴에서는 약하게 수반된 점이
생성점이며, 이는 Lemma \ref{divisors-lemma-weakass-reduced}에 따른다.
\end{proof}

\begin{lemma}
\label{divisors-lemma-reduced-normalization}
$X$를 스킴이라 하자. $X$가 기약원 없고 모든 준콤팩트 열린집합
$U \subset X$가 유한 개의 기약 성분을 갖는다고 가정하자.
그러면 정규화 사상 $\nu : X^\nu \to X$는 사상
$$
\underline{\Spec}_X(\mathcal{O}') \longrightarrow X
$$
이다. 여기서 $\mathcal{O}' \subset \mathcal{K}_X$는 유리형
함수층 안에서 $\mathcal{O}_X$의 정수적 폐포이다.
\end{lemma}

\begin{proof}
정규화 사상 $\nu : X^\nu \to X$의 정의
(Morphisms, Definition \ref{morphisms-definition-normalization}을 보라)와
위 Lemma \ref{divisors-lemma-reduced-finite-irreducible}의
$\mathcal{K}_X$에 대한 서술을 비교하면 된다.
\end{proof}

\begin{lemma}
\label{divisors-lemma-meromorphic-functions-integral-scheme}
$X$를 생성점 $\eta$를 갖는 정역 스킴이라 하자. 다음이 성립한다.
\begin{enumerate}
\item 유리형 함수층은 $X$의 함수체
(Morphisms, Definition \ref{morphisms-definition-function-field}을 보라)를
값으로 갖는 상수층과 동형이다.
\item 임의의 준연접층 $\mathcal{F}$가 $X$ 위에 주어졌을 때
층 $\mathcal{K}_X(\mathcal{F})$는 $\mathcal{F}_\eta$를 값으로
갖는 상수층과 동형이다.
\end{enumerate}
\end{lemma}

\begin{proof}
생략한다.
\end{proof}

\noindent
몇몇 경우에는 정칙 유리형 절단이 존재함을 보일 수 있다.

\begin{lemma}
\label{divisors-lemma-regular-meromorphic-section-exists}
$X$를 스킴이라 하자.
$\mathcal{L}$을 가역 $\mathcal{O}_X$-가군이라 하자.
다음 각 경우에 $\mathcal{L}$은 정칙 유리형 절단을 갖는다.
\begin{enumerate}
\item $X$는 정역 스킴이다.
\item $X$는 기약원 없고 모든 준콤팩트 열린집합은 유한 개의
기약 성분을 갖는다.
\item $X$는 국소 뇌터이고 매장점이 없다.
\end{enumerate}
\end{lemma}

\begin{proof}
(1)의 경우 $\eta \in X$를 생성점이라 하자.
Lemma \ref{divisors-lemma-meromorphic-functions-integral-scheme}에서
$\mathcal{K}_X$와 각각 $\mathcal{K}_X(\mathcal{L})$가 각각
$\kappa(\eta)$와 $\mathcal{L}_\eta$를 값으로 갖는 상수층임을 보았다.
$\dim_{\kappa(\eta)} \mathcal{L}_\eta = 1$이므로 영이 아닌 원소
$s \in \mathcal{L}_\eta$를 고를 수 있다. 명백히 $s$는
$\mathcal{L}$의 정칙 유리형 절단이다. (2)의 경우 영이 아닌
$s_\eta \in \mathcal{L}_\eta$를 모든 생성점 $\eta$에 대하여
고르자. 여기서 생성점들은 $X$의 점들이다. 이는 $\mathcal{L}_\eta$가
$1$차원 벡터공간이고 그 기초체가 $\kappa(\eta)$이므로 가능하다.
Lemma \ref{divisors-lemma-reduced-finite-irreducible}에 나오는
$\mathcal{K}_X$와 $\mathcal{K}_X(\mathcal{L})$의 서술에서 곧바로
$s = \prod s_\eta$가 $\mathcal{L}$의 정칙 유리형 절단임이 따른다.
(3)은 Lemma \ref{divisors-lemma-regular-meromorphic-section-exists-noetherian}이다.
\end{proof}





\section{베유 인자}
\label{divisors-section-Weil-divisors}

\noindent
국소 뇌터 정역 스킴에 대하여 베유 인자와 베유 인자의 유리 동치를
도입할 것이다. 스킴이 준콤팩트라고 가정하지 않으므로 베유 인자를
정의할 때 약간 주의해야 한다. 예를 들어 유리 함수는 극을 무한히
많이 가질 수 있으므로 소인자들의 무한합을 허용해야 한다. 준콤팩트
스킴에서는 통상과 같이 베유 인자가 유한합이다. 다음은 닫힌
부분스킴들의 모임이 국소 유한함을 증명할 때 자주 사용할
기초적인 보조정리이다.

\begin{lemma}
\label{divisors-lemma-components-locally-finite}
$X$를 국소 뇌터 스킴이라 하자. $Z \subset X$를 닫힌
부분스킴이라 하자. $Z$의 기약 성분들의 모임은 $X$에서
국소 유한하다.
\end{lemma}

\begin{proof}
$U \subset X$를 준콤팩트 열린 부분스킴이라 하자. 그러면 $U$는
뇌터 스킴이고, 따라서 그 바탕 위상공간은 뇌터이다
(Properties, Lemma \ref{properties-lemma-Noetherian-topology}).
그러므로 모든 부분공간이 뇌터이고 유한 개의 기약 성분을 갖는다
(Topology, Lemma \ref{topology-lemma-Noetherian}을 보라).
\end{proof}

\noindent
$Z$가 스킴 $X$의 기약 닫힌 부분집합이면 $Z$의 $X$에서의 여차원은
국소환 $\mathcal{O}_{X, \xi}$의 차원과 같음을 상기하자. 여기서
$\xi \in Z$는 생성점이다. Properties, Lemma
\ref{properties-lemma-codimension-local-ring}을 보라.

\begin{definition}
\label{divisors-definition-Weil-divisor}
$X$를 국소 뇌터 정역 스킴이라 하자.
\begin{enumerate}
\item {\it 소인자}란 정역 닫힌 부분스킴 $Z \subset X$로서
여차원이 $1$인 것을 말한다.
\item {\it 베유 인자}란 형식합 $D = \sum n_Z Z$를 말한다.
여기서 합은 $X$의 소인자들에 걸치며, 모임
$\{Z \mid n_Z \not = 0\}$은 국소 유한하다
(Topology, Definition \ref{topology-definition-locally-finite}).
\end{enumerate}
$X$ 위의 모든 베유 인자로 이루어진 군을 $\text{Div}(X)$로 쓴다.
\end{definition}

\noindent
다음 목표는 유리 함수에 수반되는 베유 인자를 정의하는 것이다.
이를 위해 소인자를 따라 유리 함수가 소멸하는 차수를 사용하며,
이는 다음과 같이 정의한다.

\begin{definition}
\label{divisors-definition-order-vanishing}
$X$를 국소 뇌터 정역 스킴이라 하자. $f \in R(X)^*$라 하자.
모든 소인자 $Z \subset X$에 대하여
{\it $f$의 $Z$를 따른 소멸 차수}를 정수
$$
\text{ord}_Z(f) = \text{ord}_{\mathcal{O}_{X, \xi}}(f)
$$
로 정의한다. 우변은 Algebra, Definition
\ref{algebra-definition-ord}의 개념이고, $\xi$는 $Z$의 생성점이다.
\end{definition}

\noindent
$f, g \in R(X)^*$에 대하여
$$
\text{ord}_Z(fg) = \text{ord}_Z(f) + \text{ord}_Z(g).
$$
임을 유의하자. 물론 $\text{ord}_Z(f) < 0$일 수도 있다.
이 경우 $f$가 $Z$를 따라 {\it 극}을 갖는다고 하고,
$-\text{ord}_Z(f) > 0$을 {\it $f$의 $Z$를 따른 극의 차수}라고 한다.
조건 $\text{ord}_Z(f) \geq 0$은 조건
$f \in \mathcal{O}_{X, \xi}$와 동치가 {\bf 아님}에 유의하는 것이
중요하다. 단, 국소환 $\mathcal{O}_{X, \xi}$가 이산 값매김환이면
두 조건은 동치이다.

\begin{lemma}
\label{divisors-lemma-divisor-locally-finite}
$X$를 국소 뇌터 정역 스킴이라 하자. $f \in R(X)^*$라 하자.
두 모임
$$
\{Z \subset X \mid Z\text{ a prime divisor with generic point }\xi
\text{ and }f\text{ not in }\mathcal{O}_{X, \xi}\}
$$
과
$$
\{Z \subset X \mid Z \text{ a prime divisor and }\text{ord}_Z(f) \not = 0\}
$$
은 $X$에서 국소 유한하다.
\end{lemma}

\begin{proof}
공집합이 아닌 열린 부분스킴 $U \subset X$가 존재하여 $f$가
$\Gamma(U, \mathcal{O}_X^*)$의 절단에 대응한다. 따라서 보조정리의
집합들에 나타날 수 있는 소인자들은 모두 $X \setminus U$의
기약 성분이다. 그러므로 Lemma
\ref{divisors-lemma-components-locally-finite}가 원하는 결과를 준다.
\end{proof}

\noindent
이 보조정리에 의해 다음 정의를 내릴 수 있다.

\begin{definition}
\label{divisors-definition-principal-divisor}
$X$를 국소 뇌터 정역 스킴이라 하자. $f \in R(X)^*$라 하자.
{\it $f$에 수반되는 주 베유 인자}란 베유 인자
$$
\text{div}(f) = \text{div}_X(f) = \sum \text{ord}_Z(f) [Z]
$$
를 말한다. 여기서 합은 소인자들에 걸치고 $\text{ord}_Z(f)$는
Definition \ref{divisors-definition-order-vanishing}에서와 같다.
Lemma \ref{divisors-lemma-divisor-locally-finite}에 의해 이는 잘 정의된다.
\end{definition}

\begin{lemma}
\label{divisors-lemma-div-additive}
$X$를 국소 뇌터 정역 스킴이라 하자. $f, g \in R(X)^*$라 하자.
그러면 $X$ 위의 베유 인자로서
$$
\text{div}_X(fg) = \text{div}_X(f) + \text{div}_X(g)
$$
이다.
\end{lemma}

\begin{proof}
$\text{ord}$ 함수들의 가법성에서 명백히 따른다.
\end{proof}

\noindent
위 보조정리에 의해 주 베유 인자들의 모임이 모든 베유 인자로
이루어진 군의 부분군임을 알 수 있다. 이로부터 다음 정의를 얻는다.

\begin{definition}
\label{divisors-definition-class-group}
$X$를 국소 뇌터 정역 스킴이라 하자. $X$의
{\it 베유 인자류군}은 베유 인자군을 주 베유 인자들의 부분군으로
나눈 몫이다. 표기는 $\text{Cl}(X)$이다.
\end{definition}

\noindent
구성에 의해 완전 복합체
\begin{equation}
\label{divisors-equation-Weil-divisor-class}
R(X)^* \xrightarrow{\text{div}} \text{Div}(X) \to \text{Cl}(X) \to 0
\end{equation}
를 얻으며, 이를 $\text{Cl}(X)$의 표시라고 생각할 수 있다.
다음 목표는 베유 인자류군을 피카르 군과 연관 짓는 것이다.





\section{가역 가군에 수반되는 베유 인자류}
\label{divisors-section-c1}

\noindent
이 절에서는 국소 뇌터 정역 스킴 위의 표준 사상
$\Pic(X) \to \text{Cl}(X)$을 정의하기 위해
Section \ref{divisors-section-Weil-divisors}와 정확히 같은 단계를 밟는다.

\medskip\noindent
$X$를 스킴이라 하자. $\mathcal{L}$을 가역
$\mathcal{O}_X$-가군이라 하자. $\xi \in X$를 점이라 하자.
$s_\xi, s'_\xi \in \mathcal{L}_\xi$가
$\mathcal{L}_\xi$를 $\mathcal{O}_{X, \xi}$-가군으로서 생성하면,
가역원 $u \in \mathcal{O}_{X, \xi}^*$가 존재하여
$s_\xi = u s'_\xi$이다. 유리형 절단층
$\mathcal{K}_X(\mathcal{L})$은 $\mathcal{L}$의 절단층이며,
그 $x$에서의 줄기는
$\mathcal{K}_{X, x} \otimes_{\mathcal{O}_{X, x}} \mathcal{L}_x$와
같다. 따라서 임의의 유리형 절단 $s$가 $\mathcal{L}$의 절단일 때,
그 $x$에서의 줄기 속 상은 $s = fs_\xi$로 쓸 수 있으며, 여기서
$f \in \mathcal{K}_{X, x}$이다. 아래에서는 이를 줄여
$f = s/s_\xi$라고 말한다. 또한 $X$가 정역 스킴이면
$\mathcal{K}_{X, x} = R(X)$는 $X$의 함수체와 같으므로
$s/s_\xi \in R(X)$이다. $s$가 정칙 유리형 절단이면 실제로
$s/s_\xi \in R(X)^*$이다. 정역 스킴에서는 정칙 유리형 절단과
영이 아닌 유리형 절단이 같은 개념이다. 끝으로 $s/s_\xi$는
$s_\xi$의 선택에 국소환 $\mathcal{O}_{X, x}$의 가역원과의
곱을 제외하면 무관하다. 이 모든 사실을 모으면 다음 정의가
잘 정의됨을 알 수 있다.

\begin{definition}
\label{divisors-definition-order-vanishing-meromorphic}
$X$를 국소 뇌터 정역 스킴이라 하자.
$\mathcal{L}$을 가역 $\mathcal{O}_X$-가군이라 하자.
$s \in \Gamma(X, \mathcal{K}_X(\mathcal{L}))$를
$\mathcal{L}$의 정칙 유리형 절단이라 하자.
모든 소인자 $Z \subset X$에 대하여
{\it $s$의 $Z$를 따른 소멸 차수}를 정수
$$
\text{ord}_{Z, \mathcal{L}}(s)
= \text{ord}_{\mathcal{O}_{X, \xi}}(s/s_\xi)
$$
로 정의한다. 우변은 Algebra, Definition
\ref{algebra-definition-ord}의 개념이고, $\xi \in Z$는
생성점이며, $s_\xi \in \mathcal{L}_\xi$는 생성원이다.
\end{definition}

\noindent
주인자의 경우와 마찬가지로 다음 보조정리를 얻는다.

\begin{lemma}
\label{divisors-lemma-divisor-meromorphic-locally-finite}
$X$를 국소 뇌터 정역 스킴이라 하자. $\mathcal{L}$을 가역
$\mathcal{O}_X$-가군이라 하자.
$s \in \mathcal{K}_X(\mathcal{L})$를 $\mathcal{L}$의 정칙
(즉, 영이 아닌) 유리형 절단이라 하자. 그러면 두 집합
$$
\{Z \subset X \mid Z \text{ a prime divisor with generic point }\xi
\text{ and }s\text{ not in }\mathcal{L}_\xi\}
$$
과
$$
\{Z \subset X \mid Z \text{ is a prime divisor and }
\text{ord}_{Z, \mathcal{L}}(s) \not = 0\}
$$
은 $X$에서 국소 유한하다.
\end{lemma}

\begin{proof}
공집합이 아닌 열린 부분스킴 $U \subset X$가 존재하여 $s$가
$\Gamma(U, \mathcal{L})$의 절단에 대응하고, 이 절단은
$\mathcal{L}$을 $U$ 위에서 생성한다. 따라서 보조정리의 집합들에
나타날 수 있는 소인자들은 모두 $X \setminus U$의 기약 성분이다.
그러므로 Lemma \ref{divisors-lemma-components-locally-finite}가 원하는
결과를 준다.
\end{proof}

\begin{lemma}
\label{divisors-lemma-divisor-meromorphic-well-defined}
$X$를 국소 뇌터 정역 스킴이라 하자.
$\mathcal{L}$을 가역 $\mathcal{O}_X$-가군이라 하자.
$s, s' \in \mathcal{K}_X(\mathcal{L})$를 $\mathcal{L}$의
영이 아닌 유리형 절단이라 하자. 그러면 $f = s/s'$는
$R(X)^*$의 원소이고
$$
\sum \text{ord}_{Z, \mathcal{L}}(s)[Z]
=
\sum \text{ord}_{Z, \mathcal{L}}(s')[Z]
+
\text{div}(f)
$$
가 베유 인자의 등식으로 성립한다.
\end{lemma}

\begin{proof}
정의들에서 명백히 따른다. Lemma
\ref{divisors-lemma-divisor-meromorphic-locally-finite}에 의해 이 합들이
실제로 베유 인자임을 유의하자.
\end{proof}

\begin{definition}
\label{divisors-definition-divisor-invertible-sheaf}
$X$를 국소 뇌터 정역 스킴이라 하자.
$\mathcal{L}$을 가역 $\mathcal{O}_X$-가군이라 하자.
\begin{enumerate}
\item 임의의 영이 아닌 유리형 절단 $s$가 $\mathcal{L}$의 절단일 때,
{\it $s$에 수반되는 베유 인자}를
$$
\text{div}_\mathcal{L}(s) =
\sum \text{ord}_{Z, \mathcal{L}}(s) [Z] \in \text{Div}(X)
$$
로 정의한다. 여기서 합은 소인자들에 걸친다.
\item {\it $\mathcal{L}$에 수반되는 베유 인자류}를
$\text{div}_\mathcal{L}(s)$의 $\text{Cl}(X)$ 속 상으로 정의한다.
여기서 $s$는 $\mathcal{L}$의 임의의 영이 아닌 유리형 절단이고
$X$ 위의
유리형 절단이다. Lemma
\ref{divisors-lemma-divisor-meromorphic-well-defined}에 의해 이는
잘 정의된다.
\end{enumerate}
\end{definition}

\noindent
예상대로 이 구성은 가역 가군에 대하여 가법적이다.

\begin{lemma}
\label{divisors-lemma-c1-additive}
$X$를 국소 뇌터 정역 스킴이라 하자.
$\mathcal{L}$, $\mathcal{N}$을 가역 $\mathcal{O}_X$-가군이라 하자.
$s$와 각각 $t$를 $\mathcal{L}$과 각각 $\mathcal{N}$의
영이 아닌 유리형 절단이라 하자. 그러면 $st$는
$\mathcal{L} \otimes \mathcal{N}$의 영이 아닌 유리형 절단이고
$$
\text{div}_{\mathcal{L} \otimes \mathcal{N}}(st)
=
\text{div}_\mathcal{L}(s) + \text{div}_\mathcal{N}(t)
$$
가 $\text{Div}(X)$에서 성립한다. 특히
$\mathcal{L} \otimes_{\mathcal{O}_X} \mathcal{N}$의 베유 인자류는
$\mathcal{L}$과 $\mathcal{N}$의 베유 인자류들의 합이다.
\end{lemma}

\begin{proof}
$s$와 각각 $t$를 $\mathcal{L}$과 각각 $\mathcal{N}$의
영이 아닌 유리형 절단이라 하자. 그러면 $st$는
$\mathcal{L} \otimes \mathcal{N}$의 영이 아닌 유리형 절단이다.
$Z \subset X$를 소인자라 하자. $\xi \in Z$를 그 생성점이라 하자.
생성원 $s_\xi \in \mathcal{L}_\xi$와
$t_\xi \in \mathcal{N}_\xi$를 고르자. 그러면 $s_\xi t_\xi$는
$(\mathcal{L} \otimes \mathcal{N})_\xi$의 생성원이다.
따라서 $st/(s_\xi t_\xi) = (s/s_\xi)(t/t_\xi)$이다. 그러므로
$$
\text{ord}_{Z, \mathcal{L} \otimes \mathcal{N}}(st)
=
\text{ord}_{Z, \mathcal{L}}(s) + \text{ord}_{Z, \mathcal{N}}(t)
$$
이며, 이는 $\text{ord}_Z$ 함수의 가법성에서 따른다.
\end{proof}

\noindent
이 방식으로 가역 가군에 그 베유 인자류를 대응시키는 아벨 군 준동형
\begin{equation}
\label{divisors-equation-c1}
\Pic(X) \longrightarrow \text{Cl}(X)
\end{equation}
을 얻는다.

\begin{lemma}
\label{divisors-lemma-normal-c1-injective}
$X$를 국소 뇌터 정역 스킴이라 하자. $X$가 정규이면 사상
(\ref{divisors-equation-c1}) $\Pic(X) \to \text{Cl}(X)$은 단사이다.
\end{lemma}

\begin{proof}
$\mathcal{L}$을 가역 $\mathcal{O}_X$-가군이라 하고, 그에 수반되는
베유 인자류가 자명하다고 하자. $s$를 $\mathcal{L}$의 정칙 유리형
절단이라 하자. 가정은
$\text{div}_\mathcal{L}(s) = \text{div}(f)$임을 뜻하며,
여기서 어떤 $f \in R(X)^*$가 존재한다.
그러면 $t = f^{-1}s$는 $\mathcal{L}$의 정칙 유리형 절단이고
$\text{div}_\mathcal{L}(t) = 0$이다. Lemma
\ref{divisors-lemma-divisor-meromorphic-well-defined}를 보라.
$t$가 $\mathcal{L}$의 자명화를 정의함을 보일 것이며, 이것으로
보조정리의 증명이 끝난다. 이를 증명할 때 $X$ 위에서 국소적으로
작업해도 된다. 따라서 $X = \Spec(A)$가 아핀이고
$\mathcal{L}$이 자명하다고 가정해도 된다. 그러면 $A$는 뇌터
정규 정역이고 $t$는 그 분수체의 원소이며,
$\text{ord}_{A_\mathfrak p}(t) = 0$이다. 이 등식은 높이
$1$인 모든 소아이디얼 $\mathfrak p$에 대하여 성립하며, 이들은
$A$의 소아이디얼이다. 목표는 $t$가 $A$의 가역원임을 보이는 것이다.
$A_\mathfrak p$는 $A$의 높이 일 소아이디얼에 대하여
이산 값매김환이므로
(Algebra, Lemma \ref{algebra-lemma-criterion-normal}), 이 조건은
$t \in A_\mathfrak p^*$임을 뜻한다. 이 명제는 모든 소아이디얼
$\mathfrak p$에 대해 성립하며, 그 높이는 $1$이다. 이로부터 Algebra, Lemma
\ref{algebra-lemma-normal-domain-intersection-localizations-height-1}에
의해 $t \in A$이고 $t^{-1} \in A$임이 따른다. 증명이 끝났다.
\end{proof}

\begin{lemma}
\label{divisors-lemma-local-rings-UFD-c1-bijective}
$X$를 국소 뇌터 정역 스킴이라 하자. 사상
(\ref{divisors-equation-c1}) $\Pic(X) \to \text{Cl}(X)$을 생각하자.
다음은 동치이다.
\begin{enumerate}
\item $X$의 국소환들은 유일 인수분해 정역이다.
\item $X$는 정규이고 $\Pic(X) \to \text{Cl}(X)$은 전사이다.
\end{enumerate}
이 경우 $\Pic(X) \to \text{Cl}(X)$은 동형이다.
\end{lemma}

\begin{proof}
(1)이 성립하면 Algebra, Lemma \ref{algebra-lemma-UFD-normal}에 의해
$X$는 정규이다. 따라서 사상 (\ref{divisors-equation-c1})은
Lemma \ref{divisors-lemma-normal-c1-injective}에 의해 단사이다. 더욱이 모든 소인자
$D \subset X$는 Lemma \ref{divisors-lemma-weil-divisor-is-cartier-UFD}에 의해
유효 카르티에 인자이다. 이 경우 표준 절단 $1_D$는
$\mathcal{O}_X(D)$의 절단이며(Definition
\ref{divisors-definition-invertible-sheaf-effective-Cartier-divisor})
$D$를 따라 정확히 소멸한다. 따라서 $D$의 류는
사상 (\ref{divisors-equation-c1})에 의한 $\mathcal{O}_X(D)$의 상이다.
그러므로 이 사상은 전사이기도 하다.

\medskip\noindent
(2)가 성립한다고 하자. 소인자 $D \subset X$를 고르자.
(\ref{divisors-equation-c1})이 전사이므로 가역층 $\mathcal{L}$,
정칙 유리형 절단 $s$, 그리고 $f \in R(X)^*$가 존재하여
$\text{div}_\mathcal{L}(s) + \text{div}(f) = [D]$이다.
다시 말해 $\text{div}_\mathcal{L}(fs) = [D]$이다.
$x \in X$라 하고 $A = \mathcal{O}_{X, x}$라 하자.
그러면 $A$는 분수체 $K = R(X)$를 갖는 뇌터 국소 정규 정역이다.
높이 $1$인 $A$의 모든 소아이디얼은 $X$ 위의 소인자에 대응하고,
모든 가역 $\mathcal{O}_X$-가군을 $\Spec(A)$에 제한하면
자명한 가역 가군을 얻는다. 따라서 높이 $1$인 모든 소아이디얼
$\mathfrak p \subset A$에 대하여 원소 $f \in K$가 존재하여
$\text{ord}_{A_\mathfrak p}(f) = 1$이고,
$\text{ord}_{A_{\mathfrak p'}}(f) = 0$이다. 뒤 등식은 높이 일인
다른 모든 소아이디얼 $\mathfrak p'$에 대하여 성립한다. 그러면 Algebra, Lemma
\ref{algebra-lemma-normal-domain-intersection-localizations-height-1}에
의해 $f \in A$이다. 같은 방식으로 논증하면 모든
$g \in \mathfrak p$는 $g = af$의 꼴임을 알 수 있으며, 여기서
어떤 $a \in A$가 존재한다. 따라서 $A$의 높이 일인 모든 소아이디얼은
주아이디얼이고, Algebra, Lemma
\ref{algebra-lemma-characterize-UFD-height-1}에 의해 $A$는
유일 인수분해 정역이다.
\end{proof}






\section{가역 가군에 관한 추가 결과}
\label{divisors-section-invertible-modules}

\noindent
이 절에서는 가역 가군의 몇 가지 성질을 논한다.

\begin{lemma}
\label{divisors-lemma-in-image-pullback}
$\varphi : X \to Y$를 스킴의 사상이라 하자.
$\mathcal{L}$을 가역 $\mathcal{O}_X$-가군이라 하자.
다음을 가정하자.
\begin{enumerate}
\item $X$는 국소 뇌터이다.
\item $Y$는 국소 뇌터 정규 정역 스킴이다.
\item $\varphi$는 정역인 (따라서 공집합이 아닌) 섬유를 갖는
평탄 사상이다.
\item $\varphi$는 준콤팩트이거나 국소 유한형이다.
\item $\mathcal{L}$을 $\varphi$의 생성 섬유에 제한하면 자명하다.
\end{enumerate}
그러면 $\mathcal{L} \cong \varphi^*\mathcal{N}$이며, 여기서
$\mathcal{O}_Y$-가군 $\mathcal{N}$은 어떤 가역 가군이다.
\end{lemma}

\begin{proof}
$\xi \in Y$를 생성점이라 하자. $X_\xi$를 $\varphi$의
$\xi$ 위 스킴론적 섬유라 하자. $\mathcal{L}_\xi$로
$\mathcal{L}$을 $X_\xi$로 당긴 것을 나타내자. 가정 (5)는
$\mathcal{L}_\xi$가 자명하다는 뜻이다. 자명화 절단
$s \in \Gamma(X_\xi, \mathcal{L}_\xi)$를 고르자.
Lemma \ref{divisors-lemma-flat-over-integral-integral-fibre}에 의해 $X$는
정역 스킴임을 관찰하자. 따라서 $s$를 $\mathcal{L}$의 정칙 유리형
절단으로 생각할 수 있다. Lemma
\ref{divisors-lemma-pullback-meromorphic-sections-defined}에 의해
$\varphi$에 대해 유리형 함수의 당김이 정의된다.
$\mathcal{N} \subset \mathcal{K}_Y$를 다음과 같은
$\mathcal{O}_Y$-가군이라 하자. 열린집합 $V \subset Y$ 위의
그 절단들은 유리형 함수 $g \in \mathcal{K}_Y(V)$ 중
$\varphi^*(g)s \in \mathcal{L}(\varphi^{-1}V)$를 만족하는 것들이다.
선험적으로 $\varphi^*(g)s$는 $\mathcal{K}_X(\mathcal{L})$의
$\varphi^{-1}V$ 위 절단이다. 우리는
$\mathcal{N}$이 가역 $\mathcal{O}_Y$-가군이고 사상
$$
\varphi^*\mathcal{N} \longrightarrow \mathcal{L},\quad
g \longmapsto gs
$$
이 동형이라고 주장한다.

\medskip\noindent
먼저 다음 상황에서 주장을 증명하자. $X$와 $Y$가 아핀이고
$\mathcal{L}$이 자명하다. $Y = \Spec(R)$, $X = \Spec(A)$이고,
$s$가 원소 $s \in A \otimes_R K$로 주어진다고 하자. 여기서
$K$는 $R$의 분수체이다. $s = a/r$로 쓸 수 있으며, 여기서
어떤 영이 아닌 $r \in R$와 $a \in A$가 존재한다.
$s$가 생성 섬유에서 $\mathcal{L}$을 생성하므로 어떤
$s' \in A \otimes_R K$가 존재하여 $ss' = 1$이다.
따라서 $s = r'/a'$로 쓸 수 있으며, 여기서 어떤 영이 아닌
$r' \in R$와 $a' \in A$가 존재한다.
$\mathfrak p_1, \ldots, \mathfrak p_n \subset R$를 $rr'$ 위의
극소 소아이디얼들이라 하자. 각 $R_{\mathfrak p_i}$는 이산
값매김환이다(Algebra, Lemmas \ref{algebra-lemma-minimal-over-1}과
\ref{algebra-lemma-criterion-normal}). 가정에 의해
$\mathfrak q_i = \mathfrak p_i A$는 소아이디얼이다. 따라서
$\mathfrak q_i A_{\mathfrak q_i}$는 하나의 원소로 생성되고,
$A_{\mathfrak q_i}$도 이산 값매김환임을 얻는다
(Algebra, Lemma \ref{algebra-lemma-characterize-dvr}). 물론
$R_{\mathfrak p_i} \to A_{\mathfrak q_i}$의 분기 지수는 $1$이다.
$e_i, e'_i \geq 0$을 $a, a'$의 $A_{\mathfrak q_i}$에서의
값매김이라 하자. 그러면 $e_i + e'_i$는 $rr'$의
$R_{\mathfrak p_i}$에서의 값매김이다. Algebra, Lemma
\ref{algebra-lemma-normal-domain-intersection-localizations-height-1}에
의해 $R$에서
$$
\mathfrak p_1^{(e_1 + e'_1)} \cap \ldots \cap \mathfrak p_i^{(e_n + e'_n)} =
(rr')
$$
이다. 다음과 같이 놓자.
$$
I = \mathfrak p_1^{(e_1)} \cap \ldots \cap \mathfrak p_i^{(e_n)}
\quad\text{and}\quad
I' = \mathfrak p_1^{(e'_1)} \cap \ldots \cap \mathfrak p_i^{(e'_n)}
$$
그러면 $II' \subset (rr')$이다. Algebra, Lemmas
\ref{algebra-lemma-symbolic-power-flat-extension}와
\ref{algebra-lemma-flat-intersect-ideals}에 의해
$$
IA =
(\mathfrak p_1^{(e_1)} \cap \ldots \cap \mathfrak p_i^{(e_n)})A =
(\mathfrak p_1A)^{(e_1)} \cap \ldots \cap (\mathfrak p_i A)^{(e_n)}
$$
임을 관찰하자. $I'A$에 대해서도 마찬가지이다. 따라서
$a \in IA$이고 $a' \in I'A$이다.
이로부터 $IA \otimes_A I'A \to rr'A$가 전사임을 얻는다.
$R \to A$의 충실한 평탄성에 의해
$I \otimes_R I' \to (rr')$도 전사이다. 따라서
$II' = (rr')$이고 $I$와 $I'$은 계수 $1$인 유한 국소 자유
가군이다. Algebra, Lemma
\ref{algebra-lemma-product-ideals-principal}을 보라.
그러므로 $R$ 위에서 자리스키 국소적으로
$I = (g)$와 $I' = (g')$로 쓸 수 있고 $gg' = rr'$이다.
그러면 $a = ug$이고 $a' = u'g'$이며, 여기서 어떤
$u, u' \in A$가 존재한다. 따라서 $u, u'$는 가역원이다. 그러므로
$R$ 위에서 자리스키 국소적으로 $s = ug/r$이고, 이 경우의 주장이
따른다.

\medskip\noindent
$y \in Y$를 점이라 하자. $x \in X$를 고르되 $y$로 사상되게 하자.
앞 문단의 결과를
$\Spec(\mathcal{O}_{X, x}) \to \Spec(\mathcal{O}_{Y, y})$에
적용할 수 있다. 이에 따라 원소 $g \in R(Y)^*$가 존재하며,
이는 $\mathcal{O}_{Y, y}^*$의 원소와의 곱을 제외하면 잘 정의되고,
$\varphi^*(g)s$가 $\mathcal{L}_x$를 생성한다. 따라서
$\varphi^*(g)s$는 $\mathcal{L}$을 근방 $U$에서 생성하며,
이는 $x$의 근방이다. $x'$을 $y$ 위에 놓인 두 번째 점이라 하고,
$g' \in R(Y)^*$가 $\varphi^*(g')s$가 $\mathcal{L}$을 열린 근방
$U'$에서 생성하도록 한다고 하자. 여기서 이 근방은 $x'$의 근방이다.
섬유가 기약이므로 점 $x''$을
$U \cap U' \cap \varphi^{-1}(\{y\})$ 안에서 고를 수 있다.
환 사상 $\mathcal{O}_{Y, y} \to \mathcal{O}_{X, x''}$에 대한
유일성에 의해 $g$와 $g'$는 $\mathcal{O}_{Y, y}^*$의 원소와의
곱만큼 다르다. 따라서 $\varphi^*(g)s$는 $\mathcal{L}$의
생성원이며, 이는 $\varphi^{-1}(y)$의 열린 근방에서 성립한다.
$Z \subset X$를 점 $z \in X$들 중 $\varphi^*(g)s$가
$\mathcal{L}_z$를 생성하지 않는 것들의 집합이라 하자.
위의 논증에 의해 $Z$는 닫혀 있고 $Z = \varphi^{-1}(T)$이며,
여기서 어떤 부분집합 $T \subset Y$가 존재하고 $y \not \in T$이다.
$T$가 닫혀 있음을 보일 수 있다면 $g$는 $\mathcal{N}$의
$\mathcal{O}_Y$-가군 생성원이며, 이는 열린 근방
$Y \setminus T$에서 성립하고 이 근방은 $y$를 포함한다. 이로써
증명이 끝난다(일부 세부사항은 생략한다).

\medskip\noindent
$\varphi$가 준콤팩트이면 Morphisms, Lemma
\ref{morphisms-lemma-fpqc-quotient-topology}에 의해 $T$는 닫혀 있다.
$\varphi$가 국소 유한형이면 Morphisms, Lemma
\ref{morphisms-lemma-fppf-open}에 의해 $\varphi$는 열린 사상이다.
그러면 $Y \setminus T$는 열린집합 $X \setminus Z$의 상이므로
열려 있다.
\end{proof}

\begin{lemma}
\label{divisors-lemma-closure-effective-cartier-divisor}
$X$를 국소 뇌터 스킴이라 하자. $U \subset X$를 열린집합이라 하고
$D \subset U$를 유효 카르티에 인자라 하자. 모든
$\mathcal{O}_{X, x}$가 유일 인수분해 정역이라고 하자. 여기서
$x \in X \setminus U$이다. 그러면 유효 카르티에 인자
$D' \subset X$가 존재하여 $D = U \cap D'$이다.
\end{lemma}

\begin{proof}
$D' \subset X$를 사상 $D \to X$의 스킴론적 상이라 하자.
$X$가 국소 뇌터이므로 사상 $D \to X$는 준콤팩트이다.
Properties, Lemma \ref{properties-lemma-immersion-into-noetherian}를
보라. 따라서 Morphisms, Lemma
\ref{morphisms-lemma-quasi-compact-scheme-theoretic-image}에 의해
$D'$의 구성은 $X$의 열린집합으로 옮기는 것과 가환한다.
그러므로 $X = \Spec(A)$가 아핀이라고 가정해도 된다.
$I \subset A$를 $D'$에 대응하는 아이디얼이라 하자.
$\mathfrak p \subset A$를 $X \setminus U$의 한 점에 대응하는
소아이디얼이라 하자. 증명을 끝내려면 $I_\mathfrak p$가 하나의
원소로 생성됨을 보이면 충분하다. Lemma
\ref{divisors-lemma-effective-Cartier-in-points}를 보라. 따라서 Morphisms,
Lemma \ref{morphisms-lemma-flat-base-change-scheme-theoretic-image}에
의해 $X$를 $\Spec(A_\mathfrak p)$로 바꾸어도 된다. 다시 말해
$X$가 국소 유일 인수분해 정역 $A$의 스펙트럼이라고 가정해도 된다.
그러면 $A$의 모든 국소환은 유일 인수분해 정역이다. 따라서
$D = \sum a_i D_i$이며, 여기서 $D_i \subset U$는 정역 유효
카르티에 인자이다. Lemma
\ref{divisors-lemma-effective-Cartier-divisor-is-a-sum}을 보라.
생성점 $\xi_i$는 $D_i$의 생성점이며, 소아이디얼
$\mathfrak p_i \subset A$에 대응하고 그 높이는 $1$이다. Lemma
\ref{divisors-lemma-effective-Cartier-codimension-1}을 보라.
그러면 $\mathfrak p_i = (f_i)$이고, 여기서 어떤 소원소
$f_i \in A$가 존재한다. 따라서 원하는 대로 $D'$는
$\prod f_i^{a_i}$에 의해 잘려 나온다.
\end{proof}

\begin{lemma}
\label{divisors-lemma-extend-invertible-module}
$X$를 국소 뇌터 스킴이라 하자. $U \subset X$를 열린집합이라 하고
$\mathcal{L}$을 가역 $\mathcal{O}_U$-가군이라 하자.
모든 $\mathcal{O}_{X, x}$가 유일 인수분해 정역이라고 하자.
여기서 $x \in X \setminus U$이다. 그러면 가역
$\mathcal{O}_X$-가군 $\mathcal{L}'$이 존재하여
$\mathcal{L} \cong \mathcal{L}'|_U$이다.
\end{lemma}

\begin{proof}
$x \in X$, $x \not \in U$를 고르자. 어떤 아핀 열린 근방
$W \subset X$가 존재하여 $\mathcal{L}|_{W \cap U}$가 $W$ 위의
가역층으로 확장됨을 보일 것이다. 층의 이어붙이기
(Sheaves, Section \ref{sheaves-section-glueing-sheaves})에 의해
이는 $\mathcal{L}$을 엄밀히 더 큰 열린집합 $U \cup W$로
확장할 수 있음을 뜻한다. $W = \Spec(A)$를 아핀 열린 근방이라 하자.
$U \cap W$가 준아핀이므로,
$\mathcal{L}|_{W \cap U}$를
$\mathcal{O}(D_1) \otimes \mathcal{O}(D_2)^{\otimes -1}$로 쓸 수 있다.
여기서 어떤 유효 카르티에 인자 $D_1, D_2 \subset W \cap U$가 존재한다.
Lemma \ref{divisors-lemma-quasi-projective-Noetherian-pic-effective-Cartier}를
보라. 그러면 Lemma \ref{divisors-lemma-closure-effective-cartier-divisor}에
의해 $D_1$과 $D_2$를 $W$의 유효 카르티에 인자로 확장할 수 있고,
이것으로 가역층의 확장을 얻는다.

\medskip\noindent
$X$가 뇌터이면(실제 응용에서 가장 자주 쓰는 경우이다) 위 결과와
뇌터 귀납법을 결합하여 증명이 끝난다. 일반적인 경우에는 다음과
같이 논증한다. 먼저 $U$ 바깥의 점들의 모든 국소환이 정역이고
$X$가 국소 뇌터이므로, $U$의 $X$에서의 폐포는 열린집합이다.
따라서 $U \subset X$가 조밀하고 스킴론적으로 조밀하다고
가정해도 된다. 이제 집합 $T$를 생각하자. 그 원소는 삼중항
$(U', \mathcal{L}', \alpha)$들이다. 여기서 $U \subset U' \subset X$는
열린 부분스킴이고, $\mathcal{L}'$은 가역
$\mathcal{O}_{U'}$-가군이며,
$\alpha : \mathcal{L}'|_U \to \mathcal{L}$는 동형이다.
$T$에 다음 규칙으로 정의되는 부분순서 $\leq$를 주자.
$(U', \mathcal{L}', \alpha) \leq (U'', \mathcal{L}'', \alpha')$인
것은 $U' \subset U''$이고 동형
$\beta : \mathcal{L}''|_{U'} \to \mathcal{L}'$이 존재하여
$\alpha$ 및 $\alpha'$와 양립하는 것과 동치이다.
$\beta$는 존재한다면 유일한데, 이는 $U \subset X$가 조밀하기
때문이다.
증명의 첫 부분에 의해 임의의 원소
$t = (U', \mathcal{L}', \alpha)$가 $T$에 속하고
$U' \not = X$를 만족하면 $t' \in T$가 존재하여 $t' > t$이다.
따라서 증명을 끝내려면 초른의 보조정리를 적용할 수 있음을 보이면
충분하다. 전순서 부분집합 $I \subset T$를 생각하자.
$i \in I$가 삼중항 $(U_i, \mathcal{L}_i, \alpha_i)$에 대응하면,
가역 가군 $\mathcal{L}'$을 $U' = \bigcup U_i$ 위에서 다음과 같이
구성할 수 있다. 열린 준콤팩트 집합 $W \subset U'$에 대하여
$W \subset U_i$인 어떤 $i$가 존재하므로
$$
\mathcal{L}'(W) = \mathcal{L}_i(W)
$$
로 놓는다. 전이 사상에는 $\beta$들을 사용한다(이들은 유일하므로
올바르게 합성된다). 이로써 가역 $\mathcal{O}$-가군
$\mathcal{L}'$이 $U'$의 준콤팩트 열린집합들의 기저 위에 정의되고, 이는
가역 가군을 정의하기에 충분하다(Sheaves, Section
\ref{sheaves-section-bases}). 세부사항은 생략한다.
\end{proof}

\begin{lemma}
\label{divisors-lemma-open-subscheme-UFD}
$R$을 유일 인수분해 정역이라 하자. 다음 공간들의 피카르 군은
자명하다.
\begin{enumerate}
\item $\Spec(R)$과 그 임의의 열린 부분스킴.
\item $\mathbf{A}^n_R = \Spec(R[x_1, \ldots, x_n])$과 그 임의의
열린 부분스킴.
\end{enumerate}
특히 아핀 $n$-공간 $\mathbf{A}^n_k$의 임의의 열린 부분스킴의
피카르 군은 자명하며, 여기서 기초체는 체 $k$이다.
\end{lemma}

\begin{proof}
Algebra, Lemma \ref{algebra-lemma-polynomial-ring-UFD}에 의해
$R$ 위의 모든 다항식환은 유일 인수분해 정역이다. 따라서
$\Pic(U)$가 자명함을 보이면 된다. 여기서
$U \subset \Spec(R)$는 (공집합이 아닌) 열린집합이다.

\medskip\noindent
뇌터 경우에는 Lemma \ref{divisors-lemma-extend-invertible-module}을 사용하여
$\Pic(\Spec(R))$이 자명함을 보이는 것으로 줄일 수 있다.
이는 다시 More on Algebra, Lemma
\ref{more-algebra-lemma-UFD-Pic-trivial}에서 증명된다.
독자는 증명의 나머지를 건너뛰어도 좋다.

\medskip\noindent
$U \subset \Spec(R)$를 공집합이 아닌 열린집합이라 하자.
영이 아닌 $f \in R$을 고르되 $D(f) \subset U$이게 하자.
$f$는 $R$의 소원소들의 곱이므로, 유한 개의 소원소
$a_1, \ldots, a_n \in R$가 존재하여 $(a_i) \not \in U$이다
(또는 동치로 $U \cap V(a_i) = \emptyset$이다).
$R$을 $R_{a_1 \ldots a_n}$으로 바꾸면 $U$가 모든 주
소아이디얼을 포함한다고 가정해도 된다. (유일 인수분해 정역의
국소화는 유일 인수분해 정역이다. 예를 들어 Algebra, Lemma
\ref{algebra-lemma-invert-prime-elements}를 보라.)

\medskip\noindent
$U$가 모든 주 소아이디얼을 포함한다고 가정하자.
$U = \Spec(R) \setminus V(I)$로 쓰자. 여기서 어떤 아이디얼
$I \subset R$가 존재한다. 영이 아닌 $f \in I$를 고르고
$f = a_1 \ldots a_n$을 소원소들의 곱으로 쓰자.
$(a_i) \not \in V(I)$이므로 $g \in I$를 고르되
$g \not \in (a_i)$가 모든 $i$에 대하여 성립하게 할 수 있다
(소아이디얼 회피).
그림으로 나타내면
$$
U' = D(f) \cup D(g) \subset U \subset \Spec(R)
$$
이다. More on Algebra, Lemma
\ref{more-algebra-lemma-UFD-Pic-trivial}에 의해 $D(f)$와
$D(g)$의 피카르 군은 자명하다. 또한 유일 인수분해 정역에 관한
초등적인 논증에 의해 사상
$R_f^* \times R_g^* \to R_{fg}^*$,
$(u, v) \mapsto uv^{-1}$은 전사이다. 따라서
$\Pic(U')$는 자명하다.

\medskip\noindent
그러므로 가역 가군 $\mathcal{L}$ 중 $U$ 위에 있고
$j : U' \to U'$에 의한 당김이 자명한 것은 자명함을 보이면 충분하다.
이를 위해 사상 $\mathcal{L} \to j_*j^*\mathcal{L}$이
동형임을 보이면 충분하다. 이를 위해 다시
$u \in U \setminus U'$라 하자. 그러면
$f, g \in \mathfrak m_u \subset \mathcal{O}_{U, u} = A$는
공통 소인수를 갖지 않는 국소 유일 인수분해 정역의 원소들이다.
소인수들을 살펴보면 이것이
$$
0 \to A \to A_f \times A_g \to A_{fg}
$$
의 완전성과 동치임을 독자가 쉽게 보일 수 있다.
$A = \mathcal{L}_u$인데, 이는 $\mathcal{L}$이 가역이기 때문이며,
$(j_*j^*\mathcal{L})_u$는 $A_f \times A_g \to A_{fg}$의
핵이므로 결론이 따른다.
\end{proof}

\begin{lemma}
\label{divisors-lemma-Pic-projective-space-UFD}
$R$을 유일 인수분해 정역이라 하자.
$\mathbf{P}^n_R$의 피카르 군은 $\mathbf{Z}$이다.
더 정확히 말하면 동형
$$
\mathbf{Z} \longrightarrow \Pic(\mathbf{P}^n_R),\quad
m \longmapsto \mathcal{O}_{\mathbf{P}^n_R}(m)
$$
이 존재한다. 특히 사영공간 $\mathbf{P}^n_k$는 체 $k$ 위에 있으며,
그 피카르 군은 $\mathbf{Z}$이다.
\end{lemma}

\begin{proof}
Lemma \ref{divisors-lemma-open-subscheme-UFD}에 의해 열린집합
$D_+(T_i) \cong \mathbf{A}^n_R$의 피카르 군들은 자명하다.
따라서 $\mathcal{L}$이 $\mathbf{P}^n_R$ 위의 가역 가군이면,
열린 덮개
$$
\mathbf{P}^n_R = \bigcup\nolimits_{i = 0, \ldots, n} D_+(T_i)
$$
에 대한 $1$-쌍대순환으로 주어지며, 그 값은 가역 함수층
$\mathcal{O}^*$에 있다.
$i \not = j$이면
$$
\mathcal{O}^*(D_+(T_i) \cap D_+(T_j)) =
\mathcal{O}^*(D_+(T_iT_j)) = \left(R[T_0, \ldots, T_n]_{(T_iT_j)}\right)^* =
R^* \times (T_i/T_j)^\mathbf{Z}
$$
임을 관찰하자. 따라서 그러한 쌍대순환 $(g_{ij})$는 가역원
$$
g_{ij} = u_{ij} (T_i/T_j)^{e_{ij}}
$$
들로 주어진다. 여기서 $u_{ij} \in R^*$이고
$e_{ij} \in \mathbf{Z}$이며 쌍대순환 조건을 만족한다.
$D_+(T_iT_jT_k)$ 위의 쌍대순환 조건은
$\#\{i, j, k\} = 3$일 때
$$
u_{ik} = u_{ij} u_{jk} \quad\text{and}\quad
e_{ik} = e_{ij} = e_{jk}
$$
임을 말해 준다. 따라서
$u_{ij} = u_{i1} u_{j1}^{-1}$은 쌍대경계이다. 그러므로 가역
가군의 모든 동형류는 어떤 $e \in \mathbf{Z}$에 대하여
$$
g_{ij} = (T_i/T_j)^e
$$
인 쌍대순환을 택함으로써 얻어진다.
$\mathcal{O}(n)$은 자명화 절단 $T_i^n$을 $D_+(T_i)$ 위에서
가지므로, $\mathcal{O}(n)$에 대응하는 쌍대순환은
$(T_i/T_j)^n$이고 증명이 끝난다.
\end{proof}






\section{정규 스킴 위의 베유 인자}
\label{divisors-section-weil-divisors-normal}

\noindent
먼저 반사적 가군의 성질들을 논한다.

\begin{lemma}
\label{divisors-lemma-reflexive-normal}
$X$를 정역 국소 뇌터 정규 스킴이라 하자.
$\mathcal{F}$와 $\mathcal{G}$가 응집 반사적
$\mathcal{O}_X$-가군이면 사상
$$
(\SheafHom_{\mathcal{O}_X}(\mathcal{F}, \mathcal{O}_X)
\otimes_{\mathcal{O}_X} \mathcal{G})^{**} \to
\SheafHom_{\mathcal{O}_X}(\mathcal{F}, \mathcal{G})
$$
은 동형이다. 규칙 $\mathcal{F}, \mathcal{G} \mapsto
(\mathcal{F} \otimes_{\mathcal{O}_X} \mathcal{G})^{**}$은
계수 $1$인 응집 반사적 $\mathcal{O}_X$-가군의 동형류 집합 위에
아벨 군 법칙을 정의한다.
\end{lemma}

\begin{proof}
엄밀히 필요하지는 않지만, 증명을 계속하기 전에 Remark
\ref{divisors-remark-tensor}를 읽기를 권한다. 열린 부분스킴
$j : U \to X$를 고르되, $X \setminus U$의 모든 기약 성분이
여차원 $\geq 2$를 $X$에서 갖고 $j^*\mathcal{F}$와
$j^*\mathcal{G}$가 유한 국소 자유 가군이 되게 하자.
Lemma \ref{divisors-lemma-reflexive-over-normal}을 보라. 사상
$$
\SheafHom_{\mathcal{O}_U}(j^*\mathcal{F}, \mathcal{O}_U)
\otimes_{\mathcal{O}_U} j^*\mathcal{G} \to
\SheafHom_{\mathcal{O}_U}(j^*\mathcal{F}, j^*\mathcal{G})
$$
은 동형이다. 실제로 이를 국소적으로 확인해도 되고, 가군들이
유한 자유이면 명백하다. 보조정리에 표시된 화살에 $j^*$를
적용하면 방금 동형임을 보인 화살을 얻는다(작은 세부사항은
생략한다). Lemma \ref{divisors-lemma-reflexive-S2-extend}와 세르의 판정법
Properties, Lemma \ref{properties-lemma-criterion-normal}에 의해
$j^*$는 $U$ 위의 응집 반사적 가군과 $X$ 위의 응집 반사적
가군 사이의 범주 동치를 정의한다. 따라서 보조정리의 화살도
동형이다. $\mathcal{F}$의 계수가 $1$이면 $j^*\mathcal{F}$는
가역 $\mathcal{O}_U$-가군이고, 반사적 가군
$\mathcal{F}^\vee = \SheafHom(\mathcal{F}, \mathcal{O}_X)$를
제한하면 그 역원을 얻는다. 앞과 같은 방식으로
$(\mathcal{F} \otimes_{\mathcal{O}_X} \mathcal{F}^\vee)^{**} = \mathcal{O}_X$
가 따른다. 이로써 보조정리의 명제에 주어진 군 법칙의 역원이
존재함을 알 수 있다.
\end{proof}

\begin{lemma}
\label{divisors-lemma-normal-class-group}
$X$를 정역 국소 뇌터 정규 스킴이라 하자. 계수 $1$인 응집 반사적
$\mathcal{O}_X$-가군들의 군은 베유 인자류군
$\text{Cl}(X)$와 동형이며, 이는 $X$의 인자류군이다.
\end{lemma}

\begin{proof}
$\mathcal{F}$를 계수 $1$인 응집 반사적 $\mathcal{O}_X$-가군이라
하자. 열린집합 $U \subset X$를 고르되, $X \setminus U$의 모든
기약 성분이 여차원 $\geq 2$를 $X$에서 갖고
$\mathcal{F}|_U$가 가역이게 하자. Lemma
\ref{divisors-lemma-reflexive-over-normal}을 보라.
$\text{Cl}(U) = \text{Cl}(X)$임을 관찰하자. 실제로 $X$의
베유 인자류군은 그 유리 함수체와 여차원 $1$인 점들 및
그 국소환들에만 의존한다.
따라서 $\mathcal{F}$의 베유 인자류를
$\mathcal{F}|_U$의 $\text{Cl}(U)$ 안 베유 인자류로 정의할 수
있다. 이것이 $U$의 선택과 무관하다는 확인은 생략한다.

\medskip\noindent
$\text{Cl}'(X)$로 계수 $1$인 응집 반사적 $\mathcal{O}_X$-가군의
동형류 집합을 나타내자. 위 구성은 군 준동형
$$
\text{Cl}'(X) \longrightarrow \text{Cl}(X)
$$
을 준다. 실제로 임의의 원소 쌍 $\mathcal{F}, \mathcal{G}$가
$\text{Cl}'(X)$에 속할 때 둘 모두에 적용되는 $U$를
고를 수 있고, 가역 가군에 그 베유 인자류를 대응시키는
(\ref{divisors-equation-c1})의 대응은 준동형이다.
$\mathcal{F}$가 이 사상의 핵에 속하면
$\mathcal{F}|_U$는 자명하고(Lemma
\ref{divisors-lemma-normal-c1-injective}), 따라서 Lemma
\ref{divisors-lemma-reflexive-S2-extend}와 세르의 판정법
Properties, Lemma \ref{properties-lemma-criterion-normal}에 의해
$\mathcal{F}$도 자명하다. 증명을 끝내려면 이 사상이 전사임을
확인하면 충분하다.

\medskip\noindent
$D = \sum n_Z Z$를 $X$ 위의 베유 인자라 하자.
열린집합 $U \subset X$가 존재하여 $X \setminus U$의 모든
기약 성분이 여차원 $\geq 2$를 $X$에서 갖고,
$Z|_U$가 $n_Z \not = 0$인 경우 유효 카르티에 인자가 됨을
주장한다. 이 주장을 증명할 때 $X$가 아핀이라고 가정해도 된다.
그러면 $D = n_1 Z_1 + \ldots + n_r Z_r$가 유한합이고
$Z_1, \ldots, Z_r$가 서로 다르다고 가정해도 된다.
$Z_i \cap Z_j$들을 $i \not = j$인 경우 제거하면
$Z_1, \ldots, Z_r$가 쌍마다 서로소라고 가정해도 된다.
이로써 하나의 소인자 $Z$가 $X$ 위에 있는 경우로 환원된다.
Properties, Lemma \ref{properties-lemma-criterion-normal}에 의해
$X$는 $(R_1)$을 만족하므로, 국소환
$\mathcal{O}_{X, \xi}$는 이산 값매김환이다. 여기서 $\xi$는
$Z$의 생성점이다.
$f \in \mathcal{O}_{X, \xi}$를 균일화원이라 하자.
$V \subset X$를 $\xi$의 열린 근방이라 하고, $f$가 어떤 원소
$f \in \mathcal{O}_X(V)$의 상이 되게 하자. $V$를 축소하여
스킴론적으로 $Z \cap V = V(f)$라고 가정해도 된다. 실제로 이는
$\xi$에서의 국소환에서 성립한다. 이 경우
$$
U = X \setminus (Z \setminus V) = (X \setminus Z) \cup V
$$
로 두면 원하는 열린집합을 얻고, 이로써 주장이 증명된다.

\medskip\noindent
$D$의 인자류가 상에 속함을 보이기 위해
$D = \sum_{n_Z < 0} n_Z Z - \sum_{n_Z > 0} (-n_Z) Z$로 쓸 수 있다.
위에서 구성한 사상의 가법성에 의해 $n_Z \leq 0$이라고
가정해도 되고 실제로 그렇게 하자. 이는 모든 소인자 $Z$에
대하여 성립한다
(독자가 원한다면 이 단계를 피할 수 있다).
$U \subset X$를 위 주장과 같이 고르자. $U$가 준콤팩트이면
$D|_U = -n_1 Z_1 - \ldots - n_r Z_r$로 쓴다. 여기서 소인자
$Z_i$들은 쌍마다 서로 다르고 $n_i > 0$이다. 그리고 가역
$\mathcal{O}_U$-가군
$$
\mathcal{L} =
\mathcal{I}_1^{n_1} \ldots \mathcal{I}_r^{n_r} \subset \mathcal{O}_U
$$
을 생각한다. 여기서 $\mathcal{I}_i$는 $Z_i$의 아이디얼층이다.
$U$의 선택과 Lemma \ref{divisors-lemma-sum-effective-Cartier-divisors}에 의해
이는 가역이다. 또한 $\text{div}_\mathcal{L}(1) = D|_U$이다.
Lemma \ref{divisors-lemma-reflexive-S2-extend}에 의해
$\mathcal{L} = \mathcal{F}|_U$이고, 여기서 어떤 계수 $1$인
응집 반사적 $\mathcal{O}_X$-가군 $\mathcal{F}$가 존재하므로
$D$가 우리 사상의 상에
속함을 얻는다.

\medskip\noindent
$U$가 준콤팩트가 아니면 위에 표시한 식으로
$\mathcal{L} \subset \mathcal{O}_U$를 국소적으로 정의한다.
이 구성이 이어붙여짐을 독자가 보이면 앞과 정확히 같은 방식으로
증명이 끝난다. 세부사항은 생략한다.
\end{proof}

\begin{lemma}
\label{divisors-lemma-structure-sheaf-Xs}
$X$를 정역 국소 뇌터 정규 스킴이라 하자.
$\mathcal{F}$를 계수 1인 응집 반사적 $\mathcal{O}_X$-가군이라 하자.
$s \in \Gamma(X, \mathcal{F})$라 하자. 다음과 같이 놓자.
$$
U = \{x \in X \mid s : \mathcal{O}_{X, x} \to \mathcal{F}_x
\text{ is an isomorphism}\}
$$
그러면 $j : U \to X$는 $X$의 열린 부분스킴이고
$$
j_*\mathcal{O}_U =
\colim (\mathcal{O}_X \xrightarrow{s} \mathcal{F}
\xrightarrow{s} \mathcal{F}^{[2]}
\xrightarrow{s} \mathcal{F}^{[3]}
\xrightarrow{s} \ldots)
$$
이다. 여기서 $\mathcal{F}^{[1]} = \mathcal{F}$이고, 귀납적으로
$\mathcal{F}^{[n + 1]} =
(\mathcal{F} \otimes_{\mathcal{O}_X} \mathcal{F}^{[n]})^{**}$이다.
\end{lemma}

\begin{proof}
Modules, Lemmas \ref{modules-lemma-finite-type-surjective-on-stalk}와
\ref{modules-lemma-finite-type-to-coherent-injective-on-stalk}에 의해
집합 $U$는 열려 있다. Properties, Lemma
\ref{properties-lemma-immersion-into-noetherian}에 의해 $j$가
준콤팩트임을 관찰하자. 마지막 명제를 증명하려면 모든 준콤팩트
열린집합 $W \subset X$에 대하여 제한 사상과 양립하는
$\mathcal{O}_X(W)$-가군의 동형
$$
\colim \Gamma(W, \mathcal{F}^{[n]})
\longrightarrow
\Gamma(U \cap W, \mathcal{O}_U)
$$
이 존재함을 보이면 충분하다. 양립성의 확인은 생략한다.
$X$를 $W$로 바꾸고 위 식을 준동형들의 말로 다시 쓰면 동형
$$
\colim \Hom_{\mathcal{O}_X}(\mathcal{O}_X, \mathcal{F}^{[n]})
\longrightarrow
\Hom_{\mathcal{O}_U}(\mathcal{O}_U, \mathcal{O}_U)
$$
을 구성하면 충분함을 알 수 있다. 열린집합 $V \subset X$를
고르되, $X \setminus V$의 모든 기약 성분이 여차원
$\geq 2$를 $X$에서 갖고 $\mathcal{F}|_V$가 가역이게 하자.
Lemma \ref{divisors-lemma-reflexive-over-normal}을 보라. 그러면 제한은
계수 $1$인 응집 반사적 가군들 중 $X$와 $V$ 위의 것들 사이,
그리고 계수 $1$인 응집 반사적 가군들 중 $U$와
$V \cap U$ 위의 것들 사이에
범주 동치를 정의한다. Lemma \ref{divisors-lemma-reflexive-S2-extend}와
세르의 판정법 Properties, Lemma
\ref{properties-lemma-criterion-normal}을 보라. 따라서 동형
$$
\colim \Gamma(V, (\mathcal{F}|_V)^{\otimes n}) \longrightarrow
\Gamma(V \cap U, \mathcal{O}_U)
$$
을 구성하면 충분하다. $\mathcal{F}|_V$는 가역이고,
$U \cap V$는 $s|_V$가 이 가역 가군을 생성하는 점들의 집합과
같으므로, 이는 Properties, Lemma
\ref{properties-lemma-invert-s-sections}의 특별한 경우이다
(사상에 대한 명시적인 식도 있다).
\end{proof}

\begin{lemma}
\label{divisors-lemma-Xs-codim-complement}
Lemma \ref{divisors-lemma-structure-sheaf-Xs}와 같은 가정과 표기를 쓰자.
$s$가 영이 아니면 $X \setminus U$의 모든 기약 성분은
여차원 $1$을 $X$에서 갖는다.
\end{lemma}

\begin{proof}
$\xi \in X$를 한 기약 성분 $Z$의 생성점이라 하자. 이 성분은
$X \setminus U$의 기약 성분이다.
$X$를 $\xi$의 열린 근방으로 바꾸면 $Z = X \setminus U$가
기약이라고 가정해도 된다.
$s : \mathcal{O}_U \to \mathcal{F}|_U$가 동형이므로
$Z$의 $X$에서의 여차원이 $\geq 2$이면
$s : \mathcal{O}_X \to \mathcal{F}$는 동형이다. 이는 Lemma
\ref{divisors-lemma-reflexive-S2-extend}와 세르의 판정법
Properties, Lemma \ref{properties-lemma-criterion-normal}에 따른다.
그러면 $Z = \emptyset$이어야 하는데, 이는 모순이다.
\end{proof}

\begin{remark}
\label{divisors-remark-structure-sheaf-Xs}
$A$를 뇌터 정규 정역이라 하자. $M$을 계수 $1$인 유한 반사적
$A$-가군이라 하자. 영이 아닌 $s \in M$을 택하자.
$\mathfrak p_1, \ldots, \mathfrak p_r$를 높이 $1$인 $A$의
소아이디얼들로서 $M/As$의 지지집합에 속하는 것들이라 하자. 그러면 Lemma
\ref{divisors-lemma-structure-sheaf-Xs}의 열린집합 $U$는 Lemma
\ref{divisors-lemma-Xs-codim-complement}에 의해
$$
U = \Spec(A) \setminus
\left(V(\mathfrak p_1) \cup \ldots \cup V(\mathfrak p_r)\right)
$$
이다. 더욱이 $M^{[n]}$이
$M \otimes_A \ldots \otimes_A M$의 반사적 포락을 나타내고
($n$개의 인자), 그러면 Lemma \ref{divisors-lemma-structure-sheaf-Xs}에 따라
$$
\Gamma(U, \mathcal{O}_U) = \colim M^{[n]}
$$
이다.
\end{remark}

\begin{lemma}
\label{divisors-lemma-affine-Xs}
Lemma \ref{divisors-lemma-structure-sheaf-Xs}와 같은 가정과 표기를 쓰자.
다음은 동치이다.
\begin{enumerate}
\item 포함 사상 $j : U \to X$는 아핀이다.
\item 모든 $x \in X \setminus U$에 대하여 어떤 $n > 0$이 존재하여
$s^n \in \mathfrak m_x \mathcal{F}^{[n]}_x$이다.
\end{enumerate}
\end{lemma}

\begin{proof}
(1)을 가정하자. $x \in X \setminus U$에 대하여 $U_x$를
$U$의 역상이라 하자. 여기서 역상은 표준 사상
$f_x : \Spec(\mathcal{O}_{X, x}) \to X$에 의한 것이다.
그 역상은 아핀이고 $x$를 포함하지 않는다. 따라서
$\mathfrak m_x \Gamma(U_x, \mathcal{O}_{U_x})$는 단위원 아이디얼이다.
특히
$$
1 = \sum f_i g_i
$$
로 쓸 수 있다. 여기서 $f_i \in \mathfrak m_x$이고
$g_i \in \Gamma(U_x, \mathcal{O}_{U_x})$이다.
Lemma \ref{divisors-lemma-structure-sheaf-Xs}에 의해
$\Gamma(U_x, \mathcal{O}_{U_x}) = \colim \mathcal{F}^{[n]}_x$이고,
전이 사상들은 $s$와의 곱으로 주어진다. 따라서 어떤 $n > 0$에
대하여
$$
s^n = \sum f_i t_i
$$
이다. 여기서 어떤 $t_i = s^ng_i \in \mathcal{F}^{[n]}_x$가
존재한다. 그러므로 (2)가 성립한다.

\medskip\noindent
역으로 (2)가 성립한다고 하자. $j$가 아핀임은 $X$ 위에서
국소적인 조건이다. Morphisms, Lemma
\ref{morphisms-lemma-characterize-affine}를 보라. 따라서
$X$가 아핀이라고 가정해도 되고 실제로 그렇게 하자. 목표는
$U$가 아핀임을 보이는 것이다. Cohomology of Schemes, Lemma
\ref{coherent-lemma-affine-if-quasi-affine}에 의해
$H^p(U, \mathcal{O}_U) = 0$임을 모든 $p > 0$에 대하여 보이면
충분하다.
Cohomology of Schemes, Lemma
\ref{coherent-lemma-quasi-coherence-higher-direct-images-application}에
의해 $H^p(U, \mathcal{O}_U) = H^0(X, R^pj_*\mathcal{O}_U)$이고,
Cohomology of Schemes, Lemma
\ref{coherent-lemma-quasi-coherence-higher-direct-images}에 의해
$R^pj_*\mathcal{O}_U$는 준연접이므로, 줄기
$(R^pj_*\mathcal{O}_U)_x$가 점 $x \in X$에서 영임을 보이면
충분하다.
밑변환 그림
$$
\xymatrix{
U_x \ar[d]_{j_x} \ar[r] & U \ar[d]^j \\
\Spec(\mathcal{O}_{X, x}) \ar[r] & X
}
$$
을 생각하자. Cohomology of Schemes, Lemma
\ref{coherent-lemma-flat-base-change-cohomology}에 의해
$(R^pj_*\mathcal{O}_U)_x = R^pj_{x, *}\mathcal{O}_{U_x}$이다.
따라서 $X$가 닫힌점 $x$를 갖는 국소 스킴이라고 가정해도 되고,
$U$가 아핀임을 보여야 한다(위의 인용 결과에 의해 이는 원하는
소멸과 동치이다). 특히 $d = \dim(X)$는 유한하다
(Algebra, Proposition \ref{algebra-proposition-dimension}).
$x \in U$이면 $U = X$이고 결과는 명백하다.
$d = 0$이고 $x \not \in U$이면 $U = \emptyset$이고 결과는
명백하다. 이제 $d > 0$이고 $x \not \in U$라고 가정하자.
$j_*\mathcal{O}_U = \colim \mathcal{F}^{[n]}$이므로 우리의 가정은
$$
1 = \sum f_i g_i
$$
로 쓸 수 있음을 뜻한다. 여기서 어떤 $n > 0$,
$f_i \in \mathfrak m_x$, $g_i \in \mathcal{O}(U)$가 존재한다.
$d$에 대한 귀납법에 의해 $D(f_i) \cap U$가 아핀임을
모든 $i$에 대하여 안다. 실제로 $X$를 $D(f_i)$로
바꾸어 방금 제시한 논증 전체를 거치면 차원이 $d$보다 엄밀히 작은
뇌터 국소환들에 이른다. 따라서 Properties, Lemma
\ref{properties-lemma-characterize-affine}에 의해 원하는 대로
$U$는 아핀이다.
\end{proof}





\section{상대 Proj}
\label{divisors-section-relative-proj}

\noindent
상대 Proj에 관한 몇 가지 결과를 다룬다. 먼저 매우 기초적인
결과들부터 보자. 상대 Proj는 언제나 밑 위에서 분리임을 상기하자.
Constructions, Lemma
\ref{constructions-lemma-relative-proj-separated}를 보라.

\begin{lemma}
\label{divisors-lemma-relative-proj-quasi-compact}
$S$를 스킴이라 하자. $\mathcal{A}$를 준연접 등급
$\mathcal{O}_S$-대수라 하자.
$p : X = \underline{\text{Proj}}_S(\mathcal{A}) \to S$를
$\mathcal{A}$의 상대 Proj라 하자. 다음 중 하나가 성립한다고 하자.
\begin{enumerate}
\item $\mathcal{A}$는 $\mathcal{A}_0$-대수의 층으로서 유한형이다.
\item $\mathcal{A}$는 $\mathcal{A}_1$에 의해
$\mathcal{A}_0$-대수로서 생성되고, $\mathcal{A}_1$은 유한형
$\mathcal{A}_0$-가군이다.
\item 유한형 준연접 $\mathcal{A}_0$-부분가군
$\mathcal{F} \subset \mathcal{A}_{+}$가 존재하여
$\mathcal{A}_{+}/\mathcal{F}\mathcal{A}$가
$\mathcal{A}/\mathcal{F}\mathcal{A}$의 국소 멱영 아이디얼층이다.
\end{enumerate}
그러면 $p$는 준콤팩트이다.
\end{lemma}

\begin{proof}
문제는 밑 위에서 국소적이다. Schemes, Lemma
\ref{schemes-lemma-quasi-compact-affine}를 보라.
따라서 $S$가 아핀이라고 가정해도 된다.
$S = \Spec(R)$이고 $\mathcal{A}$가 등급 $R$-대수 $A$에
대응한다고 하자. 그러면 $X = \text{Proj}(A)$이다.
Constructions, Section
\ref{constructions-section-relative-proj-via-glueing}을 보라.
(1)의 경우 필요하면 더 국소화한 뒤, $A$가 동차 원소
$f_1, \ldots, f_n \in A_{+}$에 의해 $A_0$ 위에서 생성된다고
가정해도 된다.
그러면 Algebra, Lemma \ref{algebra-lemma-S-plus-generated}에 의해
$A_{+} = (f_1, \ldots, f_n)$이다. (3)의 경우
$\mathcal{F} = \widetilde{M}$임을 알 수 있다. 여기서 어떤 유한형
$A_0$-가군 $M \subset A_{+}$가 존재한다.
$M = \sum A_0f_i$라 하자. 동차 성분으로의 분해를
$f_i = \sum f_{i, j}$라 하자. (3)의 조건은
$A_{+} \subset \sqrt{(f_{i, j})}$임을 뜻한다. 따라서 두 경우 모두
Constructions, Lemma
\ref{constructions-lemma-proj-quasi-compact}에 의해
$\text{Proj}(A)$가 준콤팩트임을 얻는다. 끝으로 (2)는 (1)에서
따른다.
\end{proof}

\begin{lemma}
\label{divisors-lemma-relative-proj-finite-type}
$S$를 스킴이라 하자. $\mathcal{A}$를 준연접 등급
$\mathcal{O}_S$-대수라 하자.
$p : X = \underline{\text{Proj}}_S(\mathcal{A}) \to S$를
$\mathcal{A}$의 상대 Proj라 하자. $\mathcal{A}$가
$\mathcal{O}_S$-대수의 층으로서 유한형이면 $p$는 유한형이고
$\mathcal{O}_X(d)$는 유한형 $\mathcal{O}_X$-가군이다.
\end{lemma}

\begin{proof}
가정에 의해 $p$는 준콤팩트이다. Lemma
\ref{divisors-lemma-relative-proj-quasi-compact}를 보라. 따라서 $p$가
국소 유한형임을 보이면 충분하다. 그러므로 문제는 밑과 목표 위에서
국소적이다. Morphisms, Lemma
\ref{morphisms-lemma-locally-finite-type-characterize}를 보라.
$S = \Spec(R)$이고 $\mathcal{A}$가 등급 $R$-대수 $A$에 대응한다고
하자. $S$ 위에서 더 국소화하여 $A$가 유한형 $R$-대수라고
가정해도 된다. 스킴 $X$는 환 $A_{(f)}$들의 스펙트럼을 이어붙여
구성되며, 여기서 $f \in A_{+}$는 동차 원소이다. 이들은 각각 Algebra,
Lemma \ref{algebra-lemma-dehomogenize-finite-type}의 (1)에 의해
$R$ 위에서 유한형이다. 따라서 $\text{Proj}(A)$는 $R$ 위에서
유한형이다. $\mathcal{O}_X(d)$에 관한 명제에는 Algebra, Lemma
\ref{algebra-lemma-dehomogenize-finite-type}의 (2)를 사용한다.
\end{proof}

\begin{lemma}
\label{divisors-lemma-relative-proj-universally-closed}
$S$를 스킴이라 하자. $\mathcal{A}$를 준연접 등급
$\mathcal{O}_S$-대수라 하자.
$p : X = \underline{\text{Proj}}_S(\mathcal{A}) \to S$를
$\mathcal{A}$의 상대 Proj라 하자.
$\mathcal{O}_S \to \mathcal{A}_0$가 정수적 대수 사상이고
\footnote{다시 말해 $\mathcal{O}_S$의 $\mathcal{A}_0$ 안
정수적 폐포(Morphisms, Definition
\ref{morphisms-definition-integral-closure}를 보라)가
$\mathcal{A}_0$와 같다.}
$\mathcal{A}$가 $\mathcal{A}_0$-대수로서 유한형이면 $p$는
보편 닫힌 사상이다.
\end{lemma}

\begin{proof}
문제는 밑 위에서 국소적이다. 따라서 $X = \Spec(R)$가
아핀이라고 가정해도 된다. $\mathcal{A}$를 준연접
$\mathcal{O}_X$-대수라 하되 등급 $R$-대수 $A$에 수반되게 하자.
가정은 $R \to A_0$가 정수적이고 $A$가 $A_0$ 위에서
유한형이라는 것이다. $X \to \Spec(R)$를 합성
$X \to \Spec(A_0) \to \Spec(R)$으로 쓰자.
$R \to A_0$가 정수적 환 사상이므로
$\Spec(A_0) \to \Spec(R)$은 보편 닫힌 사상이다.
Morphisms, Lemma \ref{morphisms-lemma-integral-universally-closed}를
보라. 준콤팩트 사상(Constructions, Lemma
\ref{constructions-lemma-proj-quasi-compact}를 보라)
$$
X = \text{Proj}(A) \to \Spec(A_0)
$$
은 Constructions, Lemma
\ref{constructions-lemma-proj-valuative-criterion}에 의해
값매김 판정법의 존재 부분을 만족하고, 따라서 Schemes, Proposition
\ref{schemes-proposition-characterize-universally-closed}에 의해
보편 닫힌 사상이다. 그러므로 $X \to \Spec(R)$은 보편 닫힌
사상들의 합성으로서 보편 닫힌 사상이다.
\end{proof}

\begin{lemma}
\label{divisors-lemma-relative-proj-proper}
$S$를 스킴이라 하자. $\mathcal{A}$를 준연접 등급
$\mathcal{O}_S$-대수라 하자.
$p : X = \underline{\text{Proj}}_S(\mathcal{A}) \to S$를
$\mathcal{A}$의 상대 Proj라 하자. 다음 조건들은 동치이다.
\begin{enumerate}
\item $\mathcal{A}_0$는 유한형 $\mathcal{O}_S$-가군이고,
$\mathcal{A}$는 $\mathcal{A}_0$-대수로서 유한형이다.
\item $\mathcal{A}_0$는 유한형 $\mathcal{O}_S$-가군이고,
$\mathcal{A}$는 $\mathcal{O}_S$-대수로서 유한형이다.
\end{enumerate}
이 조건들이 성립하면 $p$는 국소 사영적이고, 특히 고유하다.
\end{lemma}

\begin{proof}
$\mathcal{A}_0$가 유한형 $\mathcal{O}_S$-가군이라고 가정하자.
아핀 열린집합 $U = \Spec(R) \subset X$를 고르되
$\mathcal{A}$가 등급 $R$-대수 $A$에 대응하고 $A_0$가 유한
$R$-가군이 되게 하자. 조건 (1)은 (필요하면 $S$ 위에서 더
국소화한 뒤) $A$가 유한형 $A_0$-대수라는 뜻이고, 조건 (2)는
(필요하면 $S$ 위에서 더 국소화한 뒤) $A$가 유한형
$R$-대수라는 뜻이다. 따라서 Algebra, Lemma
\ref{algebra-lemma-compose-finite-type}에 의해 두 조건은 서로를
함의한다.

\medskip\noindent
국소 사영 사상은 고유하다. Morphisms, Lemma
\ref{morphisms-lemma-locally-projective-proper}를 보라.
따라서 이제 $S = \Spec(R)$이고 $X = \text{Proj}(A)$이며,
$A_0$가 $R$ 위에서 유한하고 $A$가 $R$ 위에서 유한형이라고
가정해도 된다. $X = \text{Proj}(A) \to \Spec(R)$이 사영 사상임을
보일 것이다. 독자에게 아래 논증을 읽는 대신 $X$에서 $R$ 위의
사영공간으로 가는 닫힌 몰입을 직접 구성하여 이를 증명해 보기를
권한다.

\medskip\noindent
Lemma \ref{divisors-lemma-relative-proj-finite-type}에 의해 $X$는
$\Spec(R)$ 위에서 유한형이다. Constructions, Lemma
\ref{constructions-lemma-ample-on-proj}에 의해
$\mathcal{O}_X(d)$는 $X$ 위에서 풍부하며, 여기서 어떤
$d \geq 1$을 택할 수 있다
(Properties, Section \ref{properties-section-ample}을 보라).
따라서 $X \to \Spec(R)$은 준사영 사상이다(Morphisms, Definition
\ref{morphisms-definition-quasi-projective}에 의해).
Morphisms, Lemma
\ref{morphisms-lemma-quasi-projective-open-projective}에 의해
$X$가 $\Spec(R)$ 위에서 사영적인 어떤 스킴의 열린 부분스킴과
동형임을 얻는다. 따라서 증명을 끝내려면
$X \to \Spec(R)$이 보편 닫힌 사상임을 보이면 충분하다
(Morphisms, Lemma
\ref{morphisms-lemma-image-proper-scheme-closed}를 사용하라).
이는 Lemma \ref{divisors-lemma-relative-proj-universally-closed}에서 따른다.
\end{proof}

\begin{lemma}
\label{divisors-lemma-relative-proj-projective}
$S$를 스킴이라 하자. $\mathcal{A}$를 준연접 등급
$\mathcal{O}_S$-대수라 하자.
$p : X = \underline{\text{Proj}}_S(\mathcal{A}) \to S$를
$\mathcal{A}$의 상대 Proj라 하자. $\mathcal{A}$가
$\mathcal{A}_1$에 의해 $\mathcal{A}_0$ 위에서 생성되고
$\mathcal{A}_1$이 유한형 $\mathcal{O}_S$-가군이면 $p$는
사영 사상이다.
\end{lemma}

\begin{proof}
구체적으로 말해 등급 $\mathcal{O}_S$-대수 사상
$$
\text{Sym}_{\mathcal{O}_X}^*(\mathcal{A}_1)
\longrightarrow
\mathcal{A}
$$
에 수반되는 사상은 닫힌 몰입 $X \to \mathbf{P}(\mathcal{A}_1)$이다.
Constructions, Lemma
\ref{constructions-lemma-surjective-generated-degree-1-map-relative-proj}를
보라.
\end{proof}

\begin{lemma}
\label{divisors-lemma-relative-proj-flat}
$S$를 스킴이라 하자. $\mathcal{A}$를 준연접 등급
$\mathcal{O}_S$-대수라 하자.
$p : X = \underline{\text{Proj}}_S(\mathcal{A}) \to S$를
$\mathcal{A}$의 상대 Proj라 하자. $\mathcal{A}_d$가 평탄
$\mathcal{O}_S$-가군이고 이것이 $d \gg 0$에 대하여 성립하면
$p$는 평탄하고
$\mathcal{O}_X(d)$는 $S$ 위에서 평탄하다.
\end{lemma}

\begin{proof}
$X$가 $S$ 위에서 평탄하다는 것은 아핀 국소적으로 다음 명제로
환원된다. $R$을 환, $A$를 등급 $R$-대수라 하고,
$A_d$가 $R$ 위에서 평탄하며 이것이 $d \gg 0$에 대하여
성립한다고 하자. 어떤 $f \in A_d$를 택하되 $d > 0$이게 하자.
그러면 $A_{(f)}$는 $R$ 위에서 평탄하다.
$A_{(f)} = \colim A_{nd}$이고 전이 사상들이 $f$와의 곱으로
주어지므로, 이는 Algebra, Lemma
\ref{algebra-lemma-colimit-flat}에서 따른다. 같은 방식으로
$\mathcal{O}_X(d)$가 $S$ 위에서 평탄함을 얻는다.
\end{proof}

\begin{lemma}
\label{divisors-lemma-relative-proj-finite-presentation}
$S$를 스킴이라 하자. $\mathcal{A}$를 준연접 등급
$\mathcal{O}_S$-대수라 하자.
$p : X = \underline{\text{Proj}}_S(\mathcal{A}) \to S$를
$\mathcal{A}$의 상대 Proj라 하자. $\mathcal{A}$가 유한 표시
$\mathcal{O}_S$-대수이면 $p$는 유한 표시 사상이고
$\mathcal{O}_X(d)$는 유한 표시 $\mathcal{O}_X$-가군이다.
\end{lemma}

\begin{proof}
아핀 국소적으로 이는 다음 명제로 환원된다. $R$을 환이라 하고
$A$를 유한 표시 등급 $R$-대수라 하자. 그러면
$\text{Proj}(A) \to \Spec(R)$은 유한 표시 사상이고
$\mathcal{O}_{\text{Proj}(A)}(d)$는 유한 표시
$\mathcal{O}_{\text{Proj}(A)}$-가군이다. 유한 표시 조건에 의해
표시
$$
A = R[X_1, \ldots, X_n]/(F_1, \ldots, F_m)
$$
를 고를 수 있다. 여기서 $R[X_1, \ldots, X_n]$은 가중치
$d_i$를 $X_i$에 주어 등급을 매긴 다항식환이고,
$F_1, \ldots, F_m$은 차수 $e_j$의 동차 다항식들이다.
$R_0 \subset R$를 다항식 $F_1, \ldots, F_m$의 계수들로 생성되는
부분환이라 하자. 다음과 같이 놓자.
$A_0 = R_0[X_1, \ldots, X_n]/(F_1, \ldots, F_m)$.
구성에 의해 $A = A_0 \otimes_{R_0} R$이다. 따라서 Constructions,
Lemma \ref{constructions-lemma-base-change-map-proj}에 의해
$X_0 = \text{Proj}(A_0)$에 대하여 결과를 증명하면 충분하며,
이는 $R_0$ 위의 결과이다. Lemma
\ref{divisors-lemma-relative-proj-finite-type}에 의해
$X_0$가 $R_0$ 위에서 유한형이고 $\mathcal{O}_{X_0}(d)$가
유한형 준연접 $\mathcal{O}_{X_0}$-가군임을 안다.
$R_0$는 뇌터이므로(유한 생성 $\mathbf{Z}$-대수이므로)
$X_0$는 $R_0$ 위에서 유한 표시이다(Morphisms, Lemma
\ref{morphisms-lemma-noetherian-finite-type-finite-presentation}).
또한 Cohomology of Schemes, Lemma
\ref{coherent-lemma-coherent-Noetherian}에 의해
$\mathcal{O}_{X_0}(d)$는 유한 표시이다. 이것으로 증명이 끝난다.
\end{proof}









\section{상대 proj의 닫힌 부분스킴}
\label{divisors-section-closed-in-relative-proj}

\noindent
상대 proj의 닫힌 부분스킴에 관한 몇 가지 보조정리를 다룬다.

\begin{lemma}
\label{divisors-lemma-closed-subscheme-proj}
$S$를 스킴이라 하자. $\mathcal{A}$를 준연접 등급
$\mathcal{O}_S$-대수라 하자.
$p : X = \underline{\text{Proj}}_S(\mathcal{A}) \to S$를
$\mathcal{A}$의 상대 Proj라 하자. $i : Z \to X$를 닫힌
부분스킴이라 하자. $\mathcal{I} \subset \mathcal{A}$로 표준 사상
$$
\mathcal{A}
\longrightarrow
\bigoplus\nolimits_{d \geq 0} p_*\left((i_*\mathcal{O}_Z)(d)\right).
$$
의 핵을 나타내자. $p$가 준콤팩트이면 동형
$Z = \underline{\text{Proj}}_S(\mathcal{A}/\mathcal{I})$이 존재한다.
\end{lemma}

\begin{proof}
Constructions, Lemma
\ref{constructions-lemma-relative-proj-separated}에 의해 사상 $p$는
분리이다. $p$가 준콤팩트이므로 $p_*$는 준연접 가군을 준연접
가군으로 보낸다. Schemes, Lemma
\ref{schemes-lemma-push-forward-quasi-coherent}를 보라.
따라서 $\mathcal{I}$는 준연접 $\mathcal{O}_S$-가군이다.
특히 $\mathcal{B} = \mathcal{A}/\mathcal{I}$는 준연접 등급
$\mathcal{O}_S$-대수이다. 함자성 사상
$Z' = \underline{\text{Proj}}_S(\mathcal{B}) \to
\underline{\text{Proj}}_S(\mathcal{A})$은 모든 곳에서 정의되고
닫힌 몰입이다. Constructions, Lemma
\ref{constructions-lemma-surjective-graded-rings-map-relative-proj}를
보라. 따라서 $Z = Z'$임을 증명하면 충분하며, 이는 $X$의
닫힌 부분스킴으로서의 등식이다.

\medskip\noindent
이제 문제는 밑 위에서 국소적이므로
$S = \Spec(R)$이고 $X = \text{Proj}(A)$라고 가정해도 된다.
여기서 $R$-대수 $A$는 등급 대수이다.
$\mathcal{I} = \widetilde{I}$라고 가정하자. 여기서
$I \subset A$는 등급 아이디얼이다. Constructions, Lemma
\ref{constructions-lemma-proj-quasi-compact}에 의해 원소
$f_0, \ldots, f_n \in A_{+}$가 존재하여
$A_{+} \subset \sqrt{(f_0, \ldots, f_n)}$, 다시 말해
$X = \bigcup D_{+}(f_i)$이다. 따라서
$Z \cap D_{+}(f_i) = Z' \cap D_{+}(f_i)$가 각 $i$에 대하여
성립함을 확인하면 충분하다.
번호를 다시 매겨 $i = 0$이라고 가정해도 된다.
$Z \cap D_{+}(f_0)$와 각각 $Z' \cap D_{+}(f_0)$가
아이디얼 $J$와 각각 $J'$에 의해 잘려 나온다고 하자. 이들은
$A_{(f_0)}$의 아이디얼이다.

\medskip\noindent
포함 $J' \subset J$. $d$를
$\deg(f_0), \ldots, \deg(f_n)$의 최소공배수라 하자.
각 꼬임 $\mathcal{O}_X(nd)$은 가역이고,
$f_i^{nd/\deg(f_i)}$에 의해 $D_{+}(f_i)$ 위에서
자명화됨을 유의하자. 또한 임의의 준연접 가군 $\mathcal{F}$가
$X$ 위에 주어질 때 곱셈 사상
$\mathcal{O}_X(nd) \otimes_{\mathcal{O}_X} \mathcal{F}(m)
\to \mathcal{F}(nd + m)$은 동형이다. Constructions, Lemma
\ref{constructions-lemma-when-invertible}을 보라.
$J'$은 $g/f_0^e$ 꼴의 원소들로 생성되는 아이디얼임을 관찰하자.
여기서 $g \in I$는 차수 $e\deg(f_0)$의 동차 원소이다
(Constructions, Lemma
\ref{constructions-lemma-surjective-graded-rings-map-proj}의
증명을 보라). 물론 $g$를 $f_0^lg$로 바꾸어도 되며, 여기서
$l$은 적당히 선택한다. 그러므로 언제나 $d | e$라고 가정할 수 있다.
$g$는 $\mathcal{O}_X(e\deg(f_0))$의 절단으로서 이를 $Z$에
제한하면 소멸하므로 $g/f_0^d$는 $J$의 원소이다.
따라서 $J' \subset J$이다.

\medskip\noindent
역으로 $g/f_0^e \in J$라고 하자. 다시 $d | e$라고 가정해도 된다.
$i \in \{1, \ldots, n\}$를 고르자. 그러면
$Z \cap D_{+}(f_i)$는 어떤 아이디얼
$J_i \subset A_{(f_i)}$에 의해 잘려 나온다. 더욱이
$$
J \cdot A_{(f_0f_i)} = J_i \cdot A_{(f_0f_i)}.
$$
우변은 $J_i$의
$f_0^{\deg(f_i)}/f_i^{\deg(f_0)}$에 관한 국소화이다.
따라서 충분히 나누어지는 어떤 $e_i \gg 0$에 대하여
$$
f_0^{e_i}g/f_i^{(e_i + e)\deg(f_0)/\deg(f_i)} \in J_i
$$
이다. 이것으로 $f_0^{\max(e_i)}g$가 $I$의 원소임을 알 수 있다.
실제로 이를 각 아핀 열린집합 $D_{+}(f_i)$에 제한하면 닫힌
부분스킴 $Z \cap D_{+}(f_i)$ 위에서 소멸한다.
따라서 $g/f_0^e \in J'$이고 원하는 대로 $J \subset J'$이다.
\end{proof}

\begin{example}
\label{divisors-example-closed-subscheme-of-proj}
$A$를 등급환이라 하자. $X = \text{Proj}(A)$ 및
$S = \Spec(A_0)$라 하자. 등급 아이디얼 $I \subset A$가 주어지면
Constructions, Lemma
\ref{constructions-lemma-surjective-graded-rings-map-proj}에 의해
닫힌 부분스킴 $V_+(I) = \text{Proj}(A/I) \to X$를 얻는다.
Lemma \ref{divisors-lemma-closed-subscheme-proj}의 결과를 옮겨 쓰면,
$X$가 준콤팩트일 때 임의의 닫힌 부분스킴 $Z$는
$V_+(I(Z))$ 꼴이다. 여기서 등급 아이디얼 $I(Z) \subset A$는
다음 규칙으로 주어진다.
$$
I(Z) = \Ker(A \longrightarrow
\bigoplus\nolimits_{n \geq 0} \Gamma(Z, \mathcal{O}_Z(n)))
$$
그러면 다음 두 가지 자연스러운 질문을 할 수 있다.
\begin{enumerate}
\item 어떤 아이디얼 $I$가 $I(Z)$ 꼴인가?
\item 연산 $I \mapsto I(V_+(I))$를 기술할 수 있는가?
\end{enumerate}
$A$가 뇌터인 경우에 이에 답할 것이다.

\medskip\noindent
먼저 $A$가 $A_1$에 의해 $A_0$ 위에서 생성된다고 가정하자.
이 경우 임의의 아이디얼 $I \subset A$에 대하여 사상
$A/I \to \bigoplus \Gamma(\text{Proj}(A/I), \mathcal{O})$의 핵은
$A/I$의 꼬임 원소들의 집합이다. Cohomology of Schemes, Proposition
\ref{coherent-proposition-coherent-modules-on-proj}을 보라.
따라서 다음을 얻는다.
$$
I(V_+(I)) = \{x \in A \mid A_n x \subset I\text{ for some }n \geq 0\}
$$
우변의 아이디얼을 때때로 $I$의 포화라고 한다.
이로써 (2)에 답했고, (1)에 대한 답은 아이디얼이 $I(Z)$ 꼴일
필요충분조건이 포화되어 있는 것, 즉 자기 자신의 포화와 같은 것이라는
것이다.

\medskip\noindent
$A$가 일반적인 뇌터 등급환이면 Cohomology of Schemes, Proposition
\ref{coherent-proposition-coherent-modules-on-proj-general}을 사용한다.
따라서 $d$를 $A$의 $A_0$ 위 생성원들의 차수의 최소공배수라 하면
다음을 얻는다.
$$
I(V_+(I)) = \{x \in A \mid (Ax)_{nd} \subset I\text{ for all }n \gg 0\}
$$
이는 $I$의 포화와 다를 수 있으며, 그 차이는 $d \not = 1$일 때
생길 수 있다. 예를 들어 $A = \mathbf{Q}[x, y]$이고
$\deg(x) = 2$이며 $\deg(y) = 3$이라고 하자. 그러면 $d = 6$이다.
$I = (y^2)$라 하자. $y \in I(V_+(I))$임을 알 수 있다. 실제로
$f \in A$인 임의의 동차 원소가 $6 | \deg(fy)$를 만족할 때
$y | f$이고 따라서 $fy \in I$이다. 그러므로
$I(V_+(I)) = (y)$이지만 $x^n y \not \in I$가 모든 $n$에 대하여
성립한다. 따라서 $I(V_+(I))$는 포화와 같지 않다.
\end{example}

\begin{lemma}
\label{divisors-lemma-equation-codim-1-in-projective-space}
$R$을 뇌터 UFD라 하자. $Z \subset \mathbf{P}^n_R$를 매장점이 없는
닫힌 부분스킴이라 하고, $Z$의 모든 기약 성분의 여차원이
$1$이라고 하자. 이 여차원은 $\mathbf{P}^n_R$에서 잰다.
그러면 아이디얼 $I(Z) \subset R[T_0, \ldots, T_n]$은
$Z$에 대응하며 주아이디얼이다.
\end{lemma}

\begin{proof}
$X = \mathbf{P}^n_R$의 국소환들은 UFD임을 관찰하자. 실제로
$X$는 $\mathbf{A}^n_R$와 동형인 아핀 조각들로 덮이고
$R[x_1, \ldots, x_n]$은 UFD이다
(Algebra, Lemma \ref{algebra-lemma-polynomial-ring-UFD}).
따라서 Lemma \ref{divisors-lemma-codimension-1-is-effective-Cartier}에 의해
$Z$는 유효 카르티에 인자이다.
$\mathcal{I} \subset \mathcal{O}_X$를 $Z$에 대응하는 준연접
아이디얼층이라 하자. 동형 $\mathcal{O}(m) \to \mathcal{I}$를
고르자. 여기서 $m \in \mathbf{Z}$이다.
Lemma \ref{divisors-lemma-Pic-projective-space-UFD}를 보라.
그러면 합성
$$
\mathcal{O}_X(m) \to \mathcal{I} \to \mathcal{O}_X
$$
은 영이 아니다. 따라서 $m \leq 0$이고, 이에 대응하는
$\mathcal{O}_X(m)^{\otimes -1} = \mathcal{O}_X(-m)$의 절단은 어떤
$F \in R[T_0, \ldots, T_n]$으로 주어지며, 그 차수는 $-m$이다.
Cohomology of Schemes, Lemma
\ref{coherent-lemma-cohomology-projective-space-over-ring}를 보라.
그러므로 $i$번째 표준 열린집합 $U_i = D_+(T_i)$ 위에서
닫힌 부분스킴 $Z \cap U_i$는 다음 아이디얼에 의해 잘려 나온다.
$$
(F(T_0/T_i, \ldots, T_n/T_i)) \subset R[T_0/T_i, \ldots, T_n/T_i]
$$
따라서 등급 아이디얼
$I(Z) = \Ker(R[T_0, \ldots, T_n] \to \bigoplus \Gamma(\mathcal{O}_Z(m)))$
의 동차 원소들은 다음을 만족하는 동차 다항식 $G$들의 집합이다.
$$
G(T_0/T_i, \ldots, T_n/T_i) \in (F(T_0/T_i, \ldots, T_n/T_i))
$$
여기서 $i = 0, \ldots, n$이다. 분모를 제거하면
다음을 만족하는 $e_i \geq 0$가 존재함을 알 수 있다.
$$
T_i^{e_i}G \in (F)
$$
여기서 $i = 0, \ldots, n$이다. $R$이 UFD이므로
$R[T_0, \ldots, T_n]$도 UFD이다. 이제 $F | T_0^{e_0}G$이고
$F | T_1^{e_1}G$이면 $F | G$이다. 실제로
$T_0^{e_0}$과 $T_1^{e_1}$에는 공통인 인수가 없다.
따라서 $I(Z) = (F)$이다.
\end{proof}

\noindent
닫힌 부분스킴이 국소적으로 유한 개의 방정식으로 잘려 나오는 경우에는
$\mathcal{A}$의 유한형 아이디얼층으로 이를 정의할 수 있다.

\begin{lemma}
\label{divisors-lemma-closed-subscheme-proj-finite}
$S$를 준콤팩트 준분리 스킴이라 하자.
$\mathcal{A}$를 준연접 등급 $\mathcal{O}_S$-대수라 하자.
$p : X = \underline{\text{Proj}}_S(\mathcal{A}) \to S$를
$\mathcal{A}$의 상대 Proj라 하자. $i : Z \to X$를 닫힌
부분스킴이라 하자. $p$가 준콤팩트이고 $i$가 유한 표시이면,
$d > 0$과 준연접 유한형 $\mathcal{O}_S$-부분가군
$\mathcal{F} \subset \mathcal{A}_d$가 존재하여
$Z = \underline{\text{Proj}}_S(\mathcal{A}/\mathcal{F}\mathcal{A})$이다.
\end{lemma}

\begin{proof}
Lemma \ref{divisors-lemma-closed-subscheme-proj}에 의해 준연접 등급
아이디얼층 $\mathcal{I} \subset \mathcal{A}$가 존재하여
$Z = \underline{\text{Proj}}(\mathcal{A}/\mathcal{I})$임을 안다.
$S$가 준콤팩트이므로 유한 아핀 열린 덮개
$S = U_1 \cup \ldots \cup U_n$을 고를 수 있다.
$U_i = \Spec(R_i)$라 하자.
$\mathcal{A}|_{U_i}$에 대응하는 등급 $R_i$-대수를 $A_i$라 하고
$\mathcal{I}|_{U_i}$에 대응하는 등급 아이디얼을
$I_i \subset A_i$라 하자.
$p^{-1}(U_i) = \text{Proj}(A_i)$가 $R_i$ 위의 스킴으로서 성립함을
유의하자. $p$가 준콤팩트이므로 유한 개의 동차 원소
$f_{i, j} \in A_{i, +}$를 골라
$p^{-1}(U_i) = \bigcup_j D_{+}(f_{i, j})$가 되게 할 수 있다.
$Z \to X$에 대한 조건은 $Z$의 아이디얼층이 $\mathcal{O}_X$에서
유한형이라는 뜻이다. Morphisms, Lemma
\ref{morphisms-lemma-closed-immersion-finite-presentation}을 보라.
따라서 유한 개의 동차 원소
$h_{i, j, k} \in I_i \cap A_{i, +}$를 찾아
$Z \cap D_{+}(f_{i, j})$의 아이디얼이 원소
$h_{i, j, k}/f_{i, j}^{e_{i, j, k}}$들로 생성되게 할 수 있다.
$d > 0$을 모든 정수 $\deg(f_{i, j})$와
$\deg(h_{i, j, k})$의 공배수로 고르자.
Properties, Lemma
\ref{properties-lemma-quasi-coherent-colimit-finite-type}에 의해
유한형 준연접 부분가군
$\mathcal{F} \subset \mathcal{I}_d$가 존재하여 모든 국소 절단
$$
h_{i, j, k}f_{i, j}^{(d - \deg(h_{i, j, k}))/\deg(f_{i, j})}
$$
이 $\mathcal{F}$의 절단이 되게 할 수 있다.
구성상 $\mathcal{F}$는 해답이다.
\end{proof}

\noindent
다음은 Lemma \ref{divisors-lemma-closed-subscheme-proj-finite}의 한 변형이며,
Lemma \ref{divisors-lemma-composition-admissible-blowups}의 증명에서 사용할 것이다.

\begin{lemma}
\label{divisors-lemma-closed-subscheme-proj-finite-type}
$S$를 준콤팩트 준분리 스킴이라 하자.
$\mathcal{A}$를 준연접 등급 $\mathcal{O}_S$-대수라 하자.
$p : X = \underline{\text{Proj}}_S(\mathcal{A}) \to S$를
$\mathcal{A}$의 상대 Proj라 하자. $i : Z \to X$를 닫힌
부분스킴이라 하자. $U \subset S$를 열린집합이라 하자.
다음을 가정하자.
\begin{enumerate}
\item $p$는 준콤팩트이다.
\item $i$는 유한 표시이다.
\item $U \cap p(i(Z)) = \emptyset$이다.
\item $U$는 준콤팩트이다.
\item $\mathcal{A}_n$은 유한형 $\mathcal{O}_S$-가군이며, 이는 모든
$n$에 대하여 성립한다.
\end{enumerate}
그러면 $d > 0$과 준연접 유한형 $\mathcal{O}_S$-부분가군
$\mathcal{F} \subset \mathcal{A}_d$가 존재하여 (a)
$Z = \underline{\text{Proj}}_S(\mathcal{A}/\mathcal{F}\mathcal{A})$이고,
(b) $\mathcal{A}_d/\mathcal{F}$의 지지는 $U$와 서로소이다.
\end{lemma}

\begin{proof}
$\mathcal{I} \subset \mathcal{A}$를
Lemma \ref{divisors-lemma-closed-subscheme-proj}에서 구성한 준연접 등급
아이디얼층이라 하자. $U_i$, $R_i$, $A_i$, $I_i$, $f_{i, j}$,
$h_{i, j, k}$ 및 $d$는
Lemma \ref{divisors-lemma-closed-subscheme-proj-finite}의 증명에서 구성한
것과 같게 하자. $U \cap p(i(Z)) = \emptyset$이므로
$\mathcal{I}_d|_U = \mathcal{A}_d|_U$임을 알 수 있다
(이는 $\mathcal{I}$를 핵으로 구성한 방식에서 따른다).
$U$가 준콤팩트이므로 유한 아핀 열린 덮개
$U = W_1 \cup \ldots \cup W_m$을 고를 수 있다.
$\mathcal{A}_d$가 유한형이므로 유한 개의 절단
$g_{t, s} \in \mathcal{A}_d(W_t)$를 찾아
$\mathcal{A}_d|_{W_t} = \mathcal{I}_d|_{W_t}$를
$\mathcal{O}_{W_t}$-가군으로서 생성하게 할 수 있다.
증명을 끝내기 위해 Properties, Lemma
\ref{properties-lemma-quasi-coherent-colimit-finite-type}에 의해
유한형 부분가군 $\mathcal{F} \subset \mathcal{I}_d$가 존재하여
모든 국소 절단
$$
h_{i, j, k}f_{i, j}^{(d - \deg(h_{i, j, k}))/\deg(f_{i, j})}
\quad\text{and}\quad
g_{t, s}
$$
이 $\mathcal{F}$의 절단이 되게 할 수 있음을 유의하자.
구성상 $\mathcal{F}$는 해답이다.
\end{proof}

\begin{lemma}
\label{divisors-lemma-conormal-sheaf-section-projective-bundle}
$X$를 스킴이라 하자. $\mathcal{E}$를 준연접
$\mathcal{O}_X$-가군이라 하자. 다음 전단사가 있다.
$$
\left\{
\begin{matrix}
\text{sections }\sigma\text{ of the } \\
\text{morphism } \mathbf{P}(\mathcal{E}) \to X
\end{matrix}
\right\}
\leftrightarrow
\left\{
\begin{matrix}
\text{surjections }\mathcal{E} \to \mathcal{L}\text{ where} \\
\mathcal{L}\text{ is an invertible }\mathcal{O}_X\text{-module}
\end{matrix}
\right\}
$$
이 경우 $\sigma$는 닫힌 몰입이고 표준 동형
$$
\Ker(\mathcal{E} \to \mathcal{L})
\otimes_{\mathcal{O}_X} \mathcal{L}^{\otimes -1}
\longrightarrow
\mathcal{C}_{\sigma(X)/\mathbf{P}(\mathcal{E})}
$$
이 존재한다. 이 전단사와 동형은 모두 밑변환과 양립한다.
\end{lemma}

\begin{proof}
$\pi : \mathbf{P}(\mathcal{E}) \to X$는 $\mathcal{E}$ 위의
대칭대수의 상대 proj임을 상기하자. Constructions, Definition
\ref{constructions-definition-projective-bundle}을 보라.
따라서 절단 $\sigma$에 대한 기술은 Constructions, Lemma
\ref{constructions-lemma-apply-relative}에 있는
$\mathbf{P}(\mathcal{E})$의 점 함자에 대한 기술에서 곧바로 따른다.
$\pi$가 분리이므로 임의의 절단은 닫힌 몰입이다
(Constructions, Lemma
\ref{constructions-lemma-relative-proj-separated} 및
Schemes, Lemma \ref{schemes-lemma-section-immersion}).
$U \subset X$를 아핀 열린집합이라 하고
$k \in \mathcal{E}(U)$와 $s \in \mathcal{E}(U)$를 국소 절단이라 하자.
$k$는 $\mathcal{L}$에서 영으로 가고, $s$는 생성원
$\overline{s}$로 간다고 하자. 이는 $\mathcal{L}$의 생성원이다.
그러면 $f = k/s$는 $\mathcal{O}_{\mathbf{P}(\mathcal{E})}$의
절단이며, 열린 근방 $D_+(s)$ 위에 정의된다. 이는 $s(U)$의
근방이고 $\pi^{-1}(U)$ 안에 있다. 또한 $k$가 $\mathcal{L}$에서
영으로 가므로 $f$는 $s(U)$의 아이디얼층의 절단이며, 이 아이디얼층은
$\pi^{-1}(U)$ 안에서 취한다. 따라서 상 $\overline{f}$를 취할 수
있는데, 이는 $f$의 상이며
$\mathcal{C}_{\sigma(X)/\mathbf{P}(\mathcal{E})}(U)$에서 취할 수 있다.
우리는 (1) $\overline{f}$의 상이 절단 $k$와 $\overline{s}$에만
의존하고 $s$의 선택에는 의존하지 않으며, (2) 이 방법으로
$U$ 위의 동형을 얻는다고 주장한다(아래를 보라).
그러나 (1)과 (2)가 확립되면, 이 구성이 밑변환
$U' \to U$와 양립함을 알 수 있다. 여기서 $U'$는 아핀이다.
따라서 이 국소 사상들은 붙고 임의의 밑변환과 양립한다.

\medskip\noindent
(1)과 (2)를 증명하기 위해 일어나는 일을 명시적으로 적자.
즉 $U = \Spec(A)$라 하고
$\mathcal{E} \to \mathcal{L}$가 $A$-가군의 사상
$M \to N$에 대응한다고 하자. 그러면
$k \in K = \Ker(M \to N)$이고 $s \in M$은 생성원
$\overline{s}$로 가며, 이는 $N$의 생성원이다.
따라서 $M = K \oplus A s$이다. 그러므로
$$
\text{Sym}(M) = \text{Sym}(K)[s]
$$
동일시 $\text{Sym}(K) \to \text{Sym}(M)_{(s)}$를 생각하자.
이는 규칙 $g \mapsto g/s^n$으로 주어진다. 여기서
$g \in \text{Sym}^n(K)$이다. 이로써 동형
$D_+(s) = \Spec(\text{Sym}(K))$을 얻으며, 이 동형 아래에서
$\sigma$는 환 준동형 $\text{Sym}(K) \to A$에 대응하며,
이 환 준동형은 $K$를 영으로 보낸다. 이 동형을 통해
$K$에 대응하는 준연접 가군은
$\mathcal{C}_{\sigma(U)/D_+(s)}$와 동일시되므로 (2)를 증명한다.
끝으로 $s' = k' + s$라고 하자. 여기서 어떤 $k' \in K$를 택한다.
그러면
$$
k/s' = (k/s) (s/s') = (k/s) (s'/s)^{-1} = (k/s) (1 + k'/s)^{-1}
$$
이 등식은 $\sigma(U)$의 열린 근방에서 성립하며, 그 근방은
$D_+(s)$ 안에 있다. 따라서 $s'/s$의 제한은 $1$이며, 이는
$\sigma(U)$ 위의 제한임을 알 수 있고,
$k/s'$가 여법선층에서 $k/s$와 같은 원소로 가는 것도 알 수 있다.
이로써 (1)이 증명된다.
\end{proof}

\section{블로업}
\label{divisors-section-blowing-up}

\noindent
블로업은 대수기하학에서 중요한 도구이다.

\begin{definition}
\label{divisors-definition-blow-up}
$X$를 스킴이라 하자. $\mathcal{I} \subset \mathcal{O}_X$를 준연접
아이디얼층이라 하고, $Z \subset X$를 $\mathcal{I}$에 대응하는 닫힌
부분스킴이라 하자. Schemes, Definition
\ref{schemes-definition-immersion}을 보라.
{\it $X$의 $Z$를 따른 블로업}, 또는
{\it $X$의 아이디얼층 $\mathcal{I}$에서의 블로업}이란 사상
$$
b :
\underline{\text{Proj}}_X
\left(\bigoplus\nolimits_{n \geq 0} \mathcal{I}^n\right)
\longrightarrow
X
$$
을 말한다. 블로업의 {\it 예외 인자}는 역상 $b^{-1}(Z)$이다.
때때로 $Z$를 블로업의 {\it 중심}이라고 한다.
\end{definition}

\noindent
뒤에서 예외 인자가 유효 카르티에 인자임을 볼 것이다.
더 나아가 블로업은 $X$ 위의 스킴 가운데 $Z$의 역상이 유효
카르티에 인자인 “가장 작은” 스킴이라는 성질로 특징지어진다.

\medskip\noindent
$b : X' \to X$가 $X$의 $Z$에서 취한 블로업이면 구조층의 꼬임을
흔히 $\mathcal{O}_{X'}(n)$으로 나타낸다. 이들은 가역
$\mathcal{O}_{X'}$-가군이고
$\mathcal{O}_{X'}(n) = \mathcal{O}_{X'}(1)^{\otimes n}$임을
유의하자. 실제로 $X'$는 준연접 등급 $\mathcal{O}_X$-대수의
상대 Proj이며, 이 대수는 차수 $1$에서 생성된다. Constructions, Lemma
\ref{constructions-lemma-apply-relative}를 보라.
$\mathcal{O}_{X'}(1)$은 $b$-상대적으로 매우 풍부함을 유의하자.
$b$가 유한형은커녕 준콤팩트일 필요도 없지만 이 명제는 성립한다.
실제로 $X'$에는 $\mathbf{P}(\mathcal{I})$로 가는 닫힌 몰입이 주어져
있다. Morphisms, Example \ref{morphisms-example-very-ample}을 보라.

\begin{lemma}
\label{divisors-lemma-blowing-up-affine}
$X$를 스킴이라 하자. $\mathcal{I} \subset \mathcal{O}_X$를 준연접
아이디얼층이라 하자. $U = \Spec(A)$를 $X$의 아핀 열린
부분스킴이라 하고, $I \subset A$를 $\mathcal{I}|_U$에 대응하는
아이디얼이라 하자. $b : X' \to X$가 $X$의
$\mathcal{I}$에서 취한 블로업이면 표준 동형
$$
b^{-1}(U) = \text{Proj}(\bigoplus\nolimits_{d \geq 0} I^d)
$$
이 존재한다. 이는 $b^{-1}(U)$를 $I$의 $A$ 안에서의 리스 대수의
동차 스펙트럼과 동일시한다. 더 나아가 $b^{-1}(U)$에는 아핀
블로업 대수 $A[\frac{I}{a}]$들의 스펙트럼으로 이루어진 아핀
열린 덮개가 있다.
\end{lemma}

\begin{proof}
첫째 명제는 상대 Proj를 붙이기로 구성하는 데서 명백하다.
Constructions, Section
\ref{constructions-section-relative-proj-via-glueing}을 보라.
$a \in I$에 대하여 $a^{(1)}$로 원소 $a$를 나타내자. 이는 차수
$1$의 원소로 본 것이며 리스 대수 $\bigoplus_{n \geq 0} I^n$ 안에 있다.
이 원소들이 리스 대수를 $A$ 위에서 생성하므로
$\text{Proj}(\bigoplus_{d \geq 0} I^d)$는 아핀 열린집합
$D_{+}(a^{(1)})$들로 덮인다. 아핀 스킴 $D_{+}(a^{(1)})$는
아핀 블로업 대수 $A' = A[\frac{I}{a}]$의 스펙트럼이다.
Algebra, Definition \ref{algebra-definition-blow-up}을 보라.
이것으로 증명이 끝난다.
\end{proof}

\begin{lemma}
\label{divisors-lemma-flat-base-change-blowing-up}
\begin{slogan}
블로업은 평탄 밑변환과 가환한다.
\end{slogan}
$X_1 \to X_2$를 스킴의 평탄 사상이라 하자.
$Z_2 \subset X_2$를 닫힌 부분스킴이라 하자.
$Z_1$을 $Z_2$의 역상이라 하되, 이 역상은 $X_1$에서 취한다.
$X'_i$를 $Z_i$의 블로업이라 하되, 이는 $X_i$에서 취한다.
그러면 스킴의
카르테시안 도표
$$
\xymatrix{
X_1' \ar[r] \ar[d] & X_2' \ar[d] \\
X_1 \ar[r] & X_2
}
$$
가 존재한다.
\end{lemma}

\begin{proof}
$\mathcal{I}_2$를 $Z_2$의 아이디얼층이라 하되, 이는 $X_2$ 안에서
취한다.
주어진 사상을 $g : X_1 \to X_2$로 나타내자. 그러면
아이디얼층 $\mathcal{I}_1$은 $Z_1$의 아이디얼층이며
$g^*\mathcal{I}_2 \to \mathcal{O}_{X_1}$의 상이다
(역상의 정의에 의한다. Schemes, Definition
\ref{schemes-definition-inverse-image-closed-subscheme}을 보라).
Constructions, Lemma \ref{constructions-lemma-relative-proj-base-change}에
의해 $X_1 \times_{X_2} X_2'$는
$\bigoplus_{n \geq 0} g^*\mathcal{I}_2^n$의 상대 Proj이다.
$g$가 평탄이므로 사상
$g^*\mathcal{I}_2^n \to \mathcal{O}_{X_1}$은 단사이고 그 상은
$\mathcal{I}_1^n$이다. 따라서
$X_1 \times_{X_2} X_2' = X_1'$임을 알 수 있다.
\end{proof}

\begin{lemma}
\label{divisors-lemma-blowing-up-gives-effective-Cartier-divisor}
$X$를 스킴이라 하자. $Z \subset X$를 닫힌 부분스킴이라 하자.
블로업 $b : X' \to X$는 $Z$에서 $X$에 대해 취한 것이며
다음 성질들을 갖는다.
\begin{enumerate}
\item $b|_{b^{-1}(X \setminus Z)} : b^{-1}(X \setminus Z) \to X \setminus Z$
는 동형이다.
\item 예외 인자 $E = b^{-1}(Z)$는 $X'$ 위의 유효 카르티에 인자이다.
\item 표준 동형 $\mathcal{O}_{X'}(-1) = \mathcal{O}_{X'}(E)$이 존재한다.
\end{enumerate}
\end{lemma}

\begin{proof}
블로업은 열린 부분스킴으로의 제한과 가환하므로
(Lemma \ref{divisors-lemma-flat-base-change-blowing-up}), 첫째 명제는
$X' = X$가 $Z = \emptyset$일 때 성립한다는 뜻일 뿐이다.
이 경우 아이디얼층 $\mathcal{I} = \mathcal{O}_X$에서 블로업을 취하므로
Constructions, Example \ref{constructions-example-trivial-proj}에서
결론이 따른다.

\medskip\noindent
둘째 명제는 $X$ 위에서 국소적이므로 $X$가 아핀이라고 가정해도 된다.
$X = \Spec(A)$ 및 $Z = \Spec(A/I)$라 하자.
Lemma \ref{divisors-lemma-blowing-up-affine}에 의해 $X'$는 아핀 블로업 대수
$A' = A[\frac{I}{a}]$들의 스펙트럼으로 덮인다.
그러면 $IA' = aA'$이고, Algebra, Lemma
\ref{algebra-lemma-affine-blowup}에 의해 $a$는 $A'$에서 영인자가 아닌
원소로 간다. $Z$의 $\Spec(A')$ 안에서의 역상은
$\Spec(A'/IA') \subset \Spec(A')$에 대응하므로 이로써 명제가 증명된다.

\medskip\noindent
표준 사상
$\psi_{univ, 1} : b^*\mathcal{I} \to \mathcal{O}_{X'}(1)$을 생각하자.
Constructions, Definition
\ref{constructions-definition-relative-proj} 뒤의 논의를 보라.
이 사상은 동형 $\mathcal{I}_E \to \mathcal{O}_{X'}(1)$을 통해
인수분해된다고 주장한다(이는 마지막 주장을 증명한다).
실제로 위에서 말한 블로업 대수에 대응하는 아핀 열린집합에서
$A' = A[\frac{I}{a}]$이고, $\psi_{univ, 1}$은 다음 $A'$-가군
사상에 대응한다.
$$
I \otimes_A A'
\longrightarrow
\left(\Big(\bigoplus\nolimits_{d \geq 0} I^d\Big)_{a^{(1)}}\right)_1
$$
여기서 $a^{(1)}$은 Algebra, Definition
\ref{algebra-definition-blow-up}에서와 같다.
이 사상이 $I \otimes_A A' \to IA' = aA'$임에 대한 확인은 생략한다.
\end{proof}

\begin{lemma}[블로업의 보편 성질]
\label{divisors-lemma-universal-property-blowing-up}
$X$를 스킴이라 하자. $Z \subset X$를 닫힌 부분스킴이라 하자.
$\mathcal{C}$를 $(\Sch/X)$의 충만한 부분범주 가운데 다음 대상들로
이루어진 것으로 두자. 즉 $Y \to X$ 중에서 $Z$의 역상이 $Y$ 위의
유효 카르티에 인자인 것들이다. 그러면 블로업
$b : X' \to X$는 $Z$에서 $X$에 대해 취한 것이며
$\mathcal{C}$의 끝 대상이다.
\end{lemma}

\begin{proof}
Lemma \ref{divisors-lemma-blowing-up-gives-effective-Cartier-divisor}에 의해
$b : X' \to X$는 $\mathcal{C}$의 대상이다.
$f : Y \to X$를 $\mathcal{C}$의 대상이라 하자.
유일한 사상 $Y \to X'$가 존재하며 이것이 $X$ 위의 사상임을
보여야 한다.
$D = f^{-1}(Z)$라 하자. $\mathcal{I} \subset \mathcal{O}_X$를
$Z$의 아이디얼층이라 하고 $\mathcal{I}_D$를 $D$의 아이디얼층이라
하자. 그러면 $f^*\mathcal{I} \to \mathcal{I}_D$는 가역
$\mathcal{O}_Y$-가군으로 가는 전사이다. 이는 사상
$\psi : \bigoplus f^*\mathcal{I}^d \to \bigoplus \mathcal{I}_D^d$로
확장되며, 이것은 등급 $\mathcal{O}_Y$-대수의 사상이다.
($\mathcal{I}_D^d = \mathcal{I}_D^{\otimes d}$임을 관찰하자.
이는 $D$가 유효 카르티에 인자이기 때문이다.)
Constructions, Section \ref{constructions-section-relative-proj}의
내용에 의해 삼중항 $(1, f : Y \to X, \psi)$는 사상
$Y \to X'$를 정의하며, 이는 $X$ 위의 사상이다. 제한
$$
Y \setminus D \longrightarrow X' \setminus b^{-1}(Z) = X \setminus Z
$$
은 유일하다. Lemma \ref{divisors-lemma-complement-effective-Cartier-divisor}에
의해 열린집합 $Y \setminus D$는 $Y$에서 스킴론적으로 조밀하다.
따라서 Morphisms, Lemma
\ref{morphisms-lemma-equality-of-morphisms}에 의해 사상
$Y \to X'$는 유일하다
(또한 Constructions, Lemma
\ref{constructions-lemma-relative-proj-separated}에 의해 $b$는
분리이다).
\end{proof}

\begin{lemma}
\label{divisors-lemma-characterize-affine-blowup}
$b : X' \to X$를 스킴 $X$의 닫힌 부분스킴 $Z$를 따른 블로업이라
하자. $U = \Spec(A)$를 $X$의 아핀 열린집합이라 하고,
$I \subset A$를 $Z \cap U$에 대응하는 아이디얼이라 하자.
$a \in I$라 하고 $x' \in X'$를 $U$의 한 점으로 가는 점이라 하자.
그러면 $x'$가 아핀 열린집합
$U' = \Spec(A[\frac{I}{a}])$의 점일 필요충분조건은
$a$의 상이 $\mathcal{O}_{X', x'}$에서 예외 인자를 잘라 내는 것이다.
\end{lemma}

\begin{proof}
$U'$ 위의 예외 인자는 $a$의 상에 의해 잘려 나오며, 이 상은
$A' = A[\frac{I}{a}]$ 안에서 취한다. 따라서
한쪽 방향은 명백하다. 역으로 $a$의 상이
$\mathcal{O}_{X', x'}$에서 $E$를 잘라 낸다고 가정하자.
$I$의 모든 원소는 $E$를 정의하는 아이디얼의 원소로 가며, 이
아이디얼은 $b^{-1}(U)$ 위에서 취한다. 따라서 $I$의 원소들은
$a$로 나누어지며, 이는 $\mathcal{O}_{X', x'}$에서 성립한다.
그러므로 $f \in I^n$에 대하여
$f = \psi(f) a^n$으로 쓸 수 있으며, 여기서 어떤
$\psi(f) \in \mathcal{O}_{X', x'}$를 택한다.
$a$가 $\mathcal{O}_{X', x'}$의 영인자가 아닌 원소로 가므로
$\psi(f)$는 이 성질에 의해 유일하게 특징지어짐을 관찰하자.
이제 다음을 정의한다.
$$
A' \longrightarrow \mathcal{O}_{X', x'},\quad
f/a^n \longmapsto \psi(f)
$$
여기서는 Algebra, Definition \ref{divisors-definition-blow-up} 뒤에 주어진
블로업 대수의 기술을 사용했다. 위에서 말한 유일성에 의해 이는
$A$-대수 준동형이다. 이로써 사상
$\Spec(\mathcal{O}_{X', x"}) \to \Spec(A') = U'$를 얻는다.
블로업의 보편 성질
(Lemma \ref{divisors-lemma-universal-property-blowing-up})에 의해 이는
$X'$ 위의 사상이고, 따라서 물론 $x' \in U'$이다.
\end{proof}

\begin{lemma}
\label{divisors-lemma-blow-up-effective-Cartier-divisor}
$X$를 스킴이라 하자. $Z \subset X$를 유효 카르티에 인자라 하자.
$X$의 $Z$에서 취한 블로업은 $X$의 항등 사상이다.
\end{lemma}

\begin{proof}
블로업의 보편 성질
(Lemma \ref{divisors-lemma-universal-property-blowing-up})에서 즉시 따른다.
\end{proof}

\begin{lemma}
\label{divisors-lemma-blow-up-reduced-scheme}
$X$를 스킴이라 하자. $\mathcal{I} \subset \mathcal{O}_X$를 준연접
아이디얼층이라 하자. $X$가 축약이면 블로업 $X'$도 축약이다.
여기서 이는 $X$의 $\mathcal{I}$에서 취한 블로업이다.
\end{lemma}

\begin{proof}
Lemma \ref{divisors-lemma-blowing-up-affine}와 Algebra, Lemma
\ref{algebra-lemma-blowup-reduced}를 결합하면 된다.
\end{proof}

\begin{lemma}
\label{divisors-lemma-blow-up-integral-scheme}
$X$를 스킴이라 하자. $\mathcal{I} \subset \mathcal{O}_X$를 영이 아닌
준연접 아이디얼층이라 하자. $X$가 정역이면 블로업 $X'$도 정역이다.
여기서 이는 $X$의 $\mathcal{I}$에서 취한 블로업이다.
\end{lemma}

\begin{proof}
Lemma \ref{divisors-lemma-blowing-up-affine}와 Algebra, Lemma
\ref{algebra-lemma-blowup-domain}을 결합하면 된다.
\end{proof}

\begin{lemma}
\label{divisors-lemma-blow-up-and-irreducible-components}
$X$를 스킴이라 하자. $Z \subset X$를 닫힌 부분스킴이라 하자.
$b : X' \to X$를 $X$의 블로업이라 하되, 이는 $Z$를 따라 취한다.
그러면
$b$는 $X'$의 기약 성분들의 일반점 집합에서 $X$의 기약 성분들의
일반점 가운데 $Z$에 속하지 않는 것들의 집합으로 가는 전단사
사상을 유도한다.
\end{lemma}

\begin{proof}
예외 인자 $E \subset X'$는 유효 카르티에 인자이고
$X' \setminus E \to X \setminus Z$는 동형이다. Lemma
\ref{divisors-lemma-blowing-up-gives-effective-Cartier-divisor}를 보라.
따라서 다음을 보이면 충분하다. 유효 카르티에 인자
$D \subset S$가 주어졌다고 하자. 여기서 스킴 $S$를 택한다.
그러면 $S$의 기약 성분의 일반점은 어느 것도
$D$에 속하지 않는다. 이를 보기 위해 $S$를 아핀 열린 덮개의
각 원소로 바꾸어도 된다. 따라서 Lemma
\ref{divisors-lemma-characterize-effective-Cartier-divisor}에 의해
$S = \Spec(A)$이고 $D = V(f)$라고 가정해도 된다. 여기서
$f \in A$는 영인자가 아니다. 이제 $f$가 어떤 극소 소아이디얼
$\mathfrak p \subset A$에도 속하지 않음을 보여야 한다.
그렇지 않다면 $f$는 $R_\mathfrak p$의 극대 아이디얼에 속하는
영인자가 아닌 원소로 갈 텐데, 이는 Algebra, Lemma
\ref{algebra-lemma-minimal-prime-reduced-ring}과 모순이다.
\end{proof}

\begin{lemma}
\label{divisors-lemma-blow-up-pullback-effective-Cartier}
$X$를 스킴이라 하자. $b : X' \to X$를 $X$의 한 닫힌
부분스킴에서 취한 블로업이라 하자. 당김 $b^{-1}D$는 모든
유효 카르티에 인자 $D \subset X$에 대하여 정의되고, $b$에 대한
유리형 함수의 당김도 정의된다
(Definitions \ref{divisors-definition-pullback-effective-Cartier-divisor} 및
\ref{divisors-definition-pullback-meromorphic-sections}).
\end{lemma}

\begin{proof}
Lemmas \ref{divisors-lemma-blowing-up-affine}와
\ref{divisors-lemma-characterize-effective-Cartier-divisor}에 의해 다음
대수적 사실로 환원된다.
$A$를 환, $I \subset A$를 아이디얼, $a \in I$를 원소,
$x \in A$를 영인자가 아닌 원소라 하자. 그러면 $x$의 상은
$A[\frac{I}{a}]$에서 영인자가 아니다.
실제로 $x (y/a^n) = 0$이라고 가정하자. 이 등식은
$A[\frac{I}{a}]$에서 성립한다. 그러면 $a^mxy = 0$이
$A$에서 어떤 $m$에 대하여 성립한다. 따라서 $a^my = 0$이다.
이는 $x$가 영인자가 아니기 때문이다. 그러므로 원하는 대로 $y/a^n$은
$A[\frac{I}{a}]$에서 영이다.
\end{proof}

\begin{lemma}
\label{divisors-lemma-blowing-up-two-ideals}
$X$를 스킴이라 하자. $\mathcal{I}, \mathcal{J} \subset \mathcal{O}_X$를
준연접 아이디얼층이라 하자. $b : X' \to X$를 $X$의
$\mathcal{I}$에서 취한 블로업이라 하자. $b' : X'' \to X'$를
$X'$의 $b^{-1}\mathcal{J} \mathcal{O}_{X'}$에서 취한 블로업이라 하자.
그러면 $X'' \to X$는 $X$의 $\mathcal{I}\mathcal{J}$에서 취한
블로업과 표준적으로 동형이다.
\end{lemma}

\begin{proof}
$E \subset X'$를 $b$의 예외 인자라 하자. 이는 Lemma
\ref{divisors-lemma-blowing-up-gives-effective-Cartier-divisor}에 의해 유효
카르티에 인자이다. 그러면 Lemma
\ref{divisors-lemma-blow-up-pullback-effective-Cartier}에 의해
$(b')^{-1}E$는 $X''$ 위의 유효 카르티에 인자이다.
$E' \subset X''$를 $b'$의 예외 인자라 하자(이 역시 유효 카르티에
인자이다). 유효 카르티에 인자 $E'' = E' + (b')^{-1}E$를 생각하자.
구성상 $E''$의 아이디얼은
$(b \circ b')^{-1}\mathcal{I} (b \circ b')^{-1}\mathcal{J} \mathcal{O}_{X''}$이다.
따라서 Lemma \ref{divisors-lemma-universal-property-blowing-up}에 의해
$X''$에서 블로업 $c : Y \to X$로 가는 표준 사상이 있다.
이는 $X$의 $\mathcal{I}\mathcal{J}$에서의 블로업이다.
역으로 $\mathcal{I}\mathcal{J}$는 가역 아이디얼로
당겨지므로 $c^{-1}\mathcal{I}\mathcal{O}_Y$가 유효 카르티에
인자를 정의함을 알 수 있다. Lemma
\ref{divisors-lemma-sum-closed-subschemes-effective-Cartier}를 보라.
따라서 Lemma \ref{divisors-lemma-universal-property-blowing-up}에 의해
사상 $c' : Y \to X'$가 $X$ 위에 존재한다.
그러면 $(c')^{-1}b^{-1}\mathcal{J}\mathcal{O}_Y = c^{-1}\mathcal{J}\mathcal{O}_Y$도
유효 카르티에 인자를 정의한다.
따라서 사상 $c'' : Y \to X''$가 $X'$ 위에 존재한다.
이 사상이 앞에서 구성한 사상 $X'' \to Y$의 역임에 대한 확인은
생략한다.
\end{proof}

\begin{lemma}
\label{divisors-lemma-blowing-up-projective}
$X$를 스킴이라 하자. $\mathcal{I} \subset \mathcal{O}_X$를 준연접
아이디얼층이라 하자. $b : X' \to X$를 $X$의 블로업이라 하되,
이는 아이디얼층 $\mathcal{I}$에서 취한다.
$\mathcal{I}$가 유한형이면 다음이 성립한다.
\begin{enumerate}
\item $b : X' \to X$는 사영 사상이다.
\item $\mathcal{O}_{X'}(1)$은 $b$-상대적으로 풍부한 가역층이다.
\end{enumerate}
\end{lemma}

\begin{proof}
등급 $\mathcal{O}_X$-대수의 전사
$$
\text{Sym}_{\mathcal{O}_X}^*(\mathcal{I})
\longrightarrow
\bigoplus\nolimits_{d \geq 0} \mathcal{I}^d
$$
는 Constructions, Lemma
\ref{constructions-lemma-surjective-generated-degree-1-map-relative-proj}를
통해 닫힌 몰입
$$
X' = \underline{\text{Proj}}_X (\bigoplus\nolimits_{d \geq 0} \mathcal{I}^d)
\longrightarrow
\mathbf{P}(\mathcal{I}).
$$
을 정의한다. 따라서 $b$는 사영이다. Morphisms, Definition
\ref{morphisms-definition-projective}을 보라.
둘째 명제는 예를 들어 상대적으로 풍부한 가역층의 특징짓기인
Morphisms, Lemma
\ref{morphisms-lemma-characterize-relatively-ample}에서 따른다.
일부 세부사항은 생략한다.
\end{proof}

\begin{lemma}
\label{divisors-lemma-composition-finite-type-blowups}
\begin{slogan}
블로업들의 합성은 블로업이다.
\end{slogan}
$X$를 준콤팩트 준분리 스킴이라 하자.
$Z \subset X$를 유한 표시인 닫힌 부분스킴이라 하자.
$b : X' \to X$를 중심이 $Z$인 블로업이라 하자.
$Z' \subset X'$를 유한 표시인 닫힌 부분스킴이라 하자.
$X'' \to X'$를 중심이 $Z'$인 블로업이라 하자.
유한 표시인 닫힌 부분스킴 $Y \subset X$가 존재하여 다음을 만족한다.
\begin{enumerate}
\item 집합론적으로 $Y = Z \cup b(Z')$이다.
\item 합성 $X'' \to X$는 $X$의 $Y$에서 취한 블로업과 동형이다.
\end{enumerate}
\end{lemma}

\begin{proof}
$Z \to X$가 유한 표시라는 조건은 $Z$가 유한형 준연접 아이디얼층
$\mathcal{I} \subset \mathcal{O}_X$에 의해 잘려 나온다는 뜻이다.
Morphisms, Lemma
\ref{morphisms-lemma-closed-immersion-finite-presentation}을 보라.
$\mathcal{A} = \bigoplus_{n \geq 0} \mathcal{I}^n$이라 써서
$X' = \underline{\text{Proj}}(\mathcal{A})$가 되게 하자.
Properties, Lemma
\ref{properties-lemma-quasi-coherent-finite-type-ideals}에 의해
$X \setminus Z$는 $X$의 준콤팩트 열린집합임을 유의하자.
$b^{-1}(X \setminus Z) \to X \setminus Z$가 동형이므로
(Lemma \ref{divisors-lemma-blowing-up-gives-effective-Cartier-divisor}),
같은 결과에 의해
$b^{-1}(X \setminus Z) \setminus Z'$는 $X'$의 준콤팩트 열린집합이다.
따라서 $U = X \setminus (Z \cup b(Z'))$는 $X$의 준콤팩트
열린집합이다. Lemma \ref{divisors-lemma-closed-subscheme-proj-finite-type}에
의해 $d > 0$과 유한형 $\mathcal{O}_X$-부분가군
$\mathcal{F} \subset \mathcal{I}^d$가 존재하여
$Z' = \underline{\text{Proj}}(\mathcal{A}/\mathcal{F}\mathcal{A})$이고
$\mathcal{I}^d/\mathcal{F}$의 지지가 $X \setminus U$에 포함된다.

\medskip\noindent
$\mathcal{F} \subset \mathcal{I}^d$는 $\mathcal{O}_X$-부분가군이므로
$\mathcal{F} \subset \mathcal{I}^d \subset \mathcal{O}_X$를
$X$ 위의 유한형 준연접 아이디얼층으로 생각할 수 있다.
혼동을 피하기 위해 이를 $\mathcal{J} \subset \mathcal{O}_X$로
나타내자. $\mathcal{I}^d / \mathcal{J}$와
$\mathcal{O}/\mathcal{I}^d$의 지지가 $X \setminus U$에 포함되므로
$V(\mathcal{J})$는 $X \setminus U$에 포함된다.
역으로 $\mathcal{J} \subset \mathcal{I}^d$이므로
$Z \subset V(\mathcal{J})$이다.
$X \setminus Z \cong X' \setminus b^{-1}(Z)$ 위에서 아이디얼층
$\mathcal{J}$는 $Z'$를 잘라 낸다(아래 표시식을 보라).
따라서 $V(\mathcal{J})$는 $Z \cup b(Z')$와 같다.
그러므로 집합론적으로
$V(\mathcal{I}\mathcal{J}) = Z \cup b(Z')$이기도 하다.
더 나아가 $\mathcal{I}\mathcal{J}$는 두 유한형 아이디얼의
곱이므로 유한형 아이디얼이다.
$X'' \to X$가 $X$의 $\mathcal{I}\mathcal{J}$에서 취한 블로업과
동형이라고 주장한다. 그러면 $Y = V(\mathcal{I}\mathcal{J})$로
놓음으로써 보조정리의 증명이 끝난다.

\medskip\noindent
먼저 $X$의 $\mathcal{I}\mathcal{J}$에서 취한 블로업은
$X'$의 $b^{-1}\mathcal{J} \mathcal{O}_{X'}$에서 취한 블로업과
같음을 상기하자. Lemma \ref{divisors-lemma-blowing-up-two-ideals}를 보라.
따라서 $X'$의 $b^{-1}\mathcal{J} \mathcal{O}_{X'}$에서 취한
블로업이 $X'$의 $Z'$에서 취한 블로업과 일치함을 보이면 충분하다.
우리는
$$
b^{-1}\mathcal{J} \mathcal{O}_{X'} = \mathcal{I}_E^d \mathcal{I}_{Z'}
$$
가 $X''$ 위의 아이디얼층으로서 성립함을 보일 것이다.
$\mathcal{I}_E^d$는 유효 카르티에 인자 $dE$를 잘라 내므로,
Lemmas \ref{divisors-lemma-blow-up-effective-Cartier-divisor}와
\ref{divisors-lemma-blowing-up-two-ideals}를 사용하면 이 등식으로 원하는
결론을 얻는다.

\medskip\noindent
표시한 아이디얼들의 등식을 보기 위해 국소적으로 작업해도 된다.
Lemma \ref{divisors-lemma-blowing-up-affine}에서와 같이 $A$, $I$,
$a \in I$라는 표기를 사용하자. 그러면 $\mathcal{F}$는
$R$-부분가군 $M \subset I^d$에 대응하며, 이것은
아이디얼 $J \subset R$로 동형으로 간다.
조건
$Z' = \underline{\text{Proj}}(\mathcal{A}/\mathcal{F}\mathcal{A})$는
$Z' \cap \Spec(A[\frac{I}{a}])$가 원소
$m/a^d$, $m \in M$들로 생성된 아이디얼에 의해 잘려 나온다는
뜻이다. 원소 $m \in M$이 함수 $f \in J$에 대응한다고 하자.
그러면 아핀 블로업 대수 $A' = A[\frac{I}{a}]$에서
$f = (a^dm)/a^d = a^d (m/a^d)$임을 알 수 있다.
따라서 등식이 성립한다.
\end{proof}

\section{엄밀 변환}
\label{divisors-section-strict-transform}

\noindent
이 절에서는 블로업 아래의 엄밀 변환을 간단히 논의한다.
$S$를 스킴이라 하고 $Z \subset S$를 닫힌 부분스킴이라 하자.
$b : S' \to S$를 $S$의 $Z$에서 취한 블로업이라 하고
$E \subset S'$로 예외 인자 $E = b^{-1}Z$를 나타내자.
이하에서는 흔히 스킴 $X$를 생각하되 이를 $S$ 위에 놓고 카르테시안 도표
$$
\xymatrix{
\text{pr}_{S'}^{-1}E \ar[r] \ar[d] &
X \times_S S' \ar[r]_-{\text{pr}_X} \ar[d]_{\text{pr}_{S'}} &
X \ar[d]^f \\
E \ar[r] & S' \ar[r] & S
}
$$
를 만든다. $E$가 유효 카르티에 인자이므로
(Lemma \ref{divisors-lemma-blowing-up-gives-effective-Cartier-divisor}),
$\text{pr}_{S'}^{-1}E \subset X \times_S S'$는 국소 주이다
(Lemma \ref{divisors-lemma-pullback-locally-principal}).
따라서 $\text{pr}_{S'}^{-1}E$의 $X \times_S S'$ 안에서의 여집합은
역콤팩트이다
(Lemma \ref{divisors-lemma-complement-locally-principal-closed-subscheme}).
그러므로 준연접 $\mathcal{O}_{X \times_S S'}$-가군
$\mathcal{G}$에 대하여 $\text{pr}_{S'}^{-1}E$에 지지된 절단들의
부분층은 준연접 부분가군이다. Properties, Lemma
\ref{properties-lemma-sections-supported-on-closed-subset}을 보라.
$\mathcal{G}$가 준연접 대수층이면, 예를 들어
$\mathcal{G} = \mathcal{O}_{X \times_S S'}$이면, 이 부분층은
$\mathcal{G}$의 아이디얼이다.

\begin{definition}
\label{divisors-definition-strict-transform}
위와 같이 $Z \subset S$와 $f : X \to S$가 주어졌다고 하자.
\begin{enumerate}
\item 준연접 $\mathcal{O}_X$-가군 $\mathcal{F}$가 주어졌을 때,
$\mathcal{F}$의 {\it 엄밀 변환}을 생각하자. 이는 $S$의
$Z$에서 취한 블로업에 대한 것으로, 몫 $\mathcal{F}'$이다.
구체적으로 $\text{pr}_X^*\mathcal{F}$를
$\text{pr}_{S'}^{-1}E$에 지지된 절단들의 부분가군으로 나눈
몫이다.
\item $X$의 {\it 엄밀 변환}은
$X' \subset X \times_S S'$인 닫힌 부분스킴이다. 이는
$\mathcal{O}_{X \times_S S'}$의 절단들 가운데
$\text{pr}_{S'}^{-1}E$에 지지된 것들로 이루어진 준연접
아이디얼에 의해 잘려 나온다.
\end{enumerate}
\end{definition}

\noindent
블로업을 따라 엄밀 변환을 취하는 것은 그 블로업에 사용한 닫힌
부분스킴에 의존함을 유의하자(사상 $S' \to S$에만 의존하는 것이
아니다). 이 개념은 $S$의 닫힌 부분스킴에 흔히 사용된다.
$X$의 엄밀 변환은 실제로 $X$의 블로업이다.

\begin{lemma}
\label{divisors-lemma-strict-transform}
Definition \ref{divisors-definition-strict-transform}의 상황에서 다음이 성립한다.
\begin{enumerate}
\item 엄밀 변환 $X'$는 $X$의 엄밀 변환이며, $X$의 닫힌 부분스킴 $f^{-1}Z$에서
취한 $X$의 블로업이다.
\item 준연접 $\mathcal{O}_X$-가군 $\mathcal{F}$에 대하여 엄밀 변환
$\mathcal{F}'$은 다음 층과 표준적으로 동형이다. 즉
$X' \to X \times_S S'$를 따라 $\mathcal{F}$의 엄밀 변환을
밀어낸 층이며, 그 엄밀 변환은 블로업 $X' \to X$에 대한 것이다.
\end{enumerate}
\end{lemma}

\begin{proof}
$X'' \to X$를 $X$의 $f^{-1}Z$에서 취한 블로업이라 하자.
블로업의 보편 성질
(Lemma \ref{divisors-lemma-universal-property-blowing-up})에 의해 가환 도표
$$
\xymatrix{
X'' \ar[r] \ar[d] & X \ar[d] \\
S' \ar[r] & S
}
$$
가 존재하고, 따라서 사상 $X'' \to X \times_S S'$가 존재한다.
그러므로 첫째 주장은 이 사상이 상이 $X'$인 닫힌 몰입이라는 것이다.
이 문제는 $X$ 위에서 국소적이다. 따라서 $X$와 $S$가 아핀이라고
가정해도 된다. $S = \Spec(A)$, $X = \Spec(B)$라 하고,
$Z$가 아이디얼 $I \subset A$에 의해 잘려 나온다고 하자.
$J = IB$라 두자. 사상
$B \otimes_A \bigoplus_{n \geq 0} I^n \to \bigoplus_{n \geq 0} J^n$은
닫힌 몰입 $X'' \to X \times_S S'$를 정의한다. Constructions,
Lemmas \ref{constructions-lemma-base-change-map-proj}와
\ref{constructions-lemma-surjective-graded-rings-generated-degree-1-map-proj}를
보라. 이 사상이 위에서 보편 성질로 구성한 사상과 같다는 확인은
생략한다.
Lemma \ref{divisors-lemma-blowing-up-affine}에 따라 $a \in I$를 고르자.
이는 아핀 열린집합 $\Spec(A[\frac{I}{a}]) \subset S'$에 대응한다.
$\Spec(A[\frac{I}{a}])$의 역상을 엄밀 변환 $X'$에서 취하자.
이는 $X$의 엄밀 변환이다. 그 역상은
다음 환의 스펙트럼이다.
$$
B' = (B \otimes_A A[\textstyle{\frac{I}{a}}])/a\text{-power-torsion}
$$
Properties, Lemma
\ref{properties-lemma-sections-supported-on-closed-subset}을 보라.
한편 $b \in J$를 $a$의 상이라 하면
$\Spec(B[\frac{J}{b}])$는 $\Spec(A[\frac{I}{a}])$의 역상이며,
이 역상은 $X''$에서 취한다.
Algebra, Lemma \ref{algebra-lemma-blowup-base-change}에 의해
열린집합 $\Spec(B[\frac{J}{b}])$는 열린 부분스킴
$\text{pr}_{S'}^{-1}(\Spec(A[\frac{I}{a}]))$로 동형으로 간다.
이 부분스킴은 $X'$의 열린 부분스킴이다.
따라서 $X'' \to X'$는 동형이다.

\medskip\noindent
위의 표기에서 $\mathcal{F}$가 $B$-가군 $N$에 대응한다고 하자.
$\mathcal{F}$의 엄밀 변환은
$B \otimes_A A[\frac{I}{a}]$-가군
$$
N' = (N \otimes_A A[\textstyle{\frac{I}{a}}])/a\text{-power-torsion}
$$
에 대응한다. Properties, Lemma
\ref{properties-lemma-sections-supported-on-closed-subset}을 보라.
$\mathcal{F}$의 엄밀 변환을 $X$의 $f^{-1}Z$에서 취한 블로업에
대하여 생각하면, 이는
$B[\frac{J}{b}]$-가군
$N \otimes_B B[\frac{J}{b}]/b\text{-power-torsion}$에 대응한다.
위와 정확히 같은 방법으로 이 두 가군이 동형임을 증명한다.
세부사항은 생략한다.
\end{proof}

\begin{lemma}
\label{divisors-lemma-strict-transform-flat}
Definition \ref{divisors-definition-strict-transform}의 상황에서 다음이 성립한다.
\begin{enumerate}
\item $X$가 $S$ 위 평탄하고 이것이 $Z$ 위에 놓이는 모든 점에서
성립하면,
$X$의 엄밀 변환은 밑변환 $X \times_S S'$와 같다.
\item $\mathcal{F}$를 준연접 $\mathcal{O}_X$-가군이라 하자.
$\mathcal{F}$가 $S$ 위 평탄하고 이것이 $Z$ 위에 놓이는 모든
점에서 성립하면, 엄밀 변환 $\mathcal{F}'$은
$\mathcal{F}$의 엄밀 변환이며 당김 $\text{pr}_X^*\mathcal{F}$와 같다.
\end{enumerate}
\end{lemma}

\begin{proof}
(2)를 증명할 것이다. 스킴의 $S$ 위 엄밀 변환의 정의에 의해
이로부터 (1)이 따르기 때문이다. 이 문제는 $X$ 위에서 국소적이다.
따라서 $S = \Spec(A)$, $X = \Spec(B)$라고 가정해도 되고,
$\mathcal{F}$가 $B$-가군 $N$에 대응한다고 가정해도 된다.
그러면 $\mathcal{F}'$은 열린집합
$\Spec(B \otimes_A A[\frac{I}{a}])$ 위에서 가군에 대응한다.
이 열린집합은 $X \times_S S'$의 열린집합이며, 그 가군은
$$
N' = (N \otimes_A A[\textstyle{\frac{I}{a}}])/a\text{-power-torsion}
$$
이다. Properties, Lemma
\ref{properties-lemma-sections-supported-on-closed-subset}을 보라.
따라서 $a$-멱 꼬임이 영임을 보여야 한다. 여기서 이는
$N \otimes_A A[\frac{I}{a}]$의 꼬임을 말한다.
$y \in N \otimes_A A[\frac{I}{a}]$이고
$a^n y = 0$이라고 하자. $\mathfrak q \subset B$가 소아이디얼이고
$a \not \in \mathfrak q$이면 $y$는
$(N \otimes_A A[\frac{I}{a}])_\mathfrak q$에서 영으로 간다.
한편 $a \in \mathfrak q$이면 $N_\mathfrak q$는 평탄
$A$-가군이고
$N_\mathfrak q \otimes_A A[\frac{I}{a}]
=(N \otimes_A A[\frac{I}{a}])_\mathfrak q$
에는 $a$-멱 꼬임이 없다
($A[\frac{I}{a}]$ 자체에 그런 꼬임이 없기 때문이다).
따라서 $y$는 이 국소화에서도 영으로 간다. Algebra, Lemma
\ref{algebra-lemma-characterize-zero-local}에 의해 $y$는 영이다.
\end{proof}

\begin{lemma}
\label{divisors-lemma-strict-transform-affine}
$S$를 스킴이라 하자. $Z \subset S$를 닫힌 부분스킴이라 하자.
$b : S' \to S$를 $Z$의 블로업이라 하되, 이는 $S$에서 취한다.
$g : X \to Y$를 $S$ 위의 스킴들의 아핀 사상이라 하자.
$\mathcal{F}$를 $X$ 위의 준연접층이라 하자.
$g' : X \times_S S' \to Y \times_S S'$를 $g$의 밑변환이라 하자.
$\mathcal{F}'$을 $\mathcal{F}$의 엄밀 변환이라 하되, 이는 $b$에
대하여 취한다.
그러면 $g'_*\mathcal{F}'$은 $g_*\mathcal{F}$의 엄밀 변환이다.
\end{lemma}

\begin{proof}
Cohomology of Schemes, Lemma
\ref{coherent-lemma-affine-base-change}에 의해
$g'_*\text{pr}_X^*\mathcal{F} = \text{pr}_Y^*g_*\mathcal{F}$임을
관찰하자. $\mathcal{K} \subset \text{pr}_X^*\mathcal{F}$를
$Z$의 역상에 지지된 절단들의 부분층이라 하자. 이 역상은
$X \times_S S'$에서 취한다.
Properties, Lemma
\ref{properties-lemma-push-sections-supported-on-closed-subset}에 의해
밀어내기 $g'_*\mathcal{K}$는
$\text{pr}_Y^*g_*\mathcal{F}$의 절단들의 부분층이며,
$Z$의 역상에 지지된다. 이 역상은 $Y \times_S S'$에서 취한다.
$g'$가 아핀이므로
(Morphisms, Lemma \ref{morphisms-lemma-base-change-affine})
$g'_*$는 완전하고, 따라서 결론을 얻는다.
\end{proof}

\begin{lemma}
\label{divisors-lemma-strict-transform-different-centers}
$S$를 스킴이라 하자. $Z \subset S$를 닫힌 부분스킴이라 하자.
$D \subset S$를 유효 카르티에 인자라 하자.
$Z' \subset S$를 $Z$와 $D$의 아이디얼층의 곱에 의해 잘려 나오는
닫힌 부분스킴이라 하자.
$S' \to S$를 $S$의 $Z$에서 취한 블로업이라 하자.
\begin{enumerate}
\item $S$의 $Z'$에서 취한 블로업은 $S' \to S$와 동형이다.
\item $f : X \to S$를 스킴의 사상이라 하고 $\mathcal{F}$를
준연접 $\mathcal{O}_X$-가군이라 하자. $\mathcal{F}$가
$f^{-1}D$에 지지된 영이 아닌 국소 절단을 갖지 않으면,
$\mathcal{F}$의 엄밀 변환을 $Z$에서 취한 블로업에 대하여
생각한 것은 $\mathcal{F}$의 엄밀 변환을 $S$의 $Z'$에서 취한
블로업에 대하여 생각한 것과 일치한다.
\end{enumerate}
\end{lemma}

\begin{proof}
첫째 명제는 Lemmas \ref{divisors-lemma-blowing-up-two-ideals}와
\ref{divisors-lemma-blow-up-effective-Cartier-divisor}를 결합하면 따른다.
Lemma \ref{divisors-lemma-blowing-up-affine}를 사용하면 둘째 명제는 다음
대수 문제로 옮겨진다. $A$를 환, $I \subset A$를 아이디얼,
$x \in A$를 영인자가 아닌 원소, $a \in I$를 원소라 하자.
$M$을 $A$-가군이라 하되 그 $x$-꼬임이 영이라고 하자.
보일 것은 $a$-멱 꼬임이
$M \otimes_A A[\frac{I}{a}]$에서 $xa$-멱 꼬임과 같다는 것이다.
그 이유는 사상 $A \to A[\frac{I}{a}]$의 핵과
여핵이 $a$-멱 꼬임이고, 따라서 $a$를 가역화하면 이 사상이
동형이 되기 때문이다. 그러므로
$M \to M \otimes_A A[\frac{I}{a}]$의 핵과 여핵도 $a$-멱
꼬임이다. 이로부터 결과가 따른다.
\end{proof}

\begin{lemma}
\label{divisors-lemma-strict-transform-composition-blowups}
$S$를 스킴이라 하자. $Z \subset S$를 닫힌 부분스킴이라 하자.
$b : S' \to S$를 중심이 $Z$인 블로업이라 하자.
$Z' \subset S'$를 닫힌 부분스킴이라 하자.
$S'' \to S'$를 중심이 $Z'$인 블로업이라 하자.
$Y \subset S$를 닫힌 부분스킴이라 하되, 집합론적으로
$Y = Z \cup b(Z')$이고 합성 $S'' \to S$가 $S$의 $Y$에서 취한
블로업과 동형이라고 하자. 이 상황에서 임의의 스킴 $X$가 주어지고
그것이 $S$ 위에 있다고 하자. 또한
$\mathcal{F} \in \QCoh(\mathcal{O}_X)$가 주어지면 다음이 성립한다.
\begin{enumerate}
\item $\mathcal{F}$의 엄밀 변환을 $S$의 $Y$에서 취한 블로업에
대하여 생각한 것은 다음 엄밀 변환과 같다. 블로업
$S'' \to S'$의 $Z'$에서 취한 엄밀 변환인데, 대상은
$\mathcal{F}$를 블로업 $S' \to S$에 대하여 엄밀 변환한 것이다.
이 마지막 블로업은 $S$의 $Z$에서 취한다.
\item $X$의 엄밀 변환을 $S$의 $Y$에서 취한 블로업에 대하여
생각한 것은 다음 엄밀 변환과 같다. 블로업
$S'' \to S'$의 $Z'$에서 취한 엄밀 변환인데, 대상은
$X$를 블로업 $S' \to S$에 대하여 엄밀 변환한 것이다.
이 마지막 블로업은 $S$의 $Z$에서 취한다.
\end{enumerate}
\end{lemma}

\begin{proof}
$\mathcal{F}'$을 $\mathcal{F}$의 엄밀 변환이라 하자. 이는 블로업
$S' \to S$에 대한 것이며, 이 블로업은 $S$의 $Z$에서 취한다.
$\mathcal{F}''$을 $\mathcal{F}'$의 엄밀 변환이라 하자. 이는 블로업
$S'' \to S'$에 대한 것이며, 이 블로업은 $S'$의 $Z'$에서 취한다.
$\mathcal{G}$를 $\mathcal{F}$의 엄밀 변환이라 하자. 이는 블로업
$S'' \to S$에 대한 것이며, 이 블로업은 $S$의 $Y$에서 취한다.
사상들에도 다음과 같이 이름을 붙이자.
$$
\xymatrix{
X \times_S S'' \ar[r]_q \ar[d]^{f''} &
X \times_S S' \ar[r]_p \ar[d]^{f'} &
X \ar[d]^f \\
S'' \ar[r] & S' \ar[r] & S
}
$$
정의에 의해 전사 $p^*\mathcal{F} \to \mathcal{F}'$과 전사
$q^*\mathcal{F}' \to \mathcal{F}''$가 있다. 이들은 $q^*$의
오른쪽 완전성에 의해 결합되어 전사
$(p \circ q)^*\mathcal{F} \to \mathcal{F}''$를 준다.
또한 전사 $(p \circ q)^*\mathcal{F} \to \mathcal{G}$가 있다.
따라서 이 두 전사의 핵이 같음을 증명하면 충분하다.

\medskip\noindent
전사 $p^*\mathcal{F} \to \mathcal{F}'$의 핵은
$(f \circ p)^{-1}Z$에 지지되므로, 이 사상은 그 여집합의 점들에서
동형이다. 따라서
$q^*p^*\mathcal{F} \to q^*\mathcal{F}'$의 핵은
$(f \circ p \circ q)^{-1}Z$에 지지된다.
$q^*\mathcal{F}' \to \mathcal{F}''$의 핵은
$(f' \circ q)^{-1}Z'$에 지지된다.
이를 결합하면 $(p \circ q)^*\mathcal{F} \to \mathcal{F}''$의
핵은
$(f \circ p \circ q)^{-1}Z \cup (f' \circ q)^{-1}Z' =
(f \circ p \circ q)^{-1}Y$에 지지된다.
$\mathcal{G}$의 구성에 의해 인수분해
$(p \circ q)^*\mathcal{F} \to \mathcal{F}'' \to \mathcal{G}$를
얻는다. 증명을 끝내려면 $\mathcal{F}''$이
$(f \circ p \circ q)^{-1}(Y) =
(f \circ p \circ q)^{-1}Z \cup (f' \circ q)^{-1}Z'$
에 지지된 영이 아닌 (국소) 절단을 갖지 않음을 보이면 충분하다.
이는 $\mathcal{F}'$에 Lemma
\ref{divisors-lemma-strict-transform-different-centers}를 적용하면 따른다.
여기서 이 층은 $X \times_S S'$ 위에 있고, 이는
$S'$ 위의 스킴이다. 닫힌 부분스킴은 $Z'$이고 유효 카르티에 인자는
$b^{-1}Z$이다.
\end{proof}

\begin{lemma}
\label{divisors-lemma-strict-transform-universally-injective}
Definition \ref{divisors-definition-strict-transform}의 상황이라 하자.
$$
0 \to \mathcal{F}_1 \to \mathcal{F}_2 \to \mathcal{F}_3 \to 0
$$
이 $X$ 위의 준연접층들의 완전열이고 임의의 밑변환
$T \to S$ 뒤에도 완전하다고 하자. 그러면 $\mathcal{F}_i'$의
엄밀 변환들을 임의의 블로업 $S' \to S$에 대하여 취해도 짧은 완전열
$0 \to \mathcal{F}'_1 \to \mathcal{F}'_2 \to \mathcal{F}'_3 \to 0$을
이룬다.
\end{lemma}

\begin{proof}
$S$와 $X$ 위에서 국소화하여 둘 다 아핀이라고 가정해도 된다.
그러면 $\mathcal{F}_i$를 $S$로 밀어내도 된다.
Lemma \ref{divisors-lemma-strict-transform-affine}을 보라.
우리의 블로업이 사상 $1 : S \to S$라고 가정해도 된다.
이는 유효 카르티에 인자 $D \subset S$에 결부된다.
그러면 대수로 옮기면 다음과 같다.
$A$가 환이고
$0 \to M_1 \to M_2 \to M_3 \to 0$이 $A$-가군들의 보편 완전열이라고
하자. $a\in A$라 하자. 그러면 열
$$
0 \to
M_1/a\text{-power torsion} \to
M_2/a\text{-power torsion} \to
M_3/a\text{-power torsion} \to 0
$$
도 완전하다. 실제로 마지막 사상의 전사성과 첫째 사상의 단사성은
즉시 알 수 있다. 문제는 가운데에서의 완전성이다.
$x \in M_2$가 $M_3/a\text{-power torsion}$에서 영으로 간다고 하자.
그러면 $y = a^n x \in M_1$이고, 이는 어떤 $n$에 대하여 성립한다.
그러면 $y$는 $M_2/a^nM_2$에서 영으로 간다.
$M_1 \to M_2$가 보편 단사이므로 $y$는
$M_1/a^nM_1$에서 영으로 간다. 따라서 $y = a^n z$이고, 여기서
어떤 $z \in M_1$을 택한다. 그러므로 $a^n(x - y) = 0$이다.
따라서 원하는 대로 $y$는 $x$의 류로 가며, 그 류는
$M_2/a\text{-power torsion}$에서 취한다.
\end{proof}

\section{허용 블로업}
\label{divisors-section-admissible-blowups}

\noindent
블로업을 조금 더 잘 제어하기 위해 다음 표준 용어를 도입한다.

\begin{definition}
\label{divisors-definition-admissible-blowup}
$X$를 스킴이라 하자. $U \subset X$를 열린 부분스킴이라 하자.
사상 $X' \to X$를 {\it $U$-허용 블로업}이라고 하는 것은 다음
조건을 만족할 때이다. 유한 표시 닫힌 몰입 $Z \to X$가 존재하고,
$Z$는 $U$와 서로소이며, $X'$는 $X$의 $Z$에서 취한 블로업과
동형이다.
\end{definition}

\noindent
$Z \to X$가 유한 표시일 필요충분조건은 아이디얼층
$\mathcal{I}_Z \subset \mathcal{O}_X$가 유한형인 것임을 상기하자.
Morphisms, Lemma
\ref{morphisms-lemma-closed-immersion-finite-presentation}을 보라.
특히 $U$-허용 블로업은 사영 사상이다.
Lemma \ref{divisors-lemma-blowing-up-projective}를 보라.
같은 사상을 낳는 중심이 여러 개일 수 있음을 유의하자.
따라서 요건은 $U$와 서로소이고 $X'$를 만들어 내는 어떤 중심이
존재한다는 것뿐이다.
끝으로 사상 $b : X' \to X$는 $U$ 위에서 동형이므로
(Lemma \ref{divisors-lemma-blowing-up-gives-effective-Cartier-divisor}),
흔히 표기를 남용하여 $U$를 $X'$의 열린 부분스킴으로도 생각한다.

\begin{lemma}
\label{divisors-lemma-composition-admissible-blowups}
\begin{slogan}
허용 블로업은 합성에 대하여 안정하다.
\end{slogan}
$X$를 준콤팩트 준분리 스킴이라 하자.
$U \subset X$를 준콤팩트 열린 부분스킴이라 하자.
$b : X' \to X$를 $U$-허용 블로업이라 하자.
$X'' \to X'$를 $U$-허용 블로업이라 하자.
그러면 합성 $X'' \to X$는 $U$-허용 블로업이다.
\end{lemma}

\begin{proof}
더 정밀한 Lemma \ref{divisors-lemma-composition-finite-type-blowups}에서
즉시 따른다.
\end{proof}

\begin{lemma}
\label{divisors-lemma-extend-admissible-blowups}
$X$를 준콤팩트 준분리 스킴이라 하자.
$U, V \subset X$를 준콤팩트 열린 부분스킴이라 하자.
$b : V' \to V$를 $U \cap V$-허용 블로업이라 하자.
그러면 $U$-허용 블로업 $X' \to X$가 존재하고, 그
$V$로의 제한은 $V'$이다.
\end{lemma}

\begin{proof}
$\mathcal{I} \subset \mathcal{O}_V$를 유한형 준연접 아이디얼층이라
하되, $V(\mathcal{I})$가 $U \cap V$와 서로소이고 $V'$가
$V$의 $\mathcal{I}$에서 취한 블로업과 동형이 되게 하자.
$\mathcal{I}' \subset \mathcal{O}_{U \cup V}$를 준연접
아이디얼층이라 하되, $U$로의 제한은 $\mathcal{O}_U$이고
$V$로의 제한은 $\mathcal{I}$가 되게 하자
(Sheaves, Section \ref{sheaves-section-glueing-sheaves}를 보라).
Properties, Lemma \ref{properties-lemma-extend}에 의해 유한형 준연접
아이디얼층 $\mathcal{J} \subset \mathcal{O}_X$가 존재하여
$U \cup V$로 제한하면 $\mathcal{I}'$이 된다.
보조정리가 따른다.
\end{proof}

\begin{lemma}
\label{divisors-lemma-dominate-admissible-blowups}
$X$를 준콤팩트 준분리 스킴이라 하자.
$U \subset X$를 준콤팩트 열린 부분스킴이라 하자.
$b_i : X_i \to X$, $i = 1, \ldots, n$을 $U$-허용 블로업들이라
하자. 다음을 만족하는 $U$-허용 블로업 $b : X' \to X$가 존재한다.
(a) $b$는 $X' \to X_i \to X$로 인수분해되며, 이는 각
$i = 1, \ldots, n$에 대하여 성립하고,
(b) 각 사상 $X' \to X_i$는 $U$-허용 블로업이다.
\end{lemma}

\begin{proof}
$\mathcal{I}_i \subset \mathcal{O}_X$를 유한형 준연접 아이디얼층이라
하되, $V(\mathcal{I}_i)$가 $U$와 서로소이고 $X_i$가
$X$의 $\mathcal{I}_i$에서 취한 블로업과 동형이 되게 하자.
$\mathcal{I} = \mathcal{I}_1 \cdot \ldots \cdot \mathcal{I}_n$이라
두고 $X'$를 $X$의 $\mathcal{I}$에서 취한 블로업이라 하자.
그러면 Lemma \ref{divisors-lemma-blowing-up-two-ideals}에 의해
$X' \to X$는 $b_i$를 통해 인수분해된다.
\end{proof}

\begin{lemma}
\label{divisors-lemma-separate-disjoint-opens-by-blowing-up}
\begin{slogan}
블로업으로 기약 성분들을 분리한다.
\end{slogan}
$X$를 준콤팩트 준분리 스킴이라 하자.
$U, V$를 $X$의 서로소인 준콤팩트 열린 부분스킴이라 하자.
그러면 $U \cup V$-허용 블로업 $b : X' \to X$가 존재하여
$X'$는 열린 부분스킴들의 서로소 합
$X' = X'_1 \amalg X'_2$이고
$b^{-1}(U) \subset X'_1$ 및 $b^{-1}(V) \subset X'_2$이다.
\end{lemma}

\begin{proof}
유한형 준연접 아이디얼층 $\mathcal{I}$와 각각 $\mathcal{J}$를
골라 $X \setminus U = V(\mathcal{I})$ 및 각각
$X \setminus V = V(\mathcal{J})$가 되게 하자. Properties, Lemma
\ref{properties-lemma-quasi-coherent-finite-type-ideals}를 보라.
그러면 집합론적으로 $V(\mathcal{I}\mathcal{J}) = X$이므로
$\mathcal{I}\mathcal{J}$는 국소 멱영 아이디얼층이다.
$\mathcal{I}$와 $\mathcal{J}$가 유한형이고 $X$가 준콤팩트이므로
어떤 $n > 0$가 존재하여
$\mathcal{I}^n \mathcal{J}^n = 0$이다.
$\mathcal{I}$를 $\mathcal{I}^n$으로, $\mathcal{J}$를
$\mathcal{J}^n$으로 바꾸며, 그렇게 하기로 한다. 따라서
$\mathcal{I} \mathcal{J} = 0$이다.
$b : X' \to X$를 $\mathcal{I} + \mathcal{J}$에서 취한 블로업이라
하자. 이는 $U \cup V$-허용이다. 실제로
$V(\mathcal{I} + \mathcal{J}) = X \setminus U \cup V$이다.
$X'$가 열린 부분스킴들의 서로소 합
$X' = X'_1 \amalg X'_2$이며
$b^{-1}\mathcal{I}|_{X'_2} = 0$과
$b^{-1}\mathcal{J}|_{X'_1} = 0$이 성립함을 보일 것이다.
그러면 보조정리가 증명된다.

\medskip\noindent
Lemma \ref{divisors-lemma-blowing-up-affine}에 있는 블로업의 기술을 사용할
것이다. $U = \Spec(A) \subset X$가 아핀 열린집합이고
$\mathcal{I}|_U$와 각각 $\mathcal{J}|_U$가 유한 생성 아이디얼
$I \subset A$와 각각 $J \subset A$에 대응한다고 하자.
그러면
$$
b^{-1}(U) = \text{Proj}(A \oplus (I + J) \oplus (I + J)^2 \oplus \ldots)
$$
이는 아핀 열린 부분집합 $A[\frac{I + J}{x}]$와
$A[\frac{I + J}{y}]$들로 덮인다. 여기서 각각
$x \in I$이고 $y \in J$이다.
$x \in I$는 $A[\frac{I + J}{x}]$에서 영인자가 아닌 원소이고
$IJ = 0$이므로 $J A[\frac{I + J}{x}] = 0$임을 알 수 있다.
$y \in J$는 $A[\frac{I + J}{y}]$에서 영인자가 아닌 원소이고
$IJ = 0$이므로 $I A[\frac{I + J}{y}] = 0$임을 알 수 있다.
더 나아가
$$
\Spec(A[\textstyle{\frac{I + J}{x}}]) \cap
\Spec(A[\textstyle{\frac{I + J}{y}}]) =
\Spec(A[\textstyle{\frac{I + J}{xy}}]) = \emptyset
$$
이다. 실제로 $xy$는 영인자가 아닌 원소이면서 동시에 영이다.
따라서 $b^{-1}(U)$는 다음 두 열린 부분스킴의 서로소 합이다.
첫째 $U_1$은 표준 열린집합
$\Spec(A[\frac{I + J}{x}])$들의 합집합으로 정의되며, 여기서
$x \in I$이다. 둘째 $U_2$는 아핀 열린집합
$\Spec(A[\frac{I + J}{y}])$들의 합집합이며, 여기서 $y \in J$이다.
$b^{-1}\mathcal{I}\mathcal{O}_{X'}$가 $U_2$ 위에서 영으로 제한되고
$b^{-1}\mathcal{I}\mathcal{O}_{X'}$가 $U_1$ 위에서 영으로
제한됨을 보았다. 이 열린 부분스킴들이 전역 열린 부분스킴
$X'_1$과 $X'_2$로 붙는다는 확인은 생략한다.
\end{proof}

\begin{lemma}
\label{divisors-lemma-blowing-up-denominators}
$X$를 국소 뇌터 스킴이라 하자.
$\mathcal{L}$을 가역 $\mathcal{O}_X$-가군이라 하자.
$s$를 $\mathcal{L}$의 정칙 유리형 절단이라 하자.
$U \subset X$를 $s$가 $\mathcal{L}$의 절단에 대응하며 이것이
$U$ 위에서 성립하는
최대 열린 부분스킴이라 하자.
분모 아이디얼에서 취한 블로업 $b : X' \to X$는 $s$의 분모
아이디얼에서 취하며
$U$-허용이다. 유효 카르티에 인자 $D \subset X'$와 동형
$$
b^*\mathcal{L} = \mathcal{O}_{X'}(D - E),
$$
이 존재한다. 여기서 $E \subset X'$는 예외 인자이며,
유리형 절단 $b^*s$는 이 동형을 통해 유리형 절단
$1_D \otimes (1_E)^{-1}$에 대응한다.
\end{lemma}

\begin{proof}
Definition
\ref{divisors-definition-regular-meromorphic-ideal-denominators}에 있는 분모
아이디얼의 정의에서 $b$가 $U$-허용 블로업임을 즉시 알 수 있다.
표기 $1_D$, $1_E$ 및 $\mathcal{O}_{X'}(D - E)$에 대해서는
Definition
\ref{divisors-definition-invertible-sheaf-effective-Cartier-divisor}를 보라.
당김 $b^*s$는 Lemmas
\ref{divisors-lemma-blow-up-pullback-effective-Cartier}와
\ref{divisors-lemma-meromorphic-sections-pullback}에 의해 정의된다.
따라서 보조정리의 명제는 의미가 있다.
마지막 주장은 다음과 같이 다시 해석할 수 있다. $b^*s$는
$b^*\mathcal{L}(E)$의 전역 정칙 절단이고 그 영점 스킴은 $D$이다.
이는 $D$를 유일하게 정의하므로 보조정리를 증명하기 위해
$X$와 $X'$ 위에서 아핀 국소적으로 작업해도 된다.
$X = \Spec(A)$가 아핀이고 $\mathcal{L} = \mathcal{O}_X$라고
가정하자. 그러면 $s$는 정칙 유리형 함수이고, 더 축소하여
$s = a'/a$라고 가정해도 되며, 여기서 $a', a \in A$는 영인자가
아닌 원소들이다.
$s$의 분모 아이디얼은
$I = \{x \in A \mid xa' \in aA\}$에 대응한다.
$X'$는 아핀 블로업 대수 $A' = A[\frac{I}{x}]$들의
스펙트럼으로 덮임을 상기하자. 여기서 $x \in I$이다
(Lemma \ref{divisors-lemma-blowing-up-affine}).
$x \in I$를 고정하고 $xa' = a a''$라고 쓰자. 여기서 어떤
$a'' \in A$를 택한다. 인자 $E \subset X'$는 $x \in A'$에 의해
잘려 나오며, 이는 $A'$의 스펙트럼 위에서 성립한다. 따라서
$1/x$는 $\mathcal{O}_{X'}(E)$의 생성원이며, 이는
$\Spec(A')$ 위에서 성립한다.
끝으로 $A'$의 전분수환에서 $a'/a = a''/x$이다.
따라서 $b^*s = a'/a$는 $\mathcal{O}_{X'}(E)$의 정칙 절단으로
제한되고, $\Spec(A')$ 위에서는 $a''/x$로 주어진다.
이것으로 증명이 끝난다.
($D \cap \Spec(A')$는 $a''$의 상에 의해 잘려 나오며, 이 상은
$A'$에서 취한다.)
\end{proof}

\section{블로업과 평탄성}
\label{divisors-section-blowup-flat}

\noindent
More on Algebra, Section \ref{more-algebra-section-blowup-flat}에서
시작한 논의를 계속한다. More on Flatness, Section
\ref{flat-section-blowup-flat}에서 더 많은 결과를 증명할 것이다.


\begin{lemma}
\label{divisors-lemma-strict-transform-blowup-fitting-ideal}
$S$를 스킴이라 하자. $\mathcal{F}$를 유한형 준연접
$\mathcal{O}_S$-가군이라 하자. $Z_k \subset S$를
$\text{Fit}_k(\mathcal{F})$에 의해 잘려 나오는 닫힌 부분스킴이라
하자. Section \ref{divisors-section-fitting-ideals}를 보라.
$S' \to S$를 $S$의 $Z_k$에서 취한 블로업이라 하고
$\mathcal{F}'$을 $\mathcal{F}$의 엄밀 변환이라 하자.
그러면 $\mathcal{F}'$은 국소적으로 $\leq k$개의 절단으로
생성될 수 있다.
\end{lemma}

\begin{proof}
$\mathcal{F}'$이 국소적으로 $\leq k$개의 절단으로 생성될 수 있을
필요충분조건은
$\text{Fit}_k(\mathcal{F}') = \mathcal{O}_{S'}$임을 상기하자.
Lemma \ref{divisors-lemma-fitting-ideal-generate-locally}를 보라.
따라서 이 보조정리는 More on Algebra, Lemma
\ref{more-algebra-lemma-blowup-fitting-ideal}를 옮긴 것이다.
\end{proof}

\begin{lemma}
\label{divisors-lemma-strict-transform-blowup-fitting-ideal-locally-free}
$S$를 스킴이라 하자. $\mathcal{F}$를 유한형 준연접
$\mathcal{O}_S$-가군이라 하자. $Z_k \subset S$를
$\text{Fit}_k(\mathcal{F})$에 의해 잘려 나오는 닫힌 부분스킴이라
하자. Section \ref{divisors-section-fitting-ideals}를 보라.
$\mathcal{F}$가 계수 $k$의 국소 자유이고 이것이 $S \setminus Z_k$ 위에서
가정하자. $S' \to S$를 $S$의 $Z_k$에서 취한 블로업이라 하고
$\mathcal{F}'$을 $\mathcal{F}$의 엄밀 변환이라 하자.
그러면 $\mathcal{F}'$은 계수 $k$의 국소 자유이다.
\end{lemma}

\begin{proof}
More on Algebra, Lemma
\ref{more-algebra-lemma-blowup-fitting-ideal-locally-free}를 옮기면 된다.
\end{proof}

\begin{lemma}
\label{divisors-lemma-blowup-fitting-ideal}
$X$를 스킴이라 하자. $\mathcal{F}$를 유한 표시
$\mathcal{O}_X$-가군이라 하자. $U \subset X$를 스킴론적으로
조밀한 열린집합이라 하고, $\mathcal{F}|_U$가 일정한 계수 $r$의
유한 국소 자유라고 하자. 그러면 다음이 성립한다.
\begin{enumerate}
\item 블로업 $b : X' \to X$는 $X$의 $r$번째 피팅 아이디얼에서
취하며, 이 아이디얼은 $\mathcal{F}$의 피팅 아이디얼이다.
이 블로업은 $U$-허용이다.
\item 엄밀 변환 $\mathcal{F}'$은 $\mathcal{F}$의 엄밀 변환이며
$b$에 대하여 취한다. 이는
계수 $r$의 국소 자유이다.
\item 핵 $\mathcal{K}$는 전사 $b^*\mathcal{F} \to \mathcal{F}'$의
핵이고, 유한 표시이며 $\mathcal{K}|_U = 0$이다.
\item $b^*\mathcal{F}$와 $\mathcal{K}$는 완전
$\mathcal{O}_{X'}$-가군이며 그 토르 차원은
$\leq 1$이다.
\end{enumerate}
\end{lemma}

\begin{proof}
Lemma \ref{divisors-lemma-fitting-ideal-of-finitely-presented}에 의해 아이디얼
$\text{Fit}_r(\mathcal{F})$은 유한형이고, Lemma
\ref{divisors-lemma-fitting-ideal-finite-locally-free}에 의해 $U$로의 제한은
$\mathcal{O}_U$와 같다. 따라서 $b : X' \to X$는 $U$-허용이다.
Definition \ref{divisors-definition-admissible-blowup}을 보라.

\medskip\noindent
Lemma \ref{divisors-lemma-fitting-ideal-finite-locally-free}에 의해
$\text{Fit}_{r - 1}(\mathcal{F})$의 $U$로의 제한은 영이고,
$U$가 스킴론적으로 조밀하므로
$\text{Fit}_{r - 1}(\mathcal{F}) = 0$이 $X$ 전체에서 성립한다.
따라서 Lemma \ref{divisors-lemma-fitting-ideal-finite-locally-free}에 의해
$\mathcal{F}$는 계수 $r$의 국소 자유이다. 이는 $r$번째 피팅
아이디얼에 의해 잘려 나오는 부분스킴의 여집합에서 성립하며,
이 피팅 아이디얼은 $\mathcal{F}$의 피팅 아이디얼이다
(이 여집합은 $U$보다 클 수 있으며, 이것이 논증에서 이 단계를
거쳐야 한 이유이다). 따라서 Lemma
\ref{divisors-lemma-strict-transform-blowup-fitting-ideal-locally-free}에 의해
엄밀 변환
$$
b^*\mathcal{F} \longrightarrow \mathcal{F}'
$$
은 계수 $r$의 국소 자유이다. 이 사상의 핵 $\mathcal{K}$는
블로업 $b$의 예외 인자에 지지되므로 $\mathcal{K}|_U = 0$이다.
끝으로 $\mathcal{F}'$은 유한 국소 자유이고 표시된 화살표는
전사이므로, $X'$ 위에서 국소적으로 $b^*\mathcal{F}$를
$\mathcal{K}$와 $\mathcal{F}'$의 직합으로 쓸 수 있다.
$b^*\mathcal{F}'$은 유한 표시이므로
(Modules, Lemma \ref{modules-lemma-pullback-finite-presentation}),
$\mathcal{K}$도 유한 표시이다.

\medskip\noindent
토르 차원에 대한 명제는 More on Algebra, Lemma
\ref{more-algebra-lemma-fitting-ideals-and-pd1}에서 따른다.
\end{proof}

\section{수정}
\label{divisors-section-modifications}

\noindent
이 절에서는 “어떤 수정을 거친 뒤에는 이러저러한 것이 성립한다”는
형태의 결과들을 모은다. 뒤에서 수정은 블로업에 의해 지배될 수 있음을
볼 것이다(More on Flatness, Lemma
\ref{flat-lemma-dominate-modification-by-blowup}).

\begin{lemma}
\label{divisors-lemma-filter-after-modification}
$X$를 정역 스킴이라 하자. $\mathcal{E}$를 유한 국소 자유
$\mathcal{O}_X$-가군이라 하자. 수정 $f : X' \to X$가 존재하여
$f^*\mathcal{E}$는 연속 몫들이 가역 $\mathcal{O}_{X'}$-가군인
여과를 갖는다.
\end{lemma}

\begin{proof}
계수 $r$에 대한 귀납법으로 증명하며, 이는 $\mathcal{E}$의 계수이다.
$r = 1$이거나 $r = 0$이면 보조정리는 명백하다. $r > 1$이라고
가정하자. $P = \mathbf{P}(\mathcal{E})$라 하고 구조 사상을
$\pi : P \to X$라 하자. Constructions, Section
\ref{constructions-section-projective-bundle}을 보라.
그러면 $\pi$는 고유이다(Lemma \ref{divisors-lemma-relative-proj-proper}).
표준 전사
$$
\pi^*\mathcal{E} \to \mathcal{O}_P(1)
$$
가 있고 그 핵은 계수 $r - 1$의 유한 국소 자유이다.
공집합이 아닌 열린 부분스킴 $U \subset X$를 골라
$\mathcal{E}|_U \cong \mathcal{O}_U^{\oplus r}$이 되게 하자.
그러면 $P_U = \pi^{-1}(U)$는 $\mathbf{P}^{r - 1}_U$와 동형이다.
특히 절단 $s : U \to P_U$가 존재하며, 이는 $\pi$의 절단이다.
$X' \subset P$를 사상 $U \to P_U \to P$의 스킴론적 상이라 하자.
그러면 $X'$는 정역이고
(Morphisms, Lemma
\ref{morphisms-lemma-scheme-theoretic-image-reduced}),
사상 $f = \pi|_{X'} : X' \to X$는 고유이며
(Morphisms, Lemmas \ref{morphisms-lemma-closed-immersion-proper}와
\ref{morphisms-lemma-composition-proper}),
$f^{-1}(U) \to U$는 동형이다. 따라서 $f$는 수정이다
(Morphisms, Definition \ref{morphisms-definition-modification}).
구성상 당김 $f^*\mathcal{E}$는 두 단계 여과를 갖는다.
그 몫은 $\mathcal{O}_P(1)|_{X'}$와 같으므로 가역이고,
그 부분 $\mathcal{E}'$은 계수 $r - 1$의 국소 자유이다.
귀납 가정에 의해 수정 $g : X'' \to X'$를 찾아
$g^*\mathcal{E}'$이 보조정리의 명제와 같은 여과를 갖게 할 수 있다.
따라서 $f \circ g : X'' \to X$가 요구되는 수정이다.
\end{proof}

\begin{lemma}
\label{divisors-lemma-extend-rational-map-after-modification}
$S$를 스킴이라 하자. $X$, $Y$를 $S$ 위의 스킴들이라 하자.
$X$가 뇌터이고 $Y$가 $S$ 위 고유라고 가정하자.
$S$-유리 사상 $f : U \to Y$가 주어지고, 이것이 $X$에서
$Y$로 가는 유리 사상이라 하자. 그러면
사상 $p : X' \to X$와 $S$-사상 $f' : X' \to Y$가 존재하여
다음을 만족한다.
\begin{enumerate}
\item $p$는 고유이고 $p^{-1}(U) \to U$는 동형이다.
\item $f'|_{p^{-1}(U)}$는 $f \circ p|_{p^{-1}(U)}$와 같다.
\end{enumerate}
\end{lemma}

\begin{proof}
포함 사상을 $j : U \to X$로 나타내자.
$X' \subset Y \times_S X$를
$(f, j) : U \to Y \times_S X$의 스킴론적 상이라 하자
(Morphisms, Definition
\ref{morphisms-definition-scheme-theoretic-image}).
사영 $g : Y \times_S X \to X$는 고유이다
(Morphisms, Lemma \ref{morphisms-lemma-base-change-proper}).
합성 $p : X' \to X$를 생각하자. 이는
$X' \to Y \times_S X$와 $g$의 합성이며 고유이다
(Morphisms, Lemmas \ref{morphisms-lemma-closed-immersion-proper}와
\ref{morphisms-lemma-composition-proper}).
$g$가 분리이고 $U \subset X$가 역콤팩트이므로
($X$가 뇌터이기 때문이다), Morphisms, Lemma
\ref{morphisms-lemma-scheme-theoretic-image-of-partial-section}에 의해
$p^{-1}(U) \to U$는 동형이다.
한편 합성 $f' : X' \to Y$를 생각하자. 이는
$X' \to Y \times_S X$와 사영 $Y \times_S X \to Y$의 합성이다.
이 사상은 $f$와 일치하며, 이는 $p^{-1}(U)$ 위에서 성립한다.
\end{proof}
